\documentclass[oneside]{book}

% ---------- Core math ----------
\usepackage{amsmath, amsthm, amsfonts, amssymb}
\usepackage{slashed}
\usepackage{graphicx}
%\usepackage{mathabx}
\usepackage{mathtools}
%\usepackage{stmaryrd}
\usepackage{bm}
% ---------- Layout & tables ----------
\usepackage[a4paper, total={6in, 8in}]{geometry}
\usepackage{booktabs}
\usepackage{comment}
\usepackage{enumitem}  % instead of enumerate
\usepackage{longtable}

% ---------- TikZ-cd (externalize diagrams) ----------
\usepackage{tikz-cd}

% ---------- Misc ----------
\usepackage{cancel}
\usepackage{epigraph}
\usepackage{listings}
\usepackage{xcolor}
\usepackage{bbm}

% ---------- Hyperref last ----------
\usepackage[colorlinks=true, linkcolor=black, urlcolor=cyan, pdfpagemode=FullScreen]{hyperref}

% ---------- Theorems ----------
\newtheorem{theorem}{Theorem}[section]
\newtheorem{example}{Example}[section]
\newtheorem{corollary}{Corollary}[section]
\newtheorem{definition}{Definition}[section]
\newtheorem{lemma}{Lemma}[section]
\newtheorem{proposition}{Proposition}[section]
\newtheorem{remark}{Remark}[section]
\newtheorem{judgement}{Judgement}[section]
\newtheorem{construction}{Construction}[section]
\newtheorem{axiom}{Axiom}[section]
\newtheorem{notation}{Notation}[section]
\newtheorem{data}{Data}[section]
\newtheorem{axiomset}{Axiom Set}[section]
\newtheorem{hypothesis}{Hypothesis}[section]
\newtheorem{assumption}{Assumption}[section]
\newtheorem{convention}{Convention}[section]
\newtheorem{observation}{Observation}[section]
\newtheorem{problem}{Problem}[section]

% ---------- Bibliography ----------
\usepackage[backend=biber]{biblatex}
\addbibresource{references.bib}

% ---------- Title / Author / Date ----------
\title{Fuzzy Domain Theory via Enriched Category Theory}
\author{Conor Sheridan}
\date{\today}

\begin{document}
    \maketitle
    \tableofcontents

    \providecommand{\Down}{\mathrm{Down}}
\providecommand{\op}{\mathrm{op}}
\providecommand{\Pre}{\mathrm{Pre}}
\providecommand{\Ord}{\mathrm{Ord}}
\providecommand{\Pred}{\mathrm{Pred}}
\providecommand{\Luk}{\mathrm{Luk}}
\providecommand{\Sub}{\mathrm{Sub}}
\providecommand{\Cut}{\mathrm{Cut}}
\providecommand{\Psh}{\mathrm{Psh}}
\providecommand{\FCut}{\operatorname{FCut}}
\providecommand{\GCut}{\Gen^{\mathsf{cut}}}
\providecommand{\Gen}{\operatorname{G}}
\providecommand{\CutSub}{\mathrm{Sub}^{\mathrm{cut}}}
\providecommand{\Sup}{\mathsf{sup}}

\providecommand{\Cat}{\mathrm{Cat}}


\providecommand{\op}{\mathrm{op}}
\providecommand{\eqdef}{\mathrel{:=}}
\providecommand{\tensor}{\otimes}
\providecommand{\meet}{\wedge}
\providecommand{\join}{\vee}
\providecommand{\bigmeet}{\bigwedge}
\providecommand{\bigjoin}{\bigvee}
\providecommand{\yon}{\mathsf y}
\providecommand{\id}{\mathrm{id}}
\providecommand{\Fix}{\mathsf{Fix}}
\providecommand{\Free}{\mathsf{Free}}
\providecommand{\Down}{\mathsf{Down}}
\providecommand{\FFrm}{\mathbf{FFrm}}
\providecommand{\FLoc}{\mathbf{FLoc}}
\providecommand{\Opens}{\mathcal O}
\providecommand{\Set}{\mathbf{Set}}
\providecommand{\Cat}{\mathbf{Cat}}
\providecommand{\CAT}{\mathbf{CAT}}
\providecommand{\Pre}{\mathbf{Pre}}
\providecommand{\Ord}{\mathbf{Ord}}
\providecommand{\Frm}{\mathbf{Frm}}
\providecommand{\Loc}{\mathbf{Loc}}
\providecommand{\disc}{\mathrm{disc}}
\providecommand{\Ob}{\mathrm{Ob}}




    \part{Fuzzy Domain Theory}
    \chapter{Fuzzy Predicates and Cut Subobjects}
\label{ch:fuzzy-predicates-and-cut-subobjects}

\section{Introduction}
\label{sec:introduction-fuzzy-predicates-and-cut-subobjects}

This chapter develops the categorical framework for fuzzy predicates and their associated cut subobjects. The guiding idea is that a system of fuzzy truth values should be treated not merely as a set of degrees, but as an enriched base. A fuzzy predicate on a fuzzy preorder \(X\) is then a contravariant truth-valued enriched functor
\[
    X^{\mathrm{op}}\longrightarrow \Lambda,
\]
and the collection of all such predicates forms the enriched presheaf object
\[
    P_\Lambda(X)=[X^{\mathrm{op}},\Lambda].
\]

The chapter has four main purposes. First, we fix the base structures used throughout the manuscript: fuzzy cosmoses in the enriched-categorical setting and fuzzy truth-value objects in the posetal, quantale-enriched setting. Second, we construct tensor products, internal homs, and the truth-value fuzzy preorder, so that fuzzy predicates can be treated as ordinary enriched presheaves. Third, we show that the predicate object \(P_\Lambda(X)\) behaves as a fuzzy power object: it carries a generic fuzzy predicate and classifies indexed fuzzy predicates by the tensor--hom adjunction. Finally, we prove a cut-subobject representation theorem, identifying fuzzy predicates with coherent systems of level cuts.

In the crisp two-valued case, this construction reduces to the familiar correspondence between lower subsets of a preorder and their characteristic maps. For standard chain-valued bases such as the G\"odel, {\L}ukasiewicz, product, and Lawvere quantales, the same formalism yields the corresponding notions of fuzzy lower predicates, cost-valued potentials, and level-cut filtrations. The point of the enriched formulation is that these examples, including non-chain truth-value bases, are treated uniformly.

\section{Fuzzy cosmoses, truth values, and fuzzy preorders}
\label{sec:truth-values-and-fuzzy-preorders}

This section fixes the categorical and order-theoretic framework used throughout the manuscript. We begin with the enriched-categorical setting: a fuzzy cosmos is an affine symmetric monoidal closed base, regarded as a categorified object of fuzzy truth values. The underlying enriched category theory follows the standard conventions of Kelly \cite[Sections~1.1--1.6]{Kelly1982BasicConcepts}; the posetal specialization is the theory of quantale-enriched categories \cite[Definitions~2.1 and~2.2]{stubbe-2014-quantaloid-enriched}.

We then specialize to the posetal case, where the enriched structure becomes a quantale of fuzzy truth values. The resulting formalism covers ordinary preorders, standard chain-valued semantics for fuzzy logic, Lawvere metric spaces, and non-chain bases within a single framework; see also Lawvere's original treatment of metric spaces as enriched categories \cite{Lawvere1973MetricSpaces} and Stubbe's survey \cite[Examples~2.3 and~2.4]{stubbe-2014-quantaloid-enriched}.

\paragraph{Set-theoretic conventions.}
We work relative to a fixed tower of Grothendieck universes
\[
    \mathbb U_0\in \mathbb U_1\in \mathbb U_2.
\]
Unless otherwise specified, \textbf{small} means \(\mathbb U_0\)-small, \textbf{large} means \(\mathbb U_1\)-small, and all constructions are formed inside \(\mathbb U_2\).

All ordinary sets underlying fuzzy preorders are assumed small. The truth-value objects used in the posetal part of the manuscript are also assumed small. Thus, when \(L\) is the underlying set of truth values, joins such as
\[
    \bigvee_{a\in L} u_a
\]
are small joins.

A categorical base \(\mathbb V\) may be large, but is always assumed locally small: its ordinary hom-sets are small. Completeness and cocompleteness of a categorical base always mean existence of small limits and small colimits.

If \(X\) is a possibly large \(\mathbb V\)-category, we write
\[
    \mathcal P(X)
\]
for the \(\mathbb V\)-category of small \(\mathbb V\)-presheaves on \(X\), that is, those presheaves
\[
    X^{\mathrm{op}}\longrightarrow \mathbb V
\]
obtained as small weighted colimits of representables. Accordingly, the Yoneda embedding is always interpreted as a functor
\[
    Y:X\longrightarrow \mathcal P(X)
\]
into small presheaves. This is the universe-bounded version of the presheaf and Yoneda formalism for enriched categories \cite[Sections~1.6 and~1.9]{Kelly1982BasicConcepts}.

Accordingly, a coend indexed by a possibly large \(\mathbb V\)-category is used only when the corresponding weight is small. In particular, for the self-enriched truth-value object \(\underline{\mathbb V}\), an expression of the form
\[
    \int^{A\in \underline{\mathbb V}} F(A)\otimes A
\]
means the weighted colimit of the identity functor
\[
    \underline{\mathbb V}\longrightarrow \mathbb V
\]
by a small presheaf \(F\in\mathcal P(\underline{\mathbb V})\).

We shall also use the tensor--hom adjunction in the following universe-bounded form. If \(X\) and \(S\) are small \(\mathbb V\)-categories and \(Y\) is a locally small possibly large \(\mathbb V\)-category for which the functor category \([X,Y]\) is defined by small ends, then there is a natural bijection
\[
    \mathbb V\text{-}\mathbf{CAT}(S\otimes X,Y)
    \cong
    \mathbb V\text{-}\mathbf{CAT}(S,[X,Y]).
\]
Here \(\mathbb V\text{-}\mathbf{CAT}\) denotes the ordinary category whose objects are the locally small \(\mathbb V\)-categories in the chosen universe and whose morphisms are \(\mathbb V\)-functors. The notation
\[
    \mathbb V\text{-}\mathbf{Cat}
\]
will be reserved for the full subcategory on small \(\mathbb V\)-categories.

\subsection{Fuzzy cosmoses and enriched categories}
\label{subsec:fuzzy-cosmoses-and-enriched-categories}

We first record the categorical setting whose posetal specialization gives the fuzzy truth-value bases used below.

\begin{definition}[Fuzzy cosmos]\label{def:fuzzy-cosmos}
    A \textbf{fuzzy cosmos} is a locally small complete and cocomplete symmetric monoidal closed category
    \[
        (\mathbb V,\otimes,I,[-,-])
    \]
    such that the monoidal unit \(I\) is terminal. We call this condition \textbf{affinity}.
\end{definition}

The terminology is adapted to the present work. The standard background notions are those of symmetric monoidal categories and closed monoidal categories in Kelly \cite[Sections~1.1, 1.4, and~1.5]{Kelly1982BasicConcepts}. Since \(\mathbb V\) is symmetric monoidal closed, for every object \(A\in\mathbb V\) the functor
\[
    A\otimes -:\mathbb V\to\mathbb V
\]
is left adjoint to \([A,-]\). Hence \(A\otimes-\) preserves all small colimits. We shall refer to this consequence as \textbf{distributivity}. In the posetal case, it becomes the quantale distributivity law
\[
    a\otimes\Big(\bigvee_i b_i\Big)
    =
    \bigvee_i(a\otimes b_i),
\]
for small joins. The affinity condition categorifies integrality: in a truth-value quantale, the monoidal unit is the top element.

\begin{definition}[Self-enrichment]\label{def:self-enrichment}
    Let \(\mathbb V\) be a fuzzy cosmos. We write
    \[
        \underline{\mathbb V}
    \]
    for the \(\mathbb V\)-enriched category obtained by self-enriching \(\mathbb V\). Its objects are the objects of \(\mathbb V\), and its hom-objects are
    \[
        \underline{\mathbb V}(A,B):=[A,B].
    \]

    Composition is induced by the closed structure, using the evaluation maps
    \[
        [B,C]\otimes[A,B]\otimes A
        \longrightarrow
        [B,C]\otimes B
        \longrightarrow
        C,
    \]
    and identities are induced by the transposes of the unit isomorphisms
    \[
        I\otimes A\cong A.
    \]
\end{definition}

This is the standard self-enrichment of a symmetric monoidal closed category \cite[Section~1.6, especially formula~(1.28)]{Kelly1982BasicConcepts}. Thus \(\underline{\mathbb V}\) denotes the categorical analogue of the object of truth values. In the posetal case, its hom-object
\[
    \underline{\mathbb V}(a,b)
\]
is the residual implication
\[
    a\multimap b.
\]

\begin{remark}[Co-Yoneda reconstruction]\label{rem:canonical-density}
    For every \(B\in\mathbb V\), the enriched co-Yoneda formula gives a canonical reconstruction
    \[
        B
        \cong
        \int^{A\in\underline{\mathbb V}}
        \underline{\mathbb V}(A,B)\otimes A
        =
        \int^{A\in\underline{\mathbb V}}
        [A,B]\otimes A.
    \]
    This is the enriched density, or co-Yoneda, formula associated to the Yoneda embedding \cite[Section~1.9; Section~3.1, especially formula~(3.17); Section~4.2, formula~(4.31); Section~5.1]{Kelly1982BasicConcepts}. By the set-theoretic conventions above, this coend is interpreted as the weighted colimit of the identity functor
    \[
        \underline{\mathbb V}\to \mathbb V
    \]
    by the small representable presheaf
    \[
        Y(B)=\underline{\mathbb V}(-,B).
    \]
\end{remark}

We now recall the corresponding enriched notion of category.

\begin{definition}[Fuzzy category]\label{def:fuzzy-category}
    Let \(\mathbb V\) be a fuzzy cosmos. A \textbf{fuzzy category}, or \textbf{small \(\mathbb V\)-category}, \(X\) consists of:
    \begin{enumerate}
        \item a small set of objects, also denoted \(X\);

        \item for each \(x,y\in X\), a hom-object
        \[
            X(x,y)\in\mathbb V;
        \]

        \item identity morphisms
        \[
            I\longrightarrow X(x,x);
        \]

        \item composition morphisms
        \[
            X(y,z)\otimes X(x,y)\longrightarrow X(x,z);
        \]
    \end{enumerate}
    satisfying the usual associativity and unit axioms.
\end{definition}

This is the definition of a \(\mathbb V\)-category \cite[Section~1.2]{Kelly1982BasicConcepts}, specialized only by the choice of the affine base \(\mathbb V\).

\begin{definition}[Fuzzy functors and opposites]
\label{def:fuzzy-functor-opposite}
    Let \(X\) and \(Y\) be fuzzy categories.
    \begin{enumerate}
        \item A \textbf{fuzzy functor}, or \textbf{\(\mathbb V\)-functor}, \(F:X\to Y\) consists of a function on objects together with morphisms
        \[
            X(x,x')\longrightarrow Y(Fx,Fx')
        \]
        compatible with identities and composition.

        \item The \textbf{opposite} fuzzy category \(X^{\mathrm{op}}\) has the same objects as \(X\) and hom-objects
        \[
            X^{\mathrm{op}}(x,y):=X(y,x),
        \]
        with composition induced using the symmetry of \(\mathbb V\).
    \end{enumerate}
\end{definition}

The definition of \(\mathbb V\)-functor is standard \cite[Section~1.2]{Kelly1982BasicConcepts}; the construction of the opposite \(\mathbb V\)-category in the symmetric case is given in Kelly \cite[formula~(1.19)]{Kelly1982BasicConcepts}.

\subsection{Fuzzy truth values as posetal bases}
\label{subsec:fuzzy-truth-values}

We now specialize the fuzzy-cosmic setting to the case where the base \(\mathbb V\) is a posetal symmetric monoidal closed category. The underlying order-theoretic structure is a commutative integral quantale; compare Stubbe's definition of quantale
\cite[Definition~2.1]{stubbe-2014-quantaloid-enriched}. We impose complete distributivity as a standing structural condition so that the induced predicate algebras are frames and the level-cut constructions below interact well with arbitrary joins and meets.

\begin{definition}[Fuzzy truth-value object]
\label{def:fuzzy-truth-value-object}
    A \textbf{fuzzy truth-value object} is a small structure
    \[
        \Lambda=(L,\leq,\wedge,\vee,\otimes,\multimap,\bot,\top)
    \]
    such that:
    \begin{enumerate}
        \item \((L,\leq,\wedge,\vee,\bot,\top)\) is a completely distributive complete lattice;

        \item \((L,\otimes,\top)\) is a commutative monoid;

        \item for all \(a,b,c\in L\),
        \begin{equation}\label{eq:residuation-law}
            a\otimes b\leq c
            \quad\Longleftrightarrow\quad
            b\leq a\multimap c.
        \end{equation}
    \end{enumerate}
    Equivalently, \(\Lambda\) is a small completely distributive complete commutative integral residuated lattice, where integrality means that the monoidal unit coincides with the top element~\(\top\).
\end{definition}

The monoidal product \(\otimes\) encodes conjunction or composition of degrees, and the operation \(a\multimap c\) is the associated residual implication. Equivalently, for each fixed \(a\in L\), residuation says that there is an adjunction
\[
    \begin{tikzcd}[column sep=large]
        L \arrow[r,shift left=1.1ex,"a\otimes -"] &
        L \arrow[l,shift left=1.1ex,"a\multimap -"']
    \end{tikzcd}
    \qquad
    a\otimes - \dashv a\multimap - .
\]
This adjunction is the order-theoretic source of quantale distributivity.

\begin{proposition}[Posetal fuzzy cosmoses and fuzzy truth values]
\label{prop:thin-fuzzy-cosmoses-and-fuzzy-truth-values}
    Every fuzzy truth-value object \(\Lambda\) determines a posetal fuzzy cosmos. Its underlying category is the thin category associated to \((L,\leq)\), its tensor product is \(\otimes\), its unit is \(\top\), and its internal hom is the residual \(\multimap\).

    Conversely, every small posetal fuzzy cosmos whose object-poset is a completely distributive complete lattice determines a fuzzy truth-value object.
\end{proposition}

\begin{proof}
    Let \(\Lambda\) be a fuzzy truth-value object. Regard \(L\) as a thin category, with a unique morphism \(a\to b\) precisely when \(a\leq b\). Completeness and cocompleteness of this thin category are exactly completeness of the lattice \(L\). The monoidal structure is given by \(\otimes\), with unit \(\top\). Affinity follows from integrality, since \(\top\) is terminal in the thin category. The residuation law \eqref{eq:residuation-law} says exactly that the tensor is closed, with internal hom \(a\multimap b\). Since the monoidal structure is commutative, the resulting posetal category is symmetric monoidal closed.

    The self-enriched truth-value object \(\underline{\Lambda}\) has object set \(L\) and hom-values
    \[
        \underline{\Lambda}(a,b)=a\multimap b.
    \]
    The canonical levels are therefore the truth values \(a\in L\) themselves, represented by the Yoneda presheaves
    \[
        Y(a)=\underline{\Lambda}(-,a).
    \]
    No level system is chosen.

    Conversely, let \(\mathbb V\) be a small posetal fuzzy cosmos whose object-poset is a completely distributive complete lattice \(L\). The affine symmetric monoidal closed structure gives a commutative monoid \((L,\otimes,\top)\), and closedness gives a residual operation \(\multimap\) satisfying \eqref{eq:residuation-law}. Thus the underlying posetal symmetric monoidal closed category determines a fuzzy truth-value object.
\end{proof}

\begin{proposition}[Tensor preserves joins]
\label{prop:tensor-preserves-joins}
    For every \(a\in L\) and every small family \((b_i)_{i\in I}\) in
    \(L\),
    \[
        a\otimes \Big(\bigvee_{i\in I} b_i\Big)
        =
        \bigvee_{i\in I}(a\otimes b_i),
        \qquad
        \Big(\bigvee_{i\in I} b_i\Big)\otimes a
        =
        \bigvee_{i\in I}(b_i\otimes a).
    \]
    Consequently, every fuzzy truth-value object determines a commutative
    unital quantale in the sense of \cite[Definition~2.1]{stubbe-2014-quantaloid-enriched}.
\end{proposition}

\begin{proof}
    Fix \(a\in L\). By \eqref{eq:residuation-law}, the map
    \[
        a\otimes - : L\to L
    \]
    is left adjoint to
    \[
        a\multimap - : L\to L.
    \]
    Since left adjoints between complete lattices preserve joins,
    \[
        a\otimes \Big(\bigvee_{i\in I} b_i\Big)
        =
        \bigvee_{i\in I}(a\otimes b_i).
    \]
    The second identity follows by commutativity of \(\otimes\).
\end{proof}

Definition~\ref{def:fuzzy-truth-value-object} includes both the standard truth-value algebras of fuzzy logic and other familiar bases for degree-valued order. The use of \([0,1]\)-valued and \(L\)-valued truth degrees goes back to Zadeh and Goguen \cite{Zadeh1965FuzzySets,Goguen1967LFuzzySets}; the connection with left-continuous \(t\)-norms, residual implications, and quantale bases is standard in mathematical fuzzy logic \cite[Example~2.4]{stubbe-2014-quantaloid-enriched} \cite[Sections~1, 2, and~7]{sep-logic-fuzzy}.

\begin{example}[The crisp quantale]\label{ex:crisp-quantale}
    Let
    \[
        L=\{0<1\},
        \qquad
        a\otimes b:=a\wedge b,
        \qquad
        a\multimap b:=(a\Rightarrow b),
    \]
    where \(\Rightarrow\) denotes ordinary two-valued implication. With
    \[
        \bot=0,
        \qquad
        \top=1,
    \]
    the structure
    \[
        \Lambda=(L,\leq,\wedge,\vee,\otimes,\multimap,\bot,\top)
    \]
    is a fuzzy truth-value object. This is the two-valued case and recovers ordinary order theory, equivalently enrichment over the two-element Boolean quantale
    \cite[Definition~2.2 and the paragraph following it]{stubbe-2014-quantaloid-enriched}.
\end{example}

\begin{example}[Standard \(t\)-norm quantales]
\label{ex:standard-tnorm-quantale}
    Let \(\ast\) be a left-continuous \(t\)-norm on \([0,1]\), and set
    \[
        L=[0,1],
        \qquad
        a\otimes b:=a\ast b,
        \qquad
        a\multimap b:=\sup\{c\in[0,1]\mid a\ast c\leq b\},
        \qquad
        \bot:=0,
        \qquad
        \top:=1.
    \]
    The complete lattice \([0,1]\) is completely distributive because it is a complete chain. For a left-continuous \(t\)-norm, each map
    \[
        a\ast-:[0,1]\to[0,1]
    \]
    preserves arbitrary suprema in the complete chain \([0,1]\): if \(s=\sup_i b_i\), then monotonicity gives
    \[
        \sup_i(a\ast b_i)\leq a\ast s,
    \]
    while left-continuity gives the reverse inequality by approximating \(s\) from below. Hence \(a\ast-\) has the displayed right adjoint, namely its residuum \cite[Example~2.4]{stubbe-2014-quantaloid-enriched}. Therefore
    \[
        \Lambda=(L,\leq,\wedge,\vee,\otimes,\multimap,\bot,\top)
    \]
    is a fuzzy truth-value object. This yields the algebraic semantics of left-continuous \(t\)-norm logics \cite[Sections~1, 2, and~7]{sep-logic-fuzzy} \cite{cintula-hajek-noguera-2011-hmfl1}.
\end{example}

\begin{example}[The G\"odel quantale]\label{ex:godel-quantale}
    Let \(L=[0,1]\) with the order, and define
    \[
        a\otimes b:=a\wedge b,
        \qquad
        a\multimap b:=
        \begin{cases}
            1,&a\leq b,\\
            b,&a>b,
        \end{cases}
        \qquad
        \bot:=0,
        \qquad
        \top:=1.
    \]
    Then
    \[
        \Lambda=(L,\leq,\wedge,\vee,\otimes,\multimap,\bot,\top)
    \]
    is a fuzzy truth-value object. This is the quantale associated to the minimum \(t\)-norm \cite[Example~2.4]{stubbe-2014-quantaloid-enriched}. Composition is governed by the weaker of the two input degrees.
\end{example}

\begin{example}[The {\L}ukasiewicz quantale]
\label{ex:lukasiewicz-quantale}
    Let \(L=[0,1]\) with the order, and define
    \[
        a\otimes b:=\max\{0,a+b-1\},
        \qquad
        a\multimap b:=\min\{1,1-a+b\},
        \qquad
        \bot:=0,
        \qquad
        \top:=1.
    \]
    Then
    \[
        \Lambda=(L,\leq,\wedge,\vee,\otimes,\multimap,\bot,\top)
    \]
    is a fuzzy truth-value object. This is the quantale associated to the
    {\L}ukasiewicz \(t\)-norm \cite[Example~2.4]{stubbe-2014-quantaloid-enriched}. Composition is bounded-additive rather than bottleneck-type.
\end{example}

\begin{example}[The product quantale]\label{ex:product-quantale}
    Let \(L=[0,1]\) with the order, and define
    \[
        a\otimes b:=ab,
        \qquad
        a\multimap b:=
        \begin{cases}
            1,&a\leq b,\\[1mm]
            \dfrac{b}{a},&a>b,
        \end{cases}
        \qquad
        \bot:=0,
        \qquad
        \top:=1.
    \]
    Then
    \[
        \Lambda=(L,\leq,\wedge,\vee,\otimes,\multimap,\bot,\top)
    \]
    is a fuzzy truth-value object. This is the quantale associated to the product \(t\)-norm \cite[Example~2.4]{stubbe-2014-quantaloid-enriched}. Composition is
    multiplicative and therefore attenuates non-maximal degrees along chains.
\end{example}

\begin{example}[The nilpotent minimum quantale]
\label{ex:nilpotent-minimum-quantale}
    Let \(L=[0,1]\) with the order, and define
    \[
        a\otimes b:=
        \begin{cases}
            a\wedge b,&a+b>1,\\
            0,&a+b\leq 1,
        \end{cases}
        \qquad
        a\multimap b:=
        \begin{cases}
            1,&a\leq b,\\
            \max\{1-a,b\},&a>b,
        \end{cases}
    \]
    with
    \[
        \bot:=0,
        \qquad
        \top:=1.
    \]
    Then
    \[
        \Lambda=(L,\leq,\wedge,\vee,\otimes,\multimap,\bot,\top)
    \]
    is a fuzzy truth-value object. The operation \(\otimes\) is the nilpotent minimum \(t\)-norm, a standard left-continuous example \cite[Example~2.4, footnote~5]{stubbe-2014-quantaloid-enriched}. Here composition is possible only above a threshold of joint support.
\end{example}

\begin{example}[Finite many-valued chains]
\label{ex:finite-many-valued-chains}
    For each integer \(n\geq 2\), let
    \[
        L_n=
        \left\{
        0,\frac{1}{n-1},\frac{2}{n-1},\dots,\frac{n-2}{n-1},1
        \right\}
        \subseteq [0,1].
    \]
    Restricting either the G\"odel operations or the {\L}ukasiewicz operations from \([0,1]\) to \(L_n\) yields a finite fuzzy truth-value object
    \[
        \Lambda_n=(L_n,\leq,\wedge,\vee,\otimes,\multimap,\bot,\top).
    \]
\end{example}

\begin{example}[The Lawvere quantale]\label{ex:lawvere-quantale}
    Let \(L=[0,\infty]\) with the reverse of the usual order:
    \[
        a\leq_L b
        \quad:\Longleftrightarrow\quad
        a\geq b
        \quad\text{in the usual order.}
    \]
    Define
    \[
        a\otimes b:=a+b,
        \qquad
        a\multimap c:=c\ominus a
        :=
        \inf_{\mathrm{usual}}\{r\in[0,\infty]\mid a+r\geq c\},
        \qquad
        \top:=0,
        \qquad
        \bot:=\infty.
    \]
    Equivalently,
    \[
        c\ominus a=
        \begin{cases}
            0,&a\geq c,\\
            c-a,&a<c<\infty,\\
            \infty,&a<\infty,\ c=\infty.
        \end{cases}
    \]
    Thus \(a\multimap c\) is the truncated difference \(c-a\), computed with respect to the usual order, with truncation at \(0\) and with the usual extended-value conventions. The reversed order on a completely distributive complete chain is again completely distributive. Hence
    \[
        \Lambda=(L,\leq_L,\wedge,\vee,\otimes,\multimap,\bot,\top)
    \]
    is a fuzzy truth-value object. Here \(\wedge\) and \(\vee\) are computed with respect to the reversed order, so \(\wedge\) is the supremum and \(\vee\) is the infimum on \([0,\infty]\). This is Lawvere's base for generalized metric spaces \cite{Lawvere1973MetricSpaces} \cite[Example~2.3]{stubbe-2014-quantaloid-enriched}.

    Under the change of variables
    \[
        p=e^{-a},
        \qquad
        e^{-\infty}=0,
    \]
    this base is isomorphic, as a commutative integral residuated quantale, to the product quantale on \([0,1]\) \cite[Example~2.4]{stubbe-2014-quantaloid-enriched}. The compatibility of tensor products is expressed by the commutative diagram
    \[
        \begin{tikzcd}[column sep=large,row sep=large]
            \left[0,\infty\right]^{\mathrm{op}}\times [0,\infty]^{\mathrm{op}}
            \arrow[r,"+"]
            \arrow[d,"e^{-(-)}\times e^{-(-)}"']
            &
            \left[0,\infty\right]^{\mathrm{op}}
            \arrow[d,"e^{-(-)}"]
            \\
            [0,1]\times[0,1]
            \arrow[r,"\cdot"']
            &
            [0,1].
        \end{tikzcd}
    \]
    Thus addition of costs corresponds to multiplication of similarities.
\end{example}

\subsection{Fuzzy preorders}
\label{subsec:fuzzy-preorders}

We now specialize Definition~\ref{def:fuzzy-category} to the posetal fuzzy cosmos determined by a fuzzy truth-value object \(\Lambda\). In the fuzzy-set literature, closely related \(L\)-valued and quantale-valued preorder structures occur throughout the work on fuzzy preorders and fuzzy topologies; see, for example, \cite{LaiZhang2006FuzzyPreorderTopology,pu-zhang-2012-gl-monoid}.

\begin{definition}[Fuzzy preorder]\label{def:fuzzy-preorder}
    A \textbf{fuzzy preorder over \(\Lambda\)} is a pair \((X,\alpha)\) consisting of a small set \(X\) and a map
    \[
        \alpha:X\times X\to L
    \]
    such that for all \(x,y,z\in X\),
    \begin{align}
        \top &\leq \alpha(x,x),\label{eq:fuzzy-refl}\\
        \alpha(x,y)\otimes \alpha(y,z)&\leq \alpha(x,z).
        \label{eq:fuzzy-trans}
    \end{align}
    We usually suppress \(\alpha\) from the notation and write \(\alpha(x,y)\) simply as \(X(x,y)\).
\end{definition}

Since \(\top\) is the greatest element of \(L\), condition \eqref{eq:fuzzy-refl} is equivalent to
\[
    X(x,x)=\top
    \qquad
    \text{for all }x\in X.
\]
Thus \(X(x,y)\) records a degree-valued comparison whose interpretation depends on the chosen base \(\Lambda\).

\begin{remark}[Posetal enrichment]\label{rem:enriched-viewpoint}
    Definition~\ref{def:fuzzy-preorder} is precisely Definition~\ref{def:fuzzy-category} in the posetal case \(\mathbb V=\Lambda\). The identity morphism
    \[
        I\longrightarrow X(x,x)
    \]
    becomes the inequality
    \[
        \top\leq X(x,x),
    \]
    and composition becomes the transitivity inequality
    \[
        X(x,y)\otimes X(y,z)\leq X(x,z).
    \]
\end{remark}

\begin{definition}[\(\Lambda\)-functors, opposites, and fuzzy orders]
\label{def:lambda-functor-opposite-separated}
    Let \(X\) and \(Y\) be fuzzy preorders over \(\Lambda\).
    \begin{enumerate}[label=\textup{(\arabic*)}]
        \item A \textbf{\(\Lambda\)-functor} \(f:X\to Y\) is a function
        \(f:X\to Y\) such that
        \[
            X(x,x')\leq Y(fx,fx')
            \qquad
            \text{for all }x,x'\in X.
        \]

        \item The \textbf{opposite} fuzzy preorder \(X^{\mathrm{op}}\)
        has the same underlying set as \(X\) and hom-values
        \[
            X^{\mathrm{op}}(x,y):=X(y,x).
        \]

        \item A fuzzy preorder \(X\) is called \textbf{separated}\footnote{This is the separatedness condition for generalized metric spaces \cite[Example~2.3]{stubbe-2014-quantaloid-enriched}.} if
        \[
            X(x,y)=\top=X(y,x)
            \quad\Longrightarrow\quad
            x=y.
        \]
    \end{enumerate}
    A \textbf{fuzzy order over \(\Lambda\)} is a separated fuzzy preorder
    over \(\Lambda\).
\end{definition}

\begin{proposition}\label{prop:fuzzy-preorders-form-category}
    Fuzzy preorders over \(\Lambda\) and \(\Lambda\)-functors form a category, denoted
    \[
        \Pre_\Lambda.
    \]
    Moreover, fuzzy orders over \(\Lambda\) and \(\Lambda\)-functors form a full subcategory
    \[
        \Ord_\Lambda\subseteq \Pre_\Lambda.
    \]
\end{proposition}

\begin{proof}
    The identity map on a fuzzy preorder is clearly a \(\Lambda\)-functor. If
    \[
        f:X\to Y,
        \qquad
        g:Y\to Z
    \]
    are \(\Lambda\)-functors, then for all \(x,x'\in X\),
    \[
        X(x,x')\leq Y(fx,fx')\leq Z(g(fx),g(fx')).
    \]
    Hence \(g\circ f:X\to Z\) is again a \(\Lambda\)-functor. Associativity of composition and the identity laws are inherited from the category of sets.

    The second statement is immediate: \(\Ord_\Lambda\) is the full subcategory of \(\Pre_\Lambda\) on the separated objects.
\end{proof}

The inclusion of fuzzy orders into fuzzy preorders will be denoted by
\[
    \begin{tikzcd}
        \Ord_\Lambda \arrow[r,hook] & \Pre_\Lambda .
    \end{tikzcd}
\]

Once the base \(\Lambda\) has been fixed, the abstract definition specializes in familiar ways.

\begin{example}[Ordinary preorders]\label{ex:ordinary-preorders}
    For \(\Lambda\) equal to the crisp quantale of Example~\ref{ex:crisp-quantale}, a fuzzy preorder over \(\Lambda\) is precisely an ordinary preorder. Indeed, if we write
    \[
        x\preceq y
        \quad:\Longleftrightarrow\quad
        X(x,y)=1,
    \]
    then \eqref{eq:fuzzy-refl} and \eqref{eq:fuzzy-trans} become reflexivity and transitivity of \(\preceq\). This is the two-element case of quantale-enriched categories \cite[Definition~2.2 and the paragraph following it]{stubbe-2014-quantaloid-enriched}.
\end{example}

\begin{example}[Lawvere metric spaces]\label{ex:lawvere-metric-spaces}
    For the Lawvere quantale of Example~\ref{ex:lawvere-quantale}, all inequalities defining fuzzy preorders are interpreted in the reversed order \(\leq_L\). Translating them back to the usual order, a fuzzy preorder is exactly a Lawvere metric space, equivalently an extended quasi-pseudometric space
    \cite{Lawvere1973MetricSpaces} \cite[Example~2.3]{stubbe-2014-quantaloid-enriched}. Writing \(d(x,y)\) for \(X(x,y)\), the axioms become
    \[
        d(x,x)=0,
        \qquad
        d(x,z)\leq d(x,y)+d(y,z)
    \]
    in the usual order on \([0,\infty]\). In this case a \(\Lambda\)-functor is exactly a nonexpansive map:
    \[
        d_Y(fx,fx')\leq d_X(x,x').
    \]
\end{example}

\begin{example}[Chain-valued fuzzy preorders]
\label{ex:chain-valued-fuzzy-preorders}
    For the G\"odel, {\L}ukasiewicz, and product quantales of Examples~\ref{ex:godel-quantale}, \ref{ex:lukasiewicz-quantale}, and \ref{ex:product-quantale}, a fuzzy preorder is a map
    \[
        X:X\times X\to [0,1]
    \]
    satisfying \(X(x,x)=1\) together with, respectively,
    \[
        X(x,z)\geq \min\{X(x,y),X(y,z)\},
    \]
    \[
        X(x,z)\geq \max\{0,\,X(x,y)+X(y,z)-1\},
    \]
    or
    \[
        X(x,z)\geq X(x,y)\,X(y,z),
    \]
    for all \(x,y,z\in X\). Thus the same definition interpolates between bottleneck, bounded-additive, and multiplicative transitivity laws. Quantale-enriched categories over left-continuous \(t\)-norms are often described as fuzzy preorders relative to the chosen \(t\)-norm \cite[Example~2.4]{stubbe-2014-quantaloid-enriched}.
\end{example}

\subsection{Non-chain fuzzy truth values}
\label{subsec:non-chain-fuzzy-truth-values}

The preceding examples are chain-valued. Since one aim of this chapter is to work uniformly over arbitrary bases, it is useful to record that genuinely non-chain examples already arise at the level of truth values. The following examples are elementary finite or product constructions; they illustrate that the theory is not intrinsically tied to linearly ordered truth values, in the spirit of \(L\)-fuzzy set theory \cite{Goguen1967LFuzzySets}.

\begin{example}[A five-element non-chain truth-value algebra]
\label{ex:finite-nonchain-truth-value-algebra}
    Let \(P=\{p,q,r\}\) be the poset with
    \[
        p<r,
        \qquad
        q<r,
    \]
    and with \(p\) and \(q\) incomparable. Let \(L=\Down(P)\) be the set of downsets of \(P\):
    \[
        L=\bigl\{\varnothing,\{p\},\{q\},\{p,q\},\{p,q,r\}\bigr\},
    \]
    ordered by inclusion. Define
    \[
        U\otimes V:=U\cap V,
        \qquad
        U\multimap V:=\bigcup\{W\in L\mid U\cap W\subseteq V\},
        \qquad
        \bot:=\varnothing,
        \qquad
        \top:=P.
    \]
    Since unions of downsets are again downsets, the right-hand side indeed lies in \(L\). Moreover, the displayed formula is precisely the Heyting implication associated to the meet operation \(U\cap V\). Since every finite distributive lattice is completely distributive, \(L\) satisfies the lattice condition in Definition~\ref{def:fuzzy-truth-value-object}. Therefore
    \[
        \Lambda=(L,\subseteq,\cap,\cup,\otimes,\multimap,\bot,\top)
    \]
    is a fuzzy truth-value object. This is not a chain, since
    \[
        \{p\}\not\subseteq \{q\},
        \qquad
        \{q\}\not\subseteq \{p\}.
    \]
\end{example}

\begin{example}[A seven-element non-chain truth-value algebra]
\label{ex:seven-element-nonchain-truth-value-algebra}
    Let \(P=\{p,q,r,s\}\) be the poset with
    \[
        p<r,\qquad p<s,\qquad q<r,\qquad q<s,
    \]
    and with \(p\) and \(q\) incomparable and \(r\) and \(s\) incomparable. Let \(L=\Down(P)\) be the set of downsets of \(P\):
    \[
        L=
        \bigl\{
            \varnothing,\{p\},\{q\},\{p,q\},
            \{p,q,r\},\{p,q,s\},\{p,q,r,s\}
        \bigr\},
    \]
    ordered by inclusion. Define
    \[
        U\otimes V:=U\cap V,
        \qquad
        U\multimap V:=\bigcup\{W\in L\mid U\cap W\subseteq V\},
        \qquad
        \bot:=\varnothing,
        \qquad
        \top:=P.
    \]
    Again, this is the Heyting implication associated to the meet operation \(U\cap V\). Since every finite distributive lattice is completely distributive, \(L\) satisfies the lattice condition in Definition~\ref{def:fuzzy-truth-value-object}. Hence
    \[
        \Lambda=(L,\subseteq,\cap,\cup,\otimes,\multimap,\bot,\top)
    \]
    is a fuzzy truth-value object. This base is not a chain: for example,
    \[
        \{p\}\not\subseteq \{q\},
        \qquad
        \{q\}\not\subseteq \{p\},
    \]
    and also
    \[
        \{p,q,r\}\not\subseteq \{p,q,s\},
        \qquad
        \{p,q,s\}\not\subseteq \{p,q,r\}.
    \]
\end{example}

\begin{example}[A non-chain product of chain-based truth-value objects]
\label{ex:coordinatewise-product-truth-value-algebra}
    Let \(L=[0,1]\times [0,1]\), ordered coordinatewise:
    \[
        (a_1,a_2)\leq (b_1,b_2)
        \quad:\Longleftrightarrow\quad
        a_1\leq b_1
        \ \text{and}\
        a_2\leq b_2.
    \]
    Define the operations coordinatewise using the {\L}ukasiewicz formulas:
    \[
        (a_1,a_2)\otimes (b_1,b_2)
        :=
        \bigl(
        \max\{0,a_1+b_1-1\},
        \max\{0,a_2+b_2-1\}
        \bigr),
    \]
    \[
        (a_1,a_2)\multimap (b_1,b_2)
        :=
        \bigl(
            \min\{1,1-a_1+b_1\},
            \min\{1,1-a_2+b_2\}
        \bigr),
    \]
    with
    \[
        \bot:=(0,0),
        \qquad
        \top:=(1,1).
    \]
    Complete distributivity is inherited by finite products, since the complete distributivity law is checked coordinatewise. Therefore
    \[
        \Lambda=(L,\leq,\wedge,\vee,\otimes,\multimap,\bot,\top)
    \]
    is a fuzzy truth-value object, where \(\wedge\) and \(\vee\) are taken coordinatewise. This is not a chain, since
    \[
        (1,0)\nleq (0,1),
        \qquad
        (0,1)\nleq (1,0).
    \]
    Thus even when one begins with standard chain-based truth-value objects, their coordinatewise combination typically produces a non-chain fuzzy truth-value object.
\end{example}

\section{Tensor products, internal homs, and truth-value objects}
\label{sec:tensor-products-internal-homs-and-the-truth-value-object}

This section records the monoidal closed structure used throughout the manuscript. We first recall the standard tensor product and internal hom for categories enriched in a fuzzy cosmos, following the construction of the symmetric monoidal closed structure on \(\mathbb V\)-categories \cite[Section~1.4 and Sections~2.2--2.3]{Kelly1982BasicConcepts}. We then specialize these constructions to fuzzy preorders over a small fuzzy truth-value object \(\Lambda\). In that specialization, the general enriched formulas become the familiar pointwise formulas for tensor products, internal homs, and the truth-value fuzzy preorder.

The final corollary isolates the object
\[
    [X^{\mathrm{op}},\Lambda],
\]
which will be studied in the next section as the fuzzy preorder of \(\Lambda\)-valued contravariant predicates on \(X\).

\subsection{Tensor products and internal homs in a fuzzy cosmos}
\label{subsec:tensor-products-and-internal-homs-in-a-fuzzy-cosmos}

Let
\[
    (\mathbb V,\otimes,I,[-,-])
\]
be a fuzzy cosmos. We write
\[
    \mathbb V\text{-}\mathbf{Cat}
\]
for the ordinary category of small \(\mathbb V\)-categories and \(\mathbb V\)-functors. Since the object sets of the enriched categories under consideration are small and \(\mathbb V\) is locally small, the ordinary hom-sets of \(\mathbb V\text{-}\mathbf{Cat}\) are small. This is the underlying ordinary category associated to the enriched functor category formalism of Kelly \cite[Section~2.2, especially formula~(2.16)]{Kelly1982BasicConcepts}.

\begin{definition}[Tensor product and unit]
\label{def:v-tensor-product-and-unit}
    Let \(X\) and \(Y\) be small \(\mathbb V\)-categories. Their \textbf{tensor product}
    \[
        X\otimes Y
    \]
    is the \(\mathbb V\)-category whose objects are pairs \((x,y)\in X\times Y\), with hom-objects
    \[
        (X\otimes Y)\bigl((x,y),(x',y')\bigr)
        :=
        X(x,x')\otimes Y(y,y').
    \]
    The identity and composition morphisms are induced from those of \(X\) and \(Y\), using the unit, associativity, and symmetry of the tensor product in \(\mathbb V\).

    The \textbf{unit \(\mathbb V\)-category}
    \[
        \mathbf 1
    \]
    has one object \(\ast\) and hom-object
    \[
        \mathbf 1(\ast,\ast)=I.
    \]
\end{definition}

This is the tensor product of enriched categories over a symmetric monoidal base \cite[Section~1.4, formulas~(1.17)--(1.18)]{Kelly1982BasicConcepts}. The use of the symmetry in the definition of composition is the middle-four interchange construction.

\begin{definition}[Enriched internal hom]
\label{def:v-internal-hom}
    Let \(X\) and \(Y\) be small \(\mathbb V\)-categories. The \textbf{internal hom}, or enriched functor category,
    \[
        [X,Y]
    \]
    is the \(\mathbb V\)-category whose objects are \(\mathbb V\)-functors
    \[
        F:X\to Y,
    \]
    and whose hom-objects are given by the end formula \cite[Section~2.2, formula~(2.10)]{Kelly1982BasicConcepts}
    \[
        [X,Y](F,G)
        :=
        \int_{x\in X}Y(Fx,Gx).
    \]
    The end exists because \(X\) is small and \(\mathbb V\) is complete. Composition and identities are induced pointwise from those of \(Y\).
\end{definition}

The composition and identity maps are those of the enriched functor
category
\cite[Section~2.2, formulas~(2.14)--(2.15)]{Kelly1982BasicConcepts}.

There is a canonical evaluation \(\mathbb V\)-functor
\[
    \operatorname{ev}:[X,Y]\otimes X\longrightarrow Y,
    \qquad
    \operatorname{ev}(F,x):=F(x),
\]
as in the enriched functor-category construction
\cite[Section~2.2, formula~(2.17)]{Kelly1982BasicConcepts}.
The tensor--hom adjunction says that every \(\mathbb V\)-functor
\[
    h:Z\otimes X\to Y
\]
has a unique transpose
\[
    \widetilde h:Z\to [X,Y],
    \qquad
    \widetilde h(z)(x):=h(z,x),
\]
equivalently making the diagram
\[
\begin{tikzcd}[column sep=large,row sep=large]
Z\otimes X
    \arrow[rr,"h"]
    \arrow[dr,"\widetilde h\otimes 1_X"']
&&
Y
\\
&
\left[X,Y\right]\otimes X
    \arrow[ur,"\operatorname{ev}"']
\end{tikzcd}
\]
commute.

\begin{proposition}[Tensor--hom adjunction in a fuzzy cosmos]
\label{prop:v-tensor-hom-adjunction}
    For small \(\mathbb V\)-categories \(X,Y,Z\), there is a natural bijection
    \[
        \mathbb V\text{-}\mathbf{Cat}(Z\otimes X,Y)
        \cong
        \mathbb V\text{-}\mathbf{Cat}(Z,[X,Y]).
    \]
    Hence the ordinary category \(\mathbb V\text{-}\mathbf{Cat}\) is symmetric monoidal closed under the tensor product of Definition~\ref{def:v-tensor-product-and-unit}.
\end{proposition}

\begin{proof}
    This is the tensor--hom adjunction for enriched categories \cite[Section~2.3, especially formula~(2.20)]{Kelly1982BasicConcepts}. Explicitly, a \(\mathbb V\)-functor
    \[
        h:Z\otimes X\to Y
    \]
    is transposed to the \(\mathbb V\)-functor
    \[
        \widetilde h:Z\to [X,Y]
    \]
    defined on objects by
    \[
        \widetilde h(z)(x):=h(z,x).
    \]
    Conversely, a \(\mathbb V\)-functor
    \[
        k:Z\to [X,Y]
    \]
    is transposed to the \(\mathbb V\)-functor
    \[
        \check k:Z\otimes X\to Y
    \]
    defined by
    \[
        \check k(z,x):=k(z)(x).
    \]
    The enriched functor axioms for these two transposes are equivalent by the universal property of the end defining \([X,Y]\). The two assignments are inverse to one another, and naturality in \(X,Y,Z\) is immediate.
\end{proof}

\begin{remark}[Universe-bounded tensor--hom use]
\label{rem:universe-bounded-tensor-hom-use}
    Later we apply the same tensor--hom construction when the codomain is a locally small possibly large \(\mathbb V\)-category, such as the self-enriched truth-value object \(\underline{\mathbb V}\) or a small-presheaf category \(\mathcal P(\underline{\mathbb V})\). This is justified by the universe-bounded convention stated in Section~\ref{sec:truth-values-and-fuzzy-preorders}: all ends are indexed by small domains, and all hom-objects of the codomain are small enough for the relevant end formulas.
\end{remark}

\subsection{The posetal specialization}
\label{subsec:posetal-specialization-of-tensor-hom}

We now specialize the preceding constructions to the posetal fuzzy cosmos determined by a small fuzzy truth-value object \(\Lambda\).

Recall that if \(X\) is a fuzzy preorder, then its opposite \(X^{\mathrm{op}}\) has the same underlying set as \(X\) and hom-values
\[
    X^{\mathrm{op}}(x,y):=X(y,x).
\]

\begin{definition}[Tensor product and unit]
\label{def:tensor-product-and-unit}
    Let \(X\) and \(Y\) be fuzzy preorders over \(\Lambda\). Their \textbf{tensor product}
    \[
        X\otimes Y
    \]
    is the fuzzy preorder with underlying set \(X\times Y\) and hom-values
    \[
        (X\otimes Y)\bigl((x,y),(x',y')\bigr)
        :=
        X(x,x')\otimes Y(y,y').
    \]
    If
    \[
        f:X\to X',
        \qquad
        g:Y\to Y'
    \]
    are \(\Lambda\)-functors, define
    \[
        f\otimes g:X\otimes Y\to X'\otimes Y'
    \]
    by
    \[
        (f\otimes g)(x,y):=\bigl(f(x),g(y)\bigr).
    \]

    The \textbf{unit fuzzy preorder}
    \[
        1
    \]
    is the one-element fuzzy preorder with unique element \(\ast\) and hom-value
    \[
        1(\ast,\ast)=\top.
    \]
\end{definition}

This is the specialization of the enriched tensor product of Definition~\ref{def:v-tensor-product-and-unit} to the quantale-enriched case \cite[Section~1.4, formulas~(1.17)--(1.18)]{Kelly1982BasicConcepts}.

\begin{proposition}
\label{prop:tensor-product-gives-symmetric-monoidal-structure}
    The tensor product of Definition~\ref{def:tensor-product-and-unit} makes
    \[
        \Pre_\Lambda
    \]
    into a symmetric monoidal category with unit object \(1\).
\end{proposition}

\begin{proof}
    This is the posetal instance of the tensor product of enriched categories \cite[Section~1.4]{Kelly1982BasicConcepts}. Directly, reflexivity follows from
    \[
        \top=\top\otimes\top
        \leq
        X(x,x)\otimes Y(y,y),
    \]
    and transitivity follows from associativity, commutativity, monotonicity of \(\otimes\), and transitivity in \(X\) and \(Y\):
    \[
        \bigl(X(x,x')\otimes Y(y,y')\bigr)
        \otimes
        \bigl(X(x',x'')\otimes Y(y',y'')\bigr)
        \leq
        X(x,x'')\otimes Y(y,y'').
    \]
    The statement for tensor products of \(\Lambda\)-functors follows from monotonicity of \(\otimes\). The associativity, unit, and symmetry isomorphisms are induced by the corresponding bijections of cartesian products of sets.
\end{proof}

\begin{lemma}[Pullbacks in \(\Pre_\Lambda\)]
\label{lem:pullbacks-in-prelambda}
    The category \(\Pre_\Lambda\) has conical pullbacks. The pullback of
    \[
        X \xrightarrow{f} Z \xleftarrow{g} Y
    \]
    has underlying set the ordinary pullback
    \[
        P=\{(x,y)\in X\times Y\mid f(x)=g(y)\},
    \]
    and hom-values
    \[
        P\bigl((x,y),(x',y')\bigr)
        :=
        X(x,x')\wedge Y(y,y').
    \]
\end{lemma}

\begin{proof}
    Reflexivity follows from reflexivity in \(X\) and \(Y\). For transitivity, use monotonicity of \(\otimes\), the inequalities
    \[
        u\wedge v\leq u,
        \qquad
        u\wedge v\leq v,
    \]
    and transitivity in \(X\) and \(Y\):
    \[
        \begin{aligned}
        &
        \bigl(X(x,x')\wedge Y(y,y')\bigr)
        \otimes
        \bigl(X(x',x'')\wedge Y(y',y'')\bigr)
        \\
        &\qquad\leq
        \bigl(X(x,x')\otimes X(x',x'')\bigr)
        \wedge
        \bigl(Y(y,y')\otimes Y(y',y'')\bigr)
        \\
        &\qquad\leq
        X(x,x'')\wedge Y(y,y'').
        \end{aligned}
    \]
    The two projections are \(\Lambda\)-functors. Given \(T\to X\) and \(T\to Y\) with equal composites into \(Z\), the induced ordinary function \(T\to P\) is a \(\Lambda\)-functor because its hom condition is exactly the conjunction of the two hom conditions for \(T\to X\) and \(T\to Y\). This gives the usual pullback universal property.
\end{proof}

\subsection{The truth-value fuzzy preorder}
\label{subsec:the-truth-value-object}

The self-enriched base \(\underline{\mathbb V}\) specializes, in the posetal case, to the fuzzy preorder whose hom-values are residual implications. This is the thin, quantale-enriched instance of the self-enrichment of a symmetric monoidal closed category \cite[Section~1.6, especially formula~(1.28)]{Kelly1982BasicConcepts}.

\begin{definition}[The truth-value fuzzy preorder]
\label{def:truth-value-object-as-fuzzy-preorder}
    By abuse of notation, we also write
    \[
        \Lambda
    \]
    for the fuzzy preorder whose underlying set is \(L\) and whose
    hom-values are
    \[
        \Lambda(a,b):=a\multimap b.
    \]
    Equivalently, this is the self-enriched \(\Lambda\)-category \(\underline{\Lambda}\).
\end{definition}

\begin{proposition}
\label{prop:truth-value-object-is-fuzzy-preorder}
    The structure \(\Lambda\) of Definition~\ref{def:truth-value-object-as-fuzzy-preorder} is a separated fuzzy preorder.
\end{proposition}

\begin{proof}
    This is the posetal form of the self-enrichment of \(\mathbb V\) \cite[Section~1.6]{Kelly1982BasicConcepts}. Directly, residuation gives
    \[
        \top\leq a\multimap a
    \]
    because \(a\otimes \top=a\leq a\). For transitivity, the inequalities
    \[
        a\otimes(a\multimap b)\leq b,
        \qquad
        b\otimes(b\multimap c)\leq c
    \]
    imply
    \[
        a\otimes\bigl((a\multimap b)\otimes(b\multimap c)\bigr)\leq c,
    \]
    hence
    \[
        (a\multimap b)\otimes(b\multimap c)\leq a\multimap c.
    \]
    If \(a\multimap b=\top=b\multimap a\), then residuation gives \(a\leq b\) and \(b\leq a\), hence \(a=b\).
\end{proof}

\subsection{Internal homs}
\label{subsec:internal-homs}

The general enriched internal hom of Definition~\ref{def:v-internal-hom} specializes to the pointwise fuzzy preorder of \(\Lambda\)-functors.

\begin{definition}[Internal hom]
\label{def:internal-hom}
    Let \(X\) and \(Y\) be fuzzy preorders over \(\Lambda\). Define
    \[
        [X,Y]
    \]
    to be the fuzzy preorder whose underlying set is the set of \(\Lambda\)-functors \(f:X\to Y\), and whose hom-values are
    \[
        [X,Y](f,g)
        :=
        \bigwedge_{x\in X}Y\bigl(f(x),g(x)\bigr).
    \]
\end{definition}

The meet in this formula is the posetal form of the enriched end
\[
    \int_{x\in X}Y(fx,gx)
\]
appearing in the enriched functor-category formula \cite[Section~2.2, formula~(2.10)]{Kelly1982BasicConcepts}. In a thin enriched base, this end is computed as the greatest lower bound of the pointwise hom-values satisfying the usual wedge conditions. These wedge conditions follow from functoriality of \(f\) and \(g\).

\begin{proposition}
\label{prop:internal-hom-is-fuzzy-preorder}
    For fuzzy preorders \(X\) and \(Y\), the structure \([X,Y]\) of Definition~\ref{def:internal-hom} is a fuzzy preorder. If \(Y\) is separated, then \([X,Y]\) is separated.
\end{proposition}

\begin{proof}
    This is the posetal instance of the enriched functor-category construction \cite[Section~2.2, formulas~(2.10), (2.14), and~(2.15)]{Kelly1982BasicConcepts}. Reflexivity and transitivity are checked pointwise:
    \[
        \top\leq Y(fx,fx)
        \quad\text{and}\quad
        [X,Y](f,g)\otimes[X,Y](g,h)\leq Y(fx,hx)
    \]
    for every \(x\in X\). Taking the meet over \(x\) gives the fuzzy preorder axioms.

    If \(Y\) is separated and
    \[
        [X,Y](f,g)=\top=[X,Y](g,f),
    \]
    then for every \(x\in X\),
    \[
        Y(fx,gx)=\top=Y(gx,fx).
    \]
    Hence \(f(x)=g(x)\) for all \(x\), and therefore \(f=g\).
\end{proof}

The enriched evaluation functor specializes to an evaluation \(\Lambda\)-functor \cite[Section~2.2, formula~(2.17)]{Kelly1982BasicConcepts}
\[
    \operatorname{ev}:[X,Y]\otimes X\to Y,
    \qquad
    \operatorname{ev}(f,x):=f(x).
\]

\begin{proposition}[Tensor--hom adjunction]
\label{prop:tensor-hom-adjunction}
    For fuzzy preorders \(X,Y,Z\), there is a natural bijection
    \[
        \Pre_\Lambda(Z\otimes X,Y)
        \cong
        \Pre_\Lambda(Z,[X,Y]).
    \]
    Consequently, the ordinary category \(\Pre_\Lambda\) is symmetric monoidal closed.
\end{proposition}

\begin{proof}
    This is Proposition~\ref{prop:v-tensor-hom-adjunction} specialized to the posetal fuzzy cosmos determined by \(\Lambda\). Explicitly, a \(\Lambda\)-functor
    \[
        h:Z\otimes X\to Y
    \]
    is sent to
    \[
        \widetilde h:Z\to [X,Y],
        \qquad
        \widetilde h(z)(x):=h(z,x),
    \]
    and a \(\Lambda\)-functor
    \[
        k:Z\to [X,Y]
    \]
    is sent to
    \[
        \check k:Z\otimes X\to Y,
        \qquad
        \check k(z,x):=k(z)(x).
    \]
    The enriched proof shows that these are mutually inverse and natural; in the posetal case the required inequalities are precisely the pointwise functoriality inequalities induced by Definition~\ref{def:internal-hom}.
\end{proof}

\begin{corollary}
\label{cor:fuzzy-orders-monoidal-closed}
    The symmetric monoidal closed structure on \(\Pre_\Lambda\) restricts to a symmetric monoidal closed structure on the full subcategory
    \[
        \Ord_\Lambda\subseteq \Pre_\Lambda .
    \]
\end{corollary}

\begin{proof}
    The unit \(1\) is separated. If \(X\) and \(Y\) are separated, then \(X\otimes Y\) is separated because integrality gives \(a\otimes b\leq a\) and \(a\otimes b\leq b\). Hence equality to \(\top\) in both directions forces equality in both coordinates. Finally, if \(Y\) is separated, then \([X,Y]\) is separated by Proposition~\ref{prop:internal-hom-is-fuzzy-preorder}. Since \(\Ord_\Lambda\) is full, the tensor--hom adjunction restricts.
\end{proof}

\begin{corollary}[Predicate transposition]
\label{cor:predicate-transposition}
    For all fuzzy preorders \(S\) and \(X\), there is a natural bijection
    \[
        \Pre_\Lambda\bigl(S,[X^{\mathrm{op}},\Lambda]\bigr)
        \cong
        \Pre_\Lambda(S\otimes X^{\mathrm{op}},\Lambda).
    \]
\end{corollary}

\begin{proof}
    This is the special case of the enriched tensor--hom adjunction \cite[Section~2.3, formula~(2.20)]{Kelly1982BasicConcepts} with \(Y=\Lambda\) and \(X\) replaced by \(X^{\mathrm{op}}\).
\end{proof}

Under this bijection, a \(\Lambda\)-functor
\[
    \Phi:S\to [X^{\mathrm{op}},\Lambda]
\]
corresponds to its uncurried form
\[
    \widehat{\Phi}:S\otimes X^{\mathrm{op}}\to\Lambda,
    \qquad
    \widehat{\Phi}(s,x):=\Phi(s)(x).
\]
Equivalently, \(\widehat{\Phi}\) is the composite
\[
    \begin{tikzcd}[column sep=large,row sep=large]
        S\otimes X^{\mathrm{op}}
            \arrow[r,"\Phi\otimes 1_{X^{\mathrm{op}}}"]
        &
        \left[X^{\mathrm{op}},\Lambda\right]\otimes X^{\mathrm{op}}
            \arrow[r,"\operatorname{ev}"]
        &
        \Lambda .
    \end{tikzcd}
\]

\begin{remark}
\label{rem:dual-covariant-construction}
    Dually, one may consider the covariant internal hom
    \[
        [X,\Lambda].
    \]
    We shall not need that construction in what follows, so we now turn to the contravariant object
    \[
        [X^{\mathrm{op}},\Lambda].
    \]
\end{remark}

\section{Fuzzy presheaves, fuzzy predicates, and predicate algebras}
\label{sec:fuzzy-presheaves-and-fuzzy-predicate-algebras}

We now specialize the internal-hom construction to contravariant truth-valued functors.

\subsection{Fuzzy presheaves in a fuzzy cosmos}
\label{subsec:fuzzy-presheaves-in-a-fuzzy-cosmos}

Let
\[
    (\mathbb V,\otimes,I,[-,-])
\]
be a fuzzy cosmos, and let \(X\) be a small \(\mathbb V\)-category.

In this subsection we use the same functor-category notation when the domain is small and the codomain is possibly large. Thus, if \(X\) is small and \(Y\) is a locally small possibly large \(\mathbb V\)-category, the \(\mathbb V\)-category \([X,Y]\) has as objects the \(\mathbb V\)-functors \(X\to Y\), with hom-objects computed by the small end
\[
    [X,Y](F,G)
    =
    \int_{x\in X}Y(Fx,Gx).
\]
This convention applies in particular to functor categories whose codomain is the self-enriched truth-value object \(\underline{\mathbb V}\) or the small-presheaf category \(\mathcal P(\underline{\mathbb V})\).

\begin{definition}[Fuzzy presheaf]
\label{def:fuzzy-presheaf-cosmos}
    A \textbf{fuzzy presheaf} on \(X\) is a fuzzy-functor
    \[
        \Phi:X^{\mathrm{op}}\to \underline{\mathbb V}.
    \]
    We write
    \[
        \Psh_{\mathbb V}(X)
        :=
        [X^{\mathrm{op}},\underline{\mathbb V}]
    \]
    for the possibly large \(\mathbb V\)-category of fuzzy presheaves on
    \(X\).
\end{definition}

This is the usual enriched presheaf construction, expressed using the self-enrichment of the base \(\mathbb V\) \cite[Sections~1.6 and~1.9]{Kelly1982BasicConcepts}. Since \(X\) is small, every \(\mathbb V\)-presheaf
\[
    X^{\mathrm{op}}\to \underline{\mathbb V}
\]
is small in the sense of the set-theoretic conventions of Section~\ref{sec:truth-values-and-fuzzy-preorders}. Indeed, by the enriched co-Yoneda formula every presheaf
\[
    \Phi:X^{\mathrm{op}}\to\underline{\mathbb V}
\]
is canonically a small weighted colimit of representables:
\[
    \Phi
    \cong
    \int^{x\in X} \Phi x\otimes y_X(x).
\]
Since \(X\) is small, this is a small colimit. Thus this notation is compatible with the earlier small-presheaf convention.

Unpacked, a fuzzy presheaf \(\Phi\) assigns to each \(x\in X\) an object
\[
    \Phi x\in\mathbb V,
\]
together with action morphisms
\[
    X(x',x)\otimes \Phi x
    \longrightarrow
    \Phi x'
\]
satisfying the evident unit and associativity axioms. Thus fuzzy presheaves are contravariant \(\mathbb V\)-valued functors.

\begin{proposition}[Presheaf hom-objects]
\label{prop:fuzzy-presheaf-hom-objects}
    For fuzzy presheaves \(\Phi,\Psi\in \Psh_{\mathbb V}(X)\), the hom-object in \(\Psh_{\mathbb V}(X)\) is
    \[
        \Psh_{\mathbb V}(X)(\Phi,\Psi)
        =
        \int_{x\in X}
        [\Phi x,\Psi x].
    \]
\end{proposition}

\begin{proof}
    This is the enriched functor-category hom formula \cite[Section~2.2, formula~(2.10)]{Kelly1982BasicConcepts}, applied to
    \[
        \Psh_{\mathbb V}(X)
        =
        [X^{\mathrm{op}},\underline{\mathbb V}],
    \]
    together with the identity
    \[
        \underline{\mathbb V}(\Phi x,\Psi x)=[\Phi x,\Psi x].
    \]
    The end is small because \(X\) is small.
\end{proof}

\begin{definition}[Representable fuzzy presheaves]
\label{def:representable-fuzzy-presheaves}
    For each \(x\in X\), the \textbf{representable fuzzy presheaf} at \(x\) is the \(\mathbb V\)-functor
    \[
        y_X(x):X^{\mathrm{op}}\to \underline{\mathbb V}
    \]
    defined by
    \[
        y_X(x)(z):=X(z,x).
    \]
    These assemble into a \(\mathbb V\)-functor
    \[
        y_X:X\to \Psh_{\mathbb V}(X),
    \]
    called the \textbf{Yoneda embedding}.
\end{definition}

This is the enriched Yoneda embedding for the self-enriched base \(\underline{\mathbb V}\) \cite[Section~1.9]{Kelly1982BasicConcepts}. The presheaf structure on \(y_X(x)\) is induced by composition in \(X\):
\[
    X(z',z)\otimes X(z,x)
    \longrightarrow
    X(z',x).
\]

\begin{proposition}[Enriched Yoneda formula]
\label{prop:enriched-yoneda-formula}
    For every small \(\mathbb V\)-category \(X\), every \(x\in X\), and every fuzzy presheaf
    \[
        \Phi:X^{\mathrm{op}}\to \underline{\mathbb V},
    \]
    there is a canonical isomorphism in \(\mathbb V\)
    \[
        \Psh_{\mathbb V}(X)\bigl(y_X(x),\Phi\bigr)
        \cong
        \Phi x.
    \]
\end{proposition}

\begin{proof}
    By Proposition~\ref{prop:fuzzy-presheaf-hom-objects},
    \[
        \Psh_{\mathbb V}(X)\bigl(y_X(x),\Phi\bigr)
        =
        \int_{z\in X}
        [X(z,x),\Phi z].
    \]
    This end is the usual enriched object of natural transformations from the representable \(X(-,x)\) to \(\Phi\). The enriched Yoneda lemma identifies it canonically with \(\Phi x\) \cite[Section~1.9]{Kelly1982BasicConcepts}.
\end{proof}

\begin{corollary}[Yoneda embedding]
\label{cor:enriched-yoneda-embedding}
    The Yoneda embedding
    \[
        y_X:X\to \Psh_{\mathbb V}(X)
    \]
    is fully faithful. Equivalently, for all \(x,x'\in X\), there is a canonical isomorphism
    \[
        X(x,x')
        \cong
        \Psh_{\mathbb V}(X)\bigl(y_X(x),y_X(x')\bigr).
    \]
\end{corollary}

\begin{proof}
    Apply Proposition~\ref{prop:enriched-yoneda-formula} to the representable fuzzy presheaf
    \[
        \Phi=y_X(x').
    \]
    Then
    \[
        \Psh_{\mathbb V}(X)\bigl(y_X(x),y_X(x')\bigr)
        \cong
        y_X(x')(x)
        =
        X(x,x').
    \]
\end{proof}

The fuzzy predicate theory below is the posetal, truth-valued specialization of this enriched presheaf picture.

\subsection{Fuzzy predicates}
\label{subsec:fuzzy-predicates}

Let \(\Lambda\) be a fuzzy truth-value object, regarded as a posetal fuzzy cosmos, and let \(X\) be a fuzzy preorder over~\(\Lambda\). Then a fuzzy predicate
\[
    X^{\mathrm{op}}\to \Lambda
\]
is exactly a \(\Lambda\)-valued contravariant predicate on \(X\). This is the posetal, quantale-enriched specialization of the enriched presheaf construction above \cite[Definition~2.2]{stubbe-2014-quantaloid-enriched}; closely related \(L\)-valued and \(GL\)-monoid-valued presheaf constructions occur in \cite{Hohle1995PresheavesGLMonoids,Hohle1998GLQuantales,HohleKubiak2011QuantaleSets}.

\begin{definition}[Fuzzy predicate]
\label{def:fuzzy-predicate}
    Let \(X\) be a fuzzy preorder over \(\Lambda\). A \textbf{fuzzy predicate} on \(X\) is a \(\Lambda\)-functor
    \[
        \phi:X^{\mathrm{op}}\to \Lambda.
    \]
    Equivalently, by Lemma~\ref{lem:fuzzy-predicate-pointwise-characterization} below, it is a map
    \[
        \phi:X\to L
    \]
    satisfying
    \begin{equation}
    \label{eq:fuzzy-predicate-law}
        X(x',x)\otimes \phi(x)\leq \phi(x')
    \end{equation}
    for all \(x,x'\in X\).

    We write
    \[
        \Pred_\Lambda(X):=[X^{\mathrm{op}},\Lambda]
    \]
    for the fuzzy preorder of fuzzy predicates on \(X\).
\end{definition}

\begin{lemma}[Pointwise characterization of fuzzy predicates]
\label{lem:fuzzy-predicate-pointwise-characterization}
    Let \(X\) be a fuzzy preorder over \(\Lambda\), and let \(\phi:X\to L\) be a map. Then \(\phi\) defines a \(\Lambda\)-functor
    \[
        X^{\mathrm{op}}\to \Lambda
    \]
    if and only if, for all \(x,x'\in X\),
    \[
        X(x',x)\otimes \phi(x)\leq \phi(x').
    \]
\end{lemma}

\begin{proof}
    A map \(\phi:X\to L\) is a \(\Lambda\)-functor \(X^{\mathrm{op}}\to \Lambda\) if and only if for all \(x,x'\in X\),
    \[
        X^{\mathrm{op}}(x,x')
        \leq
        \Lambda\bigl(\phi(x),\phi(x')\bigr).
    \]
    Since
    \[
        X^{\mathrm{op}}(x,x')=X(x',x)
        \qquad\text{and}\qquad
        \Lambda\bigl(\phi(x),\phi(x')\bigr)
        =
        \phi(x)\multimap \phi(x'),
    \]
    this is equivalent, by residuation, to
    \[
        X(x',x)\otimes \phi(x)\leq \phi(x').
    \]
\end{proof}

\begin{proposition}
\label{prop:predicates-form-fuzzy-order}
    For every fuzzy preorder \(X\), the collection \(\Pred_\Lambda(X)\) of fuzzy predicates on \(X\) carries a natural separated fuzzy preorder, with hom-values
    \[
        \Pred_\Lambda(X)(\phi,\psi)
        :=
        \bigwedge_{x\in X}\bigl(\phi(x)\multimap \psi(x)\bigr).
    \]
    Its induced crisp order is the pointwise order:
    \[
        \phi\leq \psi
        \quad:\Longleftrightarrow\quad
        (\forall x\in X)\ \phi(x)\leq \psi(x).
    \]
\end{proposition}

\begin{proof}
    This is the posetal instance of the enriched functor-category hom formula \cite[Section~2.2, formula~(2.10)]{Kelly1982BasicConcepts}, together with Proposition~\ref{prop:fuzzy-presheaf-hom-objects} and Proposition~\ref{prop:truth-value-object-is-fuzzy-preorder}. The enriched end becomes the small meet
    \[
        \bigwedge_{x\in X}
        \Lambda\bigl(\phi(x),\psi(x)\bigr)
        =
        \bigwedge_{x\in X}
        \bigl(\phi(x)\multimap \psi(x)\bigr).
    \]

    The description of the induced crisp order follows because
    \[
        \top\leq \phi(x)\multimap\psi(x)
        \quad\Longleftrightarrow\quad
        \phi(x)\leq\psi(x).
    \]
    Since the truth-value fuzzy preorder \(\Lambda\) is separated, Proposition~\ref{prop:internal-hom-is-fuzzy-preorder} implies that \(\Pred_\Lambda(X)=[X^{\mathrm{op}},\Lambda]\) is separated.
\end{proof}

Thus the crisp order underlying \(\Pred_\Lambda(X)\) is simply the order inherited pointwise from \(L\), even though the fuzzy hom-value between predicates records the stronger degree
\[
    \bigwedge_{x\in X}\bigl(\phi(x)\multimap\psi(x)\bigr).
\]

\begin{proposition}[Representable fuzzy predicates]
\label{prop:representable-fuzzy-predicates}
    For each \(x\in X\), the map
    \[
        y_X(x):X\to L,
        \qquad
        y_X(x)(z):=X(z,x),
    \]
    is a fuzzy predicate on \(X\). Thus there is a canonical \(\Lambda\)-functor
    \[
        y_X:X\to \Pred_\Lambda(X),
        \qquad
        x\mapsto y_X(x).
    \]
\end{proposition}

\begin{proof}
    This is the posetal instance of the enriched representable presheaf and Yoneda embedding \cite[Section~1.9]{Kelly1982BasicConcepts}. Directly,
    \[
        X(z',z)\otimes y_X(x)(z)
        =
        X(z',z)\otimes X(z,x)
        \leq X(z',x)
        =
        y_X(x)(z')
    \]
    by transitivity in \(X\). Hence \(y_X(x)\) is a fuzzy predicate.

    The fact that \(y_X:X\to\Pred_\Lambda(X)\) is a \(\Lambda\)-functor follows by applying transitivity:
    \[
        X(x,x')\otimes X(z,x)\leq X(z,x')
    \]
    for every \(z\), then transposing by residuation and taking the meet over \(z\).
\end{proof}

\begin{proposition}[Yoneda formula for fuzzy predicates]
\label{prop:yoneda-formula-for-fuzzy-predicates}
    For every fuzzy preorder \(X\), every \(x\in X\), and every \(\phi\in \Pred_\Lambda(X)\),
    \[
        \Pred_\Lambda(X)\bigl(y_X(x),\phi\bigr)=\phi(x).
    \]
\end{proposition}

\begin{proof}
    This is the posetal form of the enriched Yoneda lemma \cite[Section~1.9]{Kelly1982BasicConcepts}. Directly,
    \[
        \Pred_\Lambda(X)\bigl(y_X(x),\phi\bigr)
        =
        \bigwedge_{z\in X}\bigl(X(z,x)\multimap \phi(z)\bigr).
    \]
    The \(z=x\) term gives the inequality
    \[
        \Pred_\Lambda(X)\bigl(y_X(x),\phi\bigr)\leq\phi(x).
    \]
    Conversely, since \(\phi\) is a fuzzy predicate,
    \[
        X(z,x)\otimes \phi(x)\leq \phi(z),
    \]
    so residuation gives
    \[
        \phi(x)\leq X(z,x)\multimap \phi(z)
    \]
    for every \(z\). Taking the meet over \(z\) gives the reverse inequality.
\end{proof}

\begin{corollary}[Yoneda embedding]
\label{cor:yoneda-embedding-fully-faithful}
    For every fuzzy preorder \(X\), the map
    \[
        y_X:X\to \Pred_\Lambda(X)
    \]
    is fully faithful, that is,
    \[
        X(x,x')
        =
        \Pred_\Lambda(X)\bigl(y_X(x),y_X(x')\bigr)
        \qquad
        \text{for all }x,x'\in X.
    \]
    In particular, if \(X\) is separated, then \(y_X\) is injective.
\end{corollary}

\begin{proof}
    Apply Proposition~\ref{prop:yoneda-formula-for-fuzzy-predicates} to the representable fuzzy predicate
    \[
        \phi=y_X(x').
    \]
    If \(X\) is separated and \(y_X(x)=y_X(x')\), then
    \[
        \top\leq X(x,x')
        \qquad\text{and}\qquad
        \top\leq X(x',x),
    \]
    so \(x=x'\).
\end{proof}

\subsection{Fuzzy predicate algebras}
\label{subsec:fuzzy-predicate-algebras}

We now use the additional order-theoretic structure present in the posetal case. The fuzzy predicate object \(\Pred_\Lambda(X)\) carries more structure than its defining fuzzy preorder. Its induced crisp order is complete, finite meets distribute over small joins, and the truth-value object acts on it by pointwise scalar operations. Related algebraic treatments of fuzzy-valued sets, \(GL\)-monoids, and quantale-valued sets appear in \cite{Hohle1992MValuedSetsSheaves,Hohle1995PresheavesGLMonoids,Hohle1998GLQuantales,HohleKubiak2011QuantaleSets}.

\begin{proposition}
\label{prop:predicate-completeness}
    Let \(X\) be a fuzzy preorder over \(\Lambda\), and let \((\phi_i)_{i\in I}\) be a small family of fuzzy predicates on \(X\). Define maps
    \[
        \Big(\bigvee_{i\in I}\phi_i\Big)(x)
        :=
        \bigvee_{i\in I}\phi_i(x),
        \qquad
        \Big(\bigwedge_{i\in I}\phi_i\Big)(x)
        :=
        \bigwedge_{i\in I}\phi_i(x)
    \]
    for each \(x\in X\). Then both
    \[
        \bigvee_{i\in I}\phi_i
        \qquad\text{and}\qquad
        \bigwedge_{i\in I}\phi_i
    \]
    are fuzzy predicates on \(X\). Consequently, the induced crisp order of \(\Pred_\Lambda(X)\) is a complete lattice, with small joins and small meets computed pointwise.
\end{proposition}

\begin{proof}
    For joins, Proposition~\ref{prop:tensor-preserves-joins} gives
    \[
        X(x',x)\otimes
        \Big(\bigvee_{i\in I}\phi_i(x)\Big)
        =
        \bigvee_{i\in I}
        \bigl(X(x',x)\otimes \phi_i(x)\bigr)
        \leq
        \bigvee_{i\in I}\phi_i(x').
    \]

    For meets, the empty meet is the constant predicate with value \(\top\). If \(I\neq\varnothing\), then for each \(j\in I\),
    \[
        X(x',x)\otimes
        \Big(\bigwedge_{i\in I}\phi_i(x)\Big)
        \leq
        X(x',x)\otimes \phi_j(x)
        \leq
        \phi_j(x').
    \]
    Hence the left-hand side is bounded above by \(\bigwedge_i\phi_i(x')\). Thus joins and meets are computed pointwise.
\end{proof}

\begin{proposition}
\label{prop:predicate-frame}
    For every fuzzy preorder \(X\), the induced crisp order of \(\Pred_\Lambda(X)\) is a frame.
\end{proposition}

\begin{proof}
    By Proposition~\ref{prop:predicate-completeness}, small joins and meets in \(\Pred_\Lambda(X)\) are computed pointwise. Since \(L\) is completely distributive, it is in particular a frame, and finite meets distribute over small joins pointwise.
\end{proof}

This frame structure is a structure on the induced crisp order; it is not the same datum as the \(\Lambda\)-valued hom of the fuzzy preorder \(\Pred_\Lambda(X)\).

\begin{definition}[Pointwise scalar operations]
\label{def:pointwise-scalar-operations}
    Let \(X\) be a fuzzy preorder over \(\Lambda\), let \(a\in L\), and let \(\phi\in \Pred_\Lambda(X)\). Define maps
    \[
        (a\odot \phi)(x):=a\otimes \phi(x),
        \qquad
        (a\Rightarrow \phi)(x):=a\multimap \phi(x)
    \]
    for each \(x\in X\).
\end{definition}

\begin{proposition}
\label{prop:scalar-operations-preserve-predicates}
    For every \(a\in L\) and every fuzzy predicate \(\phi\) on \(X\), the maps
    \[
        a\odot \phi
        \qquad\text{and}\qquad
        a\Rightarrow \phi
    \]
    of Definition~\ref{def:pointwise-scalar-operations} are again fuzzy predicates.
\end{proposition}

\begin{proof}
    For \(a\odot \phi\), associativity, commutativity, and the predicate law for \(\phi\) give
    \[
        X(x',x)\otimes (a\odot \phi)(x)
        =
        a\otimes \bigl(X(x',x)\otimes \phi(x)\bigr)
        \leq
        a\otimes \phi(x')
        =
        (a\odot \phi)(x').
    \]

    For \(a\Rightarrow \phi\), by residuation it suffices to prove
    \[
        a\otimes
        \Bigl(X(x',x)\otimes (a\Rightarrow \phi)(x)\Bigr)
        \leq
        \phi(x').
    \]
    The left-hand side is
    \[
        X(x',x)\otimes
        \Bigl(a\otimes (a\Rightarrow \phi)(x)\Bigr)
        \leq
        X(x',x)\otimes\phi(x)
        \leq
        \phi(x'),
    \]
    so \(a\Rightarrow\phi\) is a fuzzy predicate.
\end{proof}

\begin{proposition}
\label{prop:scalar-adjunction}
    Let \(X\) be a fuzzy preorder over \(\Lambda\). For every \(a\in L\) and all \(\phi,\psi\in \Pred_\Lambda(X)\),
    \[
        a\odot \phi\leq \psi
        \quad\Longleftrightarrow\quad
        \phi\leq a\Rightarrow \psi.
    \]
\end{proposition}

\begin{proof}
    By Proposition~\ref{prop:predicates-form-fuzzy-order}, the order on \(\Pred_\Lambda(X)\) is pointwise. Therefore the result is exactly the pointwise residuation law in \(L\).
\end{proof}

Thus, for each \(a\in L\), there is an adjunction on the crisp order underlying \(\Pred_\Lambda(X)\):
\[
    \begin{tikzcd}[column sep=large]
        \Pred_\Lambda(X)
            \arrow[r,shift left=1.1ex,"a\odot -"]
        &
        \Pred_\Lambda(X)
            \arrow[l,shift left=1.1ex,"a\Rightarrow -"']
    \end{tikzcd}
    \qquad
    a\odot -\dashv a\Rightarrow - .
\]

\begin{corollary}
\label{cor:scalar-action-preserves-joins}
    For every \(a\in L\), the operation
    \[
        a\odot(-):\Pred_\Lambda(X)\to \Pred_\Lambda(X)
    \]
    preserves small joins, and the operation
    \[
        a\Rightarrow(-):\Pred_\Lambda(X)\to \Pred_\Lambda(X)
    \]
    preserves small meets.
\end{corollary}

\begin{proof}
    This is immediate from Proposition~\ref{prop:scalar-adjunction}, since left adjoints preserve joins and right adjoints preserve meets.
\end{proof}

\begin{proposition}
\label{prop:predicate-scalar-laws}
    Let \(X\) be a fuzzy preorder over \(\Lambda\). For all \(a,b\in L\) and all \(\phi\in \Pred_\Lambda(X)\),
    \[
        \top\odot \phi=\phi,
        \qquad
        (a\otimes b)\odot \phi=a\odot (b\odot \phi).
    \]
    Moreover, for every small family \((\phi_i)_{i\in I}\) in \(\Pred_\Lambda(X)\),
    \[
        a\odot \Big(\bigvee_{i\in I}\phi_i\Big)
        =
        \bigvee_{i\in I}(a\odot \phi_i).
    \]
\end{proposition}

\begin{proof}
    All three identities are pointwise consequences of the corresponding laws in \(L\). For the last identity, use Proposition~\ref{prop:tensor-preserves-joins}.
\end{proof}

\begin{proposition}[Constant fuzzy predicates]
\label{prop:constant-fuzzy-predicates}
    Let \(X\) be a fuzzy preorder over \(\Lambda\). For every \(a\in L\), the constant map
    \[
        \underline a:X\to L,
        \qquad
        \underline a(x):=a,
    \]
    is a fuzzy predicate.
\end{proposition}

\begin{proof}
    For all \(x,x'\in X\), since \(X(x',x)\leq \top\) and \(\otimes\) is monotone,
    \[
        X(x',x)\otimes a
        \leq
        \top\otimes a
        =
        a.
    \]
\end{proof}

\begin{definition}[Fuzzy predicate algebra]
\label{def:fuzzy-predicate-algebra}
    Let \(X\) be a fuzzy preorder over \(\Lambda\). The \textbf{fuzzy predicate algebra} of \(X\) is the structure carried by \(\Pred_\Lambda(X)\) consisting of:
    \begin{enumerate}
        \item its frame structure from Proposition~\ref{prop:predicate-frame};

        \item the family of pointwise scalar actions
        \[
            a\odot(-):\Pred_\Lambda(X)\to \Pred_\Lambda(X)
            \qquad
            (a\in L);
        \]

        \item the family of pointwise scalar residuals
        \[
            a\Rightarrow(-):\Pred_\Lambda(X)\to \Pred_\Lambda(X)
            \qquad
            (a\in L).
        \]
    \end{enumerate}
\end{definition}

\subsection{Principal examples}
\label{subsec:principal-examples}

We now specialize the preceding constructions to several representative truth-value objects. In each case the abstract predicate object
\[
    \Pred_\Lambda(X)=[X^{\mathrm{op}},\Lambda]
\]
admits a concrete description, and the pointwise scalar operations induced by the base become explicit. The examples below include the ordinary crisp case, the standard chain-valued bases of mathematical fuzzy logic, and the Lawvere metric base \cite{Goguen1967LFuzzySets,Lawvere1973MetricSpaces} \cite[Example~2.4]{stubbe-2014-quantaloid-enriched}\cite[Sections~1, 2, and~7]{sep-logic-fuzzy}.

\begin{proposition}[Discrete predicate object]
\label{prop:discrete-predicate-object}
    Let \(X\) be a discrete fuzzy preorder over \(\Lambda\). Then every map
    \[
        \phi:X\to L
    \]
    is a fuzzy predicate. Hence
    \[
        \Pred_\Lambda(X)\cong L^X
    \]
    as sets. Under this identification, the induced crisp order, small joins, small meets, and scalar operations
    \[
        a\odot(-),
        \qquad
        a\Rightarrow(-)
    \]
    are all computed pointwise.
\end{proposition}

\begin{proof}
    If \(x'=x\), then \(X(x,x)=\top\), so the predicate condition is immediate. If \(x'\neq x\), then \(X(x',x)=\bot\), and tensor preserves the empty join, so
    \[
        \bot\otimes \phi(x)=\bot\leq \phi(x').
    \]
    Thus every map \(X\to L\) is a fuzzy predicate. The final statement follows from the general pointwise descriptions above.
\end{proof}

In the quantale-valued setting, such full \(L\)-valued function objects are the basic powerset-like objects; compare the quantale-valued powerset perspective in \cite{ShenZhao2024PowersetMonad}.

\subsubsection{The crisp case}

The two-valued case is the ordinary posetal specialization of enriched presheaves and the Yoneda embedding \cite[Section~1.9]{Kelly1982BasicConcepts}. Let
\[
    \Lambda=\{0<1\},
    \qquad
    a\otimes b=a\wedge b.
\]
If \(X\) is a fuzzy preorder over this base, then \(X\) is just an ordinary preorder. A fuzzy predicate on \(X\) is a map
\[
    \phi:X\to \{0,1\}
\]
such that
\[
    X(x',x)\wedge \phi(x)\leq \phi(x')
    \qquad
    \text{for all }x,x'\in X.
\]

\begin{proposition}
\label{prop:crisp-predicate-algebra}
    For a preorder \(X\), the assignment
    \[
        \phi\longmapsto A_\phi:=\{x\in X\mid \phi(x)=1\}
    \]
    identifies \(\Pred_\Lambda(X)\) with the set of lower sets of~\(X\). Under this identification, the order on \(\Pred_\Lambda(X)\) is inclusion, small joins are unions, small meets are intersections, and the representable fuzzy predicate \(y_X(x)\) corresponds to the principal lower set
    \[
        \downarrow x=\{z\in X\mid z\preceq x\}.
    \]
\end{proposition}

\begin{proof}
    The fuzzy predicate condition says precisely that
    \[
        x'\preceq x
        \ \text{and}\
        x\in A_\phi
        \quad\Longrightarrow\quad
        x'\in A_\phi.
    \]
    Hence fuzzy predicates are exactly lower sets. The order, joins, meets, and representables are then identified pointwise.
\end{proof}

\begin{corollary}
\label{cor:crisp-discrete-predicate-algebra}
    If \(X\) is a discrete preorder, then
    \[
        \Pred_\Lambda(X)\cong 2^X,
    \]
    the ordinary powerset Boolean algebra. Under this identification, joins are unions, meets are intersections, and \(y_X(x)\) corresponds to the singleton \(\{x\}\).
\end{corollary}

\begin{proof}
    In a discrete preorder, every subset is lower, and \(\downarrow x=\{x\}\).
\end{proof}

\subsubsection{The G\"odel case}

The G\"odel base is the quantale associated to the minimum \(t\)-norm, one of the standard chain-valued bases in mathematical fuzzy logic \cite[Example~2.4]{stubbe-2014-quantaloid-enriched} \cite[Sections~1, 2, and~7]{sep-logic-fuzzy}. Let \(\Lambda=[0,1]\) with the G\"odel operations
\[
    a\otimes b=a\wedge b,
    \qquad
    a\multimap_G b=
    \begin{cases}
        1,&a\leq b,\\
        b,&a>b.
    \end{cases}
\]
Then a fuzzy predicate on a fuzzy preorder \(X\) is a map \(\phi:X\to [0,1]\) such that
\[
    X(x',x)\wedge \phi(x)\leq \phi(x')
    \qquad
    \text{for all }x,x'\in X.
\]

\begin{proposition}
\label{prop:godel-predicate-algebra}
    For every fuzzy preorder \(X\) over the G\"odel base, the induced crisp order on \(\Pred_\Lambda(X)\) is a frame, hence a complete Heyting algebra. Moreover, for every \(a\in [0,1]\), the constant map
    \[
        \underline{a}(x):=a
    \]
    is a fuzzy predicate, and the scalar operations
    \[
        (a\odot \phi)(x):=a\wedge \phi(x),
        \qquad
        (a\Rightarrow \phi)(x):=a\multimap_G \phi(x)
    \]
    are again fuzzy predicates. They satisfy
    \[
        a\wedge \phi\leq \psi
        \quad\Longleftrightarrow\quad
        \phi\leq a\Rightarrow \psi.
    \]
\end{proposition}

\begin{proof}
    This is the G\"odel instance of Propositions~\ref{prop:predicate-frame}, \ref{prop:constant-fuzzy-predicates}, \ref{prop:scalar-operations-preserve-predicates}, and \ref{prop:scalar-adjunction}.
\end{proof}

\begin{corollary}
\label{cor:godel-discrete-predicate-algebra}
    If \(X\) is discrete, then
    \[
        \Pred_\Lambda(X)\cong [0,1]^X
    \]
    with pointwise order and pointwise G\"odel operations
    \[
        (\phi\wedge\psi)(x):=\min\{\phi(x),\psi(x)\},
        \qquad
        (\phi\vee\psi)(x):=\max\{\phi(x),\psi(x)\},
    \]
    \[
        (\phi\to_G\psi)(x):=\phi(x)\multimap_G\psi(x).
    \]
    Under these operations \(\Pred_\Lambda(X)\) is a complete G\"odel algebra, hence in particular a complete BL-algebra.
\end{corollary}

\begin{proof}
    This follows from Proposition~\ref{prop:discrete-predicate-object}, since every pointwise operation on maps \(X\to[0,1]\) remains a fuzzy predicate when \(X\) is discrete.
\end{proof}

The natural structure here is therefore a frame equipped with an external G\"odel action. In the discrete case, one recovers the familiar complete G\"odel algebra of all \([0,1]\)-valued functions on~\(X\). What fails in general is closure under the pointwise G\"odel implication inherited from the base algebra.

\begin{proposition}
\label{prop:godel-pointwise-implication-not-closed}
    There exists a fuzzy preorder \(X\) over the G\"odel base and fuzzy predicates \(\phi,\psi\in \Pred_\Lambda(X)\) such that the pointwise G\"odel implication
    \[
        (\phi\to_G \psi)(x):=\phi(x)\multimap_G \psi(x)
    \]
    is not a fuzzy predicate. Consequently, \(\Pred_\Lambda(X)\) is not, in general, closed under pointwise G\"odel implication.
\end{proposition}

\begin{proof}
    Let \(X=\{x',x\}\) with hom-values
    \[
        X(x,x)=X(x',x')=1,
        \qquad
        X(x',x)=\tfrac12,
        \qquad
        X(x,x')=0.
    \]
    This is a fuzzy preorder. Define fuzzy predicates
    \[
        \phi(x)=0,
        \qquad
        \phi(x')=\tfrac12,
    \]
    and
    \[
        \psi(x)=0,
        \qquad
        \psi(x')=0.
    \]
    Both satisfy the predicate law. But
    \[
        (\phi\to_G \psi)(x)=1,
        \qquad
        (\phi\to_G \psi)(x')=0.
    \]
    The predicate law would require
    \[
        \tfrac12\wedge 1\leq 0,
    \]
    which is false.
\end{proof}

\begin{remark}
\label{rem:godel-frame-vs-pointwise-implication}
    This does not contradict Proposition~\ref{prop:predicate-frame}. Since \(\Pred_\Lambda(X)\) is a frame, it has an internal Heyting implication. What fails here is not the existence of implication in the frame-theoretic sense, but closure under the \textbf{pointwise} G\"odel residual inherited from the base algebra \([0,1]\). In general, these two operations do not coincide.
\end{remark}

\subsubsection{The {\L}ukasiewicz case}

The {\L}ukasiewicz base is the quantale associated to the standard {\L}ukasiewicz \(t\)-norm and residual implication, and underlies the algebraic semantics of {\L}ukasiewicz and MV-valued fuzzy logic \cite[Example~2.4]{stubbe-2014-quantaloid-enriched} \cite{cintula-hajek-noguera-2011-hmfl1}. Let \(\Lambda=[0,1]\) with the {\L}ukasiewicz operations
\[
    a\otimes b=\max\{0,a+b-1\},
    \qquad
    a\multimap_{\Luk} b=\min\{1,1-a+b\}.
\]
Then a fuzzy predicate on a fuzzy preorder \(X\) is a map \(\phi:X\to [0,1]\) such that
\[
    \max\{0,\,X(x',x)+\phi(x)-1\}\leq \phi(x')
    \qquad
    \text{for all }x,x'\in X.
\]

\begin{proposition}
\label{prop:lukasiewicz-predicate-algebra}
    For every fuzzy preorder \(X\) over the {\L}ukasiewicz base, the induced crisp order on \(\Pred_\Lambda(X)\) is a frame. Moreover, for every \(a\in [0,1]\), the constant map
    \[
        \underline{a}(x):=a
    \]
    is a fuzzy predicate, and the scalar operations
    \[
        (a\odot \phi)(x):=\max\{0,a+\phi(x)-1\},
        \qquad
        (a\Rightarrow \phi)(x):=\min\{1,1-a+\phi(x)\}
    \]
    are again fuzzy predicates. If one defines
    \[
        (a\oplus \phi)(x):=\min\{1,a+\phi(x)\},
    \]
    equivalently
    \[
        a\oplus \phi=(1-a)\Rightarrow \phi,
    \]
    then
    \[
        a\odot \phi\leq \psi
        \quad\Longleftrightarrow\quad
        \phi\leq a\Rightarrow \psi.
    \]
\end{proposition}

\begin{proof}
    This is the {\L}ukasiewicz instance of Propositions~\ref{prop:predicate-frame}, \ref{prop:constant-fuzzy-predicates}, \ref{prop:scalar-operations-preserve-predicates}, and \ref{prop:scalar-adjunction}. The formula for \(a\oplus\phi\) follows from
    \[
        (1-a)\multimap_{\Luk} u=\min\{1,a+u\}.
    \]
\end{proof}

\begin{corollary}
\label{cor:lukasiewicz-discrete-predicate-algebra}
    If \(X\) is discrete, then
    \[
        \Pred_\Lambda(X)\cong [0,1]^X
    \]
    with the pointwise {\L}ukasiewicz operations
    \[
        (\phi\odot\psi)(x):=\max\{0,\phi(x)+\psi(x)-1\},
        \qquad
        (\phi\oplus\psi)(x):=\min\{1,\phi(x)+\psi(x)\},
    \]
    \[
        (\neg\phi)(x):=1-\phi(x),
        \qquad
        (\phi\multimap_{\Luk}\psi)(x)
        :=
        \min\{1,1-\phi(x)+\psi(x)\}.
    \]
    Under these operations \(\Pred_\Lambda(X)\) is a complete MV-algebra.
\end{corollary}

\begin{proof}
    This follows from Proposition~\ref{prop:discrete-predicate-object}, since every pointwise operation on maps \(X\to[0,1]\) remains a fuzzy predicate when \(X\) is discrete.
\end{proof}

\begin{remark}
\label{rem:lukasiewicz-frame-vs-pointwise-negation}
    Again, this does not contradict the fact that \(\Pred_\Lambda(X)\) is a frame. The frame \(\Pred_\Lambda(X)\) has its own internal Heyting negation, whereas the operation considered above is the \textbf{pointwise} {\L}ukasiewicz negation inherited from the base algebra \([0,1]\). In general these are different operations, and only the former is automatically available on every frame.
\end{remark}

\begin{proposition}
\label{prop:lukasiewicz-pointwise-negation-not-closed}
    There exists a fuzzy preorder \(X\) over the {\L}ukasiewicz base and a fuzzy predicate \(\phi\in \Pred_\Lambda(X)\) such that the pointwise {\L}ukasiewicz negation
    \[
        (\neg \phi)(x):=\phi(x)\multimap_{\Luk} 0=1-\phi(x)
    \]
    is not a fuzzy predicate. Consequently, \(\Pred_\Lambda(X)\) is not, in general, closed under pointwise {\L}ukasiewicz negation.
\end{proposition}

\begin{proof}
    Let \(X=\{x',x\}\) with hom-values
    \[
        X(x,x)=X(x',x')=1,
        \qquad
        X(x',x)=\tfrac12,
        \qquad
        X(x,x')=0.
    \]
    Define
    \[
        \phi(x)=0,
        \qquad
        \phi(x')=\tfrac34.
    \]
    Then \(\phi\) is a fuzzy predicate, since
    \[
        \max\{0,\tfrac12+0-1\}=0\leq \tfrac34.
    \]
    Its pointwise negation is
    \[
        (\neg \phi)(x)=1,
        \qquad
        (\neg \phi)(x')=\tfrac14.
    \]
    The predicate law would require
    \[
        \max\{0,\tfrac12+1-1\}\leq \tfrac14,
    \]
    which is false.
\end{proof}

\subsubsection{The Lawvere case}

The Lawvere base is the quantale-enriched presentation of generalized metric spaces \cite{Lawvere1973MetricSpaces} \cite[Example~2.3]{stubbe-2014-quantaloid-enriched}. Let \(\Lambda=[0,\infty]\) with the reverse of the usual order and operations
\[
    a\otimes b=a+b,
    \qquad
    a\multimap c=c\ominus a.
\]
If \(X\) is a fuzzy preorder over this base, we write \(d(x,y):=X(x,y)\). Then a fuzzy predicate on \(X\) is a map \(\phi:X\to [0,\infty]\) such that
\[
    \phi(x')\leq d(x',x)+\phi(x)
\]
in the usual order. Thus fuzzy predicates are precisely Lawvere weights, i.e. contravariant modules into the Lawvere base in this enriched sense \cite{Lawvere1973MetricSpaces} \cite[Example~2.3]{stubbe-2014-quantaloid-enriched}. The terminology ``predicate'' reflects that they are contravariant maps into the chosen truth-value object.

\begin{proposition}
\label{prop:lawvere-predicate-algebra}
    For every Lawvere metric space \(X\), equivalently every extended quasi-pseudometric space \(X\), the induced crisp order on \(\Pred_\Lambda(X)\) is a frame. In the usual order on \([0,\infty]\), this amounts to reverse pointwise order:
    \[
        \phi\leq \psi
        \quad:\Longleftrightarrow\quad
        \phi(x)\geq \psi(x)
        \ \text{for all }x\in X.
    \]
    Accordingly, joins in \(\Pred_\Lambda(X)\) are pointwise infima in the usual order, and meets are pointwise suprema.

    Moreover, for every \(a\in [0,\infty]\), the constant map
    \[
        \underline{a}(x):=a
    \]
    is a fuzzy predicate, and the scalar operations
    \[
        (a\odot \phi)(x):=a+\phi(x),
        \qquad
        (a\Rightarrow \phi)(x):=\phi(x)\ominus a
    \]
    are again fuzzy predicates. They satisfy
    \[
        a\odot \phi\leq \psi
        \quad\Longleftrightarrow\quad
        \phi\leq a\Rightarrow \psi,
    \]
    where \(\leq\) denotes the crisp order on \(\Pred_\Lambda(X)\). Equivalently, in the usual pointwise order on \([0,\infty]\), this says
    \[
        a+\phi(x)\geq \psi(x)
        \ \text{for all }x\in X
        \quad\Longleftrightarrow\quad
        \phi(x)\geq \psi(x)\ominus a
        \ \text{for all }x\in X.
    \]
\end{proposition}

\begin{proof}
    This is the Lawvere instance of the general predicate-frame, constant-predicate, scalar-closure, and scalar-adjunction results, translated through the reversed order on \([0,\infty]\).
\end{proof}

\begin{corollary}
\label{cor:lawvere-discrete-predicate-algebra}
    If \(X\) is discrete, then
    \[
        \Pred_\Lambda(X)\cong [0,\infty]^X
    \]
    with reverse pointwise order and pointwise operations
    \[
        (\phi+\psi)(x):=\phi(x)+\psi(x),
        \qquad
        (\psi\ominus\phi)(x):=\psi(x)\ominus \phi(x).
    \]
    In particular, with reverse pointwise order and pointwise addition, \(\Pred_\Lambda(X)\) is the usual commutative unital Lawvere quantale of extended cost functions on \(X\).
\end{corollary}

\begin{proof}
    This follows from Proposition~\ref{prop:discrete-predicate-object}.
\end{proof}

\begin{proposition}
\label{prop:lawvere-pointwise-addition-not-closed}
    There exists a Lawvere metric space \(X\) and fuzzy predicates \(\phi,\psi\in \Pred_\Lambda(X)\) such that the pointwise sum
    \[
        (\phi+\psi)(x):=\phi(x)+\psi(x)
    \]
    is not a fuzzy predicate.
\end{proposition}

\begin{proof}
    Let \(X=\{x',x\}\) with
    \[
        d(x,x)=d(x',x')=0,
        \qquad
        d(x',x)=1,
        \qquad
        d(x,x')=\infty.
    \]
    Define
    \[
        \phi(x)=0,
        \qquad
        \phi(x')=1,
    \]
    and let \(\psi=\phi\). Then \(\phi\) and \(\psi\) are fuzzy predicates, since
    \[
        \phi(x')=1\leq d(x',x)+\phi(x)=1+0.
    \]
    But
    \[
        (\phi+\psi)(x)=0,
        \qquad
        (\phi+\psi)(x')=2,
    \]
    and the predicate law would require
    \[
        2\leq d(x',x)+0=1,
    \]
    which is false.
\end{proof}

\section{Fuzzy power objects and generic fuzzy predicates}
\label{sec:fuzzy-power-objects-and-generic-fuzzy-predicates}

We now regard the presheaf object as the \textbf{fuzzy power object} of \(X\):
\[
    P_{\mathbb V}(X):=\Psh_{\mathbb V}(X).
\]
The term ``power object'' is used here in the enriched classifying sense: \(P_{\mathbb V}(X)\) classifies indexed fuzzy presheaves by the tensor--hom adjunction. It is not asserted to be a subobject-classifier power object in an ambient topos.

\subsection{Fuzzy power objects in a fuzzy cosmos}
\label{subsec:power-objects-in-a-fuzzy-cosmos}

Let
\[
    (\mathbb V,\otimes,I,[-,-])
\]
be a fuzzy cosmos, and let \(X\) be a small \(\mathbb V\)-category.

\begin{definition}[Fuzzy power object in a fuzzy cosmos]
\label{def:enriched-power-object}
    The \textbf{fuzzy power object} of \(X\) is the possibly large \(\mathbb V\)-category
    \[
        P_{\mathbb V}(X)
        :=
        \Psh_{\mathbb V}(X)
        =
        [X^{\mathrm{op}},\underline{\mathbb V}].
    \]
    Thus an object of \(P_{\mathbb V}(X)\) is a fuzzy presheaf
    \[
        \Phi:X^{\mathrm{op}}\to \underline{\mathbb V}.
    \]
\end{definition}

Unpacked, an object \(\Phi\in P_{\mathbb V}(X)\) assigns an object
\[
    \Phi x\in\mathbb V
\]
to each \(x\in X\), together with action morphisms
\[
    X(x',x)\otimes \Phi x
    \longrightarrow
    \Phi x'
\]
satisfying the usual unit and associativity axioms. By Proposition~\ref{prop:fuzzy-presheaf-hom-objects}, the hom-object between two fuzzy presheaves is
\[
    P_{\mathbb V}(X)(\Phi,\Psi)
    =
    \int_{x\in X}[\Phi x,\Psi x].
\]

The fuzzy power-object construction is contravariantly functorial by reindexing.

\begin{proposition}[Reindexing of fuzzy presheaves]
\label{prop:enriched-power-object-functor}
    For every \(\mathbb V\)-functor
    \[
        F:X\to Y,
    \]
    precomposition with \(F^{\mathrm{op}}\) defines a \(\mathbb V\)-functor
    \[
        P_{\mathbb V}(F):P_{\mathbb V}(Y)\to P_{\mathbb V}(X),
        \qquad
        P_{\mathbb V}(F)(\Psi):=\Psi\circ F^{\mathrm{op}}.
    \]
    Hence \(P_{\mathbb V}\) is contravariantly functorial on the ordinary category of small \(\mathbb V\)-categories.
\end{proposition}

\begin{proof}
    This is the standard functoriality of enriched functor categories under precomposition \cite[Section~2.2]{Kelly1982BasicConcepts}. On hom-objects, precomposition is induced by the universal property of the end:
    \[
        \int_{y\in Y}[\Phi y,\Psi y]
        \longrightarrow
        \int_{x\in X}[\Phi(Fx),\Psi(Fx)].
    \]
    The identity and composition laws follow from functoriality of precomposition.
\end{proof}

The truth-value object itself is recovered as the fuzzy power object of the enriched unit.

\begin{proposition}[The fuzzy power object of the unit]
\label{prop:enriched-power-of-unit}
    There is a canonical isomorphism of \(\mathbb V\)-categories
    \[
        P_{\mathbb V}(\mathbf 1)\cong \underline{\mathbb V}.
    \]
\end{proposition}

\begin{proof}
    Since \(\mathbf 1^{\mathrm{op}}\cong \mathbf 1\), a \(\mathbb V\)-functor
    \[
        \mathbf 1^{\mathrm{op}}\to\underline{\mathbb V}
    \]
    is uniquely determined by the image of the unique object of \(\mathbf 1\). For objects \(A,B\in\mathbb V\), the hom-object is
    \[
        [\mathbf 1^{\mathrm{op}},\underline{\mathbb V}](A,B)
        =
        \int_{\ast\in\mathbf 1}\underline{\mathbb V}(A,B)
        =
        [A,B]
        =
        \underline{\mathbb V}(A,B).
    \]
\end{proof}

\subsection{The generic fuzzy presheaf}
\label{subsec:generic-fuzzy-presheaf}

A \(\mathbb V\)-functor
\[
    \Phi:S\otimes X^{\mathrm{op}}\to\underline{\mathbb V}
\]
will be called an \textbf{\(S\)-indexed fuzzy presheaf on \(X\)}. The fuzzy power object \(P_{\mathbb V}(X)\) classifies such indexed fuzzy presheaves through a universal evaluation fuzzy presheaf.

\begin{definition}[Generic fuzzy presheaf]
\label{def:generic-fuzzy-presheaf}
    Let \(X\) be a small \(\mathbb V\)-category. The \textbf{generic fuzzy presheaf} on \(X\) is the evaluation \(\mathbb V\)-functor
    \[
        \Gen_X:P_{\mathbb V}(X)\otimes X^{\mathrm{op}}
        \longrightarrow
        \underline{\mathbb V}.
    \]
    Equivalently, \(\Gen_X\) is the transpose of the identity
    \[
        1_{P_{\mathbb V}(X)}:P_{\mathbb V}(X)\to P_{\mathbb V}(X)
    \]
    under the universe-bounded tensor--hom adjunction.
\end{definition}

Thus \(\Gen_X\) is displayed as
\[
    \begin{tikzcd}[column sep=large]
        P_{\mathbb V}(X)\otimes X^{\mathrm{op}}
            \arrow[r,"\Gen_X=\operatorname{ev}"]
        &
        \underline{\mathbb V}.
    \end{tikzcd}
\]

\begin{proposition}[Evaluation formula]
\label{prop:generic-fuzzy-presheaf-is-evaluation}
    For every small \(\mathbb V\)-category \(X\), every fuzzy presheaf
    \[
        \Phi\in P_{\mathbb V}(X),
    \]
    and every \(x\in X\), the generic fuzzy presheaf is given on objects by
    \[
        \Gen_X(\Phi,x)=\Phi x.
    \]
\end{proposition}

\begin{proof}
    This is the evaluation functor associated to the internal hom
    \[
        P_{\mathbb V}(X)=[X^{\mathrm{op}},\underline{\mathbb V}].
    \]
    Under the universe-bounded tensor--hom adjunction, the transpose of the identity on \(P_{\mathbb V}(X)\) evaluates a presheaf at an object.
\end{proof}

\begin{proposition}[Naturality of the generic fuzzy presheaf]
\label{prop:enriched-naturality-generic-fuzzy-presheaf}
    For every \(\mathbb V\)-functor \(F:X\to Y\), the diagram
    \[
        \begin{tikzcd}[column sep=huge,row sep=large]
            P_{\mathbb V}(Y)\otimes X^{\mathrm{op}}
                \arrow[r,"P_{\mathbb V}(F)\otimes 1_{X^{\mathrm{op}}}"]
                \arrow[d,"1_{P_{\mathbb V}(Y)}\otimes F^{\mathrm{op}}"']
            &
            P_{\mathbb V}(X)\otimes X^{\mathrm{op}}
                \arrow[d,"\Gen_X"]
            \\
            P_{\mathbb V}(Y)\otimes Y^{\mathrm{op}}
                \arrow[r,"\Gen_Y"']
            &
            \underline{\mathbb V}
        \end{tikzcd}
    \]
    commutes.
\end{proposition}

\begin{proof}
    Both composites send an object \((\Psi,x)\) to
    \[
        \Psi(Fx).
    \]
    Equivalently, under the tensor--hom adjunction, both transpose to the reindexing functor \(P_{\mathbb V}(F)\).
\end{proof}

\begin{proposition}[Principal fuzzy presheaves]
\label{prop:principal-fuzzy-presheaves-in-cosmos}
    For every small \(\mathbb V\)-category \(X\), the Yoneda embedding
    \[
        y_X:X\to P_{\mathbb V}(X)
    \]
    satisfies
    \[
        \Gen_X(y_X(x),x')=X(x',x)
    \]
    on objects. Thus \(y_X(x)\) is the principal fuzzy presheaf generated by \(x\).
\end{proposition}

\begin{proof}
    This follows from
    \[
        y_X(x)(x')=X(x',x)
    \]
    and the evaluation formula for \(\Gen_X\).
\end{proof}

\begin{theorem}[Universal property of the fuzzy power object]
\label{thm:universal-property-enriched-power-object}
    For all small \(\mathbb V\)-categories \(S\) and \(X\), composition with the generic fuzzy presheaf induces a natural bijection
    \[
        \mathbb V\text{-}\mathbf{CAT}
        \bigl(S,P_{\mathbb V}(X)\bigr)
        \cong
        \mathbb V\text{-}\mathbf{CAT}
        \bigl(S\otimes X^{\mathrm{op}},\underline{\mathbb V}\bigr),
    \]
    given by
    \[
        F\longmapsto \Gen_X\circ (F\otimes 1_{X^{\mathrm{op}}}).
    \]
    Equivalently, \(\mathbb V\)-functors
    \[
        S\to P_{\mathbb V}(X)
    \]
    are in natural bijection with \(S\)-indexed fuzzy presheaves on \(X\).
\end{theorem}

\begin{proof}
    Since
    \[
        P_{\mathbb V}(X)=[X^{\mathrm{op}},\underline{\mathbb V}],
    \]
    this is exactly the universe-bounded tensor--hom adjunction for enriched functor categories. Explicitly, a functor
    \[
        F:S\to P_{\mathbb V}(X)
    \]
    corresponds to
    \[
        \widehat F:S\otimes X^{\mathrm{op}}\to\underline{\mathbb V},
        \qquad
        \widehat F(s,x)=F(s)(x),
    \]
    and the inverse sends \(H:S\otimes X^{\mathrm{op}}\to\underline{\mathbb V}\) to
    \[
        \widetilde H:S\to P_{\mathbb V}(X),
        \qquad
        \widetilde H(s)(x)=H(s,x).
    \]
\end{proof}

\begin{remark}
\label{rem:global-elements-of-enriched-power-object}
    Taking \(S=\mathbf 1\) in Theorem~\ref{thm:universal-property-enriched-power-object} shows that global elements
    \[
        \mathbf 1\to P_{\mathbb V}(X)
    \]
    are precisely fuzzy presheaves on \(X\).
\end{remark}

\subsection{The fuzzy power object}
\label{subsec:the-fuzzy-power-object}

We now specialize the preceding construction to the posetal fuzzy cosmos determined by a fuzzy truth-value object~\(\Lambda\). Let \(X\) be a fuzzy preorder over~\(\Lambda\).

\begin{definition}[Fuzzy power object]
\label{def:fuzzy-power-object}
    The \textbf{fuzzy power object} of \(X\) is
    \[
        P_\Lambda(X):=\Pred_\Lambda(X)=[X^{\mathrm{op}},\Lambda].
    \]
    Thus an element of \(P_\Lambda(X)\) is precisely a fuzzy predicate on \(X\), equivalently a map \(\phi:X\to L\) satisfying
    \[
        X(x',x)\otimes \phi(x)\leq \phi(x')
        \qquad
        \text{for all }x,x'\in X.
    \]
\end{definition}

Again, the term ``power object'' is used in the enriched classifying sense: \(P_\Lambda(X)\) classifies indexed fuzzy predicates by the tensor--hom adjunction. It is not meant to identify \(P_\Lambda(X)\) with a topos-theoretic power object determined by an ambient subobject classifier.

By Proposition~\ref{prop:predicates-form-fuzzy-order}, Proposition~\ref{prop:predicate-completeness}, and Proposition~\ref{prop:predicate-frame}, the object \(P_\Lambda(X)\) is a separated fuzzy preorder whose induced crisp order is the pointwise order, in fact a frame, and whose pointwise scalar operations are those of Definition~\ref{def:pointwise-scalar-operations}. We shall use this structure freely.

\begin{proposition}
\label{prop:fuzzy-power-object-functor}
    For every \(\Lambda\)-functor \(f:X\to Y\), the assignment
    \[
        P_\Lambda(f):P_\Lambda(Y)\to P_\Lambda(X),
        \qquad
        P_\Lambda(f)(\psi):=\psi\circ f,
    \]
    defines a \(\Lambda\)-functor. Hence
    \[
        P_\Lambda:\Pre_\Lambda^{\mathrm{op}}\to \Ord_\Lambda
    \]
    is a contravariant functor.
\end{proposition}

\begin{proof}
    This is the posetal instance of Proposition~\ref{prop:enriched-power-object-functor}. Directly, if \(\psi:Y^{\mathrm{op}}\to\Lambda\) is a fuzzy predicate, then \(\psi\circ f^{\mathrm{op}}\) is a fuzzy predicate on \(X\). For \(\phi,\psi\in P_\Lambda(Y)\), one has
    \[
        P_\Lambda(Y)(\phi,\psi)
        =
        \bigwedge_{y\in Y}\bigl(\phi(y)\multimap \psi(y)\bigr)
        \leq
        \bigwedge_{x\in X}
        \bigl(\phi(fx)\multimap \psi(fx)\bigr),
    \]
    so \(P_\Lambda(f)\) is a \(\Lambda\)-functor. Functoriality is ordinary functoriality of precomposition.
\end{proof}

\begin{proposition}[Reindexing preserves fuzzy predicate algebra structure]
\label{prop:reindexing-preserves-predicate-structure}
    Let \(f:X\to Y\) be a \(\Lambda\)-functor. Then the reindexing map
    \[
        P_\Lambda(f):P_\Lambda(Y)\to P_\Lambda(X)
    \]
    preserves small joins, small meets, and commutes with the pointwise scalar operations. Explicitly, for every small family \((\psi_i)_{i\in I}\) in \(P_\Lambda(Y)\), every \(a\in L\), and every \(\psi\in P_\Lambda(Y)\),
    \[
        P_\Lambda(f)\Bigl(\bigvee_{i\in I}\psi_i\Bigr)
        =
        \bigvee_{i\in I}P_\Lambda(f)(\psi_i),
        \qquad
        P_\Lambda(f)\Bigl(\bigwedge_{i\in I}\psi_i\Bigr)
        =
        \bigwedge_{i\in I}P_\Lambda(f)(\psi_i),
    \]
    \[
        P_\Lambda(f)(a\odot\psi)=a\odot P_\Lambda(f)(\psi),
        \qquad
        P_\Lambda(f)(a\Rightarrow\psi)=a\Rightarrow P_\Lambda(f)(\psi).
    \]
\end{proposition}

\begin{proof}
    All identities are pointwise after evaluating at \(x\in X\).
\end{proof}

\begin{proposition}[The fuzzy power object of the unit]
\label{prop:power-of-unit}
    There is a canonical isomorphism of fuzzy orders
    \[
        P_\Lambda(1)\cong \Lambda.
    \]
    Under this identification, a fuzzy predicate
    \[
        \phi:1^{\mathrm{op}}\to\Lambda
    \]
    corresponds to its unique value \(\phi(\ast)\in L\).
\end{proposition}

\begin{proof}
    This is the posetal instance of Proposition~\ref{prop:enriched-power-of-unit}. A predicate on \(1\) is uniquely determined by its value \(a\in L\), and the hom-value between \(a\) and \(b\) is \(a\multimap b\).
\end{proof}

\begin{corollary}[Constant fuzzy predicates by reindexing]
\label{cor:constant-predicates-by-reindexing}
    Let
    \[
        !_X:X\to 1
    \]
    be the unique \(\Lambda\)-functor. Under the identification \(P_\Lambda(1)\cong\Lambda\), the reindexing map
    \[
        P_\Lambda(!_X):\Lambda\to P_\Lambda(X)
    \]
    sends \(a\in L\) to the constant fuzzy predicate \(\underline{a}\) on
    \(X\).
\end{corollary}

\begin{proof}
    It sends the unique-value predicate \(a:1^{\mathrm{op}}\to\Lambda\) to \(a\circ !_X\), the constant predicate with value \(a\).
\end{proof}

\subsection{The generic fuzzy predicate}
\label{subsec:the-generic-fuzzy-predicate}

A \(\Lambda\)-functor
\[
    \Phi:S\otimes X^{\mathrm{op}}\to \Lambda
\]
will be called an \textbf{\(S\)-indexed fuzzy predicate on \(X\)}. Equivalently, by transposition, it is a \(\Lambda\)-functor
\[
    S\to P_\Lambda(X),
\]
that is, an \(S\)-indexed compatible family of fuzzy predicates on \(X\). The fuzzy power object \(P_\Lambda(X)\) classifies such indexed fuzzy predicates through the posetal specialization of the generic fuzzy presheaf.

\begin{definition}[Generic fuzzy predicate]
\label{def:generic-fuzzy-predicate}
    Let \(X\) be a fuzzy preorder over \(\Lambda\). The \textbf{generic fuzzy predicate} on \(X\) is the evaluation \(\Lambda\)-functor
    \[
        \Gen_X:P_\Lambda(X)\otimes X^{\mathrm{op}}\to \Lambda.
    \]
    Equivalently, \(\Gen_X\) is the transpose of the identity
    \[
        1_{P_\Lambda(X)}:P_\Lambda(X)\to P_\Lambda(X)
    \]
    under the tensor--hom adjunction.
\end{definition}

\begin{proposition}
\label{prop:generic-fuzzy-predicate-is-evaluation}
    For every fuzzy preorder \(X\), every \(\phi\in P_\Lambda(X)\), and every \(x\in X\),
    \[
        \Gen_X(\phi,x)=\phi(x).
    \]
    Thus \(\Gen_X\) is the ordinary evaluation map.
\end{proposition}

\begin{proof}
    This is the posetal instance of Proposition~\ref{prop:generic-fuzzy-presheaf-is-evaluation}.
\end{proof}

\begin{remark}[Direct functoriality of evaluation]
\label{rem:direct-functoriality-of-evaluation}
    Directly, the evaluation map \(\Gen_X(\phi,x)=\phi(x)\) is a \(\Lambda\)-functor because, for fuzzy predicates \(\phi,\psi\in P_\Lambda(X)\) and elements \(x,x'\in X\),
    \[
        P_\Lambda(X)(\phi,\psi)
        \otimes X^{\mathrm{op}}(x,x')\otimes \phi(x)
        =
        P_\Lambda(X)(\phi,\psi)
        \otimes X(x',x)\otimes \phi(x).
    \]
    Since \(\phi\) is a fuzzy predicate,
    \[
        X(x',x)\otimes \phi(x)\leq \phi(x'),
    \]
    and since
    \[
        P_\Lambda(X)(\phi,\psi)\leq \phi(x')\multimap \psi(x'),
    \]
    we obtain
    \[
        P_\Lambda(X)(\phi,\psi)
        \otimes X(x',x)\otimes \phi(x)
        \leq
        \psi(x').
    \]
    By residuation,
    \[
        P_\Lambda(X)(\phi,\psi)
        \otimes X^{\mathrm{op}}(x,x')
        \leq
        \phi(x)\multimap\psi(x'),
    \]
    which is exactly the functoriality condition for
    \[
        \Gen_X:P_\Lambda(X)\otimes X^{\mathrm{op}}\to\Lambda.
    \]
\end{remark}

\begin{proposition}[Naturality of the generic fuzzy predicate]
\label{prop:naturality-generic-predicate}
    For every \(\Lambda\)-functor \(f:X\to Y\), the diagram
    \[
        \begin{tikzcd}[column sep=huge,row sep=large]
            P_\Lambda(Y)\otimes X^{\mathrm{op}}
                \arrow[r,"P_\Lambda(f)\otimes 1_{X^{\mathrm{op}}}"]
                \arrow[d,"1_{P_\Lambda(Y)}\otimes f^{\mathrm{op}}"']
            &
            P_\Lambda(X)\otimes X^{\mathrm{op}}
                \arrow[d,"\Gen_X"]
            \\
            P_\Lambda(Y)\otimes Y^{\mathrm{op}}
                \arrow[r,"\Gen_Y"']
            &
            \Lambda
        \end{tikzcd}
    \]
    commutes. Equivalently, for every \(\psi\in P_\Lambda(Y)\) and every \(x\in X\),
    \[
        \Gen_X\bigl(P_\Lambda(f)(\psi),x\bigr)
        =
        \Gen_Y\bigl(\psi,f(x)\bigr).
    \]
\end{proposition}

\begin{proof}
    This is the posetal instance of Proposition~\ref{prop:enriched-naturality-generic-fuzzy-presheaf}; objectwise both sides are \(\psi(fx)\).
\end{proof}

\begin{proposition}[Principal fuzzy predicates]
\label{prop:principal-fuzzy-predicates}
    For every fuzzy preorder \(X\), the Yoneda map
    \[
        y_X:X\to P_\Lambda(X)
    \]
    satisfies
    \[
        \Gen_X\bigl(y_X(x),x'\bigr)=X(x',x)
        \qquad
        \text{for all }x,x'\in X.
    \]
    Thus \(y_X(x)\) is the principal fuzzy predicate generated by \(x\). If \(X\) is separated, then \(y_X\) is injective.
\end{proposition}

\begin{proof}
    This is the posetal instance of Proposition~\ref{prop:principal-fuzzy-presheaves-in-cosmos}. Injectivity follows from Corollary~\ref{cor:yoneda-embedding-fully-faithful}.
\end{proof}

\begin{theorem}[Universal property of the fuzzy power object]
\label{thm:universal-property-fuzzy-power-object}
    For all fuzzy preorders \(S\) and \(X\), composition with the generic fuzzy predicate induces a natural bijection
    \[
        \Pre_\Lambda\bigl(S,P_\Lambda(X)\bigr)
        \cong
        \Pre_\Lambda\bigl(S\otimes X^{\mathrm{op}},\Lambda\bigr),
    \]
    given by
    \[
        f\longmapsto \Gen_X\circ (f\otimes 1_{X^{\mathrm{op}}}).
    \]
    Equivalently, \(\Lambda\)-functors \(S\to P_\Lambda(X)\) are in natural bijection with \(S\)-indexed fuzzy predicates on \(X\), with inverse assignment
    \[
        \Phi\longmapsto \widetilde{\Phi},
        \qquad
        \widetilde{\Phi}(s)(x):=\Phi(s,x).
    \]
\end{theorem}

\begin{proof}
    This is the posetal instance of Theorem~\ref{thm:universal-property-enriched-power-object}, equivalently Corollary~\ref{cor:predicate-transposition}. The displayed formulas are exactly the objectwise formulas for currying and evaluation.
\end{proof}

\begin{remark}
\label{rem:global-elements-of-fuzzy-power-object}
    Taking \(S=1\) in Theorem~\ref{thm:universal-property-fuzzy-power-object} shows that global elements
    \[
        1\to P_\Lambda(X)
    \]
    are the same thing as fuzzy predicates on \(X\).
\end{remark}

\subsection{Running examples}
\label{subsec:running-examples-fuzzy-power-object}

We now record the form taken by the fuzzy power object and the generic fuzzy predicate for the main bases considered so far.

\begin{example}[The crisp power object]
\label{ex:crisp-fuzzy-power-object}
    Let \(\Lambda=\{0<1\}\), so that fuzzy preorders are ordinary preorders. By Proposition~\ref{prop:crisp-predicate-algebra},
    \[
        P_\Lambda(X)=\Pred_\Lambda(X)
    \]
    is the preorder of lower sets of \(X\), ordered by inclusion. For a lower set \(A\subseteq X\), the generic fuzzy predicate is ordinary membership:
    \[
        \Gen_X(A,x)=
        \begin{cases}
            1,&x\in A,\\
            0,&x\notin A.
        \end{cases}
    \]
    Thus monotone maps
    \[
        S\to P_\Lambda(X)
    \]
    are precisely \(S\)-indexed compatible families of lower sets of \(X\). If \(X\) is discrete, then \(P_\Lambda(X)\cong 2^X\), and \(\Gen_X\) is the ordinary element-of relation.
\end{example}

\begin{example}[The G\"odel power object]
\label{ex:godel-fuzzy-power-object}
    Let \(\Lambda=[0,1]\) with the G\"odel operations. Then
    \[
        P_\Lambda(X)=\Pred_\Lambda(X)
    \]
    consists of all maps \(\phi:X\to[0,1]\) satisfying
    \[
        X(x',x)\wedge \phi(x)\leq \phi(x')
        \qquad
        \text{for all }x,x'\in X.
    \]
    These are the lower fuzzy predicates on \(X\) in the G\"odel sense. The generic fuzzy predicate is evaluation:
    \[
        \Gen_X(\phi,x)=\phi(x).
    \]
    If \(X\) is discrete, then \(P_\Lambda(X)\cong [0,1]^X\).
\end{example}

\begin{example}[The {\L}ukasiewicz power object]
\label{ex:lukasiewicz-fuzzy-power-object}
    Let \(\Lambda=[0,1]\) with the {\L}ukasiewicz operations. Then
    \[
        P_\Lambda(X)=\Pred_\Lambda(X)
    \]
    consists of all maps \(\phi:X\to[0,1]\) satisfying
    \[
        \max\{0,\,X(x',x)+\phi(x)-1\}\leq \phi(x')
        \qquad
        \text{for all }x,x'\in X.
    \]
    These are the lower fuzzy predicates compatible with {\L}ukasiewicz transitivity. As before,
    \[
        \Gen_X(\phi,x)=\phi(x).
    \]
    If \(X\) is discrete, then \(P_\Lambda(X)\cong [0,1]^X\), now equipped with the pointwise {\L}ukasiewicz operations.
\end{example}

\begin{example}[The Lawvere power object]
\label{ex:lawvere-fuzzy-power-object}
    Let \(\Lambda=[0,\infty]\) with the reverse of the usual order, so that fuzzy preorders are Lawvere metric spaces. Writing \(d(x,y)=X(x,y)\), the fuzzy power object
    \[
        P_\Lambda(X)=\Pred_\Lambda(X)
    \]
    consists of all maps \(\phi:X\to[0,\infty]\) satisfying
    \[
        \phi(x')\leq d(x',x)+\phi(x)
        \qquad
        \text{in the usual order, for all }x,x'\in X.
    \]
    These are precisely the Lawvere weights on \(X\), read here as Lawvere-valued fuzzy predicates. The generic fuzzy predicate is again evaluation:
    \[
        \Gen_X(\phi,x)=\phi(x).
    \]
    If \(X\) is discrete, then \(P_\Lambda(X)\cong [0,\infty]^X\), the quantale of extended cost functions on \(X\) with reverse pointwise order.
\end{example}

\begin{example}[A non-chain coordinatewise power object]
\label{ex:nonchain-coordinatewise-power-object}
    Let
    \[
        \Lambda=[0,1]\times[0,1]
    \]
    with the coordinatewise {\L}ukasiewicz operations, ordered coordinatewise. If \(X\) is a fuzzy preorder over \(\Lambda\), write
    \[
        X(x,x')=\bigl(X_1(x,x'),X_2(x,x')\bigr).
    \]
    Then a fuzzy predicate
    \[
        \phi:X\to [0,1]\times[0,1]
    \]
    may be written as
    \[
        \phi(x)=\bigl(\phi_1(x),\phi_2(x)\bigr),
    \]
    and the fuzzy predicate condition is equivalent to the pair of coordinatewise conditions
    \[
        \max\{0,\,X_i(x',x)+\phi_i(x)-1\}\leq \phi_i(x')
        \qquad
        (i=1,2).
    \]
    Thus \(P_\Lambda(X)\) consists exactly of pairs of {\L}ukasiewicz fuzzy predicates, ordered and acted on coordinatewise, and the generic fuzzy predicate is
    \[
        \Gen_X\bigl((\phi_1,\phi_2),x\bigr)
        =
        \bigl(\phi_1(x),\phi_2(x)\bigr).
    \]
    This illustrates that the fuzzy power-object construction is not restricted to chain-valued bases.
\end{example}

\section{Cut subobjects and the representation theorem}
\label{sec:fuzzy-subobjects-and-the-representation-theorem}

In the preceding sections we defined the fuzzy power object of a small \(\mathbb V\)-category \(X\) by
\[
    P_{\mathbb V}(X)
    =
    [X^{\mathrm{op}},\underline{\mathbb V}],
\]
together with its generic fuzzy presheaf
\[
    \Gen_X:
    P_{\mathbb V}(X)\otimes X^{\mathrm{op}}
    \longrightarrow
    \underline{\mathbb V}.
\]
We now show that this power object represents enriched cut subobjects.

\subsection{The fuzzy cut classifier}
\label{subsec:fuzzy-cut-classifier}

Let
\[
    (\mathbb V,\otimes,I,[-,-])
\]
be a fuzzy cosmos, and let \(\underline{\mathbb V}\) be its self-enriched truth-value object. Thus the objects of \(\underline{\mathbb V}\) are the objects of \(\mathbb V\), and the hom-object from \(A\) to \(B\) is
\[
    \underline{\mathbb V}(A,B)=[A,B].
\]

\begin{definition}[Cut classifier]
\label{def:cut-classifier}
    The \textbf{cut classifier} is the self-enriched hom-functor
    \[
        \mathsf{Cut}_{\mathbb V}:
        \underline{\mathbb V}^{\mathrm{op}}
        \otimes
        \underline{\mathbb V}
        \longrightarrow
        \underline{\mathbb V}
    \]
    defined on objects by
    \[
        \mathsf{Cut}_{\mathbb V}(A,B):=[A,B].
    \]
\end{definition}

Equivalently, \(\mathsf{Cut}_{\mathbb V}\) is the Yoneda probe of the truth value \(B\) by the level \(A\).

\begin{proposition}
\label{prop:intrinsic-cut-classifier-is-functor}
    The assignment
    \[
        (A,B)\longmapsto [A,B]
    \]
    defines a \(\mathbb V\)-functor
    \[
        \mathsf{Cut}_{\mathbb V}:
        \underline{\mathbb V}^{\mathrm{op}}
        \otimes
        \underline{\mathbb V}
        \longrightarrow
        \underline{\mathbb V}.
    \]
\end{proposition}

\begin{proof}
    For objects \((A,B)\) and \((A',B')\), functoriality requires a morphism
    \[
        [A',A]\otimes [B,B']
        \longrightarrow
        [[A,B],[A',B']].
    \]
    By closedness, this is equivalently a morphism
    \[
        [A',A]\otimes [B,B']\otimes [A,B]\otimes A'
        \longrightarrow
        B'.
    \]
    Using symmetry and associativity, followed by the evaluation maps, we define it by the composite
    \[
        \begin{aligned}
            [A',A]\otimes [B,B']\otimes [A,B]\otimes A'
            &\longrightarrow
            [B,B']\otimes [A,B]\otimes [A',A]\otimes A'
            \\
            &\longrightarrow
            [B,B']\otimes [A,B]\otimes A
            \\
            &\longrightarrow
            [B,B']\otimes B
            \\
            &\longrightarrow
            B'.
        \end{aligned}
    \]
    The identity and composition axioms are the unit and associativity axioms for the self-enrichment of \(\mathbb V\).
\end{proof}

\subsection{Fuzzy cut profiles and cut nerves}
\label{subsec:raw-enriched-cut-profiles-and-cut-nerves}

Let \(X\) be a small \(\mathbb V\)-category. Let
\[
    \mathcal P(\underline{\mathbb V})
\]
denote the \(\mathbb V\)-category of small presheaves on the possibly large \(\mathbb V\)-category \(\underline{\mathbb V}\). The Yoneda embedding is
\[
    Y:
    \underline{\mathbb V}
    \longrightarrow
    \mathcal P(\underline{\mathbb V}),
    \qquad
    B\longmapsto \underline{\mathbb V}(-,B)=[-,B].
\]

\begin{definition}[Cut profile]
\label{def:cut-profile}
    A \textbf{cut profile} on \(X\) is a \(\mathbb V\)-functor
    \[
        C:X^{\mathrm{op}}\longrightarrow
        \mathcal P(\underline{\mathbb V}).
    \]
    We write
    \[
        \FCut_{\mathbb V}(X)
        :=
        [
            X^{\mathrm{op}},
            \mathcal P(\underline{\mathbb V})
        ]
    \]
    for the \(\mathbb V\)-category of raw enriched cut profiles on \(X\).
\end{definition}

Equivalently, a cut profile may be written in uncurried form as a \(\mathbb V\)-functor
\[
    C:
    \underline{\mathbb V}^{\mathrm{op}}
    \otimes
    X^{\mathrm{op}}
    \longrightarrow
    \underline{\mathbb V}
\]
such that, for every \(x\in X\), the level profile
\[
    C(-,x):
    \underline{\mathbb V}^{\mathrm{op}}
    \longrightarrow
    \underline{\mathbb V}
\]
is a small presheaf.

The hom-objects in \(\FCut_{\mathbb V}(X)\) are computed by the small end over \(X\):
\[
    \FCut_{\mathbb V}(X)(C,D)
    =
    \int_{x\in X}
    \mathcal P(\underline{\mathbb V})(C x,D x).
\]

\begin{definition}[Cut nerve]
\label{def:cut-nerve}
    Let
    \[
        \Phi:X^{\mathrm{op}}\to\underline{\mathbb V}
    \]
    be a fuzzy presheaf on \(X\). Its \textbf{cut nerve} is the cut profile
    \[
        N_X(\Phi):X^{\mathrm{op}}\to
        \mathcal P(\underline{\mathbb V})
    \]
    defined by
    \[
        N_X(\Phi)(x):=Y(\Phi x).
    \]
    Equivalently, in uncurried notation,
    \[
        N_X(\Phi)(A,x)
        =
        [A,\Phi x].
    \]
\end{definition}

Thus \(N_X(\Phi)\) is obtained by substituting \(\Phi\) into the cut classifier. Since \(N_X\) is induced by postcomposition with the Yoneda embedding \(Y\), it defines a
\(\mathbb V\)-functor
\[
    N_X:
    P_{\mathbb V}(X)
    \longrightarrow
    \FCut_{\mathbb V}(X).
\]

\begin{definition}[Generic cut]
\label{def:generic-cut}
    Let
    \[
        \Gen_X:
        P_{\mathbb V}(X)\otimes X^{\mathrm{op}}
        \longrightarrow
        \underline{\mathbb V}
    \]
    be the generic fuzzy presheaf on \(X\). The
    \textbf{generic cut} is the composite
    \[
        \begin{tikzcd}[column sep=large,row sep=large]
            \underline{\mathbb V}^{\mathrm{op}}
            \otimes
            P_{\mathbb V}(X)
            \otimes
            X^{\mathrm{op}}
            \arrow[r,"1\otimes \Gen_X"]
            \arrow[rr,bend right=18,"\GCut_X"']
            &
            \underline{\mathbb V}^{\mathrm{op}}
            \otimes
            \underline{\mathbb V}
            \arrow[r,"\mathsf{Cut}_{\mathbb V}"]
            &
            \underline{\mathbb V}.
        \end{tikzcd}
    \]
    Thus
    \[
        \GCut_X(A,\Phi,x)=[A,\Phi x].
    \]
\end{definition}

\begin{proposition}[Universality of the generic cut]
\label{prop:generic-cut-universal}
    Let
    \[
        F:S\to P_{\mathbb V}(X)
    \]
    classify an \(S\)-indexed fuzzy presheaf
    \[
        \widehat F:
        S\otimes X^{\mathrm{op}}
        \longrightarrow
        \underline{\mathbb V}.
    \]
    Then the associated \(S\)-indexed cut profile is obtained from the generic cut by substitution:
    \[
        \begin{tikzcd}[column sep=large,row sep=large]
            \underline{\mathbb V}^{\mathrm{op}}
            \otimes
            S
            \otimes
            X^{\mathrm{op}}
            \arrow[r,"1\otimes F\otimes 1"]
            \arrow[rr,bend right=18,"N_{S,X}(\widehat F)"']
            &
            \underline{\mathbb V}^{\mathrm{op}}
            \otimes
            P_{\mathbb V}(X)
            \otimes
            X^{\mathrm{op}}
            \arrow[r,"\GCut_X"]
            &
            \underline{\mathbb V}.
        \end{tikzcd}
    \]
    Objectwise,
    \[
        N_{S,X}(\widehat F)(A,s,x)
        =
        [A,\widehat F(s,x)]
        =
        [A,F(s)(x)].
    \]
\end{proposition}

\begin{proof}
    Since \(F\) classifies \(\widehat F\), one has
    \[
        \widehat F(s,x)=F(s)(x).
    \]
    Substitution into the cut classifier gives
    \[
        \mathsf{Cut}_{\mathbb V}(A,\widehat F(s,x))
        =
        [A,\widehat F(s,x)]
        =
        [A,F(s)(x)].
    \]
\end{proof}

\subsection{Realization and closed cut profiles}
\label{subsec:realization-and-closed-enriched-cut-profiles}

The Yoneda embedding
\[
    Y:\underline{\mathbb V}\to\mathcal P(\underline{\mathbb V})
\]
has a left adjoint
\[
    \Sup_{\mathbb V}:
    \mathcal P(\underline{\mathbb V})
    \longrightarrow
    \underline{\mathbb V},
\]
given by the canonical weighted colimit
\[
    \Sup_{\mathbb V}(F)
    =
    \int^{A\in\underline{\mathbb V}}
    F(A)\otimes A.
\]
This notation denotes the weighted colimit of the identity functor on \(\underline{\mathbb V}\) by the small presheaf \(F\). Thus
\[
    \Sup_{\mathbb V}\dashv Y.
\]
The counit of this adjunction is the enriched co-Yoneda isomorphism
\[
    \int^{A\in\underline{\mathbb V}}[A,B]\otimes A
    \cong
    B
\]
\cite[Section~1.9]{Kelly1982BasicConcepts}.

\begin{definition}[Realization of a cut profile]
\label{def:realization-cut-profile}
    Let
    \[
        C:X^{\mathrm{op}}\to\mathcal P(\underline{\mathbb V})
    \]
    be a cut profile. Its \textbf{realization} is the fuzzy presheaf
    \[
        R_X(C):X^{\mathrm{op}}\to\underline{\mathbb V}
    \]
    defined pointwise by
    \[
        R_X(C)(x)
        :=
        \Sup_{\mathbb V}(C x).
    \]
    Equivalently,
    \[
        R_X(C)(x)
        =
        \int^{A\in\underline{\mathbb V}}
        Cx(A)\otimes A.
    \]
\end{definition}

Since \(R_X(C)=\Sup_{\mathbb V}\circ C\), functoriality of \(R_X(C)\) follows from functoriality of \(C\) and of \(\Sup_{\mathbb V}\). Therefore realization defines a \(\mathbb V\)-functor
\[
    R_X:
    \FCut_{\mathbb V}(X)
    \longrightarrow
    P_{\mathbb V}(X).
\]

\begin{proposition}[Realization--nerve adjunction]
\label{prop:realization-nerve-adjunction-intrinsic}
    For every small \(\mathbb V\)-category \(X\), there is an adjunction of \(\mathbb V\)-categories
    \[
        R_X
        \dashv
        N_X:
        \FCut_{\mathbb V}(X)
        \rightleftarrows
        P_{\mathbb V}(X).
    \]
    Equivalently, for every cut profile \(C\) and every fuzzy presheaf \(\Phi\), there is a natural isomorphism
    \[
        P_{\mathbb V}(X)\bigl(R_X(C),\Phi\bigr)
        \cong
        \FCut_{\mathbb V}(X)\bigl(C,N_X(\Phi)\bigr).
    \]
\end{proposition}

\begin{proof}
    This is the pointwise adjunction induced by
    \[
        \Sup_{\mathbb V}\dashv Y.
    \]
    Using the enriched functor-category hom formula over the small category \(X\), we compute
    \begin{align*}
        P_{\mathbb V}(X)\bigl(R_X(C),\Phi\bigr)
        &=
        \int_{x\in X}
        \underline{\mathbb V}
        \bigl(\Sup_{\mathbb V}(C x),\Phi x\bigr)
        \\
        &\cong
        \int_{x\in X}
        \mathcal P(\underline{\mathbb V})
        \bigl(C x,Y(\Phi x)\bigr)
        \\
        &=
        \FCut_{\mathbb V}(X)\bigl(C,N_X(\Phi)\bigr).
    \end{align*}
\end{proof}

\begin{proposition}[Co-Yoneda counit]
\label{prop:intrinsic-coyoneda-counit}
    For every fuzzy presheaf
    \[
        \Phi:X^{\mathrm{op}}\to\underline{\mathbb V},
    \]
    the counit of the adjunction
    \[
        \varepsilon_\Phi:
        R_XN_X(\Phi)\longrightarrow \Phi
    \]
    is an isomorphism.
\end{proposition}

\begin{proof}
    Evaluating at \(x\in X\), we have
    \[
        R_XN_X(\Phi)(x)
        =
        \Sup_{\mathbb V}(Y(\Phi x))
        =
        \int^{A\in\underline{\mathbb V}}
        [A,\Phi x]\otimes A.
    \]
    By the enriched co-Yoneda lemma, this coend is canonically isomorphic to \(\Phi x\). These isomorphisms are natural in \(x\).
\end{proof}

\begin{corollary}
\label{cor:intrinsic-cut-nerve-fully-faithful}
    The intrinsic cut nerve functor
    \[
        N_X:
        P_{\mathbb V}(X)
        \longrightarrow
        \FCut_{\mathbb V}(X)
    \]
    is fully faithful.
\end{corollary}

\begin{proof}
    In an adjunction \(R_X\dashv N_X\), the right adjoint \(N_X\) is fully faithful precisely when the counit \(R_XN_X\to 1\) is an isomorphism. This is Proposition~\ref{prop:intrinsic-coyoneda-counit}.
\end{proof}

\begin{definition}[Closed cut profile]
\label{def:closed-cut-profile}
    A cut profile
    \[
        C:X^{\mathrm{op}}\to\mathcal P(\underline{\mathbb V})
    \]
    is called \textbf{closed} if the unit of the adjunction
    \[
        \eta_C:C\longrightarrow N_XR_X(C)
    \]
    is an isomorphism.
\end{definition}

Objectwise, the unit is the canonical map
\[
    \eta_{C,x}:
    Cx\longrightarrow
    Y(\Sup_{\mathbb V}(Cx)).
\]
Equivalently, in uncurried notation it is the map
\[
    \eta_{C,A,x}:
    C(A,x)
    \longrightarrow
    [A,R_X(C)(x)]
    =
    \left[
        A,
        \int^{B\in\underline{\mathbb V}}C(B,x)\otimes B
    \right],
\]
obtained by transposing the canonical cowedge map
\[
    C(A,x)\otimes A
    \longrightarrow
    \int^{B\in\underline{\mathbb V}}C(B,x)\otimes B.
\]
Thus \(C\) is closed precisely when, for every \(x\in X\), the small level profile
\[
    C(-,x)
\]
is represented by its realization
\[
    R_X(C)(x).
\]

\begin{definition}[Cut subobjects]
\label{def:enriched-cut-subobjects}
    The \(\mathbb V\)-category of \textbf{enriched cut subobjects} of \(X\) is the full \(\mathbb V\)-subcategory
    \[
        \CutSub_{\mathbb V}(X)
        \subseteq
        \FCut_{\mathbb V}(X)
    \]
    spanned by the closed cut profiles.
\end{definition}

\begin{theorem}[Enriched cut representation theorem]
\label{thm:cut-representation-theorem}
    For every small \(\mathbb V\)-category \(X\), the cut nerve and realization restrict to an equivalence of \(\mathbb V\)-categories
    \[
        P_{\mathbb V}(X)
        \simeq
        \CutSub_{\mathbb V}(X).
    \]
    More explicitly, the functors
    \[
        N_X:
        P_{\mathbb V}(X)
        \longrightarrow
        \CutSub_{\mathbb V}(X)
    \]
    and
    \[
        R_X:
        \CutSub_{\mathbb V}(X)
        \longrightarrow
        P_{\mathbb V}(X)
    \]
    are inverse up to the unit and counit isomorphisms of the adjunction
    \[
        R_X\dashv N_X.
    \]
\end{theorem}

\begin{proof}
    By Proposition~\ref{prop:intrinsic-coyoneda-counit}, for every fuzzy presheaf \(\Phi\), the counit
    \[
        R_XN_X(\Phi)\longrightarrow \Phi
    \]
    is an isomorphism. Hence \(N_X\) is fully faithful.

    The essential image of a fully faithful right adjoint consists precisely of those objects for which the unit is an isomorphism. In the present case, these are exactly the closed cut profiles. Restricting \(N_X\) and \(R_X\) gives the claimed equivalence.
\end{proof}

\subsection{Top-support cuts in the posetal case}
\label{subsec:top-support-cuts-posetal-case}

We now specialize the above construction to the posetal fuzzy cosmos determined by a fuzzy truth-value object
\[
    \Lambda=(L,\leq,\wedge,\vee,\otimes,\multimap,\bot,\top).
\]
Write \(\underline{\Lambda}\) for the associated \(\Lambda\)-preorder of truth values, with object set \(L\) and hom-values
\[
    \underline{\Lambda}(a,b)=a\multimap b.
\]

The cut classifier becomes the residual:
\[
    \mathsf{Cut}_\Lambda(a,b)=a\multimap b.
\]
Its top support recovers the underlying order relation on truth values.

\begin{proposition}
\label{prop:top-support-residual-classifier}
    For all \(a,b\in L\),
    \[
        \top\leq a\multimap b
        \quad\Longleftrightarrow\quad
        a\leq b.
    \]
\end{proposition}

\begin{proof}
    By residuation,
    \[
        \top\leq a\multimap b
        \quad\Longleftrightarrow\quad
        a\otimes\top\leq b.
    \]
    Since \(\Lambda\) is integral, \(a\otimes\top=a\).
\end{proof}

\begin{definition}[Principal upper cut classifiers]
\label{def:principal-upper-cut-classifiers}
    For \(a\in L\), let
    \[
        \uparrow a:=\{b\in L\mid a\leq b\},
    \]
    regarded as the induced full subpreorder of \(\underline{\Lambda}\), and let
    \[
        j_a:\uparrow a\hookrightarrow\underline{\Lambda}
    \]
    be the inclusion.
\end{definition}

The inclusion \(j_a\) is the top-support shadow of the enriched probe
\[
    b\longmapsto a\multimap b.
\]
Indeed,
\[
    b\in\uparrow a
    \quad\Longleftrightarrow\quad
    a\leq b
    \quad\Longleftrightarrow\quad
    \top\leq a\multimap b.
\]

\subsection{Level cuts of fuzzy predicates}
\label{subsec:level-cuts-of-fuzzy-predicates}

Let \(X\) be a fuzzy preorder over \(\Lambda\).

\begin{definition}[Level cuts]
\label{def:principal-upper-sets-and-level-cuts}
    Let
    \[
        \phi:X^{\mathrm{op}}\to\underline{\Lambda}
    \]
    be a fuzzy predicate on \(X\). For \(a\in L\), the \textbf{\(a\)-cut} of \(\phi\) is the induced fuzzy subpreorder
    \[
        X_{\phi,a}\hookrightarrow X
    \]
    on the underlying subset
    \[
        |X_{\phi,a}|
        :=
        \{x\in X\mid a\leq\phi(x)\}
        =
        \phi^{-1}(\uparrow a).
    \]
    We write
    \[
        m_{\phi,a}:X_{\phi,a}\hookrightarrow X
    \]
    for the inclusion.
\end{definition}

When working with induced fuzzy subpreorders of \(X\), we identify them with their underlying subsets whenever no confusion can arise.

\begin{proposition}[Level cuts as classifier pullbacks]
\label{prop:level-cuts-as-classifier-pullbacks}
    Let
    \[
        \phi:X^{\mathrm{op}}\to\underline{\Lambda}
    \]
    be a fuzzy predicate. For each \(a\in L\), the \(a\)-cut
    \[
        m_{\phi,a}:X_{\phi,a}\hookrightarrow X
    \]
    is classified by
    \[
        j_a:\uparrow a\hookrightarrow\underline{\Lambda}
    \]
    through the pullback square
    \[
        \begin{tikzcd}[column sep=large,row sep=large]
            X_{\phi,a}^{\mathrm{op}}
            \arrow[r,"\phi_a"]
            \arrow[d,hook,"m_{\phi,a}^{\mathrm{op}}"']
            &
            \uparrow a
            \arrow[d,hook,"j_a"]
            \\
            X^{\mathrm{op}}
            \arrow[r,"\phi"']
            &
            \underline{\Lambda}.
        \end{tikzcd}
    \]
\end{proposition}

\begin{proof}
    We use the conical pullbacks in \(\Pre_\Lambda\) described in Lemma~\ref{lem:pullbacks-in-prelambda}. The pullback of \(j_a\) along \(\phi\) has underlying set
    \[
        \{x\in X\mid \phi(x)\in\uparrow a\}
        =
        \{x\in X\mid a\leq\phi(x)\}.
    \]
    For \(x,y\in X_{\phi,a}\), the hom-value in the pullback is
    \[
        X^{\mathrm{op}}(x,y)
        \wedge
        \underline{\Lambda}(\phi x,\phi y)
        =
        X(y,x)\wedge(\phi(x)\multimap\phi(y)).
    \]
    Since \(\phi:X^{\mathrm{op}}\to\underline{\Lambda}\) is a \(\Lambda\)-functor,
    \[
        X^{\mathrm{op}}(x,y)
        \leq
        \phi(x)\multimap\phi(y).
    \]
    Therefore the meet is just \(X^{\mathrm{op}}(x,y)\), so the pullback carries exactly the induced full fuzzy subpreorder on the displayed subset.
\end{proof}

\begin{proposition}[Basic properties of level cuts]
\label{prop:basic-properties-cut-family}
    Let
    \[
        \phi:X^{\mathrm{op}}\to\underline{\Lambda}
    \]
    be a fuzzy predicate. Then:
    \begin{enumerate}[label=\textup{(F\arabic*)}]
        \item \(X_{\phi,\bot}=X\);

        \item if \(a\leq b\), then
        \[
            X_{\phi,b}\subseteq X_{\phi,a};
        \]

        \item for every small family \((a_i)_{i\in I}\) in \(L\),
        \[
            X_{\phi,\bigvee_{i\in I}a_i}
            =
            \bigcap_{i\in I}X_{\phi,a_i}
        \]
        as induced fuzzy subpreorders of \(X\);

        \item if \(x\in X_{\phi,a}\), then for every \(x'\in X\),
        \[
            x'\in X_{\phi,\,X(x',x)\otimes a}.
        \]
    \end{enumerate}
\end{proposition}

\begin{proof}
    The first three statements follow immediately from the definition of the cuts and the lattice order on \(L\). For the fourth, if \(x\in X_{\phi,a}\), then \(a\leq\phi(x)\). Since \(\phi\) is a fuzzy predicate,
    \[
        X(x',x)\otimes\phi(x)\leq\phi(x').
    \]
    Hence
    \[
        X(x',x)\otimes a\leq\phi(x'),
    \]
    which is exactly the asserted membership.
\end{proof}

\begin{proposition}[Recovery from level cuts]
\label{prop:recovery-classifying-predicate}
    Let
    \[
        \phi:X^{\mathrm{op}}\to\underline{\Lambda}
    \]
    be a fuzzy predicate. Then for every \(x\in X\),
    \[
        \phi(x)=\bigvee\{a\in L\mid x\in X_{\phi,a}\}.
    \]
\end{proposition}

\begin{proof}
    By definition,
    \[
        x\in X_{\phi,a}
        \quad\Longleftrightarrow\quad
        a\leq\phi(x).
    \]
    Therefore the displayed join is the join of all elements below \(\phi(x)\), hence is \(\phi(x)\).
\end{proof}

\subsection{Fuzzy cut subobjects and classifying predicates}
\label{subsec:fuzzy-subobjects-and-classifying-predicates}

The preceding properties motivate the intrinsic posetal definition. In the posetal case we now pass from enriched closed cut profiles to their top-support presentation. Thus the notation
\[
    \CutSub_\Lambda(X)
\]
below denotes the levelwise cut-family presentation of the closed enriched cut profiles, not the raw \(\Lambda\)-enriched presheaf profiles themselves. These are cut-theoretic subobjects classified by fuzzy predicates; they should not be confused with arbitrary categorical subobjects of \(X\) in \(\Pre_\Lambda\).

\begin{definition}[Fuzzy cut subobject]
\label{def:fuzzy-subobject}
    A \textbf{fuzzy cut subobject} of \(X\), or simply a \textbf{fuzzy cut subobject} in this section, is a family
    \[
        A=\bigl(m_{A,a}:A_a\hookrightarrow X\bigr)_{a\in L}
    \]
    of induced fuzzy subpreorders of \(X\) such that:
    \begin{enumerate}[label=\textup{(F\arabic*)}]
        \item \(A_\bot=X\);

        \item if \(a\leq b\), then
        \[
            A_b\subseteq A_a;
        \]

        \item for every small family \((a_i)_{i\in I}\) in \(L\),
        \[
            A_{\bigvee_{i\in I}a_i}
            =
            \bigcap_{i\in I}A_{a_i}
        \]
        as induced fuzzy subpreorders of \(X\);

        \item if \(x\in A_a\), then for every \(x'\in X\),
        \[
            x'\in A_{X(x',x)\otimes a}.
        \]
    \end{enumerate}
    We write
    \[
        \CutSub_\Lambda(X)
    \]
    for the collection of all fuzzy cut subobjects of \(X\).
\end{definition}

\begin{definition}[Classifying fuzzy predicate]
\label{def:classifying-predicate}
    Let
    \[
        A=\bigl(m_{A,a}:A_a\hookrightarrow X\bigr)_{a\in L}
        \in\CutSub_\Lambda(X).
    \]
    Define
    \[
        \chi_A:X\to L,
        \qquad
        \chi_A(x):=\bigvee\{a\in L\mid x\in A_a\}.
    \]
    We call \(\chi_A\) the \textbf{classifying fuzzy predicate} of \(A\).
\end{definition}

This formula is the top-support shadow of the enriched co-Yoneda reconstruction. The enriched nerve of a fuzzy predicate \(\phi\) is
\[
    N_X(\phi)(a,x)=a\multimap\phi(x).
\]
Passing to top support retains precisely the levels satisfying
\[
    \top\leq a\multimap\phi(x),
\]
equivalently
\[
    a\leq\phi(x).
\]
The reconstructed truth value is the join of the retained levels.

\begin{proposition}
\label{prop:coherent-cut-family-gives-predicate}
    For every fuzzy cut subobject \(A\in\CutSub_\Lambda(X)\), the map
    \[
        \chi_A:X^{\mathrm{op}}\to\underline{\Lambda}
    \]
    is a fuzzy predicate.
\end{proposition}

\begin{proof}
    Let \(x,x'\in X\), and set
    \[
        S_x:=\{a\in L\mid x\in A_a\}.
    \]
    Then
    \[
        \chi_A(x)=\bigvee S_x.
    \]
    If \(a\in S_x\), then \(x\in A_a\), so by Definition~\ref{def:fuzzy-subobject}\textup{(F4)}
    \[
        x'\in A_{X(x',x)\otimes a}.
    \]
    Hence
    \[
        X(x',x)\otimes a\leq\chi_A(x').
    \]
    Taking the join over all \(a\in S_x\), and using preservation of joins by \(X(x',x)\otimes-\), gives
    \[
        X(x',x)\otimes\chi_A(x)
        =
        \bigvee_{a\in S_x}\bigl(X(x',x)\otimes a\bigr)
        \leq
        \chi_A(x').
    \]
\end{proof}

\begin{proposition}
\label{prop:recover-cuts-from-classifying-predicate}
    Let
    \[
        A=\bigl(m_{A,a}:A_a\hookrightarrow X\bigr)_{a\in L}
        \in\CutSub_\Lambda(X).
    \]
    Then for every \(a\in L\),
    \[
        A_a
        =
        \{x\in X\mid a\leq\chi_A(x)\}
    \]
    as induced fuzzy subpreorders of \(X\). Equivalently,
    \[
        A_a=X_{\chi_A,a}.
    \]
    In particular, every fuzzy cut subobject is uniquely determined by its classifying fuzzy predicate.
\end{proposition}

\begin{proof}
    Fix \(a\in L\) and \(x\in X\).

    If \(x\in A_a\), then \(a\) belongs to the set whose join defines \(\chi_A(x)\), so \(a\leq\chi_A(x)\).

    Conversely, suppose \(a\leq\chi_A(x)\). Let
    \[
        S_x:=\{b\in L\mid x\in A_b\}.
    \]
    Since \(A_\bot=X\), the set \(S_x\) is nonempty. By definition,
    \[
        \chi_A(x)=\bigvee S_x.
    \]
    By Definition~\ref{def:fuzzy-subobject}\textup{(F3)},
    \[
        A_{\chi_A(x)}
        =
        A_{\bigvee S_x}
        =
        \bigcap_{b\in S_x}A_b.
    \]
    Since \(x\in A_b\) for every \(b\in S_x\), we have \(x\in A_{\chi_A(x)}\). Now \(a\leq\chi_A(x)\), so Definition~\ref{def:fuzzy-subobject}\textup{(F2)} gives
    \[
        A_{\chi_A(x)}\subseteq A_a.
    \]
    Hence \(x\in A_a\). Since all subpreorders are induced from \(X\), equality of underlying subsets gives equality as induced fuzzy subpreorders.
\end{proof}

\subsection{The fuzzy preorder of fuzzy cut subobjects}
\label{subsec:fuzzy-preorder-of-fuzzy-cut-subobjects}

We next put a fuzzy preorder structure on \(\CutSub_\Lambda(X)\). Instead of merely transporting the hom-values from \(P_\Lambda(X)\) by the classifying predicate, we define the hom value directly in terms of cut containment.

\begin{definition}[Fuzzy hom between cut subobjects]
\label{def:fuzzy-preorder-on-subobjects}
    For \(A,B\in\CutSub_\Lambda(X)\), define
    \[
        \CutSub_\Lambda(X)(A,B)
        :=
        \bigvee
        \left\{
            r\in L
            \ \middle|\
            \forall a\in L,\
            A_a\subseteq B_{r\otimes a}
        \right\}.
    \]
\end{definition}

Thus the degree to which \(A\) is contained in \(B\) is the greatest truth value \(r\) such that every \(a\)-cut of \(A\) is contained in the \((r\otimes a)\)-cut of \(B\).

\begin{lemma}[Cut-containment hom formula]
\label{lem:subobject-hom-equals-classifier-hom}
    For all \(A,B\in\CutSub_\Lambda(X)\),
    \[
        \CutSub_\Lambda(X)(A,B)
        =
        \bigwedge_{x\in X}
        \bigl(\chi_A(x)\multimap\chi_B(x)\bigr).
    \]
\end{lemma}

\begin{proof}
    Let
    \[
        q:=
        \bigwedge_{x\in X}
        \bigl(\chi_A(x)\multimap\chi_B(x)\bigr).
    \]
    We show that, for \(r\in L\),
    \[
        r\leq q
    \]
    if and only if
    \[
        \forall a\in L,\quad A_a\subseteq B_{r\otimes a}.
    \]

    By residuation, \(r\leq q\) is equivalent to
    \[
        r\otimes\chi_A(x)\leq\chi_B(x)
        \qquad
        \text{for every }x\in X.
    \]
    Since
    \[
        \chi_A(x)=\bigvee\{a\in L\mid x\in A_a\}
    \]
    and \(r\otimes-\) preserves joins, this is equivalent to requiring that whenever \(x\in A_a\), one has
    \[
        r\otimes a\leq\chi_B(x).
    \]
    By Proposition~\ref{prop:recover-cuts-from-classifying-predicate} applied to \(B\), this is equivalent to
    \[
        x\in B_{r\otimes a}.
    \]
    Thus \(r\leq q\) is equivalent to the displayed containment condition. Therefore the set whose join defines \(\CutSub_\Lambda(X)(A,B)\) is the principal downset of \(q\), and its join is \(q\).
\end{proof}

\begin{proposition}
\label{prop:subobjects-form-fuzzy-preorder}
    The hom-values of Definition~\ref{def:fuzzy-preorder-on-subobjects} make \(\CutSub_\Lambda(X)\) into a fuzzy preorder. Moreover, the classifying-predicate assignment
    \[
        \mathsf{Cl}_X:
        \CutSub_\Lambda(X)\longrightarrow P_\Lambda(X),
        \qquad
        A\longmapsto\chi_A
    \]
    is fully faithful.
\end{proposition}

\begin{proof}
    By Lemma~\ref{lem:subobject-hom-equals-classifier-hom},
    \[
        \CutSub_\Lambda(X)(A,B)
        =
        \bigwedge_{x\in X}
        \bigl(\chi_A(x)\multimap\chi_B(x)\bigr)
        =
        P_\Lambda(X)(\chi_A,\chi_B).
    \]
    Since \(P_\Lambda(X)\) is a fuzzy preorder, the displayed equality implies the fuzzy preorder axioms for \(\CutSub_\Lambda(X)\). It also says exactly that \(\mathsf{Cl}_X\) is fully faithful.
\end{proof}

\subsection{The fuzzy cut-subobject representation theorem}
\label{subsec:representation-theorem-and-classification}

For \(\phi\in P_\Lambda(X)\), regarded as a fuzzy predicate on \(X\), write
\[
    X_\phi
    :=
    \bigl(m_{\phi,a}:X_{\phi,a}\hookrightarrow X\bigr)_{a\in L}
\]
for its cut family.

There are two canonical assignments:
\[
    \Cut_X:P_\Lambda(X)\longrightarrow\CutSub_\Lambda(X),
    \qquad
    \phi\longmapsto X_\phi,
\]
and
\[
    \mathsf{Cl}_X:\CutSub_\Lambda(X)\longrightarrow P_\Lambda(X),
    \qquad
    A\longmapsto\chi_A.
\]

\begin{theorem}[Fuzzy cut-subobject representation theorem]
\label{thm:representation-theorem}
    For every fuzzy preorder \(X\), the assignments
    \[
        \Cut_X:P_\Lambda(X)\longrightarrow\CutSub_\Lambda(X),
        \qquad
        \mathsf{Cl}_X:\CutSub_\Lambda(X)\longrightarrow P_\Lambda(X)
    \]
    are mutually inverse isomorphisms of fuzzy preorders:
    \[
        \mathsf{Cl}_X\circ\Cut_X=1_{P_\Lambda(X)},
        \qquad
        \Cut_X\circ\mathsf{Cl}_X=1_{\CutSub_\Lambda(X)}.
    \]
    Hence
    \[
        P_\Lambda(X)\cong\CutSub_\Lambda(X)
    \]
    as fuzzy preorders.
\end{theorem}

\begin{proof}
    Let \(\phi\in P_\Lambda(X)\). By Proposition~\ref{prop:basic-properties-cut-family}, the cut family \(X_\phi\) is a fuzzy cut subobject of \(X\), so \(\Cut_X\) is well-defined on objects.

    Conversely, if \(A\in\CutSub_\Lambda(X)\), then Proposition~\ref{prop:coherent-cut-family-gives-predicate} shows that
    \[
        \chi_A\in P_\Lambda(X),
    \]
    so \(\mathsf{Cl}_X\) is well-defined on objects.

    For \(\phi\in P_\Lambda(X)\), Proposition~\ref{prop:recovery-classifying-predicate} gives
    \[
        \chi_{X_\phi}(x)=\phi(x)
    \]
    for every \(x\in X\). Hence
    \[
        \mathsf{Cl}_X(\Cut_X(\phi))=\phi.
    \]

    For \(A\in\CutSub_\Lambda(X)\), Proposition~\ref{prop:recover-cuts-from-classifying-predicate} gives
    \[
        (X_{\chi_A})_a=A_a
    \]
    for every \(a\in L\). Hence
    \[
        \Cut_X(\mathsf{Cl}_X(A))=A.
    \]

    Finally, Proposition~\ref{prop:subobjects-form-fuzzy-preorder} gives
    \[
        \CutSub_\Lambda(X)(A,B)
        =
        P_\Lambda(X)(\chi_A,\chi_B)
    \]
    for all \(A,B\in\CutSub_\Lambda(X)\). Thus \(\mathsf{Cl}_X\) is fully faithful. Since \(\Cut_X\) and \(\mathsf{Cl}_X\) are inverse on objects, they are mutually inverse isomorphisms of fuzzy preorders.
\end{proof}

\subsection{Examples of the cut-subobject representation theorem}
\label{subsec:examples-cut-subobject-representation}

We finish by spelling out the representation theorem in the principal examples. In each case, the word ``subobject'' refers to the coherent cut-subobjects of Definition~\ref{def:fuzzy-subobject}, not to arbitrary monomorphism subobjects of \(X\) in \(\Pre_\Lambda\). Thus the theorem says that fuzzy predicates on \(X\) are equivalently coherent systems of levelwise induced subpreorders of \(X\).

\begin{example}[The crisp case]
\label{ex:cut-subobjects-crisp}
    Let
    \[
        \Lambda=2=\{0<1\}
    \]
    with tensor product \(a\otimes b=a\wedge b\). A fuzzy preorder over \(\Lambda\) is an ordinary preorder. A fuzzy predicate
    \[
        \phi:X^{\mathrm{op}}\to 2
    \]
    is exactly a lower subset of \(X\). Indeed, if
    \[
        U:=\{x\in X\mid \phi(x)=1\},
    \]
    then the predicate law says that whenever \(x\in U\) and \(x'\leq x\), one has \(x'\in U\).

    The cut family of \(\phi\) has only two levels:
    \[
        X_{\phi,0}=X,
        \qquad
        X_{\phi,1}=U.
    \]
    Conversely, a fuzzy cut subobject
    \[
        A=(A_0,A_1)
    \]
    is determined by \(A_1\), since axiom \textup{(F1)} gives
    \[
        A_0=X.
    \]
    Axiom \textup{(F4)} says exactly that \(A_1\) is lower: if \(x\in A_1\) and \(x'\leq x\), then \(x'\in A_1\). Hence the classifying predicate is simply the characteristic map
    \[
        \chi_A(x)=
        \begin{cases}
            1, & x\in A_1,\\
            0, & x\notin A_1.
        \end{cases}
    \]

    The fuzzy hom between cut subobjects also reduces to ordinary inclusion. Indeed,
    \[
        \CutSub_2(X)(A,B)=1
        \quad\Longleftrightarrow\quad
        A_1\subseteq B_1.
    \]
    Thus Theorem~\ref{thm:representation-theorem} specializes to the usual correspondence between lower subsets of a preorder and their characteristic maps. If \(X\) is discrete, this is the ordinary correspondence between subsets and characteristic functions.
\end{example}

\begin{example}[The G\"odel case]
\label{ex:cut-subobjects-godel}
    Let
    \[
        \Lambda=([0,1],\leq,\wedge,\vee,\otimes,\multimap,0,1)
    \]
    be the G\"odel truth-value object, with
    \[
        a\otimes b=a\wedge b,
        \qquad
        a\multimap b=
        \begin{cases}
            1, & a\leq b,\\
            b, & a>b.
        \end{cases}
    \]
    A fuzzy predicate
    \[
        \phi:X^{\mathrm{op}}\to\underline{\Lambda}
    \]
    is a map \(\phi:X\to[0,1]\) satisfying
    \[
        X(x',x)\wedge \phi(x)\leq \phi(x')
    \]
    for all \(x,x'\in X\).

    Its \(a\)-cut is
    \[
        X_{\phi,a}
        =
        \{x\in X\mid a\leq \phi(x)\}.
    \]
    Thus a fuzzy cut subobject is a decreasing family
    \[
        (A_a)_{a\in[0,1]}
    \]
    of induced fuzzy subpreorders satisfying
    \[
        A_{\bigvee_i a_i}=\bigcap_i A_{a_i}
    \]
    and the G\"odel stability condition
    \[
        x\in A_a
        \quad\Longrightarrow\quad
        x'\in A_{\,X(x',x)\wedge a}.
    \]
    The classifying predicate is recovered by
    \[
        \chi_A(x)=\sup\{a\in[0,1]\mid x\in A_a\}.
    \]

    For \(A,B\in\CutSub_\Lambda(X)\), the hom-value is
    \[
        \CutSub_\Lambda(X)(A,B)
        =
        \sup
        \left\{
            r\in[0,1]
            \ \middle|\
            \forall a\in[0,1],\,
            A_a\subseteq B_{r\wedge a}
        \right\}.
    \]
    Equivalently, by Lemma~\ref{lem:subobject-hom-equals-classifier-hom},
    \[
        \CutSub_\Lambda(X)(A,B)
        =
        \inf_{x\in X}
        \bigl(\chi_A(x)\multimap\chi_B(x)\bigr).
    \]
    Thus \(A\) is contained in \(B\) to degree \(r\) precisely when every \(a\)-level of \(A\) is contained in the \((r\wedge a)\)-level of \(B\). In particular, containment to degree \(1\) is ordinary pointwise containment of classifying predicates:
    \[
        \CutSub_\Lambda(X)(A,B)=1
        \quad\Longleftrightarrow\quad
        \chi_A\leq\chi_B
        \quad\Longleftrightarrow\quad
        A_a\subseteq B_a
        \text{ for all }a.
    \]
\end{example}

\begin{example}[The {\L}ukasiewicz case]
\label{ex:cut-subobjects-lukasiewicz}
    Let
    \[
        \Lambda=([0,1],\leq,\wedge,\vee,\otimes,\multimap,0,1)
    \]
    be the {\L}ukasiewicz truth-value object, with
    \[
        a\otimes b=\max\{0,a+b-1\},
        \qquad
        a\multimap b=\min\{1,1-a+b\}.
    \]
    A fuzzy predicate
    \[
        \phi:X^{\mathrm{op}}\to\underline{\Lambda}
    \]
    is a map \(\phi:X\to[0,1]\) satisfying
    \[
        \max\{0,X(x',x)+\phi(x)-1\}\leq \phi(x')
    \]
    for all \(x,x'\in X\).

    The \(a\)-cut is again
    \[
        X_{\phi,a}
        =
        \{x\in X\mid a\leq\phi(x)\}.
    \]
    A fuzzy cut subobject is therefore a decreasing, join-continuous family
    \[
        (A_a)_{a\in[0,1]}
    \]
    of induced fuzzy subpreorders satisfying the {\L}ukasiewicz stability condition
    \[
        x\in A_a
        \quad\Longrightarrow\quad
        x'\in A_{\max\{0,X(x',x)+a-1\}}.
    \]
    The classifying predicate is recovered by
    \[
        \chi_A(x)=\sup\{a\in[0,1]\mid x\in A_a\}.
    \]

    For \(A,B\in\CutSub_\Lambda(X)\), the cut-containment hom is
    \[
        \CutSub_\Lambda(X)(A,B)
        =
        \sup
        \left\{
            r\in[0,1]
            \ \middle|\
            \forall a\in[0,1],\,
            A_a\subseteq B_{\max\{0,r+a-1\}}
        \right\}.
    \]
    Equivalently,
    \[
        \CutSub_\Lambda(X)(A,B)
        =
        \inf_{x\in X}
        \min\{1,1-\chi_A(x)+\chi_B(x)\}.
    \]
    Thus containment to degree \(r\) permits a uniform loss of \(1-r\) in the level parameter: an \(a\)-level point of \(A\) is required to lie in \(B\) only at level
    \[
        \max\{0,a-(1-r)\}.
    \]
    In particular, the hom-value can be written as
    \[
        \CutSub_\Lambda(X)(A,B)
        =
        1-
        \sup_{x\in X}
        \max\{0,\chi_A(x)-\chi_B(x)\}.
    \]
    Hence the {\L}ukasiewicz hom records the worst pointwise deficit of \(\chi_B\) below \(\chi_A\).
\end{example}

\begin{example}[The Lawvere case]
\label{ex:cut-subobjects-lawvere}
    Let
    \[
        \Lambda=([0,\infty],\leq_L,\wedge,\vee,+,\multimap,\infty,0)
    \]
    be the Lawvere truth-value object, where
    \[
        a\leq_L b
        \quad\Longleftrightarrow\quad
        a\geq b
        \text{ in the usual order}.
    \]
    Thus the top element is \(0\), the bottom element is \(\infty\), and the tensor product is addition.

    A fuzzy preorder over this base is a Lawvere metric space, or extended quasi-pseudometric space. Writing
    \[
        d(x,y):=X(x,y),
    \]
    the preorder axioms become
    \[
        d(x,x)=0,
        \qquad
        d(x,z)\leq d(x,y)+d(y,z)
    \]
    in the usual order on \([0,\infty]\).

    A fuzzy predicate
    \[
        \phi:X^{\mathrm{op}}\to\underline{\Lambda}
    \]
    is a map \(\phi:X\to[0,\infty]\) satisfying
    \[
        d(x',x)+\phi(x)\geq \phi(x')
    \]
    in the usual order. Equivalently,
    \[
        \phi(x')\leq \phi(x)+d(x',x).
    \]
    Thus \(\phi\) is a one-sided nonexpansive potential.

    Because the order on \(L\) is reversed, the \(a\)-cut is the usual sublevel set
    \[
        X_{\phi,a}
        =
        \{x\in X\mid a\leq_L \phi(x)\}
        =
        \{x\in X\mid \phi(x)\leq a\}.
    \]
    A fuzzy cut subobject is therefore a family of induced fuzzy subpreorders
    \[
        (A_a)_{a\in[0,\infty]}
    \]
    satisfying
    \[
        A_\infty=X,
    \]
    monotonicity in the usual threshold parameter,
    \[
        b\leq a
        \quad\Longrightarrow\quad
        A_b\subseteq A_a,
    \]
    and the closure condition
    \[
        A_{\inf_i a_i}
        =
        \bigcap_i A_{a_i}.
    \]
    The stability axiom becomes the metric expansion condition
    \[
        x\in A_a
        \quad\Longrightarrow\quad
        x'\in A_{d(x',x)+a}.
    \]
    Thus a point within cost \(a\) forces every point \(x'\) to lie within cost \(a+d(x',x)\).

    The classifying predicate of a cut family is
    \[
        \chi_A(x)
        =
        \bigvee\nolimits_L\{a\in[0,\infty]\mid x\in A_a\}.
    \]
    Since \(\bigvee_L\) is infimum in the usual order, this says
    \[
        \chi_A(x)
        =
        \inf\{a\in[0,\infty]\mid x\in A_a\}.
    \]
    So \(\chi_A(x)\) is the least usual threshold at which \(x\) enters the filtration.

    Finally, for \(A,B\in\CutSub_\Lambda(X)\), the hom-value is
    \[
        \CutSub_\Lambda(X)(A,B)
        =
        \bigvee\nolimits_L
        \left\{
            r\in[0,\infty]
            \ \middle|\
            \forall a\in[0,\infty],\,
            A_a\subseteq B_{r+a}
        \right\}.
    \]
    Since joins in \(\Lambda\) are infima in the usual order, this is the least usual cost \(r\) such that every \(a\)-sublevel of \(A\) lies inside the \((a+r)\)-sublevel of \(B\). Equivalently, written in the usual order on \([0,\infty]\),
    \[
        \CutSub_\Lambda(X)(A,B)
        =
        \sup_{x\in X}
        \max\{0,\chi_B(x)-\chi_A(x)\}.
    \]
    Thus the Lawvere hom is the directed pointwise excess of \(\chi_B\) over \(\chi_A\): it is the least additive shift \(r\) for which
    \[
        \chi_B(x)\leq \chi_A(x)+r
    \]
    for all \(x\in X\).
\end{example}
    %\chapter{Fuzzy Frames, Fuzzy Locales, and the Fuzzy Sierpiński Classifier}
\label{ch:fuzzy-frames-locales}

In the previous chapter we constructed fuzzy predicates as enriched presheaves. For a
small \(\mathbb V\)-category \(X\), the fuzzy presheaf category was
\[
    P_{\mathbb V}(X)
    =
    [X^{\op},\underline{\mathbb V}],
\]
and, after passing to the posetal specialization determined by the fuzzy truth-value
object
\[
    \Lambda=(L,\leq,\meet,\join,\tensor,\multimap,\bot_L,\top_L),
\]
the corresponding fuzzy predicate object was
\[
    P_\Lambda(X)
    =
    [X^{\op},\Lambda].
\]
This chapter gives these presheaf objects their localic interpretation.

There are two structures on \(P_\Lambda(X)\) which should be kept distinct. First,
\(P_\Lambda(X)\) is a separated \(\Lambda\)-preorder, with hom-values
\[
    P_\Lambda(X)(\phi,\psi)
    =
    \bigmeet_{x\in X}(\phi(x)\multimap\psi(x)).
\]
Second, its induced crisp order is a frame, with joins and meets computed pointwise.
In this chapter, when we call \(P_\Lambda(X)\) a frame, we mean this induced crisp
frame of predicates, equipped with the pointwise \(\Lambda\)-scalar action
\[
    (r\tensor\phi)(x)=r\tensor\phi(x).
\]
The enriched hom-values remain present, but the frame order is the induced crisp
order.

The guiding idea is that fuzzy frames are obtained by localizing fuzzy presheaves. At
the categorical level this is expressed by enriched left-exact reflective
localizations of
\[
    P_{\mathbb V}(X).
\]
At the posetal level this becomes a scalar-compatible nucleus on
\[
    P_\Lambda(X).
\]
We prove below that these two descriptions agree in the posetal specialization:
left-exact \(\Lambda\)-enriched reflective localizations of \(P_\Lambda(X)\)
correspond, up to equivalence over \(P_\Lambda(X)\), to fuzzy nuclei on
\(P_\Lambda(X)\). Fuzzy locales are then defined by formal duality.

The fuzzy Sierpiński classifier is treated by its classifying universal property. This
is an important point. In the crisp case the Sierpiński frame can be realized as the
presheaf frame on the two-element preorder
\[
    U\leq 1.
\]
For general fuzzy truth-value bases this raw presheaf construction does not, in
general, have the required free-open property. In particular, for non-idempotent
tensors such as the {\L}ukasiewicz and product tensors, the presheaf frame on the
two-point subterminal shape need not classify arbitrary fuzzy opens by
finite-meet-preserving scalar frame homomorphisms.

Accordingly, the fuzzy Sierpiński frame is defined in this chapter as the free fuzzy
frame on one fuzzy open:
\[
    \Omega_\Lambda=F_{\op}^{\Lambda}(1),
\]
with generic fuzzy open
\[
    \sigma_\Lambda\in\Omega_\Lambda.
\]
Its dual fuzzy locale
\[
    \mathbb S_\Lambda
\]
classifies fuzzy opens by the universal property
\[
    \FLoc_\Lambda(\mathcal X,\mathbb S_\Lambda)
    \cong
    |\Opens_\Lambda(\mathcal X)|.
\]
Thus the Sierpiński object is not identified with a raw presheaf object unless an
additional presentation theorem has been proved for the chosen base.

At the enriched-cosmos level, the analogous construction is bicategorical. The free
fuzzy frame on an \(S\)-indexed family of cosmos-level fuzzy opens represents a
hom-category of models:
\[
    \mathbf{FFrm}_{\mathbb V}
    \bigl(F_{\op}(S),\mathcal A\bigr)
    \simeq
    \operatorname{Sub}_{\mathcal A}(1_{\mathcal A})^S.
\]
Equivalently, if \(S_{\mathrm{disc}}\) is the discrete category on \(S\), the
right-hand side is
\[
    \mathbf{Cat}
    \bigl(
        S_{\mathrm{disc}},
        \operatorname{Sub}_{\mathcal A}(1_{\mathcal A})
    \bigr).
\]
After passing to the posetal reflection, these hom-categories become ordinary
hom-sets and the corresponding classification becomes a strict set-level universal
property.

All large categories and collections of subobjects are taken relative to the fixed
universe convention of the previous chapter. If \(\mathcal A\) is an ordinary or
enriched category with terminal object, we write
\[
    \operatorname{Sub}_{\mathcal A}(1_{\mathcal A})
\]
for the ordinary category of subobjects of \(1_{\mathcal A}\). Its objects are
monomorphisms
\[
    U\hookrightarrow 1_{\mathcal A},
\]
and its morphisms are commuting triangles over \(1_{\mathcal A}\). When only
isomorphism classes are needed, we write
\[
    \operatorname{Iso}\operatorname{Sub}_{\mathcal A}(1_{\mathcal A})
    :=
    \operatorname{Ob}\operatorname{Sub}_{\mathcal A}(1_{\mathcal A})/{\cong}.
\]
In the posetal reflection, \(\operatorname{Sub}_{A}(\top_A)\) is equivalent to the
thin category underlying the frame \(A\), and its set of objects is simply \(|A|\).

\paragraph{Finite-limit convention for this chapter.}
We fix once and for all a class \(\mathcal F\) of finite weighted limits in
\(\mathbb V\)-categories. The class \(\mathcal F\) contains the terminal-object and
binary-product weights and is the finite-limit class used in all terms such as
``left exact'', ``finite-limit preserving'', and
\(\operatorname{Lex}_{\mathbb V}\). Thus
\[
    \operatorname{Lex}_{\mathbb V}(\mathcal C,\mathcal A)
\]
denotes the ordinary category of \(\mathbb V\)-functors
\[
    \mathcal C\to\mathcal A
\]
preserving the chosen \(\mathcal F\)-limits, with enriched natural transformations.

In the posetal results below, the only finite limits used explicitly are the
terminal element and binary meets of the underlying crisp frame. Thus left exactness
in the posetal comparison theorem means preservation of \(\top\) and binary meets.

The overall progression is:
\[
\begin{tikzcd}[column sep=large]
\text{fuzzy preorders}
    \arrow[r, "{P_\Lambda}"']
&
\text{fuzzy presheaves}
    \arrow[r, "\text{localization}"']
&
\text{fuzzy frames}
    \arrow[r, "\op"']
&
\text{fuzzy locales}
    \arrow[r, "{\FLoc_\Lambda(-,\mathbb S_\Lambda)}"']
&
\text{fuzzy opens.}
\end{tikzcd}
\]

The notation \(P_{\mathbb V}(X)\) and \(P_\Lambda(X)\) is retained from the previous
chapter. The new notation introduced in this chapter is:
\[
    \FFrm_\Lambda
\]
for the category of fuzzy frames over \(\Lambda\),
\[
    \FLoc_\Lambda
\]
for the category of fuzzy locales over \(\Lambda\),
\[
    \Omega_\Lambda
\]
for the fuzzy Sierpiński frame,
\[
    \sigma_\Lambda\in\Omega_\Lambda
\]
for its generic fuzzy open, and
\[
    \mathbb S_\Lambda
\]
for the corresponding fuzzy Sierpiński locale.

\section{Fuzzy presheaves recalled}
\label{sec:fuzzy-presheaves-recalled}

We begin by recalling the presheaf constructions from the previous chapter. The
definition--theorem--proof development of these objects has already been given there.

For a small \(\mathbb V\)-category \(X\), the fuzzy presheaf category is
\[
    P_{\mathbb V}(X)
    =
    [X^{\op},\underline{\mathbb V}].
\]
This is generally a locally small, possibly large \(\mathbb V\)-category. Its objects
are fuzzy presheaves on \(X\). The Yoneda embedding is the \(\mathbb V\)-functor
\[
    \yon_X:X\to P_{\mathbb V}(X),
    \qquad
    x\longmapsto X(-,x).
\]
By the enriched Yoneda lemma,
\[
    P_{\mathbb V}(X)(\yon_Xx,\Phi)
    \cong
    \Phi x
\]
for every fuzzy presheaf \(\Phi\in P_{\mathbb V}(X)\). Equivalently, the
variance-correct Yoneda comparison says that the composite
\[
\begin{tikzcd}[column sep=large, row sep=large]
X^{\op}
    \arrow[r, "\yon_X^{\op}"]
    \arrow[dr, "\Phi"']
&
P_{\mathbb V}(X)^{\op}
    \arrow[d, "{P_{\mathbb V}(X)(-,\Phi)}"]
\\
&
\underline{\mathbb V}
\end{tikzcd}
\]
is isomorphic to \(\Phi\). These are the standard enriched presheaf and Yoneda
constructions \cite{Kelly1982BasicConcepts}.

The density formula says that every fuzzy presheaf is a weighted colimit of
representables:
\[
    \Phi
    \cong
    \Phi\star\yon_X
    \cong
    \int^{x\in X}\Phi x\tensor\yon_Xx.
\]
Thus the values of \(\Phi\) give the coefficients with which the principal presheaves
\(\yon_Xx\) are joined to reconstruct \(\Phi\). This is the usual enriched density
formula for the Yoneda embedding \cite{Kelly1982BasicConcepts}.

After passing to the posetal specialization determined by
\[
    \Lambda=(L,\leq,\meet,\join,\tensor,\multimap,\bot_L,\top_L),
\]
a fuzzy preorder \(X\) has fuzzy predicate object
\[
    P_\Lambda(X)
    =
    [X^{\op},\Lambda].
\]
An element of \(P_\Lambda(X)\) is a map
\[
    \phi:X\to L
\]
satisfying
\[
    X(x',x)\tensor\phi(x)\leq \phi(x')
\]
for all \(x,x'\in X\). Joins, meets, and scalar multiplication in the induced crisp
frame are computed pointwise:
\[
    \left(\bigjoin_i\phi_i\right)(x)
    =
    \bigjoin_i\phi_i(x),
\]
\[
    \left(\bigmeet_i\phi_i\right)(x)
    =
    \bigmeet_i\phi_i(x),
\]
and
\[
    (r\tensor\phi)(x)
    =
    r\tensor\phi(x).
\]
This is the form in which the presheaf construction will be localized below.

\section{Localizing fuzzy presheaves}
\label{sec:localizing-fuzzy-presheaves}

We now perform the localic operation. Instead of postulating an arbitrary algebraic
notion of fuzzy frame, we obtain fuzzy frames by localizing fuzzy presheaves.

At the categorical level, this means a left-exact enriched localization of
\[
    P_{\mathbb V}(X).
\]
At the posetal level, this becomes a scalar-compatible nucleus on
\[
    P_\Lambda(X).
\]
The fixed points of such a nucleus form the corresponding fuzzy frame.

\subsection{Categorical form}
\label{subsec:localizing-presheaves-categorical-form}

Let \(X\) be a small \(\mathbb V\)-category. A localization of the fuzzy presheaf
category \(P_{\mathbb V}(X)\) is given by a \(\mathbb V\)-adjunction
\[
    a\dashv i:
    P_{\mathbb V}(X)\rightleftarrows \mathcal A
\]
with \(i\) fully faithful. The associated idempotent monad is
\[
    J\eqdef ia.
\]
We require this localization to be left exact with respect to the chosen
\(\mathcal F\)-limits.

\begin{definition}[Categorical fuzzy-frame presentation]
\label{def:fuzzy-frame-presentation-categorical}
A \textbf{categorical fuzzy-frame presentation} is a \(\mathbb V\)-adjunction
\[
    a\dashv i:
    P_{\mathbb V}(X)\rightleftarrows \mathcal A
\]
between locally small, possibly large \(\mathbb V\)-categories such that:
\begin{enumerate}
    \item \(i\) is fully faithful as a \(\mathbb V\)-functor;
    \item \(\mathcal A\) has the chosen finite weighted limits;
    \item \(a\) preserves the chosen finite weighted limits.
\end{enumerate}
The fuzzy frame presented by this localization has category of opens \(\mathcal A\).
\end{definition}

The localization is displayed by
\[
\begin{tikzcd}[column sep=large, row sep=large]
P_{\mathbb V}(X)
    \arrow[r, shift left=1.1ex, "a"]
&
\mathcal A
    \arrow[l, shift left=1.1ex, "i"]
\arrow[phantom, from=1-1, to=1-2, "\dashv" description]
\end{tikzcd}
\]
and the associated idempotent monad is
\[
\begin{tikzcd}[column sep=large]
P_{\mathbb V}(X)
    \arrow[r, "a"]
&
\mathcal A
    \arrow[r, "i"]
&
P_{\mathbb V}(X).
\end{tikzcd}
\]
Weighted limits and colimits are understood in the standard enriched sense
\cite{Kelly1982BasicConcepts}.

\begin{definition}[Local object]
\label{def:local-object-categorical}
Let
\[
    J=ia:P_{\mathbb V}(X)\to P_{\mathbb V}(X)
\]
be the idempotent monad associated to a categorical fuzzy-frame presentation. A fuzzy
presheaf \(\Phi\in P_{\mathbb V}(X)\) is called \textbf{\(J\)-local} if the unit
\[
    \Phi\to J\Phi
\]
is an isomorphism.
\end{definition}

\begin{proposition}[Local objects]
\label{prop:local-objects-categorical}
Let
\[
    a\dashv i:
    P_{\mathbb V}(X)\rightleftarrows \mathcal A
\]
be a categorical fuzzy-frame presentation, and let
\[
    J=ia.
\]
Then the full subcategory of \(J\)-local objects in \(P_{\mathbb V}(X)\) is equivalent
to \(\mathcal A\).
\end{proposition}

\begin{proof}
Since \(i\) is fully faithful, the counit
\[
    aiA\to A
\]
is an isomorphism for every \(A\in\mathcal A\). Hence
\[
    J(iA)=iaiA\cong iA,
\]
so every object in the image of \(i\) is \(J\)-local.

Conversely, if \(\Phi\) is \(J\)-local, then
\[
    \Phi\cong J\Phi=ia\Phi,
\]
so \(\Phi\) is isomorphic to an object in the image of \(i\). Therefore the local
objects are precisely the objects of \(\mathcal A\), up to equivalence.
\end{proof}

The local-object subcategory sits inside the presheaf category as
\[
\begin{tikzcd}[column sep=large, row sep=large]
\mathcal A
    \arrow[r, hook, "i"]
    \arrow[dr, "\simeq"']
&
P_{\mathbb V}(X)
\\
&
\Fix(J)
    \arrow[u, hook']
\end{tikzcd}
\]
where \(\Fix(J)\) denotes the full subcategory of \(J\)-local objects.

\begin{proposition}[Closed colimits]
\label{prop:closed-colimits-categorical}
Let
\[
    a\dashv i:
    P_{\mathbb V}(X)\rightleftarrows \mathcal A
\]
be a categorical fuzzy-frame presentation. Let \(K\) be small, let
\[
    W:K^{\op}\to\underline{\mathbb V}
\]
be a weight, and let
\[
    D:K\to\mathcal A
\]
be a diagram. Whenever the weighted colimit \(W\star iD\) exists in
\(P_{\mathbb V}(X)\), the weighted colimit of \(D\) exists in \(\mathcal A\) and is
computed by
\[
    W\star_{\mathcal A}D
    \cong
    a(W\star iD).
\]
\end{proposition}

\begin{proof}
For \(A\in\mathcal A\), the enriched adjunction and the universal property of the
weighted colimit give
\[
\begin{aligned}
\mathcal A(a(W\star iD),A)
&\cong
P_{\mathbb V}(X)(W\star iD,iA)
\\
&\cong
[K^{\op},\underline{\mathbb V}]
\bigl(
W,
P_{\mathbb V}(X)(iD-,iA)
\bigr)
\\
&\cong
[K^{\op},\underline{\mathbb V}]
\bigl(
W,
\mathcal A(D-,A)
\bigr).
\end{aligned}
\]
This is the enriched universal property of the weighted colimit of \(D\) in
\(\mathcal A\).
\end{proof}

\begin{proposition}[Finite limits created by the inclusion]
\label{prop:finite-limits-created-categorical}
Let
\[
    a\dashv i:
    P_{\mathbb V}(X)\rightleftarrows \mathcal A
\]
be a categorical fuzzy-frame presentation. Let \(K\) and \(W\) belong to the chosen
finite-limit class \(\mathcal F\), and let
\[
    D:K\to\mathcal A
\]
be a diagram. Then the corresponding finite weighted limits in \(\mathcal A\) are
created by the inclusion
\[
    i:\mathcal A\to P_{\mathbb V}(X).
\]
\end{proposition}

\begin{proof}
Let
\[
    M=\{W,iD\}
\]
be the corresponding finite weighted limit in \(P_{\mathbb V}(X)\). Since \(a\)
preserves the chosen finite weighted limits,
\[
    aM
    \cong
    a\{W,iD\}
    \cong
    \{W,aiD\}
    \cong
    \{W,D\}
\]
in \(\mathcal A\).

Since \(i\) is a right adjoint, it preserves limits. Hence \(iaM\) is also a limit of
\(iD\) in \(P_{\mathbb V}(X)\). But \(M\) is already the limit of \(iD\). The unit
\[
    M\to iaM
\]
is the unique comparison map between two limiting cones over \(iD\). Since both
objects satisfy the same finite-limit universal property in \(P_{\mathbb V}(X)\),
this comparison is an isomorphism. Thus \(M\) is local, and the chosen finite limits
in \(\mathcal A\) are created by \(i\).
\end{proof}

Thus, categorically, localization has the expected frame-theoretic behavior: colimits
are formed in the ambient presheaf category and then closed, while finite limits are
inherited by the local objects.

\subsection{Posetal reflection and fuzzy nuclei}
\label{subsec:localizing-presheaves-posetal-reflection}

We now pass to the posetal specialization. The categorical localization above becomes
a closure operation on
\[
    P_\Lambda(X),
\]
compatible with finite meets and with the scalar action inherited from \(\Lambda\).
The first three axioms below are the usual nucleus axioms from frame theory
\cite{Johnstone1982}. The final axiom is the additional compatibility with fuzzy
truth-value rescaling.

\begin{definition}[Fuzzy nucleus]
\label{def:fuzzy-nucleus}
Let \(X\) be a fuzzy preorder. A \textbf{fuzzy nucleus} on \(P_\Lambda(X)\) is a map
\[
    j:P_\Lambda(X)\to P_\Lambda(X)
\]
such that, for all \(\phi,\psi\in P_\Lambda(X)\) and all \(r\in L\),
\[
    \phi\leq j(\phi),
\]
\[
    j(j(\phi))=j(\phi),
\]
\[
    j(\phi\meet\psi)=j(\phi)\meet j(\psi),
\]
and
\[
    j(r\tensor j(\phi))=j(r\tensor\phi).
\]
\end{definition}

The first three equations are equivalent to the usual nucleus axioms on the frame
\(P_\Lambda(X)\); monotonicity follows from finite-meet preservation and is recorded
below. The final equation says that localization commutes with fuzzy truth-value
rescaling after closure. This scalar-compatible form is close in spirit to nucleus
constructions in quantale theory \cite{Rosenthal1990}.

\begin{lemma}[Fuzzy nuclei are monotone]
\label{lem:fuzzy-nuclei-monotone}
Every fuzzy nucleus is monotone.
\end{lemma}

\begin{proof}
If \(\phi\leq\psi\), then
\[
    \phi=\phi\meet\psi.
\]
Therefore
\[
    j(\phi)
    =
    j(\phi\meet\psi)
    =
    j(\phi)\meet j(\psi)
    \leq
    j(\psi).
\]
\end{proof}

\begin{lemma}[Fuzzy nuclei preserve top]
\label{lem:fuzzy-nuclei-preserve-top}
Every fuzzy nucleus satisfies
\[
    j(\top_{P_\Lambda(X)})=\top_{P_\Lambda(X)}.
\]
\end{lemma}

\begin{proof}
Inflationarity gives
\[
    \top_{P_\Lambda(X)}\leq j(\top_{P_\Lambda(X)}).
\]
Since \(\top_{P_\Lambda(X)}\) is the greatest element of \(P_\Lambda(X)\), this
forces
\[
    j(\top_{P_\Lambda(X)})=\top_{P_\Lambda(X)}.
\]
\end{proof}

\begin{lemma}[Scalars as enriched tensors]
\label{lem:scalars-as-enriched-tensors}
Let \(P=P_\Lambda(X)\). For every \(r\in L\), the pointwise scalar operation
\[
    r\tensor(-):P\to P
\]
is the tensor of the \(\Lambda\)-category \(P\) by the scalar \(r\). Explicitly, for
all \(\phi,\psi\in P\),
\[
    P(r\tensor\phi,\psi)
    =
    r\multimap P(\phi,\psi).
\]
\end{lemma}

\begin{proof}
By the scalar-closure result for fuzzy predicates from the previous chapter,
\(r\tensor\phi\) is again an element of \(P\). By definition,
\[
    P(r\tensor\phi,\psi)
    =
    \bigmeet_{x\in X}
    \bigl((r\tensor\phi(x))\multimap\psi(x)\bigr).
\]
For \(q\in L\), residuation gives
\[
\begin{aligned}
q\leq P(r\tensor\phi,\psi)
&\Longleftrightarrow
q\tensor r\tensor\phi(x)\leq\psi(x)
\quad\text{for all }x
\\
&\Longleftrightarrow
q\tensor r\leq \phi(x)\multimap\psi(x)
\quad\text{for all }x
\\
&\Longleftrightarrow
q\tensor r\leq
\bigmeet_{x\in X}(\phi(x)\multimap\psi(x))
\\
&\Longleftrightarrow
q\tensor r\leq P(\phi,\psi)
\\
&\Longleftrightarrow
q\leq r\multimap P(\phi,\psi).
\end{aligned}
\]
Hence
\[
    P(r\tensor\phi,\psi)
    =
    r\multimap P(\phi,\psi).
\]
\end{proof}

\begin{lemma}[Tensors in a reflective localization]
\label{lem:tensors-in-reflective-localization}
Let
\[
    a\dashv i:
    P_\Lambda(X)\rightleftarrows\mathcal A
\]
be a \(\Lambda\)-enriched reflective localization with \(i\) fully faithful. Then
\(\mathcal A\) is tensored over \(\Lambda\). For \(r\in L\) and \(A\in\mathcal A\),
the tensor is computed by
\[
    r\tensor_{\mathcal A}A
    \cong
    a(r\tensor iA).
\]
Moreover,
\[
    a(r\tensor\phi)
    \cong
    r\tensor_{\mathcal A}a\phi
\]
for all \(\phi\in P_\Lambda(X)\).
\end{lemma}

\begin{proof}
For \(B\in\mathcal A\), using the enriched adjunction and Lemma
\ref{lem:scalars-as-enriched-tensors}, we compute
\[
\begin{aligned}
\mathcal A(a(r\tensor iA),B)
&\cong
P_\Lambda(X)(r\tensor iA,iB)
\\
&\cong
r\multimap P_\Lambda(X)(iA,iB)
\\
&\cong
r\multimap \mathcal A(A,B).
\end{aligned}
\]
Thus \(a(r\tensor iA)\) satisfies the universal property of
\(r\tensor_{\mathcal A}A\).

Taking \(A=a\phi\) and using the counit \(aia\phi\cong a\phi\), we have
\[
    r\tensor_{\mathcal A}a\phi
    \cong
    a(r\tensor ia\phi).
\]
Also, since \(a\) is an enriched left adjoint, it preserves tensors. Hence
\[
    a(r\tensor\phi)
    \cong
    r\tensor_{\mathcal A}a\phi.
\]
\end{proof}

\begin{theorem}[Posetal comparison for localizations and fuzzy nuclei]
\label{thm:posetal-comparison-localizations-nuclei}
Let \(X\) be a fuzzy preorder, and write
\[
    P=P_\Lambda(X).
\]
There is a correspondence, up to equivalence of reflective localizations over \(P\),
between:
\begin{enumerate}[label=\textup{(\arabic*)}]
    \item left-exact \(\Lambda\)-enriched reflective localizations
    \[
        a\dashv i:
        P\rightleftarrows\mathcal A
    \]
    with \(i\) fully faithful, where left exactness in the posetal reflection means
    preservation of the terminal element and binary meets in the underlying frame;

    \item fuzzy nuclei
    \[
        j:P\to P.
    \]
\end{enumerate}
Under this correspondence, a localization determines the fuzzy nucleus
\[
    j=ia,
\]
and a fuzzy nucleus \(j\) determines the reflective localization
\[
    q_j\dashv i_j:
    P\rightleftarrows P_j,
\]
where
\[
    P_j=\{\phi\in P\mid j(\phi)=\phi\},
    \qquad
    q_j(\phi)=j(\phi),
\]
and \(i_j:P_j\hookrightarrow P\) is the inclusion. Moreover,
\[
    \mathcal A\simeq P_j
\]
over \(P\).

In the posetal specialization, the underlying crisp order of \(P\) is a genuine
poset. Since \(P\) is separated as a \(\Lambda\)-preorder, any ordinary isomorphism
between objects of \(P\) is an equality. Thus, after replacing a localization by its
equivalent full subpreorder of local objects, the idempotent monad \(ia\) may be
regarded strictly as a closure operator on the underlying predicate frame.
\end{theorem}

\begin{proof}
First let
\[
    a\dashv i:
    P\rightleftarrows\mathcal A
\]
be a left-exact \(\Lambda\)-enriched reflective localization with \(i\) fully
faithful, and define
\[
    j=ia.
\]
The unit of the adjunction gives
\[
    \phi\leq ia\phi=j(\phi),
\]
so \(j\) is inflationary. Since \(i\) is fully faithful, the counit
\[
    aiA\to A
\]
is an isomorphism for every \(A\in\mathcal A\). Hence
\[
    j(j(\phi))
    =
    iaia\phi
    \cong
    ia\phi
    =
    j(\phi).
\]
Since \(P\) is separated, this isomorphism may be read as the strict equality
\[
    j(j(\phi))=j(\phi).
\]
Thus \(j\) is idempotent.

Because \(a\) preserves the finite meets of the underlying frame and \(i\), as a
right adjoint, preserves them, we have
\[
\begin{aligned}
j(\phi\meet\psi)
&=
ia(\phi\meet\psi)
\\
&\cong
i(a\phi\meet a\psi)
\\
&\cong
ia\phi\meet ia\psi
\\
&=
j(\phi)\meet j(\psi).
\end{aligned}
\]
Again using separatedness of \(P\), this comparison is a strict equality:
\[
    j(\phi\meet\psi)=j(\phi)\meet j(\psi).
\]
Similarly,
\[
    j(\top_P)=\top_P.
\]

It remains to prove scalar compatibility. By Lemma
\ref{lem:tensors-in-reflective-localization},
\[
    a(r\tensor\phi)
    \cong
    r\tensor_{\mathcal A}a\phi
    \cong
    a(r\tensor ia\phi).
\]
Therefore
\[
\begin{aligned}
j(r\tensor j(\phi))
&=
ia(r\tensor ia\phi)
\\
&\cong
ia(r\tensor\phi)
\\
&=
j(r\tensor\phi).
\end{aligned}
\]
Using separatedness once more, this gives the strict equality
\[
    j(r\tensor j(\phi))=j(r\tensor\phi).
\]
Thus \(j\) is a fuzzy nucleus.

Conversely, let \(j:P\to P\) be a fuzzy nucleus, and let
\[
    P_j=\{\phi\in P\mid j(\phi)=\phi\}
\]
be the full \(\Lambda\)-subpreorder of fixed points. Define
\[
    q_j:P\to P_j,
    \qquad
    q_j(\phi)=j(\phi),
\]
and let
\[
    i_j:P_j\hookrightarrow P
\]
be the inclusion.

First \(q_j\) is a \(\Lambda\)-functor. Let \(r\leq P(\phi,\psi)\). Equivalently,
\[
    r\tensor\phi\leq\psi.
\]
By monotonicity of \(j\),
\[
    j(r\tensor\phi)\leq j(\psi).
\]
By inflationarity applied to \(r\tensor j(\phi)\), followed by scalar compatibility,
\[
    r\tensor j(\phi)
    \leq
    j(r\tensor j(\phi))
    =
    j(r\tensor\phi).
\]
Hence
\[
    r\tensor j(\phi)\leq j(\psi),
\]
and so
\[
    r\leq P(j\phi,j\psi).
\]
Therefore
\[
    P(\phi,\psi)\leq P_j(q_j\phi,q_j\psi),
\]
so \(q_j\) is a \(\Lambda\)-functor.

We now prove that
\[
    q_j\dashv i_j.
\]
Let \(\psi\in P_j\). Since \(P_j\) is full,
\[
    P_j(q_j\phi,\psi)=P(j\phi,\psi).
\]
We prove
\[
    P(j\phi,\psi)=P(\phi,\psi).
\]
For \(r\in L\), if
\[
    r\leq P(j\phi,\psi),
\]
then
\[
    r\tensor j\phi\leq\psi.
\]
Since \(\phi\leq j\phi\), this implies
\[
    r\tensor\phi\leq\psi,
\]
so
\[
    r\leq P(\phi,\psi).
\]
Conversely, if
\[
    r\tensor\phi\leq\psi,
\]
then
\[
    j(r\tensor\phi)\leq j(\psi)=\psi.
\]
By inflationarity and scalar compatibility,
\[
    r\tensor j\phi
    \leq
    j(r\tensor j\phi)
    =
    j(r\tensor\phi)
    \leq
    \psi.
\]
Thus
\[
    r\leq P(j\phi,\psi).
\]
Therefore
\[
    P(j\phi,\psi)=P(\phi,\psi),
\]
and hence \(q_j\dashv i_j\) as a \(\Lambda\)-adjunction.

The inclusion \(i_j\) is fully faithful by definition. The reflector \(q_j\) preserves
terminal element and binary meets because
\[
    j(\top_P)=\top_P
\]
and
\[
    j(\phi\meet\psi)=j(\phi)\meet j(\psi).
\]
Thus
\[
    q_j\dashv i_j:
    P\rightleftarrows P_j
\]
is a left-exact \(\Lambda\)-enriched reflective localization.

Finally, starting from a localization and applying the construction above gives the
fixed-point subcategory of the idempotent monad \(j=ia\), which is equivalent to
\(\mathcal A\) by Proposition~\ref{prop:local-objects-categorical}. Starting from a
fuzzy nucleus \(j\), the induced monad is exactly \(i_jq_j=j\). Hence the two
constructions are inverse up to equivalence over \(P\).
\end{proof}

\subsection{Fixed points of fuzzy nuclei}
\label{subsec:fixed-points-fuzzy-nuclei}

\begin{lemma}[Intersections of fuzzy nuclei]
\label{lem:intersections-fuzzy-nuclei}
Let \((j_\alpha)_{\alpha\in I}\) be a family of fuzzy nuclei on \(P_\Lambda(X)\).
Then the pointwise meet
\[
    j(\phi)\eqdef\bigmeet_{\alpha\in I}j_\alpha(\phi)
\]
is again a fuzzy nucleus. If \(I=\varnothing\), this is the constant-top nucleus.
\end{lemma}

\begin{proof}
Inflationarity follows because \(\phi\leq j_\alpha(\phi)\) for every \(\alpha\).
Since every \(j_\alpha\) is monotone, \(j\) is monotone.

For idempotence, inflationarity gives
\[
    j(\phi)\leq j(j(\phi)).
\]
Conversely, since
\[
    j(\phi)\leq j_\alpha(\phi)
\]
for every \(\alpha\), monotonicity of \(j_\alpha\) gives
\[
    j_\alpha(j(\phi))
    \leq
    j_\alpha(j_\alpha(\phi))
    =
    j_\alpha(\phi).
\]
Taking the meet over \(\alpha\) yields
\[
    j(j(\phi))\leq j(\phi).
\]
Thus
\[
    j(j(\phi))=j(\phi).
\]

For finite meets,
\[
\begin{aligned}
j(\phi\meet\psi)
&=
\bigmeet_\alpha j_\alpha(\phi\meet\psi)
\\
&=
\bigmeet_\alpha\bigl(j_\alpha(\phi)\meet j_\alpha(\psi)\bigr)
\\
&=
\left(\bigmeet_\alpha j_\alpha(\phi)\right)
\meet
\left(\bigmeet_\alpha j_\alpha(\psi)\right)
\\
&=
j(\phi)\meet j(\psi).
\end{aligned}
\]

Finally, for scalar compatibility, inflationarity gives
\[
    r\tensor\phi\leq r\tensor j(\phi),
\]
and hence, by monotonicity of \(j\),
\[
    j(r\tensor\phi)\leq j(r\tensor j(\phi)).
\]
For the reverse inequality, since \(j(\phi)\leq j_\alpha(\phi)\), we have
\[
    r\tensor j(\phi)\leq r\tensor j_\alpha(\phi).
\]
Thus
\[
    j_\alpha(r\tensor j(\phi))
    \leq
    j_\alpha(r\tensor j_\alpha(\phi))
    =
    j_\alpha(r\tensor\phi).
\]
Taking the meet over \(\alpha\) gives
\[
    j(r\tensor j(\phi))
    \leq
    j(r\tensor\phi).
\]
Therefore
\[
    j(r\tensor j(\phi))=j(r\tensor\phi).
\]
\end{proof}

\begin{definition}[Fixed points of a fuzzy nucleus]
\label{def:fixed-points-fuzzy-nucleus}
Let \(j\) be a fuzzy nucleus on \(P_\Lambda(X)\). Define
\[
    P_\Lambda(X)_j
    \eqdef
    \{\phi\in P_\Lambda(X)\mid j(\phi)=\phi\}.
\]
For fixed points \(\phi_i,\phi,\psi\in P_\Lambda(X)_j\), define
\[
    \bigjoin\nolimits_j \phi_i
    =
    j\left(\bigjoin_i \phi_i\right),
\]
\[
    \phi\meet_j\psi
    =
    \phi\meet\psi,
\]
and
\[
    r\tensor_j\phi
    =
    j(r\tensor\phi).
\]
\end{definition}

\begin{proposition}[Fixed points of a fuzzy nucleus]
\label{prop:fixed-points-fuzzy-nucleus-frame}
Let \(j\) be a fuzzy nucleus on \(P_\Lambda(X)\). Then
\[
    P_\Lambda(X)_j
\]
is a frame. Its joins are closed joins,
\[
    \bigjoin\nolimits_j \phi_i
    =
    j\left(\bigjoin_i\phi_i\right),
\]
its finite meets are inherited from \(P_\Lambda(X)\), and it carries the closed scalar
action
\[
    r\tensor_j\phi
    =
    j(r\tensor\phi).
\]
The closed scalar action preserves arbitrary joins in each variable.
\end{proposition}

\begin{proof}
First, \(\top_{P_\Lambda(X)}\) is fixed by Lemma~\ref{lem:fuzzy-nuclei-preserve-top}.
If \(\phi,\psi\) are fixed, then
\[
    j(\phi\meet\psi)
    =
    j(\phi)\meet j(\psi)
    =
    \phi\meet\psi,
\]
so finite meets are inherited from \(P_\Lambda(X)\).

For a family \((\phi_i)\) of fixed points, define
\[
    c=j\left(\bigjoin_i\phi_i\right).
\]
Then \(c\) is fixed by idempotence. Since
\[
    \phi_i\leq \bigjoin_i\phi_i\leq j\left(\bigjoin_i\phi_i\right),
\]
the element \(c\) is an upper bound of the \(\phi_i\). If \(\psi\) is any fixed upper
bound of the \(\phi_i\), then
\[
    \bigjoin_i\phi_i\leq\psi.
\]
By monotonicity,
\[
    j\left(\bigjoin_i\phi_i\right)\leq j(\psi)=\psi.
\]
Thus \(c\) is the least fixed upper bound. Hence joins in \(P_\Lambda(X)_j\) are
closed joins. In particular, the bottom element is
\[
    \bot_j=j(\bot_{P_\Lambda(X)}).
\]

Frame distributivity follows from frame distributivity in \(P_\Lambda(X)\) and
finite-meet preservation by \(j\). If \(\phi\) and all \(\psi_i\) are fixed, then
\[
\begin{aligned}
\phi\meet_j\bigjoin\nolimits_j \psi_i
&=
\phi\meet j\left(\bigjoin_i \psi_i\right)
\\
&=
j(\phi)\meet j\left(\bigjoin_i \psi_i\right)
\\
&=
j\left(\phi\meet\bigjoin_i \psi_i\right)
\\
&=
j\left(\bigjoin_i(\phi\meet\psi_i)\right)
\\
&=
\bigjoin\nolimits_j(\phi\meet_j\psi_i).
\end{aligned}
\]

The closed scalar action is fixed by definition:
\[
    r\tensor_j\phi=j(r\tensor\phi).
\]
The scalar unit law follows from the unit law in \(P_\Lambda(X)\) and the fact that
\(\phi\) is fixed:
\[
    \top_L\tensor_j\phi
    =
    j(\top_L\tensor\phi)
    =
    j(\phi)
    =
    \phi.
\]
For associativity, use the scalar-compatibility equation:
\[
\begin{aligned}
r\tensor_j(s\tensor_j\phi)
&=
j(r\tensor j(s\tensor\phi))
\\
&=
j(r\tensor(s\tensor\phi))
\\
&=
j((r\tensor s)\tensor\phi)
\\
&=
(r\tensor s)\tensor_j\phi.
\end{aligned}
\]

It remains to verify that the closed scalar action preserves arbitrary joins. In the
frame variable,
\[
\begin{aligned}
r\tensor_j\bigjoin\nolimits_j \phi_i
&=
j\left(r\tensor j\left(\bigjoin_i\phi_i\right)\right)
\\
&=
j\left(r\tensor\bigjoin_i\phi_i\right)
\\
&=
j\left(\bigjoin_i(r\tensor\phi_i)\right).
\end{aligned}
\]
For any family \((\eta_i)\) in \(P_\Lambda(X)\), the closure identity
\[
    j\left(\bigjoin_i \eta_i\right)
    =
    j\left(\bigjoin_i j(\eta_i)\right)
\]
follows from inflationarity, monotonicity, and idempotence of \(j\). Applying this to
\(\eta_i=r\tensor\phi_i\), we get
\[
\begin{aligned}
j\left(\bigjoin_i(r\tensor\phi_i)\right)
&=
j\left(\bigjoin_i j(r\tensor\phi_i)\right)
\\
&=
\bigjoin\nolimits_j(r\tensor_j\phi_i).
\end{aligned}
\]
Therefore
\[
    r\tensor_j\bigjoin\nolimits_j \phi_i
    =
    \bigjoin\nolimits_j(r\tensor_j\phi_i).
\]

The scalar variable is similar. For a family \((r_i)\) in \(L\),
\[
\begin{aligned}
\left(\bigjoin_i r_i\right)\tensor_j\phi
&=
j\left(\left(\bigjoin_i r_i\right)\tensor\phi\right)
\\
&=
j\left(\bigjoin_i(r_i\tensor\phi)\right)
\\
&=
j\left(\bigjoin_i j(r_i\tensor\phi)\right)
\\
&=
\bigjoin\nolimits_j(r_i\tensor_j\phi).
\end{aligned}
\]
Thus the closed scalar action preserves arbitrary joins in each variable.
\end{proof}

\begin{remark}[Bottom scalars]
\label{rem:bottom-scalars}
Taking the empty join in the scalar variable gives
\[
    \bot_L\tensor_j\phi
    =
    \bot_j.
\]
Taking the empty join in the frame variable gives
\[
    r\tensor_j\bot_j
    =
    \bot_j.
\]
Thus each scalar operation preserves bottom, and the bottom scalar acts as the
constant-bottom operator. Notice that
\[
    \bot_j=j(\bot_{P_\Lambda(X)}),
\]
which need not equal the bottom element of \(P_\Lambda(X)\).
\end{remark}

\begin{lemma}[Enriched hom recovered from the scalar action]
\label{lem:enriched-hom-recovered-from-scalar-action}
Let
\[
    A=P_\Lambda(X)_j
\]
be the fixed-point frame of a fuzzy nucleus. For fixed points
\(\phi,\psi\in A\), one has
\[
    A(\phi,\psi)
    =
    \bigjoin
    \{r\in L\mid r\tensor_j\phi\leq \psi\}.
\]
Moreover this hom-value agrees with the hom inherited from \(P_\Lambda(X)\):
\[
    A(\phi,\psi)=P_\Lambda(X)(\phi,\psi).
\]
\end{lemma}

\begin{proof}
For \(r\in L\), using the tensor-hom formula in \(P_\Lambda(X)\), we have
\[
    r\leq P_\Lambda(X)(\phi,\psi)
    \quad\Longleftrightarrow\quad
    r\tensor\phi\leq \psi.
\]
Since \(\psi\) is fixed, this is equivalent to
\[
    j(r\tensor\phi)\leq \psi.
\]
But
\[
    j(r\tensor\phi)=r\tensor_j\phi.
\]
Hence
\[
    r\leq P_\Lambda(X)(\phi,\psi)
    \quad\Longleftrightarrow\quad
    r\tensor_j\phi\leq \psi.
\]
Taking the join over all such \(r\) gives the claim.
\end{proof}

\begin{definition}[\(\Lambda\)-fuzzy-frame presentation]
\label{def:fuzzy-frame-presentation}
A \textbf{\(\Lambda\)-fuzzy-frame presentation} is a pair
\[
    (X,j),
\]
where \(X\) is a fuzzy preorder and
\[
    j:P_\Lambda(X)\to P_\Lambda(X)
\]
is a fuzzy nucleus.

The frame presented by \((X,j)\) is the fixed-point frame
\[
    P_\Lambda(X)_j.
\]
\end{definition}

\begin{definition}[Fuzzy frame over \(\Lambda\)]
\label{def:fuzzy-frame}
A \textbf{fuzzy frame over \(\Lambda\)} is a frame presented by a
\(\Lambda\)-fuzzy-frame presentation. Equivalently, a fuzzy frame is a nuclear
quotient of a fuzzy presheaf frame:
\[
    P_\Lambda(X)\twoheadrightarrow P_\Lambda(X)_j.
\]
\end{definition}

Thus the phrase \emph{nuclear quotient} is frame-theoretic. Under locale duality, the
same structure is read as a fuzzy sublocale.

\begin{definition}[\(\Lambda\)-scalar frame]
\label{def:lambda-scalar-frame}
A \textbf{\(\Lambda\)-scalar frame} is a frame \(A\) equipped with an action
\[
    L\times A\to A,
    \qquad
    (r,a)\longmapsto r\tensor_A a,
\]
such that
\[
    \top_L\tensor_A a=a,
    \qquad
    (r\tensor s)\tensor_A a=r\tensor_A(s\tensor_A a),
\]
and such that the operation preserves arbitrary joins in each variable:
\[
    r\tensor_A\bigjoin_i a_i
    =
    \bigjoin_i(r\tensor_A a_i),
\]
\[
    \left(\bigjoin_i r_i\right)\tensor_A a
    =
    \bigjoin_i(r_i\tensor_A a).
\]
\end{definition}

Every fixed-point frame \(P_\Lambda(X)_j\) is a \(\Lambda\)-scalar frame by
Proposition~\ref{prop:fixed-points-fuzzy-nucleus-frame}.

\begin{definition}[Fuzzy-frame homomorphism]
\label{def:fuzzy-frame-homomorphism}
Let
\[
    P_\Lambda(X)_j
    \qquad\text{and}\qquad
    P_\Lambda(Y)_k
\]
be fuzzy frames presented by \(\Lambda\)-fuzzy-frame presentations. A
\textbf{fuzzy-frame homomorphism}
\[
    h:P_\Lambda(X)_j\to P_\Lambda(Y)_k
\]
is a frame homomorphism satisfying
\[
    h(r\tensor_j\phi)=r\tensor_k h(\phi)
\]
for all \(r\in L\) and all \(\phi\in P_\Lambda(X)_j\).
\end{definition}

\begin{corollary}[Fuzzy-frame homomorphisms are enriched functors]
\label{cor:fuzzy-frame-homomorphisms-are-enriched}
Let
\[
    h:A\to B
\]
be a fuzzy-frame homomorphism between fixed-point fuzzy frames. Then \(h\) is a
\(\Lambda\)-functor for the enriched homs recovered from the scalar actions.
\end{corollary}

\begin{proof}
Suppose
\[
    r\leq A(a,b).
\]
By Lemma~\ref{lem:enriched-hom-recovered-from-scalar-action}, this means
\[
    r\tensor_A a\leq b.
\]
Applying \(h\) and using scalar compatibility gives
\[
    r\tensor_B h(a)
    =
    h(r\tensor_A a)
    \leq
    h(b).
\]
Again by Lemma~\ref{lem:enriched-hom-recovered-from-scalar-action}, this implies
\[
    r\leq B(h(a),h(b)).
\]
Therefore
\[
    A(a,b)\leq B(h(a),h(b)),
\]
so \(h\) is a \(\Lambda\)-functor.
\end{proof}

\begin{definition}[The category of fuzzy frames]
\label{def:category-fuzzy-frames}
The category of fuzzy frames over \(\Lambda\) and fuzzy-frame homomorphisms is denoted
\[
    \FFrm_\Lambda.
\]
\end{definition}

\begin{remark}[Presentations]
\label{rem:fuzzy-frames-presentation-based}
Strictly speaking, one may keep the presentation \((X,j)\) as part of the data, or
identify presentations that present isomorphic fixed-point frames. In this chapter we
write fuzzy frames as the presented frames \(P_\Lambda(X)_j\), while remembering that
they arise from some localization of a fuzzy presheaf frame.

The definition is deliberately presentation-based. Fuzzy frames are not introduced as
arbitrary frames equipped with a scalar action; they are the fixed-point frames of
scalar-compatible nuclei on fuzzy presheaf frames.
\end{remark}

\begin{definition}[Nuclear quotient map]
\label{def:nuclear-quotient-map}
Let \(j\) be a fuzzy nucleus on \(P_\Lambda(X)\). The associated \textbf{nuclear
quotient map} is the closure map
\[
    q_j:
    P_\Lambda(X)\to P_\Lambda(X)_j,
    \qquad
    q_j(\phi)=j(\phi).
\]
\end{definition}

\begin{proposition}[The closure map is a quotient]
\label{prop:closure-map-is-quotient}
Let \(j\) be a fuzzy nucleus on \(P_\Lambda(X)\). Then
\[
    q_j:
    P_\Lambda(X)\to P_\Lambda(X)_j
\]
is a surjective fuzzy-frame homomorphism.
\end{proposition}

\begin{proof}
The map \(q_j\) is surjective because every fixed point \(\phi\in P_\Lambda(X)_j\)
satisfies
\[
    j(\phi)=\phi.
\]

It preserves arbitrary joins because joins in \(P_\Lambda(X)_j\) are closed joins:
\[
    j\left(\bigjoin_i\phi_i\right)
    =
    \bigjoin\nolimits_j j(\phi_i).
\]
For the empty family this says
\[
    q_j(\bot_{P_\Lambda(X)})=j(\bot_{P_\Lambda(X)})=\bot_j,
\]
so \(q_j\) also preserves bottom. It preserves finite meets because
\[
    j(\phi\meet\psi)=j(\phi)\meet j(\psi),
\]
and it preserves top by Lemma~\ref{lem:fuzzy-nuclei-preserve-top}.

It preserves scalars because
\[
\begin{aligned}
q_j(r\tensor\phi)
&=
j(r\tensor\phi)
\\
&=
j(r\tensor j(\phi))
\\
&=
r\tensor_j j(\phi)
\\
&=
r\tensor_j q_j(\phi).
\end{aligned}
\]
Therefore \(q_j\) is a surjective fuzzy-frame homomorphism.
\end{proof}

Conversely, every scalar-compatible quotient of a fuzzy presheaf frame determines a
fuzzy nucleus. This is the scalar-compatible version of the usual correspondence
between nuclei and quotient frames \cite{Johnstone1982}.

\begin{proposition}[Quotients determine fuzzy nuclei]
\label{prop:quotients-determine-fuzzy-nuclei}
Let \(B\) be a \(\Lambda\)-scalar frame, and let
\[
    q:P_\Lambda(X)\twoheadrightarrow B
\]
be a surjective frame homomorphism satisfying
\[
    q(r\tensor\phi)=r\tensor_B q(\phi)
\]
for all \(r\in L\) and \(\phi\in P_\Lambda(X)\). Since \(q\) preserves arbitrary
joins, it has a right adjoint
\[
    q_\ast:B\to P_\Lambda(X)
\]
defined by
\[
    q_\ast(b)
    =
    \bigjoin\{\phi\in P_\Lambda(X)\mid q(\phi)\leq b\}.
\]
Then
\[
    j_q\eqdef q_\ast q
\]
is a fuzzy nucleus on \(P_\Lambda(X)\), and the fixed-point fuzzy frame
\[
    P_\Lambda(X)_{j_q}
\]
is isomorphic to \(B\) as a fuzzy frame.
\end{proposition}

\begin{proof}
Since \(q\dashv q_\ast\), the composite
\[
    j_q=q_\ast q
\]
is inflationary. Since \(q\) is surjective, the counit
\[
    qq_\ast\leq \id_B
\]
is an equality. Indeed, if \(b=q(\phi)\), then \(\phi\leq q_\ast(b)\), so
\[
    b=q(\phi)\leq qq_\ast(b),
\]
while \(qq_\ast(b)\leq b\) follows from the counit inequality. Hence
\[
    qq_\ast=\id_B.
\]
Therefore
\[
    j_qj_q
    =
    q_\ast q q_\ast q
    =
    q_\ast q
    =
    j_q.
\]

Because \(q\) preserves finite meets and \(q_\ast\) preserves all meets,
\[
\begin{aligned}
j_q(\phi\meet\psi)
&=
q_\ast q(\phi\meet\psi)
\\
&=
q_\ast(q(\phi)\meet q(\psi))
\\
&=
q_\ast q(\phi)\meet q_\ast q(\psi)
\\
&=
j_q(\phi)\meet j_q(\psi).
\end{aligned}
\]

For scalar compatibility, use that \(q\) preserves scalars and that
\(qq_\ast=\id_B\):
\[
\begin{aligned}
j_q(r\tensor j_q(\phi))
&=
q_\ast q(r\tensor q_\ast q(\phi))
\\
&=
q_\ast\bigl(r\tensor_B q q_\ast q(\phi)\bigr)
\\
&=
q_\ast(r\tensor_B q(\phi))
\\
&=
q_\ast q(r\tensor\phi)
\\
&=
j_q(r\tensor\phi).
\end{aligned}
\]
Therefore \(j_q\) is a fuzzy nucleus.

The restriction
\[
    q:P_\Lambda(X)_{j_q}\to B
\]
is an isomorphism of frames with inverse \(q_\ast\). First, for every \(b\in B\),
\[
    j_q(q_\ast b)
    =
    q_\ast q q_\ast b
    =
    q_\ast b,
\]
so \(q_\ast b\) is fixed. Moreover,
\[
    q q_\ast(b)=b,
\]
and for \(\phi\in P_\Lambda(X)_{j_q}\),
\[
    q_\ast q(\phi)=j_q(\phi)=\phi.
\]

It remains to note that this isomorphism respects the scalar actions. The restriction
of \(q\) preserves scalars by assumption. Conversely, for \(b\in B\), write
\(b=q(\phi)\). Then
\[
\begin{aligned}
q\bigl(r\tensor_{j_q}q_\ast(b)\bigr)
&=
q\bigl(j_q(r\tensor q_\ast(b))\bigr)
\\
&=
q(r\tensor q_\ast(b))
\\
&=
r\tensor_B qq_\ast(b)
\\
&=
r\tensor_B b.
\end{aligned}
\]
Since \(q:P_\Lambda(X)_{j_q}\to B\) is an isomorphism, this implies
\[
    r\tensor_{j_q}q_\ast(b)
    =
    q_\ast(r\tensor_B b).
\]
Thus
\[
    P_\Lambda(X)_{j_q}\cong B
\]
as fuzzy frames.
\end{proof}

\subsection{Examples}
\label{subsec:localizing-presheaves-examples}

\begin{example}[The identity localization]
\label{ex:identity-localization}
For every fuzzy preorder \(X\), the identity map
\[
    \id_{P_\Lambda(X)}:
    P_\Lambda(X)\to P_\Lambda(X)
\]
is a fuzzy nucleus. Hence
\[
    P_\Lambda(X)_{\id}
    =
    P_\Lambda(X).
\]
Thus every fuzzy presheaf frame is itself a fuzzy frame.

In this case no predicates are identified or closed further. Joins, finite meets, and
scalar multiplication are the pointwise operations inherited directly from
\(\Lambda\). Localically, this is the ambient presheaf fuzzy locale
\[
    \mathcal P_\Lambda X.
\]
All other examples of fuzzy frames in this chapter may be read as sublocales of such
presheaf fuzzy locales, equivalently as nuclear quotients of their frames of opens.
\end{example}

\begin{example}[The trivial localization]
\label{ex:terminal-localization}
For every fuzzy preorder \(X\), the constant-top map
\[
    j_\top(\phi)=\top_{P_\Lambda(X)}
\]
is a fuzzy nucleus. Inflationarity and idempotence are immediate, and
\[
    j_\top(\phi\meet\psi)=\top_{P_\Lambda(X)}
    =
    \top_{P_\Lambda(X)}\meet\top_{P_\Lambda(X)}.
\]
For scalar compatibility,
\[
    j_\top(r\tensor j_\top(\phi))
    =
    \top_{P_\Lambda(X)}
    =
    j_\top(r\tensor\phi).
\]
The fixed-point frame has exactly one element:
\[
    P_\Lambda(X)_{j_\top}\cong 1.
\]
This is the trivial fuzzy frame. On the locale side it corresponds to the initial, or
empty, fuzzy sublocale of the presheaf fuzzy locale. The identity nucleus gives the
whole sublocale, while the constant-top nucleus gives the least sublocale.
\end{example}

\begin{example}[The crisp case]
\label{ex:crisp-localization-case}
Let
\[
    \Lambda=2.
\]
Then fuzzy preorders are ordinary preorders and, for an ordinary preorder \(X\),
\[
    P_2(X)
    =
    [X^{\op},2]
    \cong
    \Down(X),
\]
the frame of downsets of \(X\). A fuzzy nucleus on \(P_2(X)\) is exactly an ordinary
nucleus on the frame \(\Down(X)\), because the scalar action of \(2\) is forced:
\[
    \top\tensor U=U,
    \qquad
    \bot\tensor U=\bot.
\]
Thus the scalar-compatibility equation adds no extra condition beyond the ordinary
nucleus axioms.

Concretely, a nucleus
\[
    j:\Down(X)\to\Down(X)
\]
assigns to each downset \(U\subseteq X\) a closed downset \(j(U)\), satisfying
\[
    U\subseteq j(U),
    \qquad
    j(j(U))=j(U),
    \qquad
    j(U\cap V)=j(U)\cap j(V).
\]
The fixed points form the ordinary quotient frame
\[
    \Down(X)_j.
\]
Thus the crisp specialization recovers the standard description of frames as
nuclear quotients of downset frames.
\end{example}

\begin{theorem}[Ordinary frames as localized presheaves]
\label{thm:ordinary-frames-localized-presheaves}
Let \(A\) be a small ordinary frame. Regard the underlying poset of \(A\) as a
preorder. Then the map
\[
    q:\Down(A)\to A,
    \qquad
    q(U)=\bigjoin U,
\]
is a surjective frame homomorphism. Consequently \(A\) is isomorphic to the
fixed-point frame of the nucleus
\[
    j_A=q_\ast q
\]
on \(\Down(A)\). Explicitly,
\[
    j_A(U)=\downarrow\bigjoin U.
\]
\end{theorem}

\begin{proof}
The map \(q\) preserves arbitrary joins because
\[
\begin{aligned}
q\left(\bigcup_i U_i\right)
&=
\bigjoin\bigcup_i U_i
\\
&=
\bigjoin_i\bigjoin U_i
\\
&=
\bigjoin_i q(U_i).
\end{aligned}
\]
It preserves top because
\[
    q(A)=\bigjoin A=\top_A.
\]

It remains to check finite meets. Since
\[
    U\cap V\subseteq U
    \qquad\text{and}\qquad
    U\cap V\subseteq V,
\]
we have
\[
    q(U\cap V)\leq q(U)\meet q(V).
\]
Conversely, by frame distributivity,
\[
\begin{aligned}
q(U)\meet q(V)
&=
\left(\bigjoin U\right)\meet\left(\bigjoin V\right)
\\
&=
\bigjoin_{u\in U,\;v\in V}(u\meet v).
\end{aligned}
\]
Since \(U\) and \(V\) are downsets, for every \(u\in U\) and \(v\in V\), the element
\[
    u\meet v
\]
belongs to \(U\cap V\). Hence
\[
    q(U)\meet q(V)\leq q(U\cap V).
\]
Therefore
\[
    q(U\cap V)=q(U)\meet q(V).
\]

Finally, \(q\) is surjective because
\[
    q(\downarrow a)=a
\]
for every \(a\in A\). Therefore \(q\) is a surjective frame homomorphism. Its right
adjoint is
\[
    q_\ast(a)=\downarrow a.
\]
Hence
\[
    j_A(U)
    =
    q_\ast q(U)
    =
    \downarrow\bigjoin U.
\]
Proposition~\ref{prop:quotients-determine-fuzzy-nuclei} gives
\[
    \Down(A)_{j_A}\cong A.
\]
\end{proof}

\begin{theorem}[The crisp case recovers ordinary frames]
\label{thm:crisp-case-recovers-ordinary-frames}
When
\[
    \Lambda=2,
\]
the small fuzzy frames over \(2\) are exactly the small ordinary frames, up to
isomorphism.
\end{theorem}

\begin{proof}
Every fuzzy frame over \(2\) is, by definition, a nuclear quotient of a frame
\[
    \Down(X).
\]
By Example~\ref{ex:crisp-localization-case}, this is an ordinary nuclear quotient of
an ordinary frame, and hence is an ordinary frame.

Conversely, every small ordinary frame is a nuclear quotient of an ordinary presheaf
frame by Theorem~\ref{thm:ordinary-frames-localized-presheaves}. Therefore every
small ordinary frame is a fuzzy frame over \(2\).
\end{proof}

This recovers the usual pointfree frame-theoretic localization picture
\cite{Johnstone1982}.

\begin{example}[G\"odel-valued localizations]
\label{ex:godel-valued-localizations}
Let
\[
    \Lambda=[0,1]
\]
with G\"odel tensor
\[
    r\tensor s=r\meet s=\min(r,s).
\]
A fuzzy predicate on a fuzzy preorder \(X\) is a map
\[
    \phi:X\to[0,1]
\]
satisfying
\[
    X(x',x)\meet\phi(x)\leq \phi(x')
\]
for all \(x,x'\in X\). The scalar action is truncation by the level \(r\):
\[
    (r\tensor\phi)(x)=\min(r,\phi(x)).
\]
A G\"odel-valued fuzzy nucleus closes predicates compatibly with truncation:
\[
    j(r\meet j(\phi))=j(r\meet\phi).
\]

Thus if a predicate has already been closed, then truncating it at level \(r\) and
closing again gives the same result as truncating the original predicate and closing.
The fixed-point frame consists of the predicates satisfying the chosen closure
condition. Joins are closed pointwise suprema, finite meets are pointwise minima, and
the scalar action is closed truncation:
\[
    r\tensor_j\phi=j(r\meet\phi).
\]
This gives a localic form of G\"odel-valued fuzzy topology, obtained here by enriched
presheaf localization rather than by postulating a topology of fuzzy opens directly
\cite{HohleRodabaugh1999}.
\end{example}

\begin{example}[{\L}ukasiewicz-valued localizations]
\label{ex:lukasiewicz-valued-localizations}
Let
\[
    \Lambda=[0,1]
\]
with {\L}ukasiewicz tensor
\[
    r\tensor s=\max(0,r+s-1).
\]
A fuzzy predicate on \(X\) is a map
\[
    \phi:X\to[0,1]
\]
satisfying
\[
    X(x',x)\tensor\phi(x)\leq \phi(x').
\]
The scalar action is no longer truncation. It shifts truth degrees downward:
\[
    (r\tensor\phi)(x)=\max(0,r+\phi(x)-1).
\]
For \(r=1\), this is the identity; for \(r<1\), all values are lowered, and values
below \(1-r\) are sent to \(0\).

A {\L}ukasiewicz-valued fuzzy nucleus is therefore a closure operator satisfying
\[
    j\bigl(r\tensor j(\phi)\bigr)=j(r\tensor\phi).
\]
The fixed-point fuzzy frame has closed joins, inherited finite meets, and closed
{\L}ukasiewicz scalar action. This example emphasizes that the scalar structure is not
merely order-theoretic: the same underlying lattice \([0,1]\) produces different
fuzzy-localic behavior depending on the chosen tensor.
\end{example}

\begin{example}[Product-valued localizations]
\label{ex:product-valued-localizations}
Let
\[
    \Lambda=[0,1]
\]
with product tensor
\[
    r\tensor s=rs.
\]
Then a fuzzy predicate on \(X\) is a map
\[
    \phi:X\to[0,1]
\]
satisfying
\[
    X(x',x)\phi(x)\leq \phi(x')
\]
for all \(x,x'\in X\). Scalar multiplication is ordinary multiplication:
\[
    (r\tensor\phi)(x)=r\phi(x).
\]
A product-valued fuzzy nucleus closes these predicates compatibly with product
rescaling:
\[
    j(r\,j(\phi))=j(r\,\phi).
\]

In the fixed-point frame, multiplying a closed predicate by \(r\) may fail to remain
closed, so the closed scalar action is
\[
    r\tensor_j\phi=j(r\phi).
\]
This can be interpreted as first weakening the degree of membership by the factor
\(r\), and then imposing the local closure condition. Product-valued localizations
therefore behave differently from G\"odel-valued localizations: G\"odel scalars cut
truth values from above, whereas product scalars contract them proportionally.
\end{example}

\begin{example}[Lawvere-valued localizations]
\label{ex:lawvere-valued-localizations}
Let
\[
    \Lambda=
    ([0,\infty],\leq_L,\meet_L,\join_L,+,\ominus,\infty,0),
    \qquad
    a\leq_L b
    \Longleftrightarrow
    a\geq b
    \text{ in the usual order}.
\]
Thus the tensor is addition, the top element is \(0\), and the bottom element is
\(\infty\). Joins in \(\Lambda\) are infima in the usual order, and meets in
\(\Lambda\) are suprema in the usual order.

A fuzzy preorder is a Lawvere metric space. A Lawvere-valued fuzzy predicate is a map
\[
    \phi:X\to[0,\infty]
\]
satisfying, in the usual order,
\[
    \phi(x')\leq d(x',x)+\phi(x).
\]
Thus \(\phi(x)\) is better read as a cost, distance, or potential rather than as a
truth degree. The scalar action adds a constant:
\[
    (r\tensor\phi)(x)=r+\phi(x).
\]
Since the order is reversed, inflationarity of a nucleus
\[
    \phi\leq_L j(\phi)
\]
means
\[
    j(\phi)(x)\leq \phi(x)
\]
in the usual order. Thus a Lawvere-valued nucleus lowers costs, subject to the chosen
closure condition.

The scalar-compatibility equation becomes
\[
    j(r+j(\phi))=j(r+\phi),
\]
where the displayed equation is written in the usual arithmetic notation on
\([0,\infty]\). The fixed-point frame consists of the closed cost predicates. Closed
joins are closures of pointwise infima, while finite meets are pointwise suprema.
This is the metric analogue of passing from arbitrary presheaves to local objects
\cite{Lawvere1973MetricSpaces}.
\end{example}

\section{Fuzzy locales}
\label{sec:fuzzy-locales}

We now pass from fuzzy frames to fuzzy locales by formal duality. As in ordinary
pointfree topology, frame homomorphisms are read contravariantly as inverse-image
maps of locales \cite{Johnstone1982}.

\subsection{Fuzzy locales as duals of fuzzy frames}
\label{subsec:fuzzy-locales-duals}

\begin{definition}[Fuzzy locale]
\label{def:fuzzy-locale}
The category of \textbf{fuzzy locales over \(\Lambda\)} is
\[
    \FLoc_\Lambda
    \eqdef
    \FFrm_\Lambda^{\op}.
\]
If \(\mathcal X\) is a fuzzy locale, its fuzzy frame of opens is denoted
\[
    \Opens_\Lambda(\mathcal X).
\]
A map of fuzzy locales
\[
    f:\mathcal X\to\mathcal Y
\]
is represented by a fuzzy-frame homomorphism in the opposite direction:
\[
    f^\ast:
    \Opens_\Lambda(\mathcal Y)
    \to
    \Opens_\Lambda(\mathcal X).
\]
\end{definition}

Thus the fundamental duality is
\[
\begin{tikzcd}[column sep=large, row sep=large]
\mathcal X
    \arrow[r, "f"]
&
\mathcal Y
\\
\Opens_\Lambda(\mathcal X)
&
\Opens_\Lambda(\mathcal Y)
    \arrow[l, "f^\ast"']
\end{tikzcd}
\]
and, categorically,
\[
\begin{tikzcd}[column sep=large]
\FLoc_\Lambda
    \arrow[r, "\Opens_\Lambda"]
&
\FFrm_\Lambda^{\op}.
\end{tikzcd}
\]

\subsection{Locales presented by fuzzy presheaves}
\label{subsec:locales-presented-by-fuzzy-presheaves}

For a fuzzy preorder \(X\), the fuzzy presheaf frame \(P_\Lambda(X)\) determines a
fuzzy locale.

\begin{definition}[Presheaf fuzzy locale]
\label{def:presheaf-fuzzy-locale}
For a fuzzy preorder \(X\), the \textbf{presheaf fuzzy locale} associated to \(X\) is
the fuzzy locale
\[
    \mathcal P_\Lambda X
\]
whose fuzzy frame of opens is
\[
    \Opens_\Lambda(\mathcal P_\Lambda X)
    =
    P_\Lambda(X).
\]
Equivalently, \(\mathcal P_\Lambda X\) is the fuzzy locale corresponding to the
identity nucleus on \(P_\Lambda(X)\).
\end{definition}

\begin{proposition}[Functoriality of presheaf fuzzy locales]
\label{prop:functoriality-presheaf-fuzzy-locales}
A \(\Lambda\)-functor
\[
    f:X\to Y
\]
induces a map of fuzzy locales
\[
    \mathcal P_\Lambda X\to\mathcal P_\Lambda Y.
\]
\end{proposition}

\begin{proof}
Precomposition gives
\[
    f^\ast:P_\Lambda(Y)\to P_\Lambda(X),
    \qquad
    f^\ast(\psi)=\psi f.
\]
Explicitly, if \(\psi:Y\to L\) is a fuzzy predicate on \(Y\), then
\[
\begin{aligned}
X(x',x)\tensor(\psi f)(x)
&=
X(x',x)\tensor\psi(fx)
\\
&\leq
Y(fx',fx)\tensor\psi(fx)
\\
&\leq
\psi(fx')
\\
&=
(\psi f)(x').
\end{aligned}
\]
Thus \(f^\ast(\psi)\in P_\Lambda(X)\). Since precomposition preserves pointwise
joins, finite meets, and scalars, \(f^\ast\) is a fuzzy-frame homomorphism. Therefore
it defines a fuzzy-locale map
\[
    \mathcal P_\Lambda X\to\mathcal P_\Lambda Y.
\]
\end{proof}

The construction is covariant on fuzzy preorders:
\[
\begin{tikzcd}[column sep=large, row sep=large]
X
    \arrow[r, "f"]
    \arrow[d, "\mathcal P_\Lambda"']
&
Y
    \arrow[d, "\mathcal P_\Lambda"]
\\
\mathcal P_\Lambda X
    \arrow[r]
&
\mathcal P_\Lambda Y
\end{tikzcd}
\]
while the inverse-image map on opens is contravariant:
\[
\begin{tikzcd}[column sep=large, row sep=large]
\Opens_\Lambda(\mathcal P_\Lambda X)
&
\Opens_\Lambda(\mathcal P_\Lambda Y)
    \arrow[l, "f^\ast"']
\\
P_\Lambda(X)
    \arrow[u, equal]
&
P_\Lambda(Y)
    \arrow[l, "f^\ast"]
    \arrow[u, equal]
\end{tikzcd}
\]

\subsection{Localized presheaves as fuzzy sublocales}
\label{subsec:localized-presheaves-fuzzy-sublocales}

A fuzzy nucleus has two equivalent readings. Frame-theoretically, it is a quotient
of the fuzzy presheaf frame:
\[
    P_\Lambda(X)\twoheadrightarrow P_\Lambda(X)_j.
\]
Localically, it is a sublocale of the presheaf fuzzy locale. This is the ordinary
duality between nuclei, quotient frames, and sublocales, with the scalar compatibility
carried along \cite{Johnstone1982}.

\begin{definition}[Fuzzy sublocale presented by a nucleus]
\label{def:fuzzy-sublocale-presented-by-nucleus}
Let \(j\) be a fuzzy nucleus on \(P_\Lambda(X)\). The \textbf{fuzzy sublocale
presented by \(j\)} is the fuzzy locale
\[
    \mathcal X_j
\]
whose frame of opens is
\[
    \Opens_\Lambda(\mathcal X_j)
    =
    P_\Lambda(X)_j.
\]
\end{definition}

The quotient and sublocale readings are displayed by
\[
\begin{tikzcd}[column sep=large, row sep=large]
P_\Lambda(X)
    \arrow[r, two heads, "q_j"]
&
P_\Lambda(X)_j
\\
\mathcal P_\Lambda X
&
\mathcal X_j
    \arrow[l, hook]
\end{tikzcd}
\]
where the top row is in fuzzy frames and the bottom row is in fuzzy locales. The hook
on the locale side is justified by the fact that \(q_j\) is a quotient frame
homomorphism.

\subsection{Examples}
\label{subsec:fuzzy-locale-examples}

\begin{example}[Ordinary presheaf locales]
\label{ex:ordinary-presheaf-locales}
Let
\[
    \Lambda=2.
\]
For an ordinary preorder \(X\),
\[
    \Opens_2(\mathcal P_2X)
    =
    P_2(X)
    =
    \Down(X).
\]
Thus \(\mathcal P_2X\) is the ordinary lower presheaf locale associated to \(X\).
Nuclei on \(\Down(X)\) correspond to ordinary sublocales of this locale
\cite{Johnstone1982}.

If \(X\) is discrete, then \(\Down(X)\cong \mathcal P(X)\), so the corresponding
locale is the ordinary powerset locale on the set \(X\). If \(X\) has nontrivial
order, then the opens are downward-closed subsets, and the locale remembers the
Alexandrov topology determined by the preorder.
\end{example}

\begin{example}[G\"odel-valued presheaf locales]
\label{ex:godel-valued-presheaf-locales}
For G\"odel truth values,
\[
    \Opens_\Lambda(\mathcal P_\Lambda X)
    =
    P_\Lambda(X)
\]
is the frame of G\"odel-valued fuzzy predicates on \(X\). A fuzzy open is therefore a
monotone fuzzy lower set, where monotonicity is measured by the G\"odel tensor:
\[
    X(x',x)\meet\phi(x)\leq \phi(x').
\]
A fuzzy nucleus on this frame presents a G\"odel-valued fuzzy sublocale. Its opens are
exactly the closed G\"odel predicates. Localically, the quotient
\[
    P_\Lambda(X)\twoheadrightarrow P_\Lambda(X)_j
\]
is read as the inclusion of the fuzzy sublocale
\[
    \mathcal X_j\hookrightarrow \mathcal P_\Lambda X.
\]
\end{example}

\begin{example}[{\L}ukasiewicz-valued presheaf locales]
\label{ex:lukasiewicz-valued-presheaf-locales}
For {\L}ukasiewicz truth values,
\[
    \Opens_\Lambda(\mathcal P_\Lambda X)
    =
    P_\Lambda(X)
\]
is the frame of {\L}ukasiewicz-valued fuzzy predicates on \(X\). The enriched
lower-set condition is
\[
    \max(0,X(x',x)+\phi(x)-1)\leq \phi(x').
\]
A fuzzy nucleus presents a localized fuzzy locale whose opens are the corresponding
closed predicates. The scalar action lowers degrees by {\L}ukasiewicz multiplication
and then recloses them:
\[
    r\tensor_j\phi
    =
    j\bigl(\max(0,r+\phi-1)\bigr),
\]
where the formula is interpreted pointwise.
\end{example}

\begin{example}[Product-valued presheaf locales]
\label{ex:product-valued-presheaf-locales}
For product truth values,
\[
    \Opens_\Lambda(\mathcal P_\Lambda X)
    =
    P_\Lambda(X)
\]
is the frame of product-valued fuzzy predicates satisfying
\[
    X(x',x)\phi(x)\leq \phi(x').
\]
A fuzzy nucleus gives a product-valued fuzzy sublocale whose opens are closed under
the chosen closure operation. The scalar action multiplies fuzzy predicates by
constants and then closes:
\[
    r\tensor_j\phi=j(r\phi).
\]
Thus scalar rescaling behaves like proportional attenuation of fuzzy membership.
\end{example}

\begin{example}[Lawvere-valued presheaf locales]
\label{ex:lawvere-valued-presheaf-locales}
For a Lawvere metric space \(X\), the presheaf fuzzy locale \(\mathcal P_\Lambda X\)
has frame of opens
\[
    \Opens_\Lambda(\mathcal P_\Lambda X)=P_\Lambda(X),
\]
whose elements are maps
\[
    \phi:X\to[0,\infty]
\]
satisfying, in the usual order,
\[
    \phi(x')\leq d(x',x)+\phi(x).
\]
These are nonexpansive cost predicates. A nucleus imposes a local closure condition
on such cost predicates. Since the order is reversed, closed joins correspond to
closures of pointwise infima of costs, while finite meets correspond to pointwise
suprema. This gives a metric-localic interpretation of Lawvere presheaves.
\end{example}

\section{The fuzzy Sierpiński classifier}
\label{sec:fuzzy-sierpinski-classifier}

In ordinary locale theory, the Sierpiński locale classifies opens. On the frame
side, this says that the Sierpiński frame is freely generated by one open. In the
present enriched setting the same principle holds, but the free object is not, in
general, the raw presheaf category on the two-object subterminal shape. For
non-idempotent fuzzy tensors, such as the {\L}ukasiewicz and product tensors, that
raw two-point construction does not classify arbitrary fuzzy opens by
finite-limit-preserving fuzzy-frame morphisms.

The correct construction uses a small enriched syntactic sketch of finite fuzzy-open
formulae and then takes the associated left-exact enriched-sketch localization of its
fuzzy presheaf category. The resulting object is characterized by its classifying
universal property.

\subsection{Categorical fuzzy frames and fuzzy opens}
\label{subsec:categorical-fuzzy-frames-and-fuzzy-opens}

We recall the categorical form of fuzzy frames used in this chapter. A categorical
fuzzy frame is a left-exact reflective localization
\[
    a\dashv i:
    P_{\mathbb V}(X)
    \rightleftarrows
    \mathcal A
\]
of a fuzzy presheaf category, where \(X\) is a small \(\mathbb V\)-category,
\(i\) is fully faithful, and \(a\) preserves the chosen finite
\(\mathcal F\)-weighted limits.

A morphism of categorical fuzzy frames
\[
    H:\mathcal A\to\mathcal B
\]
is a \(\mathbb V\)-functor preserving small weighted colimits and the chosen finite
\(\mathcal F\)-weighted limits. Enriched natural transformations are the 2-cells.
We write
\[
    \mathbf{FFrm}_{\mathbb V}
\]
for the resulting 2-category: objects are categorical fuzzy frames, 1-cells are
small-colimit-preserving and finite-limit-preserving \(\mathbb V\)-functors, and
2-cells are enriched natural transformations. Thus
\[
    \mathbf{FFrm}_{\mathbb V}(\mathcal A,\mathcal B)
\]
denotes the ordinary hom-category of such morphisms and enriched natural
transformations.

A cosmos-level fuzzy open of a categorical fuzzy frame \(\mathcal A\) is a subobject
\[
    S\hookrightarrow 1_{\mathcal A}
\]
of the terminal object in the underlying ordinary category of \(\mathcal A\). We
write
\[
    \operatorname{Sub}_{\mathcal A}(1_{\mathcal A})
\]
for the ordinary category of such subobjects. Thus, if \(\mathcal A\) is enriched,
subobjects are computed in the underlying ordinary category \(\mathcal A_0\).

Equivalently, a representative object \(S\) is subterminal precisely when the
diagonal
\[
    S\longrightarrow S\times S
\]
is an isomorphism. Since morphisms of categorical fuzzy frames preserve terminal
objects and binary products, they send subterminal objects to subterminal objects.

We shall use the following tensor calculation for categorical fuzzy frames.

\begin{lemma}[Tensors in categorical fuzzy frames]
\label{lem:categorical-fuzzy-frames-are-tensored}
Let
\[
    a\dashv i:
    P_{\mathbb V}(X)\rightleftarrows\mathcal A
\]
be a categorical fuzzy-frame presentation. Then \(\mathcal A\) is tensored over
\(\mathbb V\). For \(V\in\mathbb V\) and \(A\in\mathcal A\), the tensor is computed
by
\[
    V\tensor_{\mathcal A}A
    \cong
    a(V\tensor iA).
\]
\end{lemma}

\begin{proof}
For \(B\in\mathcal A\), using the enriched adjunction and the tensor--hom
adjunction in \(P_{\mathbb V}(X)\), we compute
\[
\begin{aligned}
\mathcal A(a(V\tensor iA),B)
&\cong
P_{\mathbb V}(X)(V\tensor iA,iB)
\\
&\cong
[V,P_{\mathbb V}(X)(iA,iB)]
\\
&\cong
[V,\mathcal A(A,B)].
\end{aligned}
\]
Thus \(a(V\tensor iA)\) represents the tensor of \(A\) by \(V\) in
\(\mathcal A\).
\end{proof}

\subsection{The enriched syntactic sketch of one fuzzy open}
\label{subsec:enriched-syntactic-sketch-one-fuzzy-open}

We now construct the enriched syntactic sketch whose models are subterminal objects
together with the formal finite-limit and scalar formulae generated from them.

The construction is made relative to the universe convention fixed earlier. In
particular, the collection of scalar objects used in the sketch, the chosen class
\(\mathcal F\) of finite weights, and the generated syntax below are small in the
ambient universe in which enriched sketches are formed. When the fuzzy cosmos
\(\mathbb V\) is not small in the lowest universe, this construction is understood
in the next universe of the fixed tower, equivalently after choosing a small
skeleton of the scalar objects used by the sketch.

Let
\[
    \mathbb T_{\Omega}
\]
denote the small enriched finite-limit-and-scalar sketch freely generated by one
subterminal object. More explicitly, \(\mathbb T_{\Omega}\) is generated as follows.

First, there are two basic formal objects
\[
    1,
    \qquad
    U.
\]
The object \(1\) is declared terminal, and \(U\) is declared subterminal. Thus the
sketch contains the terminal-object data for \(1\), the canonical map
\[
    U\longrightarrow 1,
\]
and the finite-limit data expressing that the diagonal comparison
\[
    U\longrightarrow U\times U
\]
is an isomorphism.

Second, the sketch is closed under the chosen finite \(\mathcal F\)-weighted limit
operations. Thus, whenever a formal finite-limit expression
\[
    \{\varphi,D\}
\]
is built from earlier formal objects, it is again a formal object of
\(\mathbb T_{\Omega}\), together with its designated limiting cone. In particular,
binary products of formal objects are again formal objects.

Third, for every scalar object \(V\in\mathbb V\) and every formal object
\(A\in\mathbb T_{\Omega}\), there is a formal scalar object
\[
    V\diamond A.
\]
The symbol \(V\diamond A\) is syntactic at this stage. It is not required to
represent the tensor of \(A\) by \(V\) in the raw presheaf category, nor is it
required to satisfy a tensor universal property inside \(\mathbb T_{\Omega}\).

For each formal scalar object \(V\diamond A\), the sketch contains a generating
enriched morphism
\[
    \kappa_{V,A}:
    V\longrightarrow
    \mathbb T_{\Omega}(A,V\diamond A).
\]
Equivalently, \(\kappa_{V,A}\) is a formal tensor coevaluation map. Under Yoneda,
it determines a canonical comparison map in the presheaf category
\[
    \alpha_{V,A}:
    V\tensor yA
    \longrightarrow
    y(V\diamond A),
\]
where
\[
    y:\mathbb T_{\Omega}\to
    P_{\mathbb V}(\mathbb T_{\Omega})
\]
is the enriched Yoneda embedding.

The hom-objects and equations of \(\mathbb T_{\Omega}\) are generated freely by the
terminal-object data, the subterminality data for \(U\), the specified finite-limit
cones, the scalar-forming operations \(V\diamond A\), and the generating maps
\(\kappa_{V,A}\), subject only to the coherence equations for the specified
finite-limit cones and for functoriality of the displayed generating data. No tensor
universal property for \(V\diamond A\) is imposed inside
\(\mathbb T_{\Omega}\). The tensor universal property is imposed only on models, by
requiring the transposes
\[
    V\tensor_{\mathcal A}F(A)\longrightarrow F(V\diamond A)
\]
to be isomorphisms.

Equivalently, \(\mathbb T_{\Omega}\) is the free small enriched
finite-limit-and-scalar sketch generated by one subterminal object \(U\), where the
scalar operations are formal operations whose intended interpretation is by tensors
in the target category of models.

\begin{remark}[Formal scalar objects]
\label{rem:formal-scalar-objects}
One should not impose
\[
    y(V\diamond A)\cong V\tensor yA
\]
inside the raw presheaf category. In general, a scalar multiple of a representable
presheaf need not be representable.

In the posetal specialization, if such an identity held representably, then for a
scalar \(v<\top_L\) it would force
\[
    \mathbb T_{\Omega}(Z,V\diamond A)
    =
    v\tensor \mathbb T_{\Omega}(Z,A).
\]
Taking \(Z=V\diamond A\) would give
\[
    \top_L
    \leq
    \mathbb T_{\Omega}(V\diamond A,V\diamond A)
    =
    v\tensor \mathbb T_{\Omega}(V\diamond A,A)
    \leq
    v,
\]
which is impossible. Thus \(V\diamond A\) is a formal scalar formula, and the
comparison
\[
    V\tensor yA\to y(V\diamond A)
\]
is imposed only after passing to the sketch localization, and only in the correct
classifying variance.
\end{remark}

\subsection{Models of the Sierpiński sketch}
\label{subsec:models-sierpinski-sketch}

Let \(\mathcal A\) be a categorical fuzzy frame. By
Lemma~\ref{lem:categorical-fuzzy-frames-are-tensored}, \(\mathcal A\) is tensored
over \(\mathbb V\). If
\[
    a\dashv i:
    P_{\mathbb V}(X)\rightleftarrows\mathcal A
\]
is a presentation of \(\mathcal A\), then the tensor of \(A\in\mathcal A\) by
\(V\in\mathbb V\) is computed by
\[
    V\tensor_{\mathcal A}A
    \cong
    a(V\tensor iA).
\]

A \textbf{model of the Sierpiński sketch} in \(\mathcal A\) is a
\(\mathbb V\)-functor
\[
    F:\mathbb T_{\Omega}\to\mathcal A
\]
which satisfies the following conditions:
\begin{enumerate}
    \item \(F\) sends the specified finite \(\mathcal F\)-limit cones to finite
    limits in \(\mathcal A\);

    \item \(F(U)\) is subterminal;

    \item for every scalar object \(V\in\mathbb V\) and every formal object
    \(A\in\mathbb T_{\Omega}\), the image of the generating morphism
    \[
        \kappa_{V,A}:
        V\to \mathbb T_{\Omega}(A,V\diamond A)
    \]
    has transpose
    \[
        \overline{\kappa}^{F}_{V,A}:
        V\tensor_{\mathcal A}F(A)\to F(V\diamond A),
    \]
    and this transpose is an isomorphism.
\end{enumerate}
Thus scalar formulae are interpreted as actual tensors in \(\mathcal A\), and the
formal coevaluation maps are interpreted as the corresponding tensor
coevaluations.

\begin{remark}[Scalar formulae need not be subterminal]
\label{rem:scalar-formulae-not-necessarily-subterminal}
Although the generator \(U\) is required to be subterminal, the formal scalar objects
\(V\diamond A\) are not declared subterminal. In a general categorical fuzzy frame,
tensors of subterminal objects need not be subterminal. Thus the cosmos-level
classifier remembers one distinguished fuzzy open
\[
    U\hookrightarrow 1,
\]
while the scalar and finite-limit formulae are formal objects generated from it. In
the posetal reflection every object of a frame lies below the top element, so these
formula objects become ordinary fuzzy opens.
\end{remark}

We write
\[
    \operatorname{Mod}_{\Omega}(\mathcal A)
\]
for the ordinary category of such models and enriched natural transformations.

\begin{proposition}[Models are subterminal objects]
\label{prop:models-sierpinski-sketch-subterminal}
For every categorical fuzzy frame \(\mathcal A\), evaluation at \(U\) induces an
equivalence of ordinary categories
\[
    \operatorname{Mod}_{\Omega}(\mathcal A)
    \simeq
    \operatorname{Sub}_{\mathcal A}(1_{\mathcal A}).
\]
\end{proposition}

\begin{proof}
Let
\[
    F:\mathbb T_{\Omega}\to\mathcal A
\]
be a model. Since the sketch declares \(U\) subterminal, \(F(U)\) is subterminal in
\(\mathcal A\). The formal map \(U\to 1\) is sent to a map
\[
    F(U)\to F(1)\cong 1_{\mathcal A}.
\]
Thus evaluation at \(U\) gives a functor
\[
    \operatorname{Mod}_{\Omega}(\mathcal A)
    \longrightarrow
    \operatorname{Sub}_{\mathcal A}(1_{\mathcal A}).
\]

Conversely, let
\[
    S\hookrightarrow 1_{\mathcal A}
\]
be a subterminal object of \(\mathcal A\). By the initiality of the free enriched
finite-limit-and-scalar sketch generated by one subterminal object, giving a model
\[
    F:\mathbb T_{\Omega}\to\mathcal A
\]
with \(F(U)=S\) is equivalent to interpreting the generator \(U\) as the chosen
subterminal object \(S\). The values on finite-limit formulae, scalar formulae, and
the generating maps \(\kappa_{V,A}\) are then forced up to the unique coherent
isomorphisms determined by the relevant universal properties.

Explicitly, the corresponding model \(F_S\) is defined recursively by
\[
    F_S(U)=S,
    \qquad
    F_S(1)=1_{\mathcal A},
\]
by interpreting the specified finite-limit formulae as the corresponding finite
limits in \(\mathcal A\), and by interpreting scalar formulae as tensors:
\[
    F_S(V\diamond A)
    =
    V\tensor_{\mathcal A}F_S(A).
\]
The generating morphism
\[
    \kappa_{V,A}:
    V\to \mathbb T_{\Omega}(A,V\diamond A)
\]
is interpreted as the tensor coevaluation
\[
    V\to
    \mathcal A(F_S(A),V\tensor_{\mathcal A}F_S(A)).
\]

Similarly, a morphism of models
\[
    F_S\to F_T
\]
is determined by its component at the generator \(U\). The components at
finite-limit formulae are forced by the universal property of the corresponding
limits, and the components at scalar formulae are forced by tensoring with the
scalar object. The component at \(U\) is exactly a morphism
\[
    S\to T
\]
over \(1_{\mathcal A}\). Therefore evaluation at \(U\) is fully faithful and
essentially surjective, hence an equivalence:
\[
    \operatorname{Mod}_{\Omega}(\mathcal A)
    \simeq
    \operatorname{Sub}_{\mathcal A}(1_{\mathcal A}).
\]
\end{proof}

\subsection{The Sierpiński localization}
\label{subsec:sierpinski-localization}

Form the fuzzy presheaf category
\[
    P_{\mathbb V}(\mathbb T_{\Omega})
    =
    [\mathbb T_{\Omega}^{\op},\underline{\mathbb V}].
\]
Let
\[
    y:\mathbb T_{\Omega}\to P_{\mathbb V}(\mathbb T_{\Omega})
\]
be the Yoneda embedding.

We use the standard enriched finite-limit sketch localization theorem in the
following form.

\begin{proposition}[Sierpiński sketch localization]
\label{prop:sierpinski-sketch-localization}
There is a left-exact reflective localization
\[
    a_{\Omega}\dashv i_{\Omega}:
    P_{\mathbb V}(\mathbb T_{\Omega})
    \rightleftarrows
    P_{\mathbb V}(\mathbb T_{\Omega})_{J_{\Omega}}
\]
with the following universal property. For every categorical fuzzy frame
\(\mathcal A\), restriction along the composite
\[
    \mathbb T_{\Omega}
    \xrightarrow{y}
    P_{\mathbb V}(\mathbb T_{\Omega})
    \xrightarrow{a_{\Omega}}
    P_{\mathbb V}(\mathbb T_{\Omega})_{J_{\Omega}}
\]
induces a natural equivalence of ordinary categories
\[
    \mathbf{FFrm}_{\mathbb V}
    \bigl(P_{\mathbb V}(\mathbb T_{\Omega})_{J_{\Omega}},\mathcal A\bigr)
    \simeq
    \operatorname{Mod}_{\Omega}(\mathcal A).
\]
\end{proposition}

\begin{proof}
This is the enriched finite-limit sketch localization theorem applied to the small
sketch \(\mathbb T_{\Omega}\). The localization is generated by the sketch data: the
specified finite-limit cones, the subterminality equation for \(U\), and the scalar
comparison maps
\[
    \alpha_{V,A}:
    V\tensor yA
    \longrightarrow
    y(V\diamond A).
\]
The universal property says exactly that small-colimit-preserving left-exact
\(\mathbb V\)-functors out of the localized presheaf category are the same as
models of the sketch in the target categorical fuzzy frame.
\end{proof}

The localization is generated by the enriched sketch data in the following sense:
after reflection, the designated finite-limit cones become limiting cones, the
subterminality equation for \(U\) holds, and the scalar comparison maps
\[
    \alpha_{V,A}:
    V\tensor yA
    \longrightarrow
    y(V\diamond A)
\]
become isomorphisms in the correct classifying variance. Thus the localization is
characterized by its universal property for small-colimit-preserving left-exact
functors out of the localized presheaf category.

It is important not to describe the \(J_{\Omega}\)-local presheaves themselves as
presheaves satisfying the tensor equations. Orthogonality of a presheaf \(F\) to
\[
    V\tensor yA\to y(V\diamond A)
\]
would instead identify
\[
    F(V\diamond A)
    \cong
    [V,F(A)].
\]
The tensor equations are imposed in the correct variance by the universal property
for cocontinuous left-exact functors out of the localization.

\begin{definition}[Fuzzy Sierpiński frame]
\label{def:fuzzy-sierpinski-frame}
The \textbf{fuzzy Sierpiński frame} over the fuzzy cosmos \(\mathbb V\) is
\[
    \Omega_{\mathbb V}
    \eqdef
    P_{\mathbb V}(\mathbb T_{\Omega})_{J_{\Omega}}.
\]
Its \textbf{generic fuzzy open} is
\[
    \sigma_{\mathbb V}
    \eqdef
    a_{\Omega}(yU).
\]
\end{definition}

By construction,
\[
    \Omega_{\mathbb V}
\]
is a left-exact reflective localization of a fuzzy presheaf category. Hence it is a
categorical fuzzy frame in the sense of this chapter.

\begin{proposition}[The generic fuzzy open is subterminal]
\label{prop:generic-fuzzy-open-subterminal}
The object
\[
    \sigma_{\mathbb V}\in\Omega_{\mathbb V}
\]
is subterminal.
\end{proposition}

\begin{proof}
The localization \(a_{\Omega}\dashv i_{\Omega}\) is generated so that the image of
the designated subterminality data for \(U\) becomes valid in
\(\Omega_{\mathbb V}\). Since the localization also forces the specified product cone
for \(U\times U\) to become a product cone, the object
\[
    a_{\Omega}y(U\times U)
\]
is canonically isomorphic to
\[
    \sigma_{\mathbb V}\times\sigma_{\mathbb V}.
\]
Moreover, the image of the diagonal comparison
\[
    U\to U\times U
\]
becomes an isomorphism after applying \(a_{\Omega}y\). Hence
\[
    \sigma_{\mathbb V}
    \cong
    \sigma_{\mathbb V}\times\sigma_{\mathbb V}.
\]
Since \(\Omega_{\mathbb V}\) has a terminal object, this says precisely that
\(\sigma_{\mathbb V}\) is subterminal.
\end{proof}

\subsection{The classifying theorem}
\label{subsec:sierpinski-classifying-theorem}

\begin{theorem}[Fuzzy Sierpiński classification, frame form]
\label{thm:fuzzy-sierpinski-classification-frame}
For every categorical fuzzy frame \(\mathcal A\), evaluation at the generic fuzzy
open induces a natural equivalence of ordinary categories
\[
    \mathbf{FFrm}_{\mathbb V}(\Omega_{\mathbb V},\mathcal A)
    \simeq
    \operatorname{Sub}_{\mathcal A}(1_{\mathcal A}).
\]
Under this equivalence, a fuzzy-frame morphism
\[
    H:\Omega_{\mathbb V}\to\mathcal A
\]
corresponds to the subterminal object
\[
    H(\sigma_{\mathbb V})\hookrightarrow 1_{\mathcal A}.
\]
\end{theorem}

\begin{proof}
By Proposition~\ref{prop:sierpinski-sketch-localization}, restriction along
\[
    a_{\Omega}y:\mathbb T_{\Omega}\to\Omega_{\mathbb V}
\]
induces a natural equivalence
\[
    \mathbf{FFrm}_{\mathbb V}(\Omega_{\mathbb V},\mathcal A)
    \simeq
    \operatorname{Mod}_{\Omega}(\mathcal A).
\]
Under this equivalence, a morphism
\[
    H:\Omega_{\mathbb V}\to\mathcal A
\]
corresponds to the model
\[
    Ha_{\Omega}y:\mathbb T_{\Omega}\to\mathcal A.
\]
Its value at the generator \(U\) is
\[
    H(a_{\Omega}yU)
    =
    H(\sigma_{\mathbb V}).
\]

By Proposition~\ref{prop:models-sierpinski-sketch-subterminal}, evaluation at \(U\)
gives an equivalence
\[
    \operatorname{Mod}_{\Omega}(\mathcal A)
    \simeq
    \operatorname{Sub}_{\mathcal A}(1_{\mathcal A}).
\]
Combining these equivalences gives
\[
    \mathbf{FFrm}_{\mathbb V}(\Omega_{\mathbb V},\mathcal A)
    \simeq
    \operatorname{Sub}_{\mathcal A}(1_{\mathcal A}).
\]
Under the composite equivalence, \(H\) is sent to
\[
    H(\sigma_{\mathbb V})\hookrightarrow 1_{\mathcal A}.
\]
Since \(H\) preserves terminal objects and binary products,
\(H(\sigma_{\mathbb V})\) is subterminal.
\end{proof}

\begin{definition}[Characteristic morphism of a fuzzy open]
\label{def:characteristic-morphism-fuzzy-open}
Let \(\mathcal A\) be a categorical fuzzy frame, and let
\[
    S\hookrightarrow 1_{\mathcal A}
\]
be a fuzzy open. The \textbf{characteristic morphism} of \(S\) is the fuzzy-frame
morphism
\[
    \chi_S^\ast:
    \Omega_{\mathbb V}\to\mathcal A
\]
corresponding to \(S\) under the equivalence of
Theorem~\ref{thm:fuzzy-sierpinski-classification-frame}. Equivalently,
\[
    \chi_S^\ast(\sigma_{\mathbb V})
    \cong
    S
\]
as subobjects of \(1_{\mathcal A}\).
\end{definition}

\subsection{The fuzzy Sierpiński locale}
\label{subsec:fuzzy-sierpinski-locale}

Define
\[
    \mathbf{FLoc}_{\mathbb V}
    \eqdef
    \mathbf{FFrm}_{\mathbb V}^{\op}.
\]

\begin{definition}[Fuzzy Sierpiński locale]
\label{def:fuzzy-sierpinski-locale}
The \textbf{fuzzy Sierpiński locale} over \(\mathbb V\) is the fuzzy locale
\[
    \mathbb S_{\mathbb V}
\]
whose fuzzy frame of opens is
\[
    \Opens_{\mathbb V}(\mathbb S_{\mathbb V})
    =
    \Omega_{\mathbb V}.
\]
\end{definition}

\begin{theorem}[Fuzzy Sierpiński classification, localic form]
\label{thm:fuzzy-sierpinski-classification-localic}
For every categorical fuzzy locale \(\mathcal X\), there is a natural equivalence
of ordinary categories
\[
    \mathbf{FLoc}_{\mathbb V}(\mathcal X,\mathbb S_{\mathbb V})
    \simeq
    \operatorname{Sub}_{\Opens_{\mathbb V}(\mathcal X)}
    \bigl(1_{\Opens_{\mathbb V}(\mathcal X)}\bigr).
\]
Thus maps into the fuzzy Sierpiński locale classify cosmos-level fuzzy opens.
\end{theorem}

\begin{proof}
By definition of fuzzy locales,
\[
\begin{aligned}
    \mathbf{FLoc}_{\mathbb V}(\mathcal X,\mathbb S_{\mathbb V})
    &=
    \mathbf{FFrm}_{\mathbb V}
    \bigl(
        \Opens_{\mathbb V}(\mathbb S_{\mathbb V}),
        \Opens_{\mathbb V}(\mathcal X)
    \bigr)
    \\
    &=
    \mathbf{FFrm}_{\mathbb V}
    \bigl(
        \Omega_{\mathbb V},
        \Opens_{\mathbb V}(\mathcal X)
    \bigr).
\end{aligned}
\]
Applying Theorem~\ref{thm:fuzzy-sierpinski-classification-frame} to
\[
    \mathcal A=\Opens_{\mathbb V}(\mathcal X)
\]
gives
\[
    \mathbf{FFrm}_{\mathbb V}
    \bigl(
        \Omega_{\mathbb V},
        \Opens_{\mathbb V}(\mathcal X)
    \bigr)
    \simeq
    \operatorname{Sub}_{\Opens_{\mathbb V}(\mathcal X)}
    \bigl(1_{\Opens_{\mathbb V}(\mathcal X)}\bigr).
\]
Therefore
\[
    \mathbf{FLoc}_{\mathbb V}(\mathcal X,\mathbb S_{\mathbb V})
    \simeq
    \operatorname{Sub}_{\Opens_{\mathbb V}(\mathcal X)}
    \bigl(1_{\Opens_{\mathbb V}(\mathcal X)}\bigr).
\]
\end{proof}

The characteristic map of a fuzzy open
\[
    S\hookrightarrow 1_{\Opens_{\mathbb V}(\mathcal X)}
\]
is the fuzzy-locale map
\[
    \chi_S:\mathcal X\to\mathbb S_{\mathbb V}
\]
whose inverse-image morphism
\[
    \chi_S^\ast:
    \Omega_{\mathbb V}\to\Opens_{\mathbb V}(\mathcal X)
\]
satisfies
\[
    \chi_S^\ast(\sigma_{\mathbb V})\cong S.
\]

\subsection{Relation with the raw two-point presheaf}
\label{subsec:relation-raw-two-point-presheaf}

The construction above explains why the fuzzy Sierpiński frame should not be
identified, in general, with the raw presheaf frame on the two-point subterminal
shape.

Let
\[
    C=(U\leq 1)
\]
be the ordinary two-element preorder, regarded as a \(\Lambda\)-preorder by taking
\[
    C(U,1)=\top_L,
    \qquad
    C(1,U)=\bot_L,
\]
and identity homs equal to \(\top_L\). In the crisp case, the presheaf frame
\[
    P_2(C)
\]
is the ordinary Sierpiński frame. For a general fuzzy truth-value base
\[
    \Lambda=(L,\leq,\meet,\join,\tensor,\multimap,\bot_L,\top_L),
\]
however, the raw presheaf frame
\[
    P_\Lambda(C)
\]
does not have the required free finite-meet behavior among scalar-generated fuzzy
open formulae.

In the posetal reflection, an element of
\[
    P_\Lambda(C)
\]
is a pair
\[
    (u,t)\in L\times L
\]
with
\[
    t\leq u,
\]
where \(u\) is the value at \(U\) and \(t\) is the value at \(1\). The representable
corresponding to \(U\) is
\[
    \sigma=(\top_L,\bot_L),
\]
and the top predicate is
\[
    \top=(\top_L,\top_L).
\]
Every element decomposes as
\[
    (u,t)
    =
    (u\tensor\sigma)\join(t\tensor\top).
\]
Indeed,
\[
    u\tensor\sigma=(u,\bot_L),
    \qquad
    t\tensor\top=(t,t),
\]
and \(t\leq u\).

Therefore any scalar-compatible frame homomorphism
\[
    H:P_\Lambda(C)\to A
\]
into a \(\Lambda\)-scalar frame \(A\), with
\[
    H(\sigma)=a,
\]
is forced to satisfy
\[
    H(u,t)
    =
    (u\tensor_A a)\join (t\tensor_A\top_A).
\]
In the special case \(A=\Lambda\), this becomes
\[
    H_a(u,t)=(u\tensor a)\join t.
\]
For non-idempotent tensors, this forced map need not preserve binary meets.

For example, let
\[
    \Lambda=[0,1]
\]
with the {\L}ukasiewicz tensor
\[
    r\tensor s=\max(0,r+s-1).
\]
Take
\[
    a=\frac12,
    \qquad
    p=(0.8,0.1),
    \qquad
    q=(0.4,0.3).
\]
Then
\[
    p\meet q=(0.4,0.1).
\]
The forced map
\[
    H_a(u,t)=(u\tensor a)\join t
\]
satisfies
\[
    H_a(p)=0.3,
    \qquad
    H_a(q)=0.3.
\]
Thus
\[
    H_a(p)\meet H_a(q)=0.3.
\]
But
\[
    H_a(p\meet q)
    =
    H_a(0.4,0.1)
    =
    0.1.
\]
Hence
\[
    H_a(p\meet q)\neq H_a(p)\meet H_a(q).
\]
So \(H_a\) is not a frame homomorphism.

Similarly, for the product tensor
\[
    r\tensor s=rs,
\]
take the same values
\[
    a=\frac12,
    \qquad
    p=(0.8,0.1),
    \qquad
    q=(0.4,0.3).
\]
Then
\[
    H_a(u,t)=\max(au,t).
\]
Thus
\[
    H_a(p)=\max(0.4,0.1)=0.4,
    \qquad
    H_a(q)=\max(0.2,0.3)=0.3,
\]
so
\[
    H_a(p)\meet H_a(q)=0.3.
\]
But
\[
    H_a(p\meet q)
    =
    H_a(0.4,0.1)
    =
    \max(0.2,0.1)
    =
    0.2.
\]
Hence
\[
    H_a(p\meet q)\neq H_a(p)\meet H_a(q).
\]
Therefore the raw two-point presheaf frame does not classify arbitrary fuzzy opens
for these non-idempotent bases.

The enriched syntactic site \(\mathbb T_{\Omega}\) corrects this defect by adding
formal scalar formulae and their finite-limit combinations before freely adjoining
presheaves and then imposing the sketch localization. Thus the fuzzy Sierpiński
frame is
\[
    \Omega_{\mathbb V}
    =
    P_{\mathbb V}(\mathbb T_{\Omega})_{J_{\Omega}},
\]
not the raw presheaf object on the two-point shape.

\begin{corollary}[Crisp specialization]
\label{cor:crisp-sierpinski-specialization}
When
\[
    \Lambda=2,
\]
the fuzzy Sierpiński frame \(\Omega_2\) is isomorphic to the ordinary three-element
Sierpiński frame. Equivalently, it is isomorphic to the downset frame of the
two-element preorder
\[
    U\leq 1.
\]
\end{corollary}

\begin{proof}
For \(\Lambda=2\), the scalar action is forced:
\[
    1\tensor a=a,
    \qquad
    0\tensor a=0.
\]
Thus scalar compatibility adds no structure beyond the ordinary finite-meet and
join structure. The free fuzzy frame on one fuzzy open is therefore the ordinary
free frame on one open, namely the Sierpiński frame. This is the downset frame of
the preorder \(U\leq 1\), whose generator is the open \(\{U\}\), and which has three
elements:
\[
    \varnothing,
    \qquad
    \{U\},
    \qquad
    \{U,1\}.
\]
\end{proof}

\subsection{Posetal reflection}
\label{subsec:sierpinski-posetal-reflection}

When the fuzzy cosmos is the thin cosmos associated to
\[
    \Lambda=(L,\leq,\meet,\join,\tensor,\multimap,\bot_L,\top_L),
\]
we write
\[
    \Omega_\Lambda,
    \qquad
    \sigma_\Lambda,
    \qquad
    \mathbb S_\Lambda
\]
for the corresponding posetal reflections of
\[
    \Omega_{\mathbb V},
    \qquad
    \sigma_{\mathbb V},
    \qquad
    \mathbb S_{\mathbb V}.
\]

The enriched syntactic sketch has a transparent posetal shadow. In the posetal
comparison where the only finite limits used explicitly are terminal objects and
binary meets, its formal objects are finite fuzzy-open formulae generated by
\[
    \top,
    \qquad
    \sigma,
    \qquad
    r\tensor\varphi \quad (r\in L),
    \qquad
    \varphi\meet\psi.
\]
Here \(\sigma\) denotes the formal generator corresponding to \(U\). Arbitrary joins
are not part of the finite syntactic sketch; they are freely adjoined by passing to
the presheaf frame and then imposing the sketch localization.

The fuzzy hom-value
\[
    \mathbb T_{\Omega}(\varphi,\psi)
\]
is the degree to which the sequent
\[
    \varphi\Rightarrow\psi
\]
is derivable from the finite-limit, subterminality, and scalar rules of the sketch.

The posetal fuzzy Sierpiński frame is
\[
    \Omega_\Lambda
    =
    P_\Lambda(\mathbb T_{\Omega})_{J_{\Omega}},
\]
with generic fuzzy open
\[
    \sigma_\Lambda
    =
    a_{\Omega}(y\sigma).
\]
For every fuzzy frame \(A\) over \(\Lambda\), the classification theorem specializes
to the set-level universal property
\[
    \FFrm_\Lambda(\Omega_\Lambda,A)
    \cong
    |A|.
\]
Here the right-hand side is the underlying set of the frame \(A\), because in the
posetal reflection every element \(a\in A\) determines a unique subobject
\[
    a\hookrightarrow \top_A,
\]
and every subobject of \(\top_A\) arises this way. The element corresponding to a
fuzzy-frame homomorphism
\[
    H:\Omega_\Lambda\to A
\]
is
\[
    H(\sigma_\Lambda)\in |A|.
\]

Dualizing, for every fuzzy locale \(\mathcal X\),
\[
    \FLoc_\Lambda(\mathcal X,\mathbb S_\Lambda)
    \cong
    |\Opens_\Lambda(\mathcal X)|.
\]
Thus maps into the fuzzy Sierpiński locale classify fuzzy opens exactly as in the
ordinary localic case, but the construction of the classifier uses the enriched
syntactic site rather than the naive two-point presheaf object.

    %\chapter{The Fuzzy Ideal Monad, Fuzzy Domains, and Fuzzy Scott Opens}
\label{ch:fuzzy-ideal-monad-domains-scott}

In the previous chapters we developed fuzzy predicates as enriched presheaves and
then localized their predicate frames to obtain fuzzy frames and fuzzy locales.  We
now apply the same predicate construction to \(X^{\op}\).  Thus the upper fuzzy
predicate object of a small \(\mathbb V\)-category \(X\) is
\[
    P_{\mathbb V}(X^{\op})
    =
    [X,\underline{\mathbb V}],
\]
and, in the posetal specialization determined by
\[
    \Lambda=(L,\leq,\meet,\join,\tensor,\multimap,\bot_L,\top_L),
\]
it is
\[
    P_\Lambda(X^{\op})
    =
    [X,\Lambda].
\]
Its elements are upper fuzzy predicates on \(X\).

The central construction of this chapter is the character adjunction
\[
    \mathcal P_{\mathbb V}^{\uparrow}
    \dashv
    \operatorname{Char}_{\mathbb V}.
\]
Here \(\mathcal P_{\mathbb V}^{\uparrow}X\) is the upper fuzzy Alexandrov locale
whose frame of opens is \([X,\underline{\mathbb V}]\), while
\(\operatorname{Char}_{\mathbb V}\) sends a fuzzy locale to its
\(\mathbb V\)-category of fuzzy-frame morphisms into the truth-value fuzzy frame
\(\underline{\mathbb V}\).  Thus characters do not target the fuzzy Sierpiński
frame.  The fuzzy Sierpiński frame classifies fuzzy opens; the truth-value frame
receives fuzzy characters.

The induced monad on fuzzy preorders is the fuzzy ideal monad:
\[
    T_{\mathbb V}X
    \simeq
    \operatorname{Idl}_{\mathbb V}(X),
\]
where \(\operatorname{Idl}_{\mathbb V}(X)\) is the \(\mathbb V\)-category of
finite-flat weights
\[
    X^{\op}\to\underline{\mathbb V}.
\]
In the posetal specialization,
\[
    T_\Lambda X
    \cong
    \operatorname{Idl}_\Lambda(X),
\]
whose elements are fuzzy ideals of \(X\).

A weak fuzzy domain is a fuzzy preorder equipped with a chosen operation
\[
    \alpha:\operatorname{Idl}_\Lambda(X)\to X
\]
assigning candidate joins to fuzzy ideals.  A fuzzy domain is a strict algebra for
the fuzzy ideal monad.  In the crisp case \(\Lambda=2\), this recovers the ordinary
ideal-completion monad, and strict algebras on separated preorders are precisely
dcpos.

The same monad also contains the algebraic and continuous layers of domain theory.
Free fuzzy domains
\[
    \operatorname{Idl}_\Lambda(B)
\]
are the presented fuzzy algebraic domains.  Fuzzy continuous domains are fuzzy
domains for which the ideal-join map
\[
    \alpha:\operatorname{Idl}_\Lambda(X)\to X
\]
admits an enriched left adjoint
\[
    \beta:X\to\operatorname{Idl}_\Lambda(X).
\]
In the crisp case, \(\beta(x)\) is the ideal of elements way-below \(x\).

Finally, fuzzy Scott opens are upper fuzzy predicates preserving fuzzy ideal joins.
For a weak fuzzy domain
\[
    \alpha:\operatorname{Idl}_\Lambda(X)\to X,
\]
a fuzzy Scott open is an upper fuzzy predicate
\[
    U:X\to L
\]
satisfying
\[
    U(\alpha(I))
    =
    \bigjoin_{x\in X} I(x)\tensor U(x)
\]
for every fuzzy ideal \(I\).

\begin{remark}[Universe and doctrine conventions]
\label{rem:ideal-monad-universe-and-doctrine-convention}
We use the same universe convention as in the preceding chapters.  If
\(\mathcal A\) is not small, then
\[
    [\mathcal A,\underline{\mathbb V}]
    \qquad\text{and}\qquad
    \operatorname{Char}_{\mathbb V}(\mathcal A)
\]
are understood in the appropriate larger universe.  In particular, for small
\(X\), the category
\[
    T_{\mathbb V}X
    =
    \operatorname{Char}_{\mathbb V}(P_{\mathbb V}(X^{\op}))
\]
is taken in the universe where the corresponding category of finite-flat weights
exists.

At the categorical-cosmos level, \(T_{\mathbb V}\) is therefore a
universe-relative monad on the chosen \(\mathbb V\)-CAT of locally small
\(\mathbb V\)-categories for which the finite-flat-weight category exists in the
ambient universe.  In the small posetal specialization it restricts to an honest
monad
\[
    T_\Lambda:\Pre_\Lambda\to\Pre_\Lambda.
\]

Throughout the chapter, ``finite-flat'' means flat with respect to the chosen
finite-limit doctrine used in the definition of categorical fuzzy frames.  In the
posetal scalar-frame specialization, this doctrine consists of the terminal element
and binary meets.  Thus finite-flatness specializes to preservation of top and
binary meets.
\end{remark}

\begin{remark}[Strictification of the ideal monad]
\label{rem:strictification-ideal-monad}
The enriched character construction gives
\[
    T_{\mathbb V}X
    =
    \operatorname{Char}_{\mathbb V}(P_{\mathbb V}(X^{\op}))
    \simeq
    \operatorname{Idl}_{\mathbb V}(X).
\]
After the finite-flat-weight representation is established, we freely replace this
monad by the transported equivalent monad whose value on \(X\) is
\(\operatorname{Idl}_{\mathbb V}(X)\), with unit and multiplication given by the
Yoneda and flattening formulas below.  In the posetal specialization this
identification is an isomorphism of fuzzy preorders.
\end{remark}

\section{Upper fuzzy predicates and upper fuzzy Alexandrov locales}
\label{sec:upper-fuzzy-predicates}

\subsection{The enriched construction}
\label{subsec:upper-predicates-enriched}

\begin{definition}[Upper fuzzy predicates]
\label{def:upper-fuzzy-predicates-enriched}
Let \(X\) be a small \(\mathbb V\)-category.  The \textbf{upper fuzzy predicate
category} of \(X\) is
\[
    P_{\mathbb V}(X^{\op})
    =
    [X,\underline{\mathbb V}].
\]
Its objects are \(\mathbb V\)-functors
\[
    U:X\to\underline{\mathbb V}.
\]
\end{definition}

\begin{definition}[Upper fuzzy Alexandrov locale]
\label{def:upper-fuzzy-alexandrov-locale}
The \textbf{upper fuzzy Alexandrov locale} associated to \(X\) is the fuzzy locale
\[
    \mathcal P_{\mathbb V}^{\uparrow}X
\]
whose fuzzy frame of opens is
\[
    \operatorname{Opens}_{\mathbb V}(\mathcal P_{\mathbb V}^{\uparrow}X)
    =
    P_{\mathbb V}(X^{\op})
    =
    [X,\underline{\mathbb V}].
\]
Since \(P_{\mathbb V}(X^{\op})\) is itself a fuzzy presheaf category, it is a
categorical fuzzy frame via the identity left-exact localization
\[
    1\dashv 1:
    P_{\mathbb V}(X^{\op})
    \rightleftarrows
    P_{\mathbb V}(X^{\op}).
\]
Thus it defines a fuzzy locale by duality.
\end{definition}

\begin{proposition}[Functoriality of upper fuzzy predicates]
\label{prop:upper-fuzzy-predicates-functorial}
Let
\[
    f:X\to Y
\]
be a \(\mathbb V\)-functor.  Precomposition defines a fuzzy-frame morphism
\[
    f^\ast:
    P_{\mathbb V}(Y^{\op})
    =
    [Y,\underline{\mathbb V}]
    \longrightarrow
    [X,\underline{\mathbb V}]
    =
    P_{\mathbb V}(X^{\op}),
    \qquad
    f^\ast(V)=Vf.
\]
Consequently there is a fuzzy-locale map
\[
    \mathcal P_{\mathbb V}^{\uparrow}X
    \longrightarrow
    \mathcal P_{\mathbb V}^{\uparrow}Y,
\]
and hence a covariant functor
\[
    \mathcal P_{\mathbb V}^{\uparrow}:
    \mathbb V\text{-}\mathbf{Cat}
    \longrightarrow
    \mathbf{FLoc}_{\mathbb V}.
\]
\end{proposition}

\begin{proof}
All fuzzy-frame operations in the predicate frame are computed pointwise.
Precomposition preserves them.  Passing to frame opposites gives the displayed
locale map.
\end{proof}

\subsection{The posetal specialization}
\label{subsec:upper-predicates-posetal}

\begin{definition}[Upper fuzzy predicate over \(\Lambda\)]
\label{def:upper-fuzzy-predicate-posetal}
Let \(X\) be a fuzzy preorder over \(\Lambda\).  An \textbf{upper fuzzy predicate}
on \(X\) is a map
\[
    U:X\to L
\]
satisfying
\[
    X(x,x')\tensor U(x)\leq U(x')
\]
for all \(x,x'\in X\).
\end{definition}

\begin{proposition}[The upper Alexandrov fuzzy frame]
\label{prop:upper-alexandrov-fuzzy-frame}
The upper fuzzy predicates on \(X\) form the fuzzy frame
\[
    P_\Lambda(X^{\op})
    =
    [X,\Lambda].
\]
Its hom-values are
\[
    P_\Lambda(X^{\op})(U,V)
    =
    \bigmeet_{x\in X}
    \bigl(U(x)\multimap V(x)\bigr),
\]
and its induced crisp order is pointwise.  Joins, meets, and scalar multiplication
are also computed pointwise:
\[
    \left(\bigjoin_i U_i\right)(x)=\bigjoin_i U_i(x),
    \qquad
    \left(\bigmeet_i U_i\right)(x)=\bigmeet_i U_i(x),
    \qquad
    (r\tensor U)(x)=r\tensor U(x).
\]
\end{proposition}

\begin{proof}
This is the predicate-frame construction from the previous chapters applied to
\(X^{\op}\).
\end{proof}

\begin{example}[The crisp case]
\label{ex:upper-predicates-crisp}
For \(\Lambda=2\), a fuzzy preorder is an ordinary preorder and
\[
    P_2(X^{\op})
    \cong
    \operatorname{Up}(X).
\]
Indeed, the upper-predicate law becomes
\[
    x\leq x'
    \quad\text{and}\quad
    U(x)
    \quad\Longrightarrow\quad
    U(x').
\]
Later, Theorem~\ref{thm:scott-opens-free-fuzzy-domain} identifies this frame with
the Scott frame of the free dcpo \(\operatorname{Idl}(X)\).
\end{example}

\begin{example}[Standard truth-value objects]
\label{ex:standard-upper-fuzzy-predicates}
For the G\"odel tensor \(r\tensor s=r\meet s\), an upper fuzzy predicate satisfies
\[
    X(x,x')\meet U(x)\leq U(x').
\]
For the {\L}ukasiewicz tensor \(r\tensor s=\max(0,r+s-1)\), it satisfies
\[
    \max(0,X(x,x')+U(x)-1)\leq U(x').
\]
For the product tensor \(r\tensor s=rs\), it satisfies
\[
    X(x,x')\,U(x)\leq U(x').
\]
For the Lawvere quantale \([0,\infty]\), with the Lawvere order
\(a\leq_\Lambda b\) iff \(a\geq b\) in the usual order, tensor \(+\), top \(0\),
and bottom \(\infty\), the upper-predicate law becomes
\[
    U(x')\leq d(x,x')+U(x)
\]
in the usual numerical order.
\end{example}

\section{Fuzzy characters}
\label{sec:fuzzy-characters}

\subsection{The enriched construction}
\label{subsec:fuzzy-characters-enriched}

The truth-value object \(\underline{\mathbb V}\) is itself a categorical fuzzy
frame.  Indeed, if \(\mathbf 1\) denotes the terminal small \(\mathbb V\)-category,
then
\[
    P_{\mathbb V}(\mathbf 1)
    =
    [\mathbf 1^{\op},\underline{\mathbb V}]
    \simeq
    \underline{\mathbb V}.
\]
Thus \(\underline{\mathbb V}\) is the presheaf fuzzy frame on \(\mathbf 1\).

\begin{definition}[Fuzzy character]
\label{def:fuzzy-character-enriched}
Let \(\mathcal A\) be a categorical fuzzy frame.  Its \textbf{fuzzy character
category} is the full \(\mathbb V\)-subcategory
\[
    \operatorname{Char}_{\mathbb V}(\mathcal A)
    \subseteq
    [\mathcal A,\underline{\mathbb V}]
\]
spanned by those \(\mathbb V\)-functors
\[
    \chi:\mathcal A\to\underline{\mathbb V}
\]
that are fuzzy-frame morphisms, that is, that preserve colimits, the chosen finite
limits, and scalars.
\end{definition}

\begin{definition}[Characters of a fuzzy locale]
\label{def:fuzzy-characters-locale}
If \(\mathcal X\) is a fuzzy locale, define
\[
    \operatorname{Char}_{\mathbb V}(\mathcal X)
    \coloneqq
    \operatorname{Char}_{\mathbb V}
    (\operatorname{Opens}_{\mathbb V}(\mathcal X)).
\]
\end{definition}

\begin{remark}[Presentation-wise characters]
\label{rem:character-universe-convention}
If a fuzzy frame is presented as a reflective localization
\[
    a\dashv i:
    P_{\mathbb V}(X)\rightleftarrows\mathcal A,
\]
then characters of \(\mathcal A\) may be described as characters of
\(P_{\mathbb V}(X)\) satisfying the descent condition imposed by the localization.
For upper predicate frames, this descent-free case becomes the finite-flat-weight
description of Section~\ref{sec:fuzzy-ideals-flat-weights}.
\end{remark}

\subsection{The posetal specialization}
\label{subsec:fuzzy-characters-posetal}

\begin{definition}[Fuzzy character over \(\Lambda\)]
\label{def:fuzzy-character-posetal}
Let \(A\in\FFrm_\Lambda\).  Its \textbf{fuzzy character preorder}
\[
    \operatorname{Char}_\Lambda(A)
\]
is the full fuzzy subpreorder of \([A,\Lambda]\) spanned by the scalar-frame
homomorphisms
\[
    \chi:A\to\Lambda.
\]
Equivalently, its underlying set is
\[
    |\operatorname{Char}_\Lambda(A)|
    =
    \FFrm_\Lambda(A,\Lambda).
\]
Thus an element of \(\operatorname{Char}_\Lambda(A)\) is a map preserving arbitrary
joins, finite meets, and scalar multiplication:
\[
    \chi(r\tensor_A a)=r\tensor\chi(a).
\]
\end{definition}

\begin{proposition}[The fuzzy preorder of characters]
\label{prop:fuzzy-character-preorder}
The fuzzy preorder on \(\operatorname{Char}_\Lambda(A)\) is inherited from the
pointwise functor preorder:
\[
    \operatorname{Char}_\Lambda(A)(\chi,\psi)
    =
    \bigmeet_{a\in A}
    \bigl(\chi(a)\multimap\psi(a)\bigr).
\]
Its induced crisp order is pointwise.
\end{proposition}

\begin{proof}
The character preorder is a full fuzzy subpreorder of \([A,\Lambda]\).
\end{proof}

\begin{definition}[Characters of a fuzzy locale over \(\Lambda\)]
\label{def:fuzzy-character-locale-posetal}
For a fuzzy locale \(\mathcal X\), define
\[
    \operatorname{Char}_\Lambda(\mathcal X)
    \coloneqq
    \operatorname{Char}_\Lambda(\Opens_\Lambda(\mathcal X)).
\]
\end{definition}

\begin{proposition}[Characters of the truth-value frame]
\label{prop:characters-truth-value-frame}
There is exactly one scalar-frame character
\[
    \Lambda\to\Lambda,
\]
namely the identity.  Hence
\[
    \operatorname{Char}_\Lambda(\Lambda)
\]
is the terminal fuzzy preorder.
\end{proposition}

\begin{proof}
Let \(\chi:\Lambda\to\Lambda\) be a scalar-frame homomorphism.  Since \(\chi\)
preserves top,
\[
    \chi(\top_L)=\top_L.
\]
For every \(r\in L\),
\[
    \chi(r)
    =
    \chi(r\tensor\top_L)
    =
    r\tensor\chi(\top_L)
    =
    r\tensor\top_L
    =
    r.
\]
Thus \(\chi=1_\Lambda\).
\end{proof}

\begin{example}[The crisp case]
\label{ex:fuzzy-characters-crisp}
For \(\Lambda=2\), a character of a frame \(A\) is a frame homomorphism
\[
    A\to 2,
\]
equivalently a completely prime filter of \(A\).  Thus, in the crisp localic case,
characters are points.
\end{example}

\section{The upper fuzzy character adjunction}
\label{sec:upper-fuzzy-character-adjunction}

\subsection{The enriched adjunction}
\label{subsec:upper-character-adjunction-enriched}

\begin{theorem}[Upper fuzzy character adjunction]
\label{thm:upper-fuzzy-character-adjunction-enriched}
Let \(X\) be a small \(\mathbb V\)-category and let \(\mathcal X\) be a fuzzy
locale.  There is a natural equivalence
\[
    \mathbf{FLoc}_{\mathbb V}
    (\mathcal P_{\mathbb V}^{\uparrow}X,\mathcal X)
    \simeq
    \mathbb V\text{-}\mathbf{Cat}
    \bigl(
        X,
        \operatorname{Char}_{\mathbb V}(\mathcal X)
    \bigr).
\]
Equivalently, writing
\[
    A=\operatorname{Opens}_{\mathbb V}(\mathcal X),
\]
there is a natural equivalence
\[
    \mathbf{FFrm}_{\mathbb V}
    \bigl(
        A,
        P_{\mathbb V}(X^{\op})
    \bigr)
    \simeq
    \mathbb V\text{-}\mathbf{Cat}
    \bigl(
        X,
        \operatorname{Char}_{\mathbb V}(A)
    \bigr).
\]
\end{theorem}

\begin{proof}
A fuzzy-locale map
\[
    \mathcal P_{\mathbb V}^{\uparrow}X\to\mathcal X
\]
is frame-dually a fuzzy-frame morphism
\[
    h:A\to [X,\underline{\mathbb V}].
\]
By the enriched tensor--hom adjunction and the symmetry of \(\tensor\), this is
equivalently a \(\mathbb V\)-functor
\[
    \widehat h:X\to [A,\underline{\mathbb V}],
    \qquad
    \widehat h(x)(a)=h(a)(x).
\]
Explicitly, this is the composite transposition
\[
    \mathbb V\text{-}\mathbf{CAT}(A,[X,\underline{\mathbb V}])
    \cong
    \mathbb V\text{-}\mathbf{CAT}(A\tensor X,\underline{\mathbb V})
    \cong
    \mathbb V\text{-}\mathbf{CAT}(X,[A,\underline{\mathbb V}]).
\]

Since evaluation at \(x\) preserves the pointwise fuzzy-frame structure on
\([X,\underline{\mathbb V}]\), each
\[
    \widehat h(x):A\to\underline{\mathbb V}
\]
is a fuzzy character of \(A\).  Hence \(\widehat h\) factors through
\[
    \operatorname{Char}_{\mathbb V}(A)
    \subseteq
    [A,\underline{\mathbb V}].
\]

Conversely, a \(\mathbb V\)-functor
\[
    k:X\to\operatorname{Char}_{\mathbb V}(A)
\]
defines
\[
    \widetilde k:A\to[X,\underline{\mathbb V}],
    \qquad
    \widetilde k(a)(x)=k(x)(a).
\]
For each \(a\), this is an upper predicate on \(X\); for each \(x\), the map
\(a\mapsto k(x)(a)\) preserves the fuzzy-frame structure.  Therefore
\(\widetilde k\) is a fuzzy-frame morphism.  The two constructions are inverse by
the enriched tensor--hom adjunction, and naturality is inherited from the
naturality of that adjunction.
\end{proof}

\subsection{The posetal specialization}
\label{subsec:upper-character-adjunction-posetal}

\begin{theorem}[Posetal upper fuzzy character adjunction]
\label{thm:upper-fuzzy-character-adjunction-posetal}
Let \(X\) be a fuzzy preorder and let \(A\in\FFrm_\Lambda\).  There is a natural
isomorphism
\[
    \FFrm_\Lambda(A,P_\Lambda(X^{\op}))
    \cong
    \Pre_\Lambda(X,\operatorname{Char}_\Lambda(A)).
\]
\end{theorem}

\begin{proof}
For a fuzzy-frame homomorphism
\[
    h:A\to P_\Lambda(X^{\op}),
\]
define
\[
    \widehat h(x)(a)=h(a)(x).
\]
Each \(\widehat h(x)\) is a character because evaluation at \(x\) preserves the
pointwise scalar-frame structure.  Moreover, for all \(x,x'\in X\) and \(a\in A\),
the upper-predicate law for \(h(a)\) gives
\[
    X(x,x')\tensor h(a)(x)\leq h(a)(x'),
\]
or equivalently
\[
    X(x,x')\leq h(a)(x)\multimap h(a)(x').
\]
Taking the meet over \(a\) shows that
\[
    X(x,x')
    \leq
    \operatorname{Char}_\Lambda(A)(\widehat h(x),\widehat h(x')).
\]
Thus \(\widehat h\) is a \(\Lambda\)-functor.

Conversely, a \(\Lambda\)-functor
\[
    k:X\to\operatorname{Char}_\Lambda(A)
\]
defines
\[
    \widetilde k(a)(x)=k(x)(a).
\]
Functoriality of \(k\) gives the upper-predicate law in \(x\), while the character
condition on each \(k(x)\) gives preservation of joins, finite meets, and scalars
in \(a\).  Hence \(\widetilde k\) is a fuzzy-frame homomorphism.  The two
assignments are inverse.
\end{proof}

\begin{example}[The crisp case]
\label{ex:upper-character-adjunction-crisp}
For \(\Lambda=2\), the adjunction becomes
\[
    \Frm(A,\operatorname{Up}(X))
    \cong
    \Pre(X,\operatorname{Char}_2(A)).
\]
Thus frame maps into an upper Alexandrov frame are the same as monotone
point-valued families indexed by \(X\).
\end{example}

\begin{example}[Discrete fuzzy preorders]
\label{ex:upper-character-adjunction-discrete}
If \(X\) is discrete, then
\[
    P_\Lambda(X^{\op})\cong L^X.
\]
A fuzzy-frame homomorphism
\[
    A\to L^X
\]
is precisely an \(X\)-indexed family of characters \(A\to\Lambda\).
\end{example}

\section{Fuzzy ideals as finite-flat weights}
\label{sec:fuzzy-ideals-flat-weights}

\subsection{The enriched construction}
\label{subsec:fuzzy-ideals-enriched}

\begin{definition}[Finite-flat weight]
\label{def:finite-flat-weight}
Let \(X\) be a small \(\mathbb V\)-category.  A \textbf{finite-flat weight} on
\(X\) is a weight
\[
    I:X^{\op}\to\underline{\mathbb V}
\]
such that the weighted-colimit functional
\[
    \chi_I:[X,\underline{\mathbb V}]\to\underline{\mathbb V},
    \qquad
    \chi_I(U)=\int^{x\in X} I(x)\tensor U(x),
\]
preserves the chosen finite limits of the fuzzy-frame doctrine.  Equivalently,
\(I\)-weighted colimits in \(\underline{\mathbb V}\) commute with those finite
limits.
\end{definition}

\begin{definition}[Fuzzy ideal]
\label{def:fuzzy-ideal-enriched}
A \textbf{fuzzy ideal} of a small \(\mathbb V\)-category \(X\) is a finite-flat
weight
\[
    I:X^{\op}\to\underline{\mathbb V}.
\]
We write
\[
    \operatorname{Idl}_{\mathbb V}(X)
\]
for the full \(\mathbb V\)-subcategory of
\[
    [X^{\op},\underline{\mathbb V}]
\]
spanned by the finite-flat weights.
\end{definition}

\begin{proposition}[Flat-weight representation of cocontinuous functionals]
\label{prop:flat-weight-representation-of-characters}
Let \(X\) be small.  Cocontinuous \(\mathbb V\)-functors
\[
    [X,\underline{\mathbb V}]\to\underline{\mathbb V}
\]
are, up to canonical isomorphism, precisely the weighted-colimit functionals
\[
    U\longmapsto \int^{x\in X} I(x)\tensor U(x)
\]
for weights
\[
    I:X^{\op}\to\underline{\mathbb V}.
\]
Such a functional preserves the chosen finite limits precisely when \(I\) is
finite-flat.
\end{proposition}

\begin{proof}
Let
\[
    \chi:[X,\underline{\mathbb V}]\to\underline{\mathbb V}
\]
be cocontinuous.  Define
\[
    I_\chi(x)=\chi(X(x,-)).
\]
The assignment
\[
    x\longmapsto X(x,-)
\]
is the Yoneda embedding of \(X^{\op}\) into \([X,\underline{\mathbb V}]\).  Since
\(\chi\) is a \(\mathbb V\)-functor, \(I_\chi\) is a weight
\[
    I_\chi:X^{\op}\to\underline{\mathbb V}.
\]

For every
\[
    U:X\to\underline{\mathbb V},
\]
the enriched co-Yoneda formula in the copresheaf category
\([X,\underline{\mathbb V}]\) gives
\[
    U
    \cong
    \int^{x\in X} U(x)\tensor X(x,-).
\]
Applying the cocontinuous functor \(\chi\), we obtain
\[
\begin{aligned}
    \chi(U)
    &\cong
    \chi\left(\int^{x\in X} U(x)\tensor X(x,-)\right)\\
    &\cong
    \int^{x\in X} U(x)\tensor\chi(X(x,-))\\
    &=
    \int^{x\in X} U(x)\tensor I_\chi(x)\\
    &\cong
    \int^{x\in X} I_\chi(x)\tensor U(x).
\end{aligned}
\]
Thus \(\chi\) is represented by \(I_\chi\).

Conversely, every weight \(I:X^{\op}\to\underline{\mathbb V}\) determines the
weighted-colimit functional
\[
    \chi_I(U)=\int^{x\in X}I(x)\tensor U(x).
\]
This functional is cocontinuous because it is built from tensors and small
colimits in \(\underline{\mathbb V}\), and tensors preserve colimits in each
variable.  The two constructions are inverse by the same co-Yoneda calculation,
since
\[
    \chi_I(X(x,-))
    \cong
    \int^{y\in X}I(y)\tensor X(x,y)
    \cong
    I(x).
\]

Finally, preservation of the chosen finite limits by \(\chi_I\) is exactly the
definition of finite-flatness of the weight \(I\).
\end{proof}

\begin{theorem}[Characters of upper predicate frames]
\label{thm:characters-upper-predicate-frames-flat}
For every small \(\mathbb V\)-category \(X\), there is a canonical equivalence
\[
    \operatorname{Char}_{\mathbb V}(P_{\mathbb V}(X^{\op}))
    \simeq
    \operatorname{Idl}_{\mathbb V}(X).
\]
Under this equivalence, a fuzzy ideal
\[
    I:X^{\op}\to\underline{\mathbb V}
\]
corresponds to the character
\[
    \chi_I(U)
    =
    \int^{x\in X}I(x)\tensor U(x).
\]
\end{theorem}

\begin{proof}
Since
\[
    P_{\mathbb V}(X^{\op})=[X,\underline{\mathbb V}],
\]
a character is a cocontinuous, finite-limit-preserving, scalar-preserving
\(\mathbb V\)-functor
\[
    [X,\underline{\mathbb V}]\to\underline{\mathbb V}.
\]
By Proposition~\ref{prop:flat-weight-representation-of-characters}, the
cocontinuous finite-limit-preserving functors are exactly the finite-flat
weighted-colimit functionals
\[
    U\longmapsto\int^{x\in X}I(x)\tensor U(x).
\]
Scalar preservation is automatic for these functionals because tensoring
distributes over the coends defining the weighted colimits:
\[
\begin{aligned}
    \chi_I(r\tensor U)
    &=
    \int^{x\in X} I(x)\tensor r\tensor U(x)\\
    &\cong
    r\tensor
    \int^{x\in X}I(x)\tensor U(x)\\
    &=
    r\tensor\chi_I(U).
\end{aligned}
\]
Thus characters of the upper predicate frame are precisely finite-flat weights on
\(X\).
\end{proof}

\subsection{The posetal specialization}
\label{subsec:fuzzy-ideals-posetal}

\begin{definition}[Fuzzy ideal over \(\Lambda\)]
\label{def:fuzzy-ideal-posetal}
Let \(X\) be a fuzzy preorder over \(\Lambda\).  A \textbf{fuzzy ideal} of \(X\) is
a lower fuzzy predicate
\[
    I:X^{\op}\to\Lambda
\]
such that
\[
    \chi_I(U)
    =
    \bigjoin_{x\in X} I(x)\tensor U(x)
\]
preserves top and binary meets as a map
\[
    P_\Lambda(X^{\op})\to\Lambda.
\]
Equivalently, \(I\) satisfies the lower-predicate law
\[
    X(x',x)\tensor I(x)\leq I(x')
\]
and the finite-flatness equations
\[
    \bigjoin_{x\in X}I(x)=\top_L
\]
and
\[
    \bigjoin_{x\in X}I(x)\tensor\bigl(U(x)\meet V(x)\bigr)
    =
    \left(\bigjoin_{x\in X}I(x)\tensor U(x)\right)
    \meet
    \left(\bigjoin_{x\in X}I(x)\tensor V(x)\right)
\]
for all upper fuzzy predicates \(U,V\).
\end{definition}

\begin{definition}[The fuzzy preorder of fuzzy ideals]
\label{def:fuzzy-ideal-preorder-posetal}
For a fuzzy preorder \(X\), define
\[
    \operatorname{Idl}_\Lambda(X)
\]
to be the full fuzzy subpreorder of \(P_\Lambda(X)\) spanned by the fuzzy ideals.
Thus, for fuzzy ideals \(I,J\),
\[
    \operatorname{Idl}_\Lambda(X)(I,J)
    =
    \bigmeet_{x\in X}
    \bigl(I(x)\multimap J(x)\bigr).
\]
\end{definition}

\begin{theorem}[Characters of upper fuzzy Alexandrov frames]
\label{thm:characters-upper-alexandrov-fuzzy-ideals}
For every fuzzy preorder \(X\), there is an isomorphism of fuzzy preorders
\[
    \operatorname{Char}_\Lambda(P_\Lambda(X^{\op}))
    \cong
    \operatorname{Idl}_\Lambda(X).
\]
Under this isomorphism, a fuzzy ideal \(I\) corresponds to the character
\[
    \chi_I(U)
    =
    \bigjoin_{x\in X}I(x)\tensor U(x).
\]
\end{theorem}

\begin{proof}
The coend formula in
Theorem~\ref{thm:characters-upper-predicate-frames-flat} specializes to the
displayed join formula, and finite-flatness specializes to preservation of top and
binary meets.

It remains to identify hom-values.  If
\[
    r\leq \bigmeet_x (I(x)\multimap J(x)),
\]
then \(r\tensor I(x)\leq J(x)\) for every \(x\), and therefore
\[
    r\tensor \chi_I(U)
    =
    r\tensor \bigjoin_x I(x)\tensor U(x)
    \leq
    \bigjoin_x J(x)\tensor U(x)
    =
    \chi_J(U)
\]
for every upper predicate \(U\).  Hence
\[
    r\leq
    \bigmeet_U(\chi_I(U)\multimap\chi_J(U)).
\]

Conversely, evaluate at the upper representable \(X(x,-)\).  The enriched Yoneda
calculation gives
\[
    \chi_I(X(x,-))
    =
    \bigjoin_y I(y)\tensor X(x,y)
    =
    I(x),
\]
and similarly for \(J\).  Hence the character hom is bounded above by
\[
    \bigmeet_x(I(x)\multimap J(x)).
\]
Thus the hom-values agree.
\end{proof}

\begin{example}[The crisp case]
\label{ex:fuzzy-ideals-crisp}
For \(\Lambda=2\), a fuzzy ideal is an ordinary ideal
\[
    I\subseteq X.
\]
Thus \(I\) is a lower set, is inhabited, and is directed.  Therefore
\[
    \operatorname{Char}_2(\operatorname{Up}(X))
    \cong
    \operatorname{Idl}(X).
\]
\end{example}

\begin{example}[Standard truth-value objects]
\label{ex:standard-fuzzy-ideals}
For G\"odel truth values,
\[
    \chi_I(U)=\bigjoin_x \min(I(x),U(x)).
\]
For {\L}ukasiewicz truth values,
\[
    \chi_I(U)
    =
    \bigjoin_x \max(0,I(x)+U(x)-1).
\]
For product truth values,
\[
    \chi_I(U)=\bigjoin_x I(x)U(x).
\]
For the Lawvere quantale, joins in \(\Lambda\) are usual infima and tensor is
addition, so
\[
    \chi_I(U)
    =
    \inf_x\bigl(I(x)+U(x)\bigr).
\]
\end{example}

\section{The fuzzy ideal monad}
\label{sec:fuzzy-ideal-monad}

\subsection{The enriched monad}
\label{subsec:fuzzy-ideal-monad-enriched}

\begin{definition}[Fuzzy ideal monad]
\label{def:fuzzy-ideal-monad-enriched}
The \textbf{fuzzy ideal monad} is the monad
\[
    T_{\mathbb V}
    =
    \operatorname{Char}_{\mathbb V}\mathcal P_{\mathbb V}^{\uparrow}
\]
induced by the adjunction
\[
    \mathcal P_{\mathbb V}^{\uparrow}
    \dashv
    \operatorname{Char}_{\mathbb V}.
\]
Thus
\[
    T_{\mathbb V}X
    =
    \operatorname{Char}_{\mathbb V}(P_{\mathbb V}(X^{\op}))
    \simeq
    \operatorname{Idl}_{\mathbb V}(X).
\]
After the strictification convention of
Remark~\ref{rem:strictification-ideal-monad}, we write simply
\[
    T_{\mathbb V}X=\operatorname{Idl}_{\mathbb V}(X).
\]
\end{definition}

\begin{proposition}[Unit and multiplication]
\label{prop:fuzzy-ideal-monad-unit-multiplication-enriched}
Under the finite-flat-weight description, the unit
\[
    \eta_X:X\to\operatorname{Idl}_{\mathbb V}(X)
\]
is the Yoneda embedding into finite-flat weights:
\[
    \eta_X(x)=X(-,x).
\]
The multiplication
\[
    \mu_X:
    \operatorname{Idl}_{\mathbb V}(\operatorname{Idl}_{\mathbb V}(X))
    \to
    \operatorname{Idl}_{\mathbb V}(X)
\]
is flattening:
\[
    \mu_X(\mathcal I)(x)
    =
    \int^{I\in\operatorname{Idl}_{\mathbb V}(X)}
    \mathcal I(I)\tensor I(x).
\]
\end{proposition}

\begin{proof}
The unit of the character adjunction sends \(x\) to evaluation at \(x\):
\[
    U\longmapsto U(x).
\]
The corresponding weight evaluates covariant representables by
\[
    X(y,-)(x)=X(y,x),
\]
so \(\eta_X(x)=X(-,x)\).  The multiplication is the monad multiplication induced by
the adjunction, transported through the finite-flat-weight representation; it is
the weighted colimit of an ideal of ideals.
\end{proof}

\subsection{The posetal specialization}
\label{subsec:fuzzy-ideal-monad-posetal}

\begin{definition}[The posetal fuzzy ideal monad]
\label{def:fuzzy-ideal-monad-posetal}
In the posetal specialization,
\[
    T_\Lambda(X)=\operatorname{Idl}_\Lambda(X).
\]
The unit is
\[
    \eta_X:X\to\operatorname{Idl}_\Lambda(X),
    \qquad
    \eta_X(x)(y)=X(y,x),
\]
and the multiplication is
\[
    \mu_X:
    \operatorname{Idl}_\Lambda(\operatorname{Idl}_\Lambda(X))
    \to
    \operatorname{Idl}_\Lambda(X),
\]
where
\[
    \mu_X(\mathcal I)(x)
    =
    \bigjoin_{I\in\operatorname{Idl}_\Lambda(X)}
    \mathcal I(I)\tensor I(x).
\]
\end{definition}

\begin{proposition}[Functorial action of the fuzzy ideal monad]
\label{prop:fuzzy-ideal-monad-functorial-action}
Let
\[
    f:X\to Y
\]
be a \(\Lambda\)-functor.  Then
\[
    T_\Lambda(f):
    \operatorname{Idl}_\Lambda(X)\to\operatorname{Idl}_\Lambda(Y)
\]
is given by
\[
    T_\Lambda(f)(I)(y)
    =
    \bigjoin_{x\in X}I(x)\tensor Y(y,f(x)).
\]
This is the fuzzy lower closure of the direct image of \(I\) along \(f\).
\end{proposition}

\begin{proof}
This is the posetal specialization of enriched left Kan extension of a weight
along \(f\).  Since it is the functorial action of the monad induced by the
character adjunction, it preserves fuzzy ideals and is functorial in \(f\).
\end{proof}

\begin{example}[The crisp ideal monad]
\label{ex:crisp-ideal-monad}
For \(\Lambda=2\),
\[
    T_2(X)=\operatorname{Idl}(X).
\]
The unit sends
\[
    x\longmapsto\downarrow x,
\]
multiplication sends an ideal of ideals to its union,
\[
    \mu_X(\mathcal I)=\bigcup_{I\in\mathcal I}I,
\]
and for a monotone map \(f:X\to Y\),
\[
    T_2(f)(I)=\downarrow f[I].
\]
Thus \(T_2\) is the ordinary ideal-completion monad.
\end{example}

\begin{example}[Standard direct-image formulas]
\label{ex:standard-ideal-direct-image}
For G\"odel truth values,
\[
    T_\Lambda(f)(I)(y)
    =
    \bigjoin_x \min(I(x),Y(y,f(x))).
\]
For {\L}ukasiewicz truth values,
\[
    T_\Lambda(f)(I)(y)
    =
    \bigjoin_x \max(0,I(x)+Y(y,f(x))-1).
\]
For product truth values,
\[
    T_\Lambda(f)(I)(y)
    =
    \bigjoin_x I(x)Y(y,f(x)).
\]
For the Lawvere quantale,
\[
    T_\Lambda(f)(I)(y)
    =
    \inf_x\bigl(I(x)+d_Y(y,f(x))\bigr).
\]
\end{example}

\section{Weak fuzzy domains and fuzzy domains}
\label{sec:weak-fuzzy-domains-fuzzy-domains}

\subsection{The enriched notion}
\label{subsec:weak-fuzzy-domains-enriched}

\begin{definition}[Weak fuzzy domain]
\label{def:weak-fuzzy-domain-enriched}
A \textbf{weak fuzzy domain} is a small \(\mathbb V\)-category \(X\) equipped with
a chosen \(\mathbb V\)-functor
\[
    \alpha:T_{\mathbb V}X\to X.
\]
Using
\[
    T_{\mathbb V}X=\operatorname{Idl}_{\mathbb V}(X),
\]
we regard \(\alpha\) as assigning to each fuzzy ideal
\[
    I:X^{\op}\to\underline{\mathbb V}
\]
a chosen candidate fuzzy ideal-join
\[
    \alpha(I)\in X.
\]
\end{definition}

\begin{definition}[Fuzzy domain]
\label{def:fuzzy-domain-enriched}
A \textbf{fuzzy domain} is a weak fuzzy domain whose structure map
\[
    \alpha:T_{\mathbb V}X\to X
\]
is a strict algebra for the fuzzy ideal monad:
\[
    \alpha\eta_X=1_X,
    \qquad
    \alpha\mu_X=\alpha T_{\mathbb V}(\alpha).
\]
\end{definition}

\begin{remark}
Weak fuzzy domains carry chosen candidate ideal-joins.  Strict fuzzy domains are
those for which these choices are coherent with principal ideals and iterated
fuzzy ideals.
\end{remark}

\subsection{The posetal specialization}
\label{subsec:weak-fuzzy-domains-posetal}

\begin{definition}[Weak fuzzy domain over \(\Lambda\)]
\label{def:weak-fuzzy-domain-posetal}
A \textbf{weak fuzzy domain over \(\Lambda\)} is a fuzzy preorder \(X\) equipped
with a \(\Lambda\)-functor
\[
    \alpha:\operatorname{Idl}_\Lambda(X)\to X.
\]
\end{definition}

\begin{definition}[Fuzzy domain over \(\Lambda\)]
\label{def:fuzzy-domain-posetal}
A \textbf{fuzzy domain over \(\Lambda\)} is a weak fuzzy domain satisfying
\[
    \alpha(\eta_X x)=x
\]
and
\[
    \alpha(\mu_X\mathcal I)=\alpha(T_\Lambda\alpha(\mathcal I)).
\]
Equivalently,
\[
    \alpha(X(-,x))=x
\]
and
\[
\alpha\left(
    x\longmapsto
    \bigjoin_{I\in\operatorname{Idl}_\Lambda(X)}
    \mathcal I(I)\tensor I(x)
\right)
=
\alpha\left(
    x\longmapsto
    \bigjoin_{I\in\operatorname{Idl}_\Lambda(X)}
    \mathcal I(I)\tensor X(x,\alpha(I))
\right).
\]
\end{definition}

\begin{theorem}[Crisp algebras are dcpos]
\label{thm:crisp-algebras-are-dcpos}
For \(\Lambda=2\), strict algebras for the ideal monad on separated preorders are
precisely dcpos, where dcpos are understood to have suprema of nonempty directed
subsets.
\end{theorem}

\begin{proof}
Let
\[
    \alpha:\operatorname{Idl}(X)\to X
\]
be a strict algebra.  The unit law says
\[
    \alpha(\downarrow x)=x.
\]
If \(I\in\operatorname{Idl}(X)\) and \(x\in I\), then
\[
    \downarrow x\subseteq I,
\]
so by monotonicity
\[
    x=\alpha(\downarrow x)\leq \alpha(I).
\]
Thus \(\alpha(I)\) is an upper bound of \(I\).  If \(y\) is any upper bound of
\(I\), then
\[
    I\subseteq\downarrow y,
\]
and again by monotonicity,
\[
    \alpha(I)\leq\alpha(\downarrow y)=y.
\]
Hence \(\alpha(I)=\bigjoin I\).  Every nonempty directed subset generates an ideal
with the same supremum, so \(X\) is a dcpo when separated.

Conversely, if \(X\) is a dcpo, define
\[
    \alpha(I)=\bigjoin I.
\]
The unit law is immediate, and the multiplication law says that the supremum of a
directed family of directed suprema is the supremum of its union.  Thus \(X\) is a
strict algebra.
\end{proof}

\begin{example}[G\"odel, product, and Lawvere weak fuzzy domains]
\label{ex:weak-fuzzy-domain-standard}
A G\"odel-valued weak fuzzy domain assigns candidate joins to finite-flat lower
weights governed by the minimum tensor.  A product-valued weak fuzzy domain assigns
candidate joins to finite-flat lower weights with multiplicative membership.  For
the Lawvere quantale, a weak fuzzy domain assigns colimits to finite-flat Lawvere
weights, giving a metric-enriched domain-like completion structure.
\end{example}

\section{The category of fuzzy domains}
\label{sec:category-fuzzy-domains}

We now record the categorical structure carried by fuzzy domains.  The point is
not to add extra structure to a fuzzy domain, but to observe that the
finite-flat doctrine already used to define fuzzy ideals also determines the
colimits, morphisms, tensor products, and internal homs of fuzzy domains.

Throughout this section, finite-flatness is understood relative to the
finite-limit doctrine fixed at the beginning of the chapter.  Thus, in each
predicate category \([X,\underline{\mathbb V}]\), the specified finite limits are
computed pointwise.  In the posetal scalar-frame specialization, this means
preservation of the terminal predicate and binary meets.

\begin{definition}[The category of fuzzy domains]
\label{def:category-fuzzy-domains}
A morphism of fuzzy domains
\[
    f:(X,\alpha_X)\longrightarrow (Y,\alpha_Y)
\]
is a \(\mathbb V\)-functor
\[
    f:X\to Y
\]
such that
\[
    f\alpha_X=\alpha_YT_{\mathbb V}(f).
\]
The ordinary category of fuzzy domains and fuzzy-domain morphisms is denoted
\[
    \mathbf{FDom}_{\mathbb V}.
\]
\end{definition}

Thus \(\mathbf{FDom}_{\mathbb V}\) is the Eilenberg--Moore category of strict
algebras for the fuzzy ideal monad.  The rest of this section identifies these
strict algebras as finite-flat-cocomplete \(\mathbb V\)-categories and then applies
Kelly's tensor-product theorem for saturated commutative classes of weights.

\subsection{Finite-flat colimits in fuzzy domains}
\label{subsec:finite-flat-colimits-fdom}

\begin{lemma}[Representables are finite-flat]
\label{lem:representables-finite-flat-fdom}
Let \(X\) be a small \(\mathbb V\)-category.  For every \(x\in X\), the
representable weight
\[
    X(-,x):X^{\op}\to\underline{\mathbb V}
\]
is finite-flat.
\end{lemma}

\begin{proof}
The colimit functional associated to \(X(-,x)\) is
\[
    \chi_{X(-,x)}(U)
    =
    \int^{y\in X}X(y,x)\tensor U(y).
\]
By enriched co-Yoneda,
\[
    \chi_{X(-,x)}(U)\cong U(x).
\]
Thus \(\chi_{X(-,x)}\) is evaluation at \(x\).  Since the specified finite
limits in \([X,\underline{\mathbb V}]\) are computed pointwise, evaluation at
\(x\) preserves them.  Hence \(X(-,x)\) is finite-flat.
\end{proof}

\begin{lemma}[Finite-flat direct images]
\label{lem:finite-flat-direct-images-fdom}
Let
\[
    D:K\to X
\]
be a \(\mathbb V\)-functor, and let
\[
    W:K^{\op}\to\underline{\mathbb V}
\]
be finite-flat.  Define the lower direct image
\[
    D_!W:X^{\op}\to\underline{\mathbb V}
\]
by
\[
    (D_!W)(x)
    =
    \int^{k\in K}W(k)\tensor X(x,Dk).
\]
Then \(D_!W\) is finite-flat.
\end{lemma}

\begin{proof}
Let
\[
    U:X\to\underline{\mathbb V}
\]
be an upper predicate on \(X\).  Then
\[
\begin{aligned}
    \chi_{D_!W}(U)
    &=
    \int^{x\in X}(D_!W)(x)\tensor U(x)
    \\
    &=
    \int^{x\in X}
        \left(
            \int^{k\in K}W(k)\tensor X(x,Dk)
        \right)\tensor U(x)
    \\
    &\cong
    \int^{k\in K}
        W(k)\tensor
        \left(
            \int^{x\in X}X(x,Dk)\tensor U(x)
        \right)
    \\
    &\cong
    \int^{k\in K}W(k)\tensor U(Dk)
    \\
    &=
    \chi_W(UD).
\end{aligned}
\]
The first isomorphism is Fubini for coends, and the second is enriched
co-Yoneda.

Precomposition
\[
    D^*:[X,\underline{\mathbb V}]\to[K,\underline{\mathbb V}],
    \qquad
    U\longmapsto UD,
\]
preserves the specified finite limits because they are computed pointwise.
Since \(W\) is finite-flat, \(\chi_W\) preserves those finite limits.  Hence
\[
    \chi_{D_!W}\cong \chi_WD^*
\]
preserves them, so \(D_!W\) is finite-flat.
\end{proof}

\begin{proposition}[Finite-flat saturation]
\label{prop:finite-flat-saturation-fdom}
Finite-flat weights are closed under finite-flat weighted colimits in presheaf
categories.

Explicitly, let
\[
    W:K^{\op}\to\underline{\mathbb V}
\]
be finite-flat, and let
\[
    H:K\to\operatorname{Idl}_{\mathbb V}(X)
\]
be a \(\mathbb V\)-functor.  Define
\[
    J:X^{\op}\to\underline{\mathbb V}
\]
pointwise by
\[
    J(x)
    =
    \int^{k\in K}W(k)\tensor Hk(x).
\]
Then \(J\) is finite-flat.
\end{proposition}

\begin{proof}
For every upper predicate
\[
    U:X\to\underline{\mathbb V},
\]
we have
\[
\begin{aligned}
    \chi_J(U)
    &=
    \int^{x\in X}J(x)\tensor U(x)
    \\
    &=
    \int^{x\in X}
        \left(
            \int^{k\in K}W(k)\tensor Hk(x)
        \right)\tensor U(x)
    \\
    &\cong
    \int^{k\in K}
        W(k)\tensor
        \left(
            \int^{x\in X}Hk(x)\tensor U(x)
        \right)
    \\
    &=
    \int^{k\in K}W(k)\tensor \chi_{Hk}(U).
\end{aligned}
\]
Thus
\[
    \chi_J(U)\cong\chi_W(\overline H(U)),
\]
where
\[
    \overline H:[X,\underline{\mathbb V}]\to[K,\underline{\mathbb V}]
\]
is defined by
\[
    \overline H(U)(k)=\chi_{Hk}(U).
\]

We first check that \(\overline H(U)\) is an upper predicate on \(K\).  Since
\(H\) is a \(\mathbb V\)-functor, for every \(k,k'\in K\) there is a morphism
\[
    K(k,k')\longrightarrow
    \operatorname{Idl}_{\mathbb V}(X)(Hk,Hk').
\]
Functoriality of the coend construction gives, naturally in \(U\),
\[
    \operatorname{Idl}_{\mathbb V}(X)(Hk,Hk')
    \longrightarrow
    \underline{\mathbb V}(\chi_{Hk}(U),\chi_{Hk'}(U)).
\]
Composing these maps gives
\[
    K(k,k')\longrightarrow
    \underline{\mathbb V}(\overline H(U)(k),\overline H(U)(k')),
\]
so \(\overline H(U)\) is indeed a \(\mathbb V\)-functor
\[
    K\to\underline{\mathbb V}.
\]

Now let \((U_i)\) be a finite diagram of upper predicates whose specified limit
exists in the finite-limit doctrine.  Since each \(Hk\) is finite-flat,
\[
    \chi_{Hk}(\lim_i U_i)
    \cong
    \lim_i\chi_{Hk}(U_i).
\]
Since specified finite limits in \([K,\underline{\mathbb V}]\) are computed
pointwise, this gives
\[
    \overline H(\lim_i U_i)\cong \lim_i\overline H(U_i).
\]
Thus \(\overline H\) preserves the specified finite limits.  Since \(W\) is
finite-flat, \(\chi_W\) preserves them.  Hence
\[
    \chi_J\cong\chi_W\overline H
\]
preserves the specified finite limits, so \(J\) is finite-flat.
\end{proof}

\begin{corollary}[Flattening is finite-flat]
\label{cor:flattening-finite-flat-fdom}
If
\[
    \mathcal I:
    \operatorname{Idl}_{\mathbb V}(X)^{\op}\to\underline{\mathbb V}
\]
is finite-flat, then the flattened weight
\[
    \mu_X(\mathcal I)(x)
    =
    \int^{I\in\operatorname{Idl}_{\mathbb V}(X)}
        \mathcal I(I)\tensor I(x)
\]
is finite-flat.
\end{corollary}

\begin{proof}
Apply Proposition~\ref{prop:finite-flat-saturation-fdom} to the identity diagram
\[
    1_{\operatorname{Idl}_{\mathbb V}(X)}:
    \operatorname{Idl}_{\mathbb V}(X)\to\operatorname{Idl}_{\mathbb V}(X).
\]
\end{proof}

\begin{proposition}[The KZ comparison for the ideal monad]
\label{prop:ideal-monad-kz-fdom}
For every small \(\mathbb V\)-category \(X\), there is an enriched natural
transformation
\[
    T_{\mathbb V}\eta_X\Longrightarrow \eta_{T_{\mathbb V}X}.
\]
Equivalently, for every finite-flat weight
\[
    I:X^{\op}\to\underline{\mathbb V},
\]
there is a canonical morphism of weights
\[
    (T_{\mathbb V}\eta_X)(I)
    \longrightarrow
    \operatorname{Idl}_{\mathbb V}(X)(-,I).
\]
Together with the monad unit and multiplication, these comparisons make the ideal
monad a KZ, or lax-idempotent, monad.
\end{proposition}

\begin{proof}
Let
\[
    I,J\in\operatorname{Idl}_{\mathbb V}(X).
\]
The value of \((T_{\mathbb V}\eta_X)(I)\) at \(J\) is
\[
    (T_{\mathbb V}\eta_X)(I)(J)
    =
    \int^{x\in X}
        I(x)\tensor
        \operatorname{Idl}_{\mathbb V}(X)(J,X(-,x)).
\]
By enriched Yoneda,
\[
    I(x)
    \cong
    [X^{\op},\underline{\mathbb V}](X(-,x),I).
\]
Composition in the presheaf category gives morphisms
\[
    \operatorname{Idl}_{\mathbb V}(X)(J,X(-,x))
    \tensor
    [X^{\op},\underline{\mathbb V}](X(-,x),I)
    \longrightarrow
    \operatorname{Idl}_{\mathbb V}(X)(J,I).
\]
Tensoring with the Yoneda identification and integrating over \(x\) gives
\[
    (T_{\mathbb V}\eta_X)(I)(J)
    \longrightarrow
    \operatorname{Idl}_{\mathbb V}(X)(J,I).
\]
This is natural in \(J\), and therefore defines a morphism of weights
\[
    (T_{\mathbb V}\eta_X)(I)
    \longrightarrow
    \operatorname{Idl}_{\mathbb V}(X)(-,I).
\]
It is also enriched-natural in \(I\), since it is induced by composition in
\([X^{\op},\underline{\mathbb V}]\).

We now record the KZ consequences explicitly.  The monad laws give
\[
    \mu_X\,T_{\mathbb V}\eta_X=1_{T_{\mathbb V}X},
    \qquad
    \mu_X\,\eta_{T_{\mathbb V}X}=1_{T_{\mathbb V}X}.
\]
The comparison just constructed gives
\[
    T_{\mathbb V}\eta_X\Longrightarrow \eta_{T_{\mathbb V}X}.
\]
Composing with \(\mu_X\) and using the first monad law gives the identity on
\(T_{\mathbb V}X\).  Equivalently, the comparison exhibits
\[
    \mu_X\dashv \eta_{T_{\mathbb V}X}
\]
in the enriched order of endomorphisms.  The remaining triangle inequality is the
statement that every finite-flat ideal of finite-flat ideals is contained in the
principal ideal generated by its flattening:
\[
    1_{T_{\mathbb V}T_{\mathbb V}X}
    \Longrightarrow
    \eta_{T_{\mathbb V}X}\mu_X.
\]
On a finite-flat weight
\[
    \mathcal I:
    \operatorname{Idl}_{\mathbb V}(X)^{\op}\to\underline{\mathbb V},
\]
this is the morphism
\[
    \mathcal I(J)
    \longrightarrow
    \operatorname{Idl}_{\mathbb V}(X)(J,\mu_X\mathcal I)
\]
obtained by composing
\[
    \mathcal I(J)\tensor J(x)
    \longrightarrow
    \int^{I}\mathcal I(I)\tensor I(x)
    =
    \mu_X(\mathcal I)(x)
\]
and transposing in the presheaf hom.  Naturality and the triangle identities are
the associativity and unitality of composition in
\([X^{\op},\underline{\mathbb V}]\), together with Fubini for the coends defining
flattening.  Hence the ideal monad is KZ.
\end{proof}

\begin{theorem}[The algebra map is left adjoint to Yoneda]
\label{thm:algebra-map-left-adjoint-y-fdom}
Let
\[
    \alpha:\operatorname{Idl}_{\mathbb V}(X)\to X
\]
be a strict algebra structure for the fuzzy ideal monad.  Then
\[
    \alpha\dashv \eta_X
\]
with identity counit
\[
    \alpha\eta_X=1_X.
\]
Consequently, enriched-naturally in
\(I\in\operatorname{Idl}_{\mathbb V}(X)\) and \(x\in X\), there is an isomorphism
\[
    X(\alpha I,x)
    \cong
    \operatorname{Idl}_{\mathbb V}(X)(I,X(-,x)).
\]
\end{theorem}

\begin{proof}
The algebra unit law gives
\[
    \alpha\eta_X=1_X,
\]
so the counit of the desired adjunction is the identity.

It remains to construct the unit
\[
    1_{\operatorname{Idl}_{\mathbb V}(X)}
    \longrightarrow
    \eta_X\alpha.
\]
By Proposition~\ref{prop:ideal-monad-kz-fdom}, we have
\[
    T_{\mathbb V}\eta_X\Longrightarrow \eta_{T_{\mathbb V}X}.
\]
Applying \(T_{\mathbb V}\alpha\) gives
\[
    T_{\mathbb V}\alpha\;T_{\mathbb V}\eta_X
    \longrightarrow
    T_{\mathbb V}\alpha\;\eta_{T_{\mathbb V}X}.
\]
The left-hand side is
\[
    T_{\mathbb V}(\alpha\eta_X)=T_{\mathbb V}(1_X)=1_{\operatorname{Idl}_{\mathbb V}(X)}.
\]
By naturality of \(\eta\), the right-hand side is
\[
    \eta_X\alpha.
\]
Thus we obtain
\[
    1_{\operatorname{Idl}_{\mathbb V}(X)}
    \longrightarrow
    \eta_X\alpha.
\]
The KZ coherence supplies the triangle identities.  Therefore
\[
    \alpha\dashv\eta_X.
\]
The displayed hom-isomorphism is the enriched hom-form of this adjunction.
\end{proof}

\begin{theorem}[Fuzzy domains and finite-flat colimits]
\label{thm:fuzzy-domains-ff-colimits}
Let \(X\) be a small \(\mathbb V\)-category.  To give \(X\) the structure of a
fuzzy domain is equivalently to give a strict coherent choice of all finite-flat
weighted colimits in \(X\).

Under this equivalence, the algebra map
\[
    \alpha_X:\operatorname{Idl}_{\mathbb V}(X)\to X
\]
sends a finite-flat weight \(I:X^{\op}\to\underline{\mathbb V}\) to the
\(I\)-weighted colimit of the identity diagram:
\[
    \alpha_X(I)=I*1_X.
\]
More generally, if
\[
    D:K\to X
\]
is a diagram and
\[
    W:K^{\op}\to\underline{\mathbb V}
\]
is finite-flat, then
\[
    W*D=\alpha_X(D_!W),
\]
where
\[
    (D_!W)(x)
    =
    \int^{k\in K}W(k)\tensor X(x,Dk).
\]

A \(\mathbb V\)-functor between fuzzy domains is a fuzzy-domain morphism if and
only if it preserves finite-flat weighted colimits.
\end{theorem}

\begin{proof}
Suppose first that
\[
    \alpha_X:\operatorname{Idl}_{\mathbb V}(X)\to X
\]
is a strict algebra.  By
Theorem~\ref{thm:algebra-map-left-adjoint-y-fdom},
\[
    \alpha_X\dashv\eta_X.
\]
Hence, for every finite-flat weight \(I\) and every \(x\in X\),
\[
\begin{aligned}
    X(\alpha_XI,x)
    &\cong
    \operatorname{Idl}_{\mathbb V}(X)(I,\eta_Xx)
    \\
    &=
    [X^{\op},\underline{\mathbb V}](I,X(-,x)).
\end{aligned}
\]
This is the universal property of the \(I\)-weighted colimit of the identity
diagram.  Thus
\[
    \alpha_X(I)=I*1_X.
\]

Now let \(D:K\to X\), and let \(W:K^{\op}\to\underline{\mathbb V}\) be
finite-flat.  By Lemma~\ref{lem:finite-flat-direct-images-fdom}, \(D_!W\) is
finite-flat.  For every \(x\in X\),
\[
\begin{aligned}
    X(\alpha_X(D_!W),x)
    &\cong
    [X^{\op},\underline{\mathbb V}](D_!W,X(-,x))
    \\
    &\cong
    [K^{\op},\underline{\mathbb V}](W,X(D-,x)).
\end{aligned}
\]
The second isomorphism is the weighted left-Kan-extension adjunction defining
\(D_!W\).  This is the universal property of the weighted colimit \(W*D\).
Therefore
\[
    W*D=\alpha_X(D_!W)
\]
under the strictification convention.

Conversely, suppose \(X\) is equipped with a strict coherent choice of all
finite-flat weighted colimits.  Define
\[
    \alpha_X(I)=I*1_X.
\]
The representing isomorphisms
\[
    X(\alpha_XI,x)
    \cong
    [X^{\op},\underline{\mathbb V}](I,X(-,x))
\]
are enriched-natural in \(I\) and \(x\).  Hence \(I\mapsto \alpha_XI\) extends
uniquely to a \(\mathbb V\)-functor
\[
    \alpha_X:\operatorname{Idl}_{\mathbb V}(X)\to X
\]
left adjoint to the Yoneda embedding.

For a representable weight,
\[
    X(-,x)*1_X=x,
\]
so
\[
    \alpha_X\eta_X=1_X.
\]
For a finite-flat weight
\[
    \mathcal I:
    \operatorname{Idl}_{\mathbb V}(X)^{\op}\to\underline{\mathbb V},
\]
the left-hand side of the multiplication law is
\[
    \alpha_X\mu_X(\mathcal I)=\mu_X(\mathcal I)*1_X.
\]
The right-hand side is
\[
    \alpha_XT_{\mathbb V}(\alpha_X)(\mathcal I)
    =
    (T_{\mathbb V}\alpha_X)(\mathcal I)*1_X.
\]
The weight \(T_{\mathbb V}(\alpha_X)(\mathcal I)\) is the lower direct image of
\(\mathcal I\) along
\[
    \alpha_X:\operatorname{Idl}_{\mathbb V}(X)\to X.
\]
Hence its colimit against \(1_X\) is the same weighted colimit as
\(\mathcal I\) against the diagram \(\alpha_X\):
\[
    (T_{\mathbb V}\alpha_X)(\mathcal I)*1_X
    \cong
    \mathcal I*\alpha_X.
\]
On the other hand,
\[
    \mu_X(\mathcal I)(x)
    =
    \int^{I}\mathcal I(I)\tensor I(x)
\]
is the flattened weight computing the same iterated finite-flat colimit.
Strict coherence of the chosen finite-flat colimits gives
\[
    \mu_X(\mathcal I)*1_X
    =
    \mathcal I*(I\mapsto I*1_X)
    =
    \mathcal I*\alpha_X.
\]
Therefore
\[
    \alpha_X\mu_X=\alpha_XT_{\mathbb V}(\alpha_X),
\]
so \(\alpha_X\) is a strict algebra.

Finally, let
\[
    f:(X,\alpha_X)\to(Y,\alpha_Y)
\]
be a \(\mathbb V\)-functor.  If \(f\) is a fuzzy-domain morphism, then for every
finite-flat \(I:X^{\op}\to\underline{\mathbb V}\),
\[
    f(I*1_X)
    =
    f\alpha_X(I)
    =
    \alpha_YT_{\mathbb V}(f)(I)
    =
    T_{\mathbb V}(f)(I)*1_Y.
\]
Thus \(f\) preserves finite-flat colimits of identity diagrams, and hence, by
the direct-image formula, all finite-flat weighted colimits.

Conversely, if \(f\) preserves finite-flat weighted colimits, then for every
finite-flat \(I\),
\[
    f\alpha_X(I)
    =
    f(I*1_X)
    =
    T_{\mathbb V}(f)(I)*1_Y
    =
    \alpha_YT_{\mathbb V}(f)(I).
\]
Therefore
\[
    f\alpha_X=\alpha_YT_{\mathbb V}(f),
\]
so \(f\) is a fuzzy-domain morphism.
\end{proof}

\subsection{Free fuzzy domains and limits}
\label{subsec:free-fuzzy-domains-limits}

\begin{theorem}[Free fuzzy domains]
\label{thm:free-fuzzy-domains}
The forgetful functor
\[
    U:\mathbf{FDom}_{\mathbb V}\to \mathbb V\text{-}\mathbf{Cat}
\]
has a left adjoint
\[
    F:\mathbb V\text{-}\mathbf{Cat}\to\mathbf{FDom}_{\mathbb V},
    \qquad
    F(X)=\operatorname{Idl}_{\mathbb V}(X),
\]
where \(F(X)\) has algebra structure
\[
    \mu_X:
    \operatorname{Idl}_{\mathbb V}(\operatorname{Idl}_{\mathbb V}(X))
    \to
    \operatorname{Idl}_{\mathbb V}(X).
\]
Thus, for every small \(\mathbb V\)-category \(X\) and every fuzzy domain \(A\),
there is a natural bijection
\[
    \mathbf{FDom}_{\mathbb V}(\operatorname{Idl}_{\mathbb V}(X),A)
    \cong
    \mathbb V\text{-}\mathbf{Cat}(X,UA).
\]
\end{theorem}

\begin{proof}
This is the free-algebra adjunction of the monad \(T_{\mathbb V}\).  Explicitly,
if
\[
    f:X\to UA
\]
is a \(\mathbb V\)-functor and
\[
    \alpha_A:\operatorname{Idl}_{\mathbb V}(A)\to A
\]
is the algebra map of \(A\), define
\[
    \overline f:
    \operatorname{Idl}_{\mathbb V}(X)\to A
\]
by
\[
    \overline f(I)
    =
    \alpha_A(T_{\mathbb V}f(I)).
\]
Equivalently,
\[
    \overline f(I)=I*f.
\]
By Theorem~\ref{thm:fuzzy-domains-ff-colimits}, \(\overline f\) preserves
finite-flat weighted colimits, so it is a fuzzy-domain morphism.  Moreover,
\[
    \overline f(\eta_Xx)
    =
    \alpha_A(\eta_A(fx))
    =
    fx.
\]
Thus
\[
    \overline f\eta_X=f.
\]

Conversely, if
\[
    g:\operatorname{Idl}_{\mathbb V}(X)\to A
\]
is a fuzzy-domain morphism, then \(g\) preserves finite-flat weighted colimits.
For every finite-flat weight \(I:X^{\op}\to\underline{\mathbb V}\), the object
\(I\in\operatorname{Idl}_{\mathbb V}(X)\) is the \(I\)-weighted colimit of the
Yoneda diagram:
\[
    I\cong I*\eta_X.
\]
Therefore
\[
    g(I)
    \cong
    g(I*\eta_X)
    \cong
    I*(g\eta_X).
\]
Hence
\[
    g=\overline{g\eta_X}
\]
under the strictification convention.  The two constructions are inverse and
natural.
\end{proof}

\begin{corollary}[Monadicity]
\label{cor:fdom-monadicity}
The forgetful functor
\[
    U:\mathbf{FDom}_{\mathbb V}\to \mathbb V\text{-}\mathbf{Cat}
\]
is monadic, and the induced monad is the fuzzy ideal monad
\[
    T_{\mathbb V}.
\]
Thus
\[
    \mathbf{FDom}_{\mathbb V}=(\mathbb V\text{-}\mathbf{Cat})^{T_{\mathbb V}}.
\]
\end{corollary}

\begin{proof}
By definition, \(\mathbf{FDom}_{\mathbb V}\) is the Eilenberg--Moore category of
strict algebras for \(T_{\mathbb V}\).  The adjunction of
Theorem~\ref{thm:free-fuzzy-domains} has underlying monad \(T_{\mathbb V}\), with
Yoneda unit and flattening multiplication.
\end{proof}

\begin{theorem}[Small conical limits]
\label{thm:fdom-limits}
The category \(\mathbf{FDom}_{\mathbb V}\) has all small conical limits whose
underlying limits exist in
\[
    \mathbb V\text{-}\mathbf{Cat}.
\]
These limits are created by the forgetful functor
\[
    U:\mathbf{FDom}_{\mathbb V}\to \mathbb V\text{-}\mathbf{Cat}.
\]
\end{theorem}

\begin{proof}
Eilenberg--Moore categories have all limits that exist in the base category, and
the forgetful functor creates them.

Explicitly, let
\[
    D:J\to\mathbf{FDom}_{\mathbb V}
\]
be a small diagram, and suppose that
\[
    L=\lim_{j\in J}UD_j
\]
exists in \(\mathbb V\text{-}\mathbf{Cat}\), with projections
\[
    p_j:L\to UD_j.
\]
There is a unique algebra structure
\[
    \alpha_L:T_{\mathbb V}L\to L
\]
such that
\[
    p_j\alpha_L=\alpha_jT_{\mathbb V}(p_j)
\]
for every \(j\).  With this algebra structure, \(L\) satisfies the universal
property of the limit in \(\mathbf{FDom}_{\mathbb V}\).
\end{proof}

\subsection{Tensor products and internal homs}
\label{subsec:tensor-products-internal-homs-fdom}

\begin{definition}[External product of weights]
\label{def:external-product-weights-fdom}
Let
\[
    W:K^{\op}\to\underline{\mathbb V},
    \qquad
    Q:J^{\op}\to\underline{\mathbb V}
\]
be weights.  Their external product is the weight
\[
    W\boxtimes Q:(K\tensor J)^{\op}\to\underline{\mathbb V}
\]
defined by
\[
    (W\boxtimes Q)(k,j)=W(k)\tensor Q(j).
\]
\end{definition}

\begin{proposition}[External products of finite-flat weights]
\label{prop:external-products-finite-flat-fdom}
If \(W\) and \(Q\) are finite-flat, then
\[
    W\boxtimes Q
\]
is finite-flat.
\end{proposition}

\begin{proof}
Let
\[
    E:K\tensor J\to\underline{\mathbb V}
\]
be an upper predicate.  Then
\[
\begin{aligned}
    \chi_{W\boxtimes Q}(E)
    &=
    \int^{(k,j)\in K\tensor J}
        W(k)\tensor Q(j)\tensor E(k,j)
    \\
    &\cong
    \int^{k\in K}
        W(k)\tensor
        \left(
            \int^{j\in J}Q(j)\tensor E(k,j)
        \right).
\end{aligned}
\]
For each \(k\), the inner expression is
\[
    \chi_Q(E(k,-)).
\]

Since \(E:K\tensor J\to\underline{\mathbb V}\) is a \(\mathbb V\)-functor, the
assignment
\[
    k\longmapsto E(k,-)
\]
defines a \(\mathbb V\)-functor
\[
    K\to[J,\underline{\mathbb V}].
\]
Composing with
\[
    \chi_Q:[J,\underline{\mathbb V}]\to\underline{\mathbb V}
\]
gives an upper predicate
\[
    k\longmapsto \chi_Q(E(k,-))
\]
on \(K\).  Since \(Q\) is finite-flat, this construction preserves the specified
finite limits in the \(J\)-variable.  Since \(W\) is finite-flat, applying
\(\chi_W\) preserves the specified finite limits in the remaining \(K\)-variable.
Hence \(\chi_{W\boxtimes Q}\) preserves the specified finite limits, so
\(W\boxtimes Q\) is finite-flat.
\end{proof}

\begin{proposition}[Finite-flat Fubini]
\label{prop:finite-flat-fubini-fdom}
Let
\[
    W:K^{\op}\to\underline{\mathbb V},
    \qquad
    Q:J^{\op}\to\underline{\mathbb V}
\]
be finite-flat.  For every diagram
\[
    E:K\tensor J\to C
\]
in a fuzzy domain \(C\), there are canonical isomorphisms
\[
    W*\bigl(k\mapsto Q*(j\mapsto E(k,j))\bigr)
    \cong
    (W\boxtimes Q)*E
    \cong
    Q*\bigl(j\mapsto W*(k\mapsto E(k,j))\bigr).
\]
\end{proposition}

\begin{proof}
These are the Fubini isomorphisms for enriched weighted colimits.  By
Proposition~\ref{prop:external-products-finite-flat-fdom}, the external product
\(W\boxtimes Q\) is finite-flat.  Since \(C\) is a fuzzy domain, all three
finite-flat weighted colimits appearing above exist.
\end{proof}

\begin{definition}[Internal hom fuzzy domain]
\label{def:internal-hom-fdom}
For fuzzy domains \(B\) and \(C\), define
\[
    [B,C]_{\mathrm{ff}}
\]
to be the full \(\mathbb V\)-subcategory of the enriched functor category
\([B,C]\) spanned by the finite-flat-colimit-preserving \(\mathbb V\)-functors
\[
    B\to C.
\]
\end{definition}

\begin{proposition}[Internal homs are fuzzy domains]
\label{prop:internal-homs-fdom}
For fuzzy domains \(B\) and \(C\), the \(\mathbb V\)-category
\[
    [B,C]_{\mathrm{ff}}
\]
is a fuzzy domain.  Its finite-flat weighted colimits are computed pointwise in
\(C\).
\end{proposition}

\begin{proof}
Let
\[
    W:K^{\op}\to\underline{\mathbb V}
\]
be finite-flat, and let
\[
    H:K\to[B,C]_{\mathrm{ff}}
\]
be a diagram.  Since \(C\) is a fuzzy domain, the pointwise formula
\[
    (W*H)(b)=W*(k\mapsto Hk(b))
\]
defines the \(W\)-weighted colimit of \(H\) in the ambient functor category
\([B,C]\).

It remains to show that \(W*H\) preserves finite-flat weighted colimits in
\(B\).  Let
\[
    Q:J^{\op}\to\underline{\mathbb V}
\]
be finite-flat, and let
\[
    D:J\to B
\]
be a diagram.  Then
\[
\begin{aligned}
    (W*H)(Q*D)
    &\cong
    W*\bigl(k\mapsto Hk(Q*D)\bigr)
    \\
    &\cong
    W*\bigl(k\mapsto Q*(j\mapsto Hk(Dj))\bigr)
    \\
    &\cong
    Q*\bigl(j\mapsto W*(k\mapsto Hk(Dj))\bigr)
    \\
    &\cong
    Q*((W*H)D).
\end{aligned}
\]
The second isomorphism uses that each \(Hk:B\to C\) preserves finite-flat
weighted colimits.  The third isomorphism is
Proposition~\ref{prop:finite-flat-fubini-fdom}.  Hence \(W*H\) preserves
finite-flat weighted colimits in \(B\).

Therefore \([B,C]_{\mathrm{ff}}\) is closed under finite-flat weighted colimits
inside \([B,C]\).  Since \([B,C]_{\mathrm{ff}}\) is full in \([B,C]\), the
ambient weighted-colimit universal property restricts to
\([B,C]_{\mathrm{ff}}\).  By
Theorem~\ref{thm:fuzzy-domains-ff-colimits},
\([B,C]_{\mathrm{ff}}\) is a fuzzy domain.
\end{proof}

\begin{definition}[Finite-flat bimorphism]
\label{def:finite-flat-bimorphism-fdom}
Let \(A,B,C\) be fuzzy domains.  A finite-flat bimorphism
\[
    M:A\tensor B\to C
\]
is a \(\mathbb V\)-functor such that:
\begin{enumerate}
    \item for each \(a\in A\), the partial functor
    \[
        M(a,-):B\to C
    \]
    preserves finite-flat weighted colimits;
    \item for each \(b\in B\), the partial functor
    \[
        M(-,b):A\to C
    \]
    preserves finite-flat weighted colimits.
\end{enumerate}
We write
\[
    \operatorname{Bimor}_{\mathrm{ff}}(A,B;C)
\]
for the set of finite-flat bimorphisms from \(A,B\) to \(C\).
\end{definition}

\begin{proposition}[Currying finite-flat bimorphisms]
\label{prop:currying-bimorphisms-fdom}
For fuzzy domains \(A,B,C\), there is a natural bijection
\[
    \operatorname{Bimor}_{\mathrm{ff}}(A,B;C)
    \cong
    \mathbf{FDom}_{\mathbb V}(A,[B,C]_{\mathrm{ff}}).
\]
\end{proposition}

\begin{proof}
Let
\[
    M:A\tensor B\to C
\]
be a finite-flat bimorphism.  By enriched currying, \(M\) corresponds to a
\(\mathbb V\)-functor
\[
    \widehat M:A\to[B,C],
    \qquad
    \widehat M(a)(b)=M(a,b).
\]
Since \(M\) preserves finite-flat weighted colimits in the \(B\)-variable, each
\[
    \widehat M(a):B\to C
\]
lies in \([B,C]_{\mathrm{ff}}\).  Thus \(\widehat M\) factors through
\[
    [B,C]_{\mathrm{ff}}\subseteq[B,C].
\]
Since \(M\) preserves finite-flat weighted colimits in the \(A\)-variable and
finite-flat colimits in \([B,C]_{\mathrm{ff}}\) are computed pointwise,
\[
    \widehat M:A\to[B,C]_{\mathrm{ff}}
\]
preserves finite-flat weighted colimits.  Hence \(\widehat M\) is a
fuzzy-domain morphism.

Conversely, let
\[
    H:A\to[B,C]_{\mathrm{ff}}
\]
be a fuzzy-domain morphism.  Uncurrying gives
\[
    \widetilde H:A\tensor B\to C,
    \qquad
    \widetilde H(a,b)=H(a)(b).
\]
For each \(a\), the partial functor
\[
    \widetilde H(a,-)=H(a)
\]
preserves finite-flat weighted colimits because \(H(a)\) is an object of
\([B,C]_{\mathrm{ff}}\).  For each \(b\), the partial functor
\[
    \widetilde H(-,b):A\to C
\]
preserves finite-flat weighted colimits because \(H\) does and colimits in
\([B,C]_{\mathrm{ff}}\) are computed pointwise.  Hence \(\widetilde H\) is a
finite-flat bimorphism.

The two constructions are inverse by the enriched tensor--hom adjunction.
\end{proof}

\begin{theorem}[Kelly tensor product for fuzzy domains]
\label{thm:kelly-tensor-product-fdom}
For fuzzy domains \(A\) and \(B\), there is a fuzzy domain
\[
    A\boxtimes_{\mathrm{ff}}B
\]
and a finite-flat bimorphism
\[
    \kappa_{A,B}:A\tensor B\to A\boxtimes_{\mathrm{ff}}B
\]
such that, for every fuzzy domain \(C\), composition with \(\kappa_{A,B}\)
induces a natural bijection
\[
    \mathbf{FDom}_{\mathbb V}(A\boxtimes_{\mathrm{ff}}B,C)
    \cong
    \operatorname{Bimor}_{\mathrm{ff}}(A,B;C).
\]
\end{theorem}

\begin{proof}
By Lemma~\ref{lem:representables-finite-flat-fdom},
Lemma~\ref{lem:finite-flat-direct-images-fdom}, and
Proposition~\ref{prop:finite-flat-saturation-fdom}, the assignment
\[
    X\longmapsto \operatorname{Idl}_{\mathbb V}(X)
\]
contains the representables, is stable under lower direct image along
\(\mathbb V\)-functors, and is closed under finite-flat weighted colimits in
presheaf categories.  Hence the finite-flat weights form a saturated doctrine.

Moreover, Proposition~\ref{prop:external-products-finite-flat-fdom} and
Proposition~\ref{prop:finite-flat-fubini-fdom} show that the finite-flat
doctrine is commutative: finite-flat weighted colimits commute with finite-flat
weighted colimits.

Kelly's tensor-product theorem for categories equipped with a saturated
commutative class of specified weighted colimits therefore applies to the
finite-flat doctrine.  The representing object for separately
finite-flat-colimit-preserving \(\mathbb V\)-functors
\[
    A\tensor B\to C
\]
is denoted
\[
    A\boxtimes_{\mathrm{ff}}B.
\]
Its universal map
\[
    \kappa_{A,B}:A\tensor B\to A\boxtimes_{\mathrm{ff}}B
\]
is a finite-flat bimorphism, and the displayed bijection is exactly the
representing universal property.
\end{proof}

\begin{corollary}[Tensor product of free fuzzy domains]
\label{cor:tensor-free-fdom}
For small \(\mathbb V\)-categories \(X\) and \(Y\), there is a canonical
equivalence
\[
    \operatorname{Idl}_{\mathbb V}(X)
    \boxtimes_{\mathrm{ff}}
    \operatorname{Idl}_{\mathbb V}(Y)
    \simeq
    \operatorname{Idl}_{\mathbb V}(X\tensor Y).
\]
\end{corollary}

\begin{proof}
For every fuzzy domain \(C\), there are natural bijections
\[
\begin{aligned}
\mathbf{FDom}_{\mathbb V}
\bigl(
    \operatorname{Idl}_{\mathbb V}(X)
    \boxtimes_{\mathrm{ff}}
    \operatorname{Idl}_{\mathbb V}(Y),
    C
\bigr)
&\cong
\operatorname{Bimor}_{\mathrm{ff}}
\bigl(
    \operatorname{Idl}_{\mathbb V}(X),
    \operatorname{Idl}_{\mathbb V}(Y);
    C
\bigr)
\\
&\cong
\mathbb V\text{-}\mathbf{Cat}(X\tensor Y,UC)
\\
&\cong
\mathbf{FDom}_{\mathbb V}
\bigl(
    \operatorname{Idl}_{\mathbb V}(X\tensor Y),
    C
\bigr).
\end{aligned}
\]
The middle bijection sends a bimorphism to its restriction along
\[
    \eta_X\tensor\eta_Y:
    X\tensor Y\to
    \operatorname{Idl}_{\mathbb V}(X)\tensor\operatorname{Idl}_{\mathbb V}(Y).
\]
Conversely, a \(\mathbb V\)-functor
\[
    m:X\tensor Y\to UC
\]
extends uniquely to the finite-flat bimorphism
\[
    \widetilde m:
    \operatorname{Idl}_{\mathbb V}(X)\tensor\operatorname{Idl}_{\mathbb V}(Y)\to C
\]
defined by
\[
    \widetilde m(I,J)
    =
    I*\bigl(x\mapsto J*(y\mapsto m(x,y))\bigr).
\]
By Proposition~\ref{prop:external-products-finite-flat-fdom} and
Proposition~\ref{prop:finite-flat-fubini-fdom}, this is equivalently
\[
    \widetilde m(I,J)\cong (I\boxtimes J)*m.
\]
Thus \(\widetilde m\) preserves finite-flat weighted colimits separately in
\(I\) and \(J\), and it is the unique such extension of \(m\).  By
representability, the displayed equivalence follows.
\end{proof}

\begin{corollary}[Tensor unit]
\label{cor:tensor-unit-fdom}
Let
\[
    \mathbf 1
\]
be the terminal small \(\mathbb V\)-category.  The tensor unit for fuzzy domains is
\[
    \operatorname{Idl}_{\mathbb V}(\mathbf 1).
\]
\end{corollary}

\begin{proof}
For every fuzzy domain \(C\),
\[
\begin{aligned}
    \mathbf{FDom}_{\mathbb V}
    \bigl(
        \operatorname{Idl}_{\mathbb V}(\mathbf 1)
        \boxtimes_{\mathrm{ff}}A,
        C
    \bigr)
    &\cong
    \operatorname{Bimor}_{\mathrm{ff}}
    \bigl(
        \operatorname{Idl}_{\mathbb V}(\mathbf 1),
        A;
        C
    \bigr)
    \\
    &\cong
    \mathbf{FDom}_{\mathbb V}(A,C).
\end{aligned}
\]
The second bijection uses that \(\operatorname{Idl}_{\mathbb V}(\mathbf 1)\) is
freely generated as a fuzzy domain by the unique object of \(\mathbf 1\).  Hence
\[
    \operatorname{Idl}_{\mathbb V}(\mathbf 1)\boxtimes_{\mathrm{ff}}A
    \simeq
    A.
\]
The right unit law is analogous.
\end{proof}

\begin{theorem}[Symmetric monoidal closure of fuzzy domains]
\label{thm:fdom-smc}
The category
\[
    \mathbf{FDom}_{\mathbb V}
\]
is symmetric monoidal closed.  Its tensor product is
\[
    A\boxtimes_{\mathrm{ff}}B,
\]
its tensor unit is
\[
    \operatorname{Idl}_{\mathbb V}(\mathbf 1),
\]
and its internal hom is
\[
    [B,C]_{\mathrm{ff}}.
\]
The closed structure is the natural bijection
\[
    \mathbf{FDom}_{\mathbb V}(A\boxtimes_{\mathrm{ff}}B,C)
    \cong
    \mathbf{FDom}_{\mathbb V}(A,[B,C]_{\mathrm{ff}}).
\]
\end{theorem}

\begin{proof}
By Theorem~\ref{thm:kelly-tensor-product-fdom},
\[
    \mathbf{FDom}_{\mathbb V}(A\boxtimes_{\mathrm{ff}}B,C)
    \cong
    \operatorname{Bimor}_{\mathrm{ff}}(A,B;C).
\]
By Proposition~\ref{prop:currying-bimorphisms-fdom},
\[
    \operatorname{Bimor}_{\mathrm{ff}}(A,B;C)
    \cong
    \mathbf{FDom}_{\mathbb V}(A,[B,C]_{\mathrm{ff}}).
\]
Composing these bijections gives the closed structure.

The symmetry of \(\boxtimes_{\mathrm{ff}}\) follows from the symmetry of
\[
    A\tensor B\cong B\tensor A
\]
and from the symmetric nature of the condition of preserving finite-flat
weighted colimits separately in the two variables.

Associativity follows by representability.  For every fuzzy domain \(D\), both
\[
    (A\boxtimes_{\mathrm{ff}}B)\boxtimes_{\mathrm{ff}}C
\]
and
\[
    A\boxtimes_{\mathrm{ff}}(B\boxtimes_{\mathrm{ff}}C)
\]
represent trilinear \(\mathbb V\)-functors
\[
    A\tensor B\tensor C\to D
\]
that preserve finite-flat weighted colimits separately in each variable.  Hence
they are canonically equivalent.

The tensor unit is \(\operatorname{Idl}_{\mathbb V}(\mathbf 1)\) by
Corollary~\ref{cor:tensor-unit-fdom}.  Therefore \(\mathbf{FDom}_{\mathbb V}\) is
symmetric monoidal closed.
\end{proof}

\begin{corollary}[Evaluation]
\label{cor:evaluation-fdom}
For fuzzy domains \(B\) and \(C\), there is an evaluation morphism
\[
    [B,C]_{\mathrm{ff}}\boxtimes_{\mathrm{ff}}B
    \longrightarrow
    C.
\]
On generators it is given by
\[
    (f,b)\longmapsto f(b).
\]
\end{corollary}

\begin{proof}
The ordinary enriched evaluation
\[
    [B,C]\tensor B\to C
\]
restricts to a finite-flat bimorphism
\[
    [B,C]_{\mathrm{ff}}\tensor B\to C.
\]
It is finite-flat-colimit-preserving in the \(B\)-variable because every object
of \([B,C]_{\mathrm{ff}}\) preserves finite-flat weighted colimits.  It is
finite-flat-colimit-preserving in the \([B,C]_{\mathrm{ff}}\)-variable because
finite-flat colimits there are computed pointwise in \(C\).  By the universal
property of \(\boxtimes_{\mathrm{ff}}\), this bimorphism induces the displayed
fuzzy-domain morphism.
\end{proof}

\subsection{Posetal and crisp specializations}
\label{subsec:posetal-crisp-fdom-category}

In the posetal specialization, a fuzzy domain over \(\Lambda\) is a fuzzy
preorder \(X\) equipped with a strict algebra
\[
    \alpha_X:\operatorname{Idl}_{\Lambda}(X)\to X.
\]
The map \(\alpha_X\) sends a fuzzy ideal \(I\) to its fuzzy ideal-join.

If
\[
    D:K\to X
\]
is a \(\Lambda\)-functor and
\[
    W:K^{\op}\to\Lambda
\]
is a fuzzy ideal of \(K\), that is, a finite-flat lower fuzzy predicate on
\(K\), then
\[
    W*D=\alpha_X(D_!W),
\]
where
\[
    (D_!W)(x)
    =
    \bigjoin_{k\in K}W(k)\tensor X(x,Dk).
\]
A morphism of fuzzy domains
\[
    f:X\to Y
\]
is a \(\Lambda\)-functor satisfying
\[
    f\alpha_X=\alpha_YT_{\Lambda}(f).
\]
Equivalently, for every fuzzy ideal \(I\) of \(X\),
\[
    f(\alpha_X(I))
    =
    \alpha_Y(T_{\Lambda}f(I)).
\]
Here the direct image \(T_{\Lambda}f(I)\) is given by
\[
    (T_{\Lambda}f)(I)(y)
    =
    \bigjoin_{x\in X}I(x)\tensor Y(y,fx).
\]

For fuzzy domains \(B\) and \(C\), the internal hom
\[
    [B,C]_{\mathrm{ff}}
\]
has as objects the fuzzy-domain morphisms
\[
    B\to C.
\]
Its hom-values are inherited from the enriched functor preorder:
\[
    [B,C]_{\mathrm{ff}}(f,g)
    =
    \bigmeet_{b\in B}C(fb,gb).
\]
Finite-flat joins in \([B,C]_{\mathrm{ff}}\) are computed pointwise in \(C\).

For \(\Lambda=2\), finite-flat weights are ordinary ideals.  On separated
preorders, fuzzy domains are precisely dcpos, and fuzzy-domain morphisms are
precisely Scott-continuous maps.  Hence
\[
    \mathbf{FDom}_{2}^{\mathrm{sep}}
    \simeq
    \mathbf{DCPO},
\]
where \(\mathbf{DCPO}\) denotes the category of dcpos and Scott-continuous maps,
small in the chosen universe.

In this crisp separated case, finite-flat bimorphisms are monotone maps
\[
    A\times B\to C
\]
between dcpos that are Scott-continuous in each variable separately.  Such maps
are Scott-continuous for the product dcpo.  Indeed, let
\[
    D\subseteq A\times B
\]
be nonempty and directed.  Then \(\pi_A D\) and \(\pi_BD\) are nonempty directed
subsets, and
\[
    \bigjoin D
    =
    \left(
        \bigjoin\pi_A D,
        \bigjoin\pi_BD
    \right).
\]
If
\[
    f:A\times B\to C
\]
is separately Scott-continuous, then
\[
\begin{aligned}
    f\left(\bigjoin D\right)
    &=
    f\left(
        \bigjoin\pi_A D,
        \bigjoin\pi_B D
    \right)
    \\
    &=
    \bigjoin_{a\in\pi_A D}
        f\left(a,\bigjoin\pi_BD\right)
    \\
    &=
    \bigjoin_{a\in\pi_A D}
        \bigjoin_{b\in\pi_BD}
            f(a,b).
\end{aligned}
\]
Since \(D\) is directed, every pair
\[
    (a,b)\in\pi_A D\times\pi_BD
\]
is bounded above by some element of \(D\).  By monotonicity,
\[
    \bigjoin_{a\in\pi_A D,\;b\in\pi_BD}f(a,b)
    =
    \bigjoin_{(a,b)\in D}f(a,b).
\]
Hence
\[
    f\left(\bigjoin D\right)
    =
    \bigjoin_{(a,b)\in D}f(a,b).
\]
Thus \(f\) is Scott-continuous on the product dcpo.

Consequently, in the separated crisp specialization, the Kelly tensor product
\(\boxtimes_{\mathrm{ff}}\) agrees with the cartesian product of dcpos, and the
internal hom agrees with the dcpo of Scott-continuous maps ordered pointwise.
Hence, under the equivalence
\[
    \mathbf{FDom}_{2}^{\mathrm{sep}}
    \simeq
    \mathbf{DCPO},
\]
the symmetric monoidal closed structure constructed above specializes to the
standard cartesian closed structure on the category of dcpos and Scott-continuous
maps.

\section{The Kleisli category and fuzzy approximable maps}
\label{sec:kleisli-fuzzy-ideal-monad}

\subsection{The enriched Kleisli category}
\label{subsec:kleisli-enriched}

\begin{definition}[Fuzzy Kleisli morphism]
\label{def:fuzzy-kleisli-morphism-enriched}
A \textbf{fuzzy Kleisli morphism}
\[
    X\rightsquigarrow Y
\]
is a \(\mathbb V\)-functor
\[
    X\to T_{\mathbb V}Y
    =
    \operatorname{Idl}_{\mathbb V}(Y).
\]
Equivalently, it is a profunctor-like object
\[
    R:X\otimes Y^{\op}\to\underline{\mathbb V}
\]
whose fibre
\[
    R(x,-):Y^{\op}\to\underline{\mathbb V}
\]
is a fuzzy ideal of \(Y\), functorially in \(x\).
\end{definition}

\begin{proposition}[Kleisli composition]
\label{prop:kleisli-composition-enriched}
If
\[
    R:X\rightsquigarrow Y
    \qquad\text{and}\qquad
    S:Y\rightsquigarrow Z
\]
are fuzzy Kleisli morphisms, their composite is
\[
    (S\circ R)(x,z)
    =
    \int^{y\in Y}R(x,y)\tensor S(y,z).
\]
The identity Kleisli morphism on \(X\) is
\[
    1_X(x,y)=X(y,x).
\]
\end{proposition}

\begin{proof}
This is Kleisli composition for the monad \(T_{\mathbb V}\), written under the
finite-flat-weight representation.  The identity is the unit
\(\eta_X(x)=X(-,x)\).
\end{proof}

\begin{proposition}[Extension to free fuzzy domains]
\label{prop:kleisli-extension-free-domain}
A fuzzy Kleisli morphism
\[
    R:X\rightsquigarrow Y
\]
extends to a morphism between free fuzzy domains
\[
    \overline R:
    \operatorname{Idl}_{\mathbb V}(X)
    \to
    \operatorname{Idl}_{\mathbb V}(Y)
\]
given by
\[
    \overline R(I)(y)
    =
    \int^{x\in X} I(x)\tensor R(x,y).
\]
\end{proposition}

\begin{proof}
This is the Kleisli extension
\[
    T_{\mathbb V}X
    \xrightarrow{T_{\mathbb V}R}
    T_{\mathbb V}T_{\mathbb V}Y
    \xrightarrow{\mu_Y}
    T_{\mathbb V}Y.
\]
The displayed formula is the resulting convolution.
\end{proof}

\subsection{The posetal specialization}
\label{subsec:kleisli-posetal}

\begin{definition}[Fuzzy approximable map]
\label{def:fuzzy-approximable-map}
For fuzzy preorders \(X,Y\), a \textbf{fuzzy approximable map}
\[
    R:X\rightsquigarrow Y
\]
is a map
\[
    R:X\times Y\to L
\]
such that each fibre
\[
    R(x,-):Y^{\op}\to\Lambda
\]
is a fuzzy ideal of \(Y\), and the assignment
\[
    x\longmapsto R(x,-)
\]
is a \(\Lambda\)-functor
\[
    X\to\operatorname{Idl}_\Lambda(Y).
\]
\end{definition}

\begin{proposition}[Posetal Kleisli formulas]
\label{prop:posetal-kleisli-formulas}
For fuzzy approximable maps
\[
    R:X\rightsquigarrow Y,
    \qquad
    S:Y\rightsquigarrow Z,
\]
composition is
\[
    (S\circ R)(x,z)
    =
    \bigjoin_{y\in Y}R(x,y)\tensor S(y,z).
\]
The identity is
\[
    1_X(x,y)=X(y,x).
\]
The extension of \(R\) to free fuzzy domains is
\[
    \overline R(I)(y)
    =
    \bigjoin_{x\in X}I(x)\tensor R(x,y).
\]
\end{proposition}

\begin{proof}
These are the posetal specializations of
Propositions~\ref{prop:kleisli-composition-enriched} and
\ref{prop:kleisli-extension-free-domain}.
\end{proof}

\begin{example}[The crisp Kleisli category]
\label{ex:crisp-kleisli}
For \(\Lambda=2\), a Kleisli morphism
\[
    X\rightsquigarrow Y
\]
is a monotone map
\[
    X\to\operatorname{Idl}(Y).
\]
Equivalently, it is an ideal-valued relation
\[
    R\subseteq X\times Y
\]
such that each fibre
\[
    R[x]=\{y\in Y\mid xRy\}
\]
is an ideal of \(Y\), and
\[
    x\leq x'
    \quad\Longrightarrow\quad
    R[x]\subseteq R[x'].
\]
Composition is
\[
    (S\circ R)[x]
    =
    \bigcup_{y\in R[x]}S[y],
\]
and the Kleisli extension is
\[
    \overline r(I)=\bigcup_{x\in I}r(x).
\]
These are the classical approximable maps between bases.
\end{example}

\begin{example}[Lawvere Kleisli morphisms]
\label{ex:lawvere-kleisli}
For the Lawvere quantale, Kleisli morphisms are flat-weight-valued metric
relations.  Composition is infimal convolution:
\[
    (S\circ R)(x,z)
    =
    \inf_{y\in Y}\bigl(R(x,y)+S(y,z)\bigr).
\]
\end{example}

\section{Fuzzy algebraic and fuzzy continuous domains}
\label{sec:fuzzy-algebraic-continuous-domains}

\subsection{Fuzzy algebraic domains}
\label{subsec:fuzzy-algebraic-domains}

\begin{definition}[Presented fuzzy algebraic domain]
\label{def:presented-fuzzy-algebraic-domain}
A \textbf{presented fuzzy algebraic domain} is a free fuzzy domain
\[
    \operatorname{Idl}_{\mathbb V}(B)
\]
for a small \(\mathbb V\)-category \(B\).  The objects of \(B\) are regarded as
basic elements, and the principal ideals
\[
    \eta_B(b)=B(-,b)
\]
are their basic compact approximants.
\end{definition}

\begin{proposition}[Algebra structure on a free fuzzy domain]
\label{prop:free-fuzzy-domain-algebra-structure}
The algebra structure on
\[
    \operatorname{Idl}_{\mathbb V}(B)
\]
is the monad multiplication
\[
    \mu_B:
    \operatorname{Idl}_{\mathbb V}(\operatorname{Idl}_{\mathbb V}(B))
    \to
    \operatorname{Idl}_{\mathbb V}(B).
\]
In the posetal case,
\[
    \mu_B(\mathcal I)(b)
    =
    \bigjoin_{I\in\operatorname{Idl}_\Lambda(B)}
    \mathcal I(I)\tensor I(b).
\]
\end{proposition}

\begin{proof}
This is the free-algebra construction for a monad, with multiplication given by
flattening.
\end{proof}

\begin{example}[Crisp algebraic domains]
\label{ex:crisp-algebraic-domains}
For \(\Lambda=2\),
\[
    \operatorname{Idl}(B)
\]
is an algebraic dcpo with compact elements the principal ideals
\[
    \downarrow b.
\]
Every ideal is the directed supremum of its principal subideals:
\[
    I=\bigjoin_{b\in I}\downarrow b.
\]
Thus free algebras for the crisp ideal monad are algebraic dcpos equipped with a
chosen basis of compact generators.
\end{example}

\subsection{Fuzzy continuous domains}
\label{subsec:fuzzy-continuous-domains}

\begin{definition}[Fuzzy continuous domain]
\label{def:fuzzy-continuous-domain-enriched}
Let \(X\) be a fuzzy domain with algebra structure
\[
    \alpha:\operatorname{Idl}_{\mathbb V}(X)\to X.
\]
We say that \(X\) is a \textbf{fuzzy continuous domain} if \(\alpha\) admits an
enriched left adjoint
\[
    \beta:X\to\operatorname{Idl}_{\mathbb V}(X),
    \qquad
    \beta\dashv\alpha.
\]
Thus
\[
    \operatorname{Idl}_{\mathbb V}(X)(\beta x,I)
    \cong
    X(x,\alpha I).
\]
The object \(\beta(x)\) is the fuzzy ideal of approximants to \(x\).
\end{definition}

\begin{definition}[Fuzzy way-below degree]
\label{def:fuzzy-way-below-degree}
In the posetal specialization, the associated \textbf{fuzzy way-below degree} is
\[
    \operatorname{wb}_X(y,x)
    =
    \beta(x)(y).
\]
\end{definition}

\begin{proposition}[Adjunction and approximation]
\label{prop:fuzzy-continuity-adjunction}
Let \(X\) be a separated fuzzy domain over \(\Lambda\), with algebra map
\[
    \alpha:\operatorname{Idl}_\Lambda(X)\to X.
\]
If
\[
    \beta:X\to\operatorname{Idl}_\Lambda(X)
\]
satisfies
\[
    \beta\dashv\alpha,
\]
then
\[
    \alpha\beta=1_X
\]
in the induced crisp order.  Hence every element is recovered as the fuzzy join of
its approximants.
\end{proposition}

\begin{proof}
The unit of the adjunction gives
\[
    1_X\leq \alpha\beta.
\]
The counit gives
\[
    \beta\alpha\leq 1_{\operatorname{Idl}_\Lambda(X)}.
\]
Apply the counit inequality to the principal ideal \(\eta_X(x)\).  Since
\[
    \alpha\eta_X(x)=x,
\]
we obtain
\[
    \beta(x)\leq\eta_X(x).
\]
Applying \(\alpha\) gives
\[
    \alpha\beta(x)\leq\alpha\eta_X(x)=x.
\]
Together with \(x\leq\alpha\beta(x)\), this gives equality in the induced crisp
order.  Since \(X\) is separated, this is equality of objects.
\end{proof}

\begin{example}[Crisp continuous domains]
\label{ex:crisp-continuous-domains}
Let \(D\) be a dcpo with ideal-supremum map
\[
    \bigjoin:\operatorname{Idl}(D)\to D.
\]
Then \(D\) is continuous precisely when this map has a left adjoint
\[
    \beta:D\to\operatorname{Idl}(D).
\]
The left adjoint is
\[
    \beta(x)=\{y\in D\mid y\ll x\},
\]
and the adjunction says
\[
    \beta(x)\subseteq I
    \quad\Longleftrightarrow\quad
    x\leq \bigjoin I.
\]
\end{example}

\begin{example}[Split Kleisli idempotents]
\label{ex:split-kleisli-idempotents-continuous}
In the crisp case, the map
\[
    \beta:D\to\operatorname{Idl}(D),
    \qquad
    \beta(x)=\{y\mid y\ll x\},
\]
is a Kleisli endomorphism
\[
    D\rightsquigarrow D.
\]
Interpolation of the way-below relation gives
\[
    \beta\circ_{\operatorname{Kl}}\beta=\beta.
\]
Thus continuous dcpos appear as splittings of Kleisli idempotents, equivalently as
retracts of algebraic dcpos.
\end{example}

\begin{example}[Lawvere-valued approximation]
\label{ex:lawvere-fuzzy-continuous}
For the Lawvere quantale, a fuzzy continuity structure assigns to each point \(x\)
a finite-flat weight
\[
    \beta(x):X^{\op}\to[0,\infty].
\]
The value \(\beta(x)(y)\) is the cost by which \(y\) approximates \(x\).  The
adjunction
\[
    \beta\dashv\alpha
\]
says that approximation into \(x\) is controlled exactly by comparison with
finite-flat-weight colimits.
\end{example}

\section{Fuzzy Scott opens and fuzzy Scott locales}
\label{sec:fuzzy-scott-opens}

\subsection{The enriched construction}
\label{subsec:fuzzy-scott-opens-enriched}

\begin{definition}[Fuzzy Scott open]
\label{def:fuzzy-scott-open-enriched}
Let \(X\) be a weak fuzzy domain with structure map
\[
    \alpha:T_{\mathbb V}X\to X.
\]
A \textbf{fuzzy Scott open} on \(X\) is an upper fuzzy predicate
\[
    U:X\to\underline{\mathbb V}
\]
such that, for every fuzzy ideal \(I\),
\[
    U(\alpha(I))
    \cong
    \int^{x\in X} I(x)\tensor U(x).
\]
\end{definition}

\begin{definition}[Enriched fuzzy Scott frame]
\label{def:enriched-fuzzy-scott-frame}
When the required equalizer exists in categorical fuzzy frames, define
\[
    \operatorname{Scott}_{\mathbb V}(X)
\]
to be the equalizer
\[
\begin{tikzcd}[column sep=large]
\operatorname{Scott}_{\mathbb V}(X)
    \arrow[r, hook]
&
P_{\mathbb V}(X^{\op})
    \arrow[r, shift left=.8ex, "{U\mapsto U\alpha}"]
    \arrow[r, shift right=.8ex, "{U\mapsto (I\mapsto \int^x I(x)\tensor U(x))}"']
&
P_{\mathbb V}(\operatorname{Idl}_{\mathbb V}(X)^{\op}).
\end{tikzcd}
\]
The second map is a fuzzy-frame morphism because, for each finite-flat weight \(I\),
the functional
\[
    U\longmapsto I*U
\]
is a character, and functoriality in \(I\) gives an upper predicate on
\(\operatorname{Idl}_{\mathbb V}(X)\).
\end{definition}

\begin{definition}[Fuzzy Scott locale]
\label{def:fuzzy-scott-locale-enriched}
The \textbf{fuzzy Scott locale} of \(X\) is the fuzzy locale
\[
    \operatorname{ScottLoc}_{\mathbb V}(X)
\]
whose frame of opens is
\[
    \operatorname{Opens}_{\mathbb V}(\operatorname{ScottLoc}_{\mathbb V}(X))
    =
    \operatorname{Scott}_{\mathbb V}(X).
\]
The inclusion
\[
    \operatorname{Scott}_{\mathbb V}(X)
    \hookrightarrow
    P_{\mathbb V}(X^{\op})
\]
defines a fuzzy-locale map
\[
    \mathcal P_{\mathbb V}^{\uparrow}X
    \longrightarrow
    \operatorname{ScottLoc}_{\mathbb V}(X).
\]
\end{definition}

\subsection{The posetal specialization}
\label{subsec:fuzzy-scott-opens-posetal}

\begin{definition}[Fuzzy Scott open over \(\Lambda\)]
\label{def:fuzzy-scott-open-posetal}
Let \(X\) be a weak fuzzy domain over \(\Lambda\), with structure map
\[
    \alpha:\operatorname{Idl}_\Lambda(X)\to X.
\]
A \textbf{fuzzy Scott open} is an upper fuzzy predicate
\[
    U:X\to L
\]
satisfying
\[
    U(\alpha(I))
    =
    \bigjoin_{x\in X} I(x)\tensor U(x)
\]
for every fuzzy ideal \(I\).
\end{definition}

\begin{definition}[The fuzzy Scott frame]
\label{def:fuzzy-scott-frame-posetal}
The \textbf{fuzzy Scott frame} of \(X\) is the full fuzzy subpreorder
\[
    \operatorname{Scott}_\Lambda(X)
    \hookrightarrow
    P_\Lambda(X^{\op})
\]
spanned by the fuzzy Scott opens.  Equivalently, it is the equalizer
\[
\begin{tikzcd}[column sep=large]
\operatorname{Scott}_\Lambda(X)
    \arrow[r, hook]
&
P_\Lambda(X^{\op})
    \arrow[r, shift left=.8ex, "{U\mapsto U\alpha}"]
    \arrow[r, shift right=.8ex, "{\Theta}"']
&
P_\Lambda(\operatorname{Idl}_\Lambda(X)^{\op}),
\end{tikzcd}
\]
where
\[
    \Theta(U)(I)
    =
    \bigjoin_{x\in X} I(x)\tensor U(x).
\]
\end{definition}

\begin{proposition}[Fuzzy Scott opens form a fuzzy frame]
\label{prop:fuzzy-scott-opens-frame}
Let \(X\) be a weak fuzzy domain over \(\Lambda\).  Then
\[
    \operatorname{Scott}_\Lambda(X)
\]
is a fuzzy frame, and the inclusion
\[
    \operatorname{Scott}_\Lambda(X)\hookrightarrow P_\Lambda(X^{\op})
\]
is a fuzzy-frame homomorphism.
\end{proposition}

\begin{proof}
First, \(\Theta(U)\) is an upper fuzzy predicate on
\(\operatorname{Idl}_\Lambda(X)\).  If
\[
    r\leq \operatorname{Idl}_\Lambda(X)(I,J),
\]
then \(r\tensor I(x)\leq J(x)\) for every \(x\), and hence
\[
    r\tensor \Theta(U)(I)
    =
    r\tensor \bigjoin_x I(x)\tensor U(x)
    \leq
    \bigjoin_x J(x)\tensor U(x)
    =
    \Theta(U)(J).
\]
Thus
\[
    \Theta(U)\in
    P_\Lambda(\operatorname{Idl}_\Lambda(X)^{\op}).
\]

For each fuzzy ideal \(I\), the map
\[
    U\longmapsto \Theta(U)(I)
\]
is the character \(\chi_I\), so it preserves arbitrary joins, finite meets, and
scalars.  Since fuzzy-frame structure in
\[
    P_\Lambda(\operatorname{Idl}_\Lambda(X)^{\op})
\]
is computed pointwise in \(I\), the map
\[
    \Theta:
    P_\Lambda(X^{\op})
    \to
    P_\Lambda(\operatorname{Idl}_\Lambda(X)^{\op})
\]
preserves arbitrary joins, finite meets, and scalars.  The map
\(U\mapsto U\alpha\) is precomposition by \(\alpha\), so it preserves the same
structure.  Therefore their equalizer is closed under arbitrary joins, finite
meets, and scalars, and the inclusion is a fuzzy-frame homomorphism.
\end{proof}

\begin{definition}[Fuzzy Scott locale over \(\Lambda\)]
\label{def:fuzzy-scott-locale-posetal}
For a weak fuzzy domain \(X\), the \textbf{fuzzy Scott locale} of \(X\) is the
fuzzy locale
\[
    \operatorname{ScottLoc}_\Lambda(X)
\]
whose frame of opens is
\[
    \Opens_\Lambda(\operatorname{ScottLoc}_\Lambda(X))
    =
    \operatorname{Scott}_\Lambda(X).
\]
The inclusion
\[
    \operatorname{Scott}_\Lambda(X)\hookrightarrow P_\Lambda(X^{\op})
\]
gives a fuzzy-locale map
\[
    \mathcal P_\Lambda^\uparrow X
    \longrightarrow
    \operatorname{ScottLoc}_\Lambda(X).
\]
\end{definition}

\begin{theorem}[Scott opens of a free fuzzy domain]
\label{thm:scott-opens-free-fuzzy-domain}
For every fuzzy preorder \(B\), there is a canonical isomorphism of fuzzy frames
\[
    P_\Lambda(B^{\op})
    \cong
    \operatorname{Scott}_\Lambda(\operatorname{Idl}_\Lambda(B)).
\]
An upper fuzzy predicate
\[
    U:B\to L
\]
corresponds to the fuzzy Scott open
\[
    \widehat U:\operatorname{Idl}_\Lambda(B)\to L
\]
defined by
\[
    \widehat U(I)
    =
    \bigjoin_{b\in B}I(b)\tensor U(b).
\]
\end{theorem}

\begin{proof}
Let \(U\in P_\Lambda(B^{\op})\).  Define
\[
    \widehat U(I)
    =
    \bigjoin_b I(b)\tensor U(b).
\]
The calculation in Proposition~\ref{prop:fuzzy-scott-opens-frame} shows that
\(\widehat U\) is an upper fuzzy predicate on
\(\operatorname{Idl}_\Lambda(B)\).

For
\[
    \mathcal I\in
    \operatorname{Idl}_\Lambda(\operatorname{Idl}_\Lambda(B)),
\]
we compute
\[
\begin{aligned}
    \widehat U(\mu_B(\mathcal I))
    &=
    \bigjoin_b \mu_B(\mathcal I)(b)\tensor U(b)\\
    &=
    \bigjoin_b
    \left(
        \bigjoin_I \mathcal I(I)\tensor I(b)
    \right)\tensor U(b)\\
    &=
    \bigjoin_I \mathcal I(I)\tensor
    \left(
        \bigjoin_b I(b)\tensor U(b)
    \right)\\
    &=
    \bigjoin_I \mathcal I(I)\tensor \widehat U(I).
\end{aligned}
\]
Thus \(\widehat U\) is Scott open.

Conversely, let
\[
    W\in\operatorname{Scott}_\Lambda(\operatorname{Idl}_\Lambda(B)).
\]
Restrict along the unit and define
\[
    U_W(b)=W(\eta_B b).
\]
Then \(U_W\) is an upper fuzzy predicate on \(B\).  For \(U\), Yoneda gives
\[
    \widehat U(\eta_B b)
    =
    \bigjoin_c B(c,b)\tensor U(c)
    =
    U(b),
\]
so \(U_{\widehat U}=U\).

It remains to show that \(\widehat{U_W}=W\).  For
\(I\in\operatorname{Idl}_\Lambda(B)\), the upper-predicate law for \(W\) gives
\[
    \operatorname{Idl}_\Lambda(B)(\eta_B b,I)
    \tensor W(\eta_B b)
    \leq
    W(I).
\]
By Yoneda,
\[
    \operatorname{Idl}_\Lambda(B)(\eta_B b,I)=I(b).
\]
Hence
\[
    \widehat{U_W}(I)
    =
    \bigjoin_b I(b)\tensor W(\eta_B b)
    \leq
    W(I).
\]

For the converse, use
\[
    \mu_B T_\Lambda(\eta_B)=1_{\operatorname{Idl}_\Lambda(B)}.
\]
Since \(W\) is Scott open,
\[
\begin{aligned}
    W(I)
    &=
    W(\mu_B T_\Lambda(\eta_B)(I))\\
    &=
    \bigjoin_J T_\Lambda(\eta_B)(I)(J)\tensor W(J).
\end{aligned}
\]
Moreover,
\[
    T_\Lambda(\eta_B)(I)(J)
    =
    \bigjoin_b I(b)\tensor
    \operatorname{Idl}_\Lambda(B)(J,\eta_B b).
\]
Using the upper-predicate law for \(W\),
\[
    \operatorname{Idl}_\Lambda(B)(J,\eta_B b)\tensor W(J)
    \leq
    W(\eta_B b).
\]
Therefore
\[
\begin{aligned}
    W(I)
    &\leq
    \bigjoin_J \bigjoin_b
    I(b)\tensor
    \operatorname{Idl}_\Lambda(B)(J,\eta_B b)\tensor W(J)\\
    &\leq
    \bigjoin_b I(b)\tensor W(\eta_B b)\\
    &=
    \widehat{U_W}(I).
\end{aligned}
\]
Thus \(W=\widehat{U_W}\).  The two constructions are inverse and preserve the
fuzzy-frame structure because all operations are computed pointwise and each
evaluation functional is a character.
\end{proof}

\begin{example}[The crisp Scott condition]
\label{ex:crisp-scott-opens}
For \(\Lambda=2\), if \(X\) is a dcpo with algebra
\[
    \alpha:\operatorname{Idl}(X)\to X,
    \qquad
    I\mapsto\bigjoin I,
\]
then a fuzzy Scott open is an upper set \(U\subseteq X\) satisfying
\[
    \bigjoin I\in U
    \quad\Longleftrightarrow\quad
    I\cap U\neq\varnothing
\]
for every ideal \(I\).  This is the ordinary Scott-open condition expressed using
ideals.
\end{example}

\begin{example}[Upper sets as Scott opens of ideal completions]
\label{ex:upper-sets-as-scott-opens-ideal-completion}
For a preorder \(B\),
\[
    \operatorname{Scott}(\operatorname{Idl}(B))
    \cong
    \operatorname{Up}(B).
\]
The correspondence sends an upper set \(U\subseteq B\) to
\[
    \widehat U
    =
    \{I\in\operatorname{Idl}(B)\mid I\cap U\neq\varnothing\}.
\]
Thus the upper Alexandrov frame of \(B\) is the Scott frame of the free dcpo
generated by \(B\).
\end{example}

\begin{example}[Standard fuzzy Scott formulas]
\label{ex:standard-fuzzy-scott-opens}
For G\"odel truth values, the fuzzy Scott condition is
\[
    U(\alpha(I))
    =
    \bigjoin_x \min(I(x),U(x)).
\]
For {\L}ukasiewicz truth values, it is
\[
    U(\alpha(I))
    =
    \bigjoin_x \max(0,I(x)+U(x)-1).
\]
For product truth values, it is
\[
    U(\alpha(I))
    =
    \bigjoin_x I(x)U(x).
\]
For the Lawvere quantale, it becomes
\[
    U(\alpha(I))
    =
    \inf_x\bigl(I(x)+U(x)\bigr).
\]
\end{example}

\section{The Alexandrov character comonad}
\label{sec:alexandrov-character-comonad}

\subsection{The enriched comonad}
\label{subsec:alexandrov-character-comonad-enriched}

\begin{definition}[Alexandrov character comonad]
\label{def:alexandrov-character-comonad-enriched}
The \textbf{Alexandrov character comonad} is the comonad
\[
    G_{\mathbb V}
    =
    \mathcal P_{\mathbb V}^{\uparrow}\operatorname{Char}_{\mathbb V}
\]
on fuzzy locales induced by
\[
    \mathcal P_{\mathbb V}^{\uparrow}
    \dashv
    \operatorname{Char}_{\mathbb V}.
\]
For a fuzzy locale \(\mathcal X\),
\[
    G_{\mathbb V}\mathcal X
    =
    \mathcal P_{\mathbb V}^{\uparrow}
    \bigl(
        \operatorname{Char}_{\mathbb V}(\mathcal X)
    \bigr),
\]
so
\[
    \operatorname{Opens}_{\mathbb V}(G_{\mathbb V}\mathcal X)
    =
    P_{\mathbb V}
    \bigl(
        \operatorname{Char}_{\mathbb V}(\mathcal X)^{\op}
    \bigr).
\]
\end{definition}

\begin{proposition}[Counit of the character comonad]
\label{prop:alexandrov-character-comonad-counit}
The counit
\[
    G_{\mathbb V}\mathcal X\to\mathcal X
\]
is frame-dual to the evaluation fuzzy-frame morphism
\[
    \operatorname{Opens}_{\mathbb V}(\mathcal X)
    \longrightarrow
    P_{\mathbb V}
    \bigl(
        \operatorname{Char}_{\mathbb V}(\mathcal X)^{\op}
    \bigr),
\]
given by
\[
    a\longmapsto \widehat a,
    \qquad
    \widehat a(\chi)=\chi(a).
\]
\end{proposition}

\begin{proof}
This is the counit of the adjunction, evaluated at \(\mathcal X\).  On frames it is
the transpose of the identity functor on
\(\operatorname{Char}_{\mathbb V}(\mathcal X)\), namely evaluation.
\end{proof}

\begin{definition}[Coalgebra for the Alexandrov character comonad]
\label{def:alexandrov-character-coalgebra}
A \textbf{\(G_{\mathbb V}\)-coalgebra} is a fuzzy locale \(\mathcal X\) equipped
with a fuzzy-locale map
\[
    \mathcal X\to G_{\mathbb V}\mathcal X
\]
satisfying the usual comonad coalgebra laws.  Frame-dually, this says that
\[
    \operatorname{Opens}_{\mathbb V}(\mathcal X)
\]
is a comonad-coherent retract of the upper fuzzy predicate frame on its character
object.
\end{definition}

\subsection{The posetal specialization}
\label{subsec:alexandrov-character-comonad-posetal}

\begin{proposition}[The posetal character comonad]
\label{prop:posetal-character-comonad}
For a fuzzy locale \(\mathcal X\),
\[
    \Opens_\Lambda(G_\Lambda\mathcal X)
    =
    P_\Lambda
    \bigl(
        \operatorname{Char}_\Lambda(\Opens_\Lambda(\mathcal X))^{\op}
    \bigr).
\]
The counit is frame-dual to the evaluation map
\[
    a\longmapsto \widehat a,
    \qquad
    \widehat a(\chi)=\chi(a).
\]
A \(G_\Lambda\)-coalgebra is frame-dually a fuzzy-frame homomorphism
\[
    P_\Lambda
    \bigl(
        \operatorname{Char}_\Lambda(\Opens_\Lambda(\mathcal X))^{\op}
    \bigr)
    \to
    \Opens_\Lambda(\mathcal X)
\]
compatible with the counit and comultiplication.
\end{proposition}

\begin{proof}
This is the posetal specialization of the enriched comonad and its counit.
\end{proof}

\begin{example}[The crisp case]
\label{ex:crisp-character-comonad}
For \(\Lambda=2\), the character object of a locale is its preorder of points,
ordered by specialization according to the convention
\[
    p\leq q
    \quad\Longleftrightarrow\quad
    \forall a,\;p(a)\leq q(a).
\]
Thus
\[
    G_2\mathcal X
\]
is the upper Alexandrov locale of the specialization preorder of points of
\(\mathcal X\).  The counit is frame-dual to
\[
    \Opens(\mathcal X)
    \to
    \operatorname{Up}(\operatorname{Pt}(\mathcal X)),
    \qquad
    a\longmapsto \{p\mid p(a)=\top\}.
\]
\end{example}

\begin{example}[Presheaf fuzzy locales]
\label{ex:presheaf-fuzzy-locale-character-comonad}
If
\[
    \mathcal X=\mathcal P_\Lambda^{\uparrow}X,
\]
then
\[
    \operatorname{Char}_\Lambda(\Opens_\Lambda(\mathcal X))
    =
    \operatorname{Char}_\Lambda(P_\Lambda(X^{\op}))
    \cong
    \operatorname{Idl}_\Lambda(X).
\]
Therefore
\[
    G_\Lambda(\mathcal P_\Lambda^{\uparrow}X)
\]
has opens given by upper fuzzy predicates on
\(\operatorname{Idl}_\Lambda(X)\).
\end{example}

\section{The co-Kleisli category}
\label{sec:cokleisli-character-comonad}

\subsection{The enriched co-Kleisli category}
\label{subsec:cokleisli-enriched}

\begin{theorem}[Co-Kleisli morphisms]
\label{thm:cokleisli-morphisms-enriched}
For fuzzy locales \(\mathcal X,\mathcal Y\),
\[
    \operatorname{CoKl}(G_{\mathbb V})(\mathcal X,\mathcal Y)
    =
    \mathbf{FLoc}_{\mathbb V}(G_{\mathbb V}\mathcal X,\mathcal Y)
\]
is naturally equivalent to
\[
    \mathbb V\text{-}\mathbf{Cat}
    \bigl(
        \operatorname{Char}_{\mathbb V}(\mathcal X),
        \operatorname{Char}_{\mathbb V}(\mathcal Y)
    \bigr).
\]
Thus a co-Kleisli morphism
\[
    \mathcal X\to\mathcal Y
\]
is equivalently a \(\mathbb V\)-functor between the character objects.
\end{theorem}

\begin{proof}
By definition,
\[
    G_{\mathbb V}\mathcal X
    =
    \mathcal P_{\mathbb V}^{\uparrow}
    \operatorname{Char}_{\mathbb V}(\mathcal X).
\]
The upper fuzzy character adjunction gives
\[
\begin{aligned}
    \mathbf{FLoc}_{\mathbb V}(G_{\mathbb V}\mathcal X,\mathcal Y)
    &=
    \mathbf{FLoc}_{\mathbb V}
    \bigl(
        \mathcal P_{\mathbb V}^{\uparrow}
        \operatorname{Char}_{\mathbb V}(\mathcal X),
        \mathcal Y
    \bigr)\\
    &\simeq
    \mathbb V\text{-}\mathbf{Cat}
    \bigl(
        \operatorname{Char}_{\mathbb V}(\mathcal X),
        \operatorname{Char}_{\mathbb V}(\mathcal Y)
    \bigr).
\end{aligned}
\]
\end{proof}

\subsection{The posetal specialization}
\label{subsec:cokleisli-posetal}

\begin{corollary}[Posetal co-Kleisli morphisms]
\label{cor:posetal-cokleisli-morphisms}
In the posetal specialization, a co-Kleisli morphism
\[
    \mathcal X\to\mathcal Y
\]
is a \(\Lambda\)-functor
\[
    \operatorname{Char}_\Lambda(\Opens_\Lambda(\mathcal X))
    \to
    \operatorname{Char}_\Lambda(\Opens_\Lambda(\mathcal Y)).
\]
Such a map need not arise from a fuzzy-locale map
\[
    \mathcal X\to\mathcal Y.
\]
\end{corollary}

\begin{proof}
This is the posetal specialization of
Theorem~\ref{thm:cokleisli-morphisms-enriched}.
\end{proof}

\begin{example}[The crisp case]
\label{ex:crisp-cokleisli}
For \(\Lambda=2\), co-Kleisli morphisms are monotone maps between the specialization
preorders of points:
\[
    \operatorname{Pt}(\mathcal X)
    \to
    \operatorname{Pt}(\mathcal Y).
\]
These need not be continuous maps of locales or spaces.
\end{example}

\begin{example}[Presheaf fuzzy locales]
\label{ex:presheaf-fuzzy-locale-cokleisli}
Between upper presheaf fuzzy locales, co-Kleisli morphisms are fuzzy functors
between fuzzy ideal completions:
\[
    \operatorname{Idl}_\Lambda(X)\to\operatorname{Idl}_\Lambda(Y).
\]
\end{example}

\begin{example}[Lawvere co-Kleisli maps]
\label{ex:lawvere-cokleisli}
For the Lawvere quantale, co-Kleisli morphisms are nonexpansive maps between metric
spaces of finite-flat weights.
\end{example}

\section{Summary}
\label{sec:summary-fuzzy-ideal-monad-domains-scott}

The upper fuzzy predicate frame of \(X\) is
\[
    P_{\mathbb V}(X^{\op})
    =
    [X,\underline{\mathbb V}],
\]
and in the posetal specialization it is
\[
    P_\Lambda(X^{\op})
    =
    [X,\Lambda].
\]
The character adjunction
\[
    \mathcal P_{\mathbb V}^{\uparrow}
    \dashv
    \operatorname{Char}_{\mathbb V}
\]
induces the fuzzy ideal monad
\[
    T_\Lambda(X)
    =
    \operatorname{Idl}_\Lambda(X),
\]
whose unit is the principal-ideal embedding and whose multiplication is fuzzy
flattening.

Weak fuzzy domains carry chosen candidate fuzzy ideal-joins, while fuzzy domains
are strict algebras for this monad.  Free fuzzy domains
\[
    \operatorname{Idl}_\Lambda(B)
\]
are the presented fuzzy algebraic domains.  Fuzzy continuous domains are those for
which
\[
    \alpha:\operatorname{Idl}_\Lambda(X)\to X
\]
has an enriched left adjoint
\[
    \beta:X\to\operatorname{Idl}_\Lambda(X),
\]
with fuzzy way-below degree
\[
    \operatorname{wb}_X(y,x)=\beta(x)(y).
\]

Kleisli morphisms for the fuzzy ideal monad are fuzzy ideal-valued maps
\[
    X\rightsquigarrow Y,
\]
or fuzzy approximable maps.  Fuzzy Scott opens of a weak fuzzy domain
\[
    \alpha:\operatorname{Idl}_\Lambda(X)\to X
\]
are exactly the upper fuzzy predicates \(U\) satisfying
\[
    U(\alpha(I))
    =
    \bigjoin_x I(x)\tensor U(x).
\]
For \(\Lambda=2\), these constructions recover ordinary ideals, ideal completion,
dcpos, algebraic dcpos, continuous dcpos, approximable maps, and ordinary Scott
opens.
    %\chapter{The Fuzzy Distributor Equipment and the Fuzzy Way-Below Relation}
\label{ch:fuzzy-distributors-way-below}

In the previous chapter we constructed the fuzzy ideal monad
\[
    T_{\mathbb V}X
    =
    \operatorname{Idl}_{\mathbb V}(X)
\]
from the upper fuzzy character adjunction. Its strict algebras were called fuzzy
domains. Continuous fuzzy domains were defined by the existence of an enriched
left adjoint
\[
    \beta_X:X\to \operatorname{Idl}_{\mathbb V}(X)
\]
to the ideal-algebra map
\[
    \alpha_X:\operatorname{Idl}_{\mathbb V}(X)\to X.
\]
Thus
\[
    \beta_X\dashv \alpha_X.
\]

The basic algebra and continuity data are summarized by the diagrams
\[
\begin{tikzcd}[column sep=huge,row sep=large]
X
    \ar[r,"\eta_X"]
    \ar[dr,"1_X"']
&
\operatorname{Idl}_{\mathbb V}(X)
    \ar[d,"\alpha_X"]
\\
&
X
\end{tikzcd}
\qquad
\begin{tikzcd}[column sep=huge,row sep=large]
X
    \ar[r,bend left=35,"\beta_X"]
&
\operatorname{Idl}_{\mathbb V}(X)
    \ar[l,bend left=35,"\alpha_X"]
\end{tikzcd}
\]
where the first triangle expresses the strict unit law
\[
    \alpha_X\eta_X=1_X.
\]
In the posetal specialization this gave the fuzzy way-below degree
\[
    \operatorname{wb}_X(a,b)=\beta_X(b)(a).
\]

The purpose of this chapter is to explain the equipment-theoretic origin of this
formula. The way-below relation is not added as further structure. It is
constructed from the ideal-algebra map and the principal-ideal embedding by a
right lifting in the fuzzy distributor equipment.

For a fuzzy domain \(X\), the available maps are
\[
    \eta_X:X\to \operatorname{Idl}_{\mathbb V}(X),
    \qquad
    \eta_X(x)=X(-,x),
\]
and
\[
    \alpha_X:\operatorname{Idl}_{\mathbb V}(X)\to X.
\]
The conjoint of \(\alpha_X\) and the companion of \(\eta_X\) are distributors
with the same source and the same target,
\[
    X\nrightarrow \operatorname{Idl}_{\mathbb V}(X).
\]
They fit into the following equipment diagram:
\[
\begin{tikzcd}[column sep=huge,row sep=large]
&
\operatorname{Idl}_{\mathbb V}(X)
\\
X
    \ar[ur,dashed,bend left=13,"\alpha_X^\ast"]
    \ar[ur,dashed,bend right=13,"{(\eta_X)_\ast}"']
\end{tikzcd}
\]
where
\[
    \alpha_X^\ast(b,I)=X(b,\alpha_X I),
    \qquad
    (\eta_X)_\ast(a,I)
    =
    \operatorname{Idl}_{\mathbb V}(X)(\eta_X a,I).
\]

The fuzzy way-below distributor is defined by the right lifting
\[
    \ll_X
    \coloneqq
    [\alpha_X^\ast,(\eta_X)_\ast]
    :
    X\nrightarrow X.
\]
Equivalently, in the enriched distributor bicategory it represents the object of
distributors \(\Theta:X\nrightarrow X\) equipped with a comparison
\[
    \alpha_X^\ast\circ\Theta\to(\eta_X)_\ast.
\]
In the locally posetal reflection, it is the largest distributor \(\Theta\) for
which
\[
    \alpha_X^\ast\circ\Theta\leq(\eta_X)_\ast.
\]

Diagrammatically, a comparison
\[
    \alpha_X^\ast\circ\Theta\to(\eta_X)_\ast
\]
has the form
\[
\begin{tikzcd}[column sep=huge,row sep=large]
X
    \ar[r,dashed,"\Theta"]
    \ar[rr,bend right=18,dashed,"{(\eta_X)_\ast}"']
&
X
    \ar[r,dashed,"\alpha_X^\ast"]
&
\operatorname{Idl}_{\mathbb V}(X)
    \ar[phantom,from=1-2,to=1-3,"\Downarrow" description]
\end{tikzcd}
\]
and the universal property of \(\ll_X\) says that such comparisons are
represented by comparisons
\[
    \Theta\to \ll_X.
\]

Unwinding the right lifting gives
\[
    \ll_X(a,b)
    \cong
    \int_{I\in \operatorname{Idl}_{\mathbb V}(X)}
    \left[
        X(b,\alpha_X I),
        \operatorname{Idl}_{\mathbb V}(X)(\eta_X a,I)
    \right].
\]
By enriched Yoneda,
\[
    \operatorname{Idl}_{\mathbb V}(X)(\eta_X a,I)
    \cong
    I(a),
\]
and hence
\[
    \ll_X(a,b)
    \cong
    \int_{I\in \operatorname{Idl}_{\mathbb V}(X)}
    \left[
        X(b,\alpha_X I),
        I(a)
    \right].
\]
In the posetal reflection this becomes
\[
    \ll_X(a,b)
    =
    \bigmeet_{I\in\operatorname{Idl}_\Lambda(X)}
    \left(
        X(b,\alpha_X I)\multimap I(a)
    \right).
\]

If \(X\) is continuous, with
\[
    \beta_X\dashv \alpha_X,
\]
then this right-lifted distributor is represented by the approximant map:
\[
    \ll_X(a,b)
    \cong
    \operatorname{Idl}_{\mathbb V}(X)(\eta_X a,\beta_X b)
    \cong
    \beta_X(b)(a).
\]
Thus the way-below degree of the previous chapter is the value of a right-lifted
fuzzy distributor.

Throughout this chapter
\[
    (\mathbb V,\tensor,I,[-,-])
\]
is the fuzzy cosmos fixed in the previous chapters. We write
\[
    \underline{\mathbb V}
\]
for its self-enriched category, and we use the same universe conventions as
before: all ends and coends are formed in a universe in which their indexing
fuzzy category is small and the corresponding limits or colimits in
\(\mathbb V\) exist.

All distributor categories below are understood in this universe-bounded sense.
Thus, when \(\operatorname{Idl}_{\mathbb V}(X)\) is not small in the original
universe, we regard it as small in the next universe and form the corresponding
distributor homs, ends, and coends there.

We write
\[
    \Phi:X\nrightarrow Y
\]
for a fuzzy distributor from \(X\) to \(Y\). The arrow \(\nrightarrow\) denotes a
proarrow, not a negated function arrow.

\section{Inherited data from the fuzzy ideal monad}
\label{sec:inherited-data-from-ideal-monad}

This section records the data from the previous chapter that will be used below.
No new structure is imposed on fuzzy domains in this chapter.

\subsection{The enriched construction}
\label{subsec:inherited-data-enriched}

For a small fuzzy category \(X\), the fuzzy ideal object is the full
\(\mathbb V\)-subcategory
\[
    \operatorname{Idl}_{\mathbb V}(X)
    \subseteq
    [X^{\op},\underline{\mathbb V}]
\]
spanned by finite-flat weights. The unit of the fuzzy ideal monad is the
principal-weight embedding
\[
    \eta_X:X\to \operatorname{Idl}_{\mathbb V}(X),
    \qquad
    \eta_X(x)=X(-,x).
\]
This is the restricted Yoneda embedding:
\[
\begin{tikzcd}[column sep=huge,row sep=large]
X
    \ar[rr,"y_X"]
    \ar[dr,swap,"\eta_X"]
&&
{[X^{\op},\underline{\mathbb V}]}
\\
&
{\operatorname{Idl}_{\mathbb V}(X)}
    \ar[from=2-2,to=1-3,hook,"i_X"']
\end{tikzcd}
\]

The multiplication is flattening:
\[
    \mu_X:
    \operatorname{Idl}_{\mathbb V}(\operatorname{Idl}_{\mathbb V}(X))
    \to
    \operatorname{Idl}_{\mathbb V}(X),
\]
with value
\[
    \mu_X(\mathcal I)(x)
    =
    \int^{I\in \operatorname{Idl}_{\mathbb V}(X)}
        \mathcal I(I)\tensor I(x).
\]

The two unit laws are
\[
    \mu_X\eta_{\operatorname{Idl}_{\mathbb V}(X)}
    =
    1_{\operatorname{Idl}_{\mathbb V}(X)},
    \qquad
    \mu_X\operatorname{Idl}_{\mathbb V}(\eta_X)
    =
    1_{\operatorname{Idl}_{\mathbb V}(X)}.
\]
The first of these is displayed by
\[
\begin{tikzcd}[column sep=huge,row sep=large]
\operatorname{Idl}_{\mathbb V}(X)
    \ar[r,"\eta_{\operatorname{Idl}_{\mathbb V}(X)}"]
    \ar[dr,"1_{\operatorname{Idl}_{\mathbb V}(X)}"']
&
\operatorname{Idl}_{\mathbb V}(\operatorname{Idl}_{\mathbb V}(X))
    \ar[d,"\mu_X"]
\\
&
\operatorname{Idl}_{\mathbb V}(X).
\end{tikzcd}
\]
The second is displayed by
\[
\begin{tikzcd}[column sep=huge,row sep=large]
\operatorname{Idl}_{\mathbb V}(X)
    \ar[r,"\operatorname{Idl}_{\mathbb V}(\eta_X)"]
    \ar[dr,"1_{\operatorname{Idl}_{\mathbb V}(X)}"']
&
\operatorname{Idl}_{\mathbb V}(\operatorname{Idl}_{\mathbb V}(X))
    \ar[d,"\mu_X"]
\\
&
\operatorname{Idl}_{\mathbb V}(X).
\end{tikzcd}
\]
The associativity law is
\[
    \mu_X\operatorname{Idl}_{\mathbb V}(\mu_X)
    =
    \mu_X\mu_{\operatorname{Idl}_{\mathbb V}(X)},
\]
displayed by
\[
\begin{tikzcd}[column sep=huge,row sep=large]
\operatorname{Idl}_{\mathbb V}
    (\operatorname{Idl}_{\mathbb V}
    (\operatorname{Idl}_{\mathbb V}(X)))
    \ar[r,"\operatorname{Idl}_{\mathbb V}(\mu_X)"]
    \ar[d,"\mu_{\operatorname{Idl}_{\mathbb V}(X)}"']
&
\operatorname{Idl}_{\mathbb V}
    (\operatorname{Idl}_{\mathbb V}(X))
    \ar[d,"\mu_X"]
\\
\operatorname{Idl}_{\mathbb V}
    (\operatorname{Idl}_{\mathbb V}(X))
    \ar[r,"\mu_X"']
&
\operatorname{Idl}_{\mathbb V}(X).
\end{tikzcd}
\]

A fuzzy domain is a strict algebra for this monad. Thus it is a small fuzzy
category \(X\) equipped with a fuzzy functor
\[
    \alpha_X:
    \operatorname{Idl}_{\mathbb V}(X)
    \to
    X
\]
such that
\[
    \alpha_X\eta_X=1_X
\]
and
\[
    \alpha_X\mu_X
    =
    \alpha_X\operatorname{Idl}_{\mathbb V}(\alpha_X).
\]
The algebra laws are encoded by the diagrams
\[
\begin{tikzcd}[column sep=huge,row sep=large]
X
    \ar[r,"\eta_X"]
    \ar[dr,"1_X"']
&
\operatorname{Idl}_{\mathbb V}(X)
    \ar[d,"\alpha_X"]
\\
&
X
\end{tikzcd}
\]
and
\[
\begin{tikzcd}[column sep=huge,row sep=large]
\operatorname{Idl}_{\mathbb V}
    (\operatorname{Idl}_{\mathbb V}(X))
    \ar[r,"\mu_X"]
    \ar[d,"\operatorname{Idl}_{\mathbb V}(\alpha_X)"']
&
\operatorname{Idl}_{\mathbb V}(X)
    \ar[d,"\alpha_X"]
\\
\operatorname{Idl}_{\mathbb V}(X)
    \ar[r,"\alpha_X"']
&
X.
\end{tikzcd}
\]
The map \(\alpha_X\) evaluates a finite-flat weight by its chosen ideal colimit.

A fuzzy domain \(X\) is continuous when
\[
    \alpha_X:
    \operatorname{Idl}_{\mathbb V}(X)
    \to
    X
\]
has an enriched left adjoint
\[
    \beta_X:X\to \operatorname{Idl}_{\mathbb V}(X),
    \qquad
    \beta_X\dashv\alpha_X.
\]
The object \(\beta_X(b)\) is the finite-flat weight of approximants of \(b\). The
adjunction supplies unit and counit maps
\[
    1_X\to \alpha_X\beta_X,
    \qquad
    \beta_X\alpha_X\to 1_{\operatorname{Idl}_{\mathbb V}(X)}.
\]
These may be represented by
\[
\begin{tikzcd}[column sep=huge,row sep=large]
X
    \ar[r,"\beta_X"]
    \ar[dr,"1_X"']
&
\operatorname{Idl}_{\mathbb V}(X)
    \ar[d,"\alpha_X"]
\\
&
X
\ar[ul,phantom,"\scriptstyle\eta^{\beta\dashv\alpha}" description]
\end{tikzcd}
\qquad
\begin{tikzcd}[column sep=huge,row sep=large]
\operatorname{Idl}_{\mathbb V}(X)
    \ar[r,"\alpha_X"]
    \ar[dr,"1_{\operatorname{Idl}_{\mathbb V}(X)}"']
&
X
    \ar[d,"\beta_X"]
\\
&
\operatorname{Idl}_{\mathbb V}(X)
\ar[ul,phantom,"\scriptstyle\varepsilon^{\beta\dashv\alpha}" description]
\end{tikzcd}
\]

\subsection{The posetal reflection}
\label{subsec:inherited-data-posetal}

Let
\[
    \Lambda=(L,\leq,\meet,\join,\tensor,\multimap,\bot_L,\top_L)
\]
be a fuzzy truth-value object. For a fuzzy preorder \(X\), the object
\[
    \operatorname{Idl}_\Lambda(X)
\]
is a full subpreorder of
\[
    [X^{\op},\Lambda].
\]
Its elements are the finite-flat fuzzy ideals
\[
    I:X^{\op}\to \Lambda.
\]
The principal ideal at \(x\) is
\[
    \eta_X(x)(y)=X(y,x).
\]
A fuzzy domain over \(\Lambda\) has an algebra map
\[
    \alpha_X:
    \operatorname{Idl}_\Lambda(X)
    \to
    X
\]
satisfying the strict algebra laws
\[
    \alpha_X\eta_X=1_X
\]
and
\[
    \alpha_X\mu_X
    =
    \alpha_X\operatorname{Idl}_\Lambda(\alpha_X).
\]
These are strict algebra laws for the transported ideal monad. In the posetal
specialization they are literal equalities of \(\Lambda\)-functors, equivalently
objectwise equalities compatible with the fuzzy preorder structure.

If \(X\) is continuous, then there is a \(\Lambda\)-adjunction
\[
    \beta_X\dashv\alpha_X,
\]
and the previous chapter defined the fuzzy way-below degree by
\[
    \operatorname{wb}_X(a,b)=\beta_X(b)(a).
\]
The goal of the present chapter is to reconstruct this value as a right lifting.

\subsection{Examples}
\label{subsec:inherited-data-examples}

\begin{example}[The crisp case]
\label{ex:inherited-data-crisp}
For
\[
    \Lambda=2,
\]
fuzzy preorders are ordinary preorders and fuzzy ideals are ordinary ideals,
namely directed lower sets. The principal ideal embedding is
\[
    \eta_X(x)=\downarrow x,
\]
and a strict algebra
\[
    \alpha_X:\operatorname{Idl}(X)\to X
\]
assigns to each ideal its supremum:
\[
    \alpha_X(I)=\bigvee I.
\]
Thus separated strict algebras recover dcpos. If \(X\) is continuous, then
\[
    \beta_X(b)=\{a\in X\mid a\ll b\}.
\]
The crisp data fit the diagram
\[
\begin{tikzcd}[column sep=huge,row sep=large]
X
    \ar[r,"x\mapsto\downarrow x"]
    \ar[dr,"1_X"']
&
\operatorname{Idl}(X)
    \ar[d,"I\mapsto\bigvee I"]
\\
&
X.
\end{tikzcd}
\]
\end{example}

\begin{example}[G\"odel truth values]
\label{ex:inherited-data-godel}
For G\"odel truth values,
\[
    r\tensor s=r\meet s.
\]
A fuzzy ideal is a finite-flat lower predicate
\[
    I:X^{\op}\to [0,1],
\]
and \(\alpha_X(I)\) is the chosen G\"odel-valued ideal colimit of \(I\).
\end{example}

\begin{example}[{\L}ukasiewicz truth values]
\label{ex:inherited-data-lukasiewicz}
For {\L}ukasiewicz truth values,
\[
    r\tensor s=\max(0,r+s-1),
    \qquad
    r\multimap s=\min(1,1-r+s).
\]
Finite-flat fuzzy ideals are lower predicates whose flatness is computed using
this bounded-additive tensor.
\end{example}

\begin{example}[Product truth values]
\label{ex:inherited-data-product}
For the product tensor,
\[
    r\tensor s=rs,
\]
and
\[
    r\multimap s
    =
    \begin{cases}
    1,& r\leq s,\\
    s/r,& r>s.
    \end{cases}
\]
The ideal monad is the finite-flat completion for multiplicative fuzzy weights.
\end{example}

\begin{example}[Lawvere truth values]
\label{ex:inherited-data-lawvere}
For the Lawvere truth-value object, the order is
\[
    a\leq_L b
    \quad\Longleftrightarrow\quad
    a\geq b
\]
in the usual numerical order, the tensor is addition, the top element is \(0\),
and the bottom element is \(\infty\). The residual is
\[
    a\multimap c=c\ominus a.
\]
Fuzzy ideals are finite-flat cost weights. The algebra map
\[
    \alpha_X:\operatorname{Idl}_\Lambda(X)\to X
\]
evaluates such weights by their Lawvere-enriched weighted colimits.
\end{example}

\section{Fuzzy distributors}
\label{sec:fuzzy-distributors}

We now construct the horizontal arrows of the fuzzy distributor equipment.

\subsection{The enriched construction}
\label{subsec:fuzzy-distributors-enriched}

\begin{definition}[Fuzzy distributor]
\label{def:fuzzy-distributor}
Let \(X\) and \(Y\) be small fuzzy categories. A \textbf{fuzzy distributor}
\[
    \Phi:X\nrightarrow Y
\]
is a fuzzy functor
\[
    \Phi:X^{\op}\tensor Y\to \underline{\mathbb V}.
\]
We write its value at \((x,y)\) as
\[
    \Phi(x,y).
\]
\end{definition}

Thus a fuzzy distributor is contravariant in its first variable and covariant in
its second variable. This variance is indicated by the diagram
\[
\begin{tikzcd}[column sep=large,row sep=large]
X^{\op}\tensor Y
    \ar[rr,"\Phi"]
&&
\underline{\mathbb V}.
\end{tikzcd}
\]

\begin{definition}[Composition of fuzzy distributors]
\label{def:fuzzy-distributor-composition}
Let
\[
    \Phi:X\nrightarrow Y,
    \qquad
    \Psi:Y\nrightarrow Z
\]
be fuzzy distributors. Their composite
\[
    \Psi\circ\Phi:X\nrightarrow Z
\]
is defined by
\[
    (\Psi\circ\Phi)(x,z)
    =
    \int^{y\in Y}
    \Phi(x,y)\tensor \Psi(y,z).
\]
The identity distributor on \(X\) is
\[
    1_X:X\nrightarrow X,
    \qquad
    1_X(x,x')=X(x,x').
\]
\end{definition}

The composite is represented by the coend-convolution diagram
\[
\begin{tikzcd}[column sep=huge,row sep=large]
X
    \ar[r,dashed,"\Phi"]
    \ar[rr,bend right=18,dashed,"\Psi\circ\Phi"']
&
Y
    \ar[r,dashed,"\Psi"]
&
Z.
\end{tikzcd}
\]
The identity distributor is the horizontal shadow of the hom-object:
\[
\begin{tikzcd}[column sep=large,row sep=large]
X
    \ar[r,dashed,"1_X"]
&
X.
\end{tikzcd}
\]

\begin{proposition}[The fuzzy distributor bicategory]
\label{prop:fuzzy-distributor-bicategory}
Small fuzzy categories in a fixed universe, fuzzy distributors, and enriched
distributor transformations form a bicategory
\[
    \operatorname{Dist}_{\mathbb V}.
\]
For small fuzzy categories \(X\) and \(Y\), the hom \(\mathbb V\)-category of
distributors from \(X\) to \(Y\) is
\[
    \operatorname{Dist}_{\mathbb V}(X,Y)
    =
    [X^{\op}\tensor Y,\underline{\mathbb V}].
\]
Composition is the coend convolution of
Definition~\ref{def:fuzzy-distributor-composition}, and identities are the
hom-distributors.
\end{proposition}

\begin{proof}
The coend formula defines a fuzzy functor
\[
    X^{\op}\tensor Z\to \underline{\mathbb V}
\]
by the functoriality of enriched coends and the functoriality of \(\tensor\).
Associativity follows from Fubini for enriched coends together with associativity
of the tensor product. In diagrammatic form, both routes below evaluate to the
same iterated coend:
\[
\begin{tikzcd}[column sep=huge,row sep=large]
X
    \ar[r,dashed,"\Phi"]
    \ar[rr,bend left=22,dashed,"\Psi\circ\Phi"]
    \ar[rrr,bend right=34,dashed,"{(\Xi\circ\Psi)\circ\Phi}"']
    \ar[rrr,bend left=34,dashed,"{\Xi\circ(\Psi\circ\Phi)}"]
&
Y
    \ar[r,dashed,"\Psi"]
    \ar[rr,bend right=22,swap,dashed,"\Xi\circ\Psi"]
&
Z
    \ar[r,dashed,"\Xi"]
&
W.
\end{tikzcd}
\]
More explicitly,
\[
\begin{aligned}
    \Xi\circ(\Psi\circ\Phi)
    &\cong
    \int^{z}
    \left(
        \int^{y}\Phi(-,y)\tensor\Psi(y,z)
    \right)
    \tensor \Xi(z,-)
\\
    &\cong
    \int^{y,z}
        \Phi(-,y)\tensor\Psi(y,z)\tensor\Xi(z,-)
\\
    &\cong
    (\Xi\circ\Psi)\circ\Phi.
\end{aligned}
\]
The unit laws are the enriched co-Yoneda formulas
\[
    \int^{y\in Y}
        Y(y,y')\tensor \Phi(x,y)
    \cong
    \Phi(x,y')
\]
and
\[
    \int^{x'\in X}
        X(x,x')\tensor \Phi(x',y)
    \cong
    \Phi(x,y).
\]
These are expressed by
\[
\begin{tikzcd}[column sep=huge,row sep=large]
X
    \ar[r,dashed,"\Phi"]
    \ar[rr,bend right=18,dashed,"\Phi"']
&
Y
    \ar[r,dashed,"1_Y"]
&
Y
\end{tikzcd}
\qquad
\begin{tikzcd}[column sep=huge,row sep=large]
X
    \ar[r,dashed,"1_X"]
    \ar[rr,bend right=18,dashed,"\Phi"']
&
X
    \ar[r,dashed,"\Phi"]
&
Y.
\end{tikzcd}
\]
\end{proof}

Together with fuzzy functors as vertical arrows and the companion--conjoint
assignments below, this bicategory underlies the fuzzy distributor equipment.

A cell in the equipment
\[
\begin{tikzcd}[column sep=huge,row sep=large]
X
    \ar[d,"f"']
    \ar[r,dashed,"\Phi"]
&
Y
    \ar[d,"g"]
\\
X'
    \ar[r,dashed,"\Psi"']
&
Y'
\ar[phantom,from=1-2,to=2-1,"\Downarrow" description]
\end{tikzcd}
\]
is a distributor transformation
\[
    \Phi(x,y)\to \Psi(fx,gy),
\]
natural in \(x\in X\) and \(y\in Y\). In the locally posetal reflection, this is
the pointwise inequality
\[
    \Phi(x,y)\leq \Psi(fx,gy).
\]

\subsection{The posetal reflection}
\label{subsec:fuzzy-distributors-posetal}

In the posetal specialization determined by
\[
    \Lambda=(L,\leq,\meet,\join,\tensor,\multimap,\bot_L,\top_L),
\]
a fuzzy distributor
\[
    \Phi:X\nrightarrow Y
\]
is a map
\[
    \Phi:X\times Y\to L
\]
satisfying
\[
    X(x',x)\tensor \Phi(x,y)\tensor Y(y,y')
    \leq
    \Phi(x',y')
\]
for all \(x,x'\in X\) and \(y,y'\in Y\). Thus \(\Phi\) is lower in the first
variable and upper in the second variable.

Composition is
\[
    (\Psi\circ\Phi)(x,z)
    =
    \bigjoin_{y\in Y}
    \Phi(x,y)\tensor\Psi(y,z),
\]
and the identity distributor is
\[
    1_X(x,x')=X(x,x').
\]

\subsection{Examples}
\label{subsec:fuzzy-distributors-examples}

\begin{example}[The crisp case]
\label{ex:fuzzy-distributors-crisp}
For \(\Lambda=2\), fuzzy categories are preorders and a distributor
\[
    \Phi:X\nrightarrow Y
\]
is a relation
\[
    \Phi\subseteq X\times Y
\]
such that
\[
    x'\leq_X x,\quad \Phi(x,y),\quad y\leq_Y y'
    \quad\Longrightarrow\quad
    \Phi(x',y').
\]
Composition is ordinary relational composition:
\[
    (\Psi\circ\Phi)(x,z)
    \Longleftrightarrow
    \exists y\in Y,\; \Phi(x,y)\wedge \Psi(y,z).
\]
\end{example}

\begin{example}[G\"odel-valued distributors]
\label{ex:fuzzy-distributors-godel}
For G\"odel truth values,
\[
    (\Psi\circ\Phi)(x,z)
    =
    \bigjoin_{y\in Y}
    \min(\Phi(x,y),\Psi(y,z)).
\]
\end{example}

\begin{example}[{\L}ukasiewicz-valued distributors]
\label{ex:fuzzy-distributors-lukasiewicz}
For {\L}ukasiewicz truth values,
\[
    (\Psi\circ\Phi)(x,z)
    =
    \bigjoin_{y\in Y}
    \max(0,\Phi(x,y)+\Psi(y,z)-1).
\]
\end{example}

\begin{example}[Product-valued distributors]
\label{ex:fuzzy-distributors-product}
For product truth values,
\[
    (\Psi\circ\Phi)(x,z)
    =
    \bigjoin_{y\in Y}
    \Phi(x,y)\Psi(y,z).
\]
\end{example}

\begin{example}[Lawvere-valued distributors]
\label{ex:fuzzy-distributors-lawvere}
For the Lawvere quantale, joins in the truth-value order are infima in the usual
numerical order, and tensor is addition. Hence distributor composition is
infimal convolution:
\[
    (\Psi\circ\Phi)(x,z)
    =
    \inf_{y\in Y}
    \bigl(\Phi(x,y)+\Psi(y,z)\bigr).
\]
\end{example}

\section{Companions and conjoints}
\label{sec:companions-conjoints}

Fuzzy functors determine distinguished fuzzy distributors. These are the
companion and conjoint constructions. They place vertical arrows inside the
horizontal distributor equipment.

\subsection{The enriched construction}
\label{subsec:companions-conjoints-enriched}

\begin{definition}[Companion and conjoint]
\label{def:companion-conjoint}
Let
\[
    f:X\to Y
\]
be a fuzzy functor. Its \textbf{companion} is the distributor
\[
    f_\ast:X\nrightarrow Y
\]
defined by
\[
    f_\ast(x,y)=Y(fx,y).
\]
Its \textbf{conjoint} is the distributor
\[
    f^\ast:Y\nrightarrow X
\]
defined by
\[
    f^\ast(y,x)=Y(y,fx).
\]
\end{definition}

The companion and conjoint sit on opposite sides of \(f\):
\[
\begin{tikzcd}[column sep=huge,row sep=large]
X
    \ar[d,"f"']
    \ar[r,dashed,"f_\ast"]
&
Y
    \ar[d,"1_Y"]
\\
Y
    \ar[r,dashed,"1_Y"']
&
Y
\end{tikzcd}
\qquad
\begin{tikzcd}[column sep=huge,row sep=large]
Y
    \ar[d,"1_Y"']
    \ar[r,dashed,"f^\ast"]
&
X
    \ar[d,"f"]
\\
Y
    \ar[r,dashed,"1_Y"']
&
Y.
\end{tikzcd}
\]
Equivalently,
\[
\begin{tikzcd}[column sep=huge,row sep=large]
X
    \ar[r,bend left=14,dashed,"f_\ast"]
&
Y
    \ar[l,bend left=14,dashed,"f^\ast"]
\end{tikzcd}
\]
is an adjoint pair in the distributor bicategory.

\begin{proposition}[Companion--conjoint adjunction]
\label{prop:companion-conjoint-adjunction}
For every fuzzy functor
\[
    f:X\to Y,
\]
there is an adjunction in the distributor bicategory
\[
    f_\ast\dashv f^\ast.
\]
\end{proposition}

\begin{proof}
The unit is a distributor transformation
\[
    1_X\to f^\ast\circ f_\ast.
\]
At \((x,x')\), the target is
\[
    (f^\ast\circ f_\ast)(x,x')
    =
    \int^{y\in Y}
        Y(fx,y)\tensor Y(y,fx').
\]
The required map is obtained by composing the functoriality map
\[
    X(x,x')\to Y(fx,fx')
\]
with the coend coprojection at \(y=fx'\):
\[
    Y(fx,fx')
    \cong
    Y(fx,fx')\tensor I
    \to
    Y(fx,fx')\tensor Y(fx',fx')
    \to
    \int^{y\in Y}
        Y(fx,y)\tensor Y(y,fx').
\]

The counit is a distributor transformation
\[
    f_\ast\circ f^\ast\to 1_Y.
\]
At \((y,y')\), it is induced by composition in \(Y\):
\[
    \int^{x\in X}
        Y(y,fx)\tensor Y(fx,y')
    \to
    Y(y,y').
\]
The triangle identities follow from associativity and unitality of composition in
\(Y\), equivalently from the standard companion--conjoint construction in the
enriched distributor equipment.
\end{proof}

\begin{proposition}[Functoriality of companions and conjoints]
\label{prop:companion-conjoint-functoriality}
For fuzzy functors
\[
    X\xrightarrow{f}Y\xrightarrow{g}Z,
\]
there are canonical isomorphisms of distributors
\[
    (gf)_\ast
    \cong
    g_\ast\circ f_\ast
\]
and
\[
    (gf)^\ast
    \cong
    f^\ast\circ g^\ast.
\]
Moreover,
\[
    (1_X)_\ast\cong 1_X,
    \qquad
    (1_X)^\ast\cong 1_X.
\]
\end{proposition}

\begin{proof}
For companions,
\[
\begin{aligned}
    (g_\ast\circ f_\ast)(x,z)
    &=
    \int^{y\in Y}
        Y(fx,y)\tensor Z(gy,z)
\\
    &\cong
    Z(gfx,z)
\\
    &=
    (gf)_\ast(x,z),
\end{aligned}
\]
by enriched Yoneda. Diagrammatically:
\[
\begin{tikzcd}[column sep=huge,row sep=large]
X
    \ar[r,dashed,"f_\ast"]
    \ar[rr,bend right=18,dashed,"{(gf)_\ast}"']
&
Y
    \ar[r,dashed,"g_\ast"]
&
Z.
\end{tikzcd}
\]
For conjoints,
\[
\begin{aligned}
    (f^\ast\circ g^\ast)(z,x)
    &=
    \int^{y\in Y}
        Z(z,gy)\tensor Y(y,fx)
\\
    &\cong
    Z(z,gfx)
\\
    &=
    (gf)^\ast(z,x).
\end{aligned}
\]
This is represented by the reverse diagram
\[
\begin{tikzcd}[column sep=huge,row sep=large]
Z
    \ar[r,dashed,"g^\ast"]
    \ar[rr,bend right=18,dashed,"{(gf)^\ast}"']
&
Y
    \ar[r,dashed,"f^\ast"]
&
X.
\end{tikzcd}
\]
The identity statements follow directly from the definitions.
\end{proof}

\subsection{The posetal reflection}
\label{subsec:companions-conjoints-posetal}

For a \(\Lambda\)-functor
\[
    f:X\to Y,
\]
the companion and conjoint are
\[
    f_\ast(x,y)=Y(fx,y)
\]
and
\[
    f^\ast(y,x)=Y(y,fx).
\]
The adjunction
\[
    f_\ast\dashv f^\ast
\]
amounts to the distributor inequalities
\[
    1_X\leq f^\ast\circ f_\ast
\]
and
\[
    f_\ast\circ f^\ast\leq 1_Y.
\]
The first inequality follows from functoriality of \(f\), and the second follows
from transitivity in \(Y\).

\subsection{Examples}
\label{subsec:companions-conjoints-examples}

\begin{example}[The crisp case]
\label{ex:companions-crisp}
For a monotone map
\[
    f:X\to Y
\]
between preorders,
\[
    f_\ast(x,y)
    \Longleftrightarrow
    f(x)\leq_Y y,
\]
and
\[
    f^\ast(y,x)
    \Longleftrightarrow
    y\leq_Y f(x).
\]
\end{example}

\begin{example}[Standard fuzzy truth values]
\label{ex:companions-standard}
For G\"odel, {\L}ukasiewicz, and product truth values, the formulas remain
\[
    f_\ast(x,y)=Y(fx,y),
    \qquad
    f^\ast(y,x)=Y(y,fx).
\]
The tensor and residual of the chosen truth-value object determine the distributor
composition and the adjunction inequalities.
\end{example}

\begin{example}[Lawvere truth values]
\label{ex:companions-lawvere}
For the Lawvere quantale, a fuzzy functor is a nonexpansive map. The companion
and conjoint of
\[
    f:X\to Y
\]
are
\[
    f_\ast(x,y)=d_Y(fx,y),
    \qquad
    f^\ast(y,x)=d_Y(y,fx).
\]
Thus companions and conjoints record the forward and backward distance kernels
induced by \(f\).
\end{example}

\section{Right liftings of fuzzy distributors}
\label{sec:right-liftings-fuzzy-distributors}

Right liftings are the categorical operation that will produce the fuzzy
way-below distributor.

\subsection{The enriched construction}
\label{subsec:right-liftings-enriched}

Let
\[
    \Phi:A\nrightarrow C,
    \qquad
    \Psi:B\nrightarrow C
\]
be fuzzy distributors with common target \(C\).

\begin{definition}[Right lifting]
\label{def:right-lifting}
The \textbf{right lifting} of \(\Psi\) through \(\Phi\) is the fuzzy distributor
\[
    [\Phi,\Psi]:B\nrightarrow A
\]
defined by
\[
    [\Phi,\Psi](b,a)
    =
    \int_{c\in C}
    [\Phi(a,c),\Psi(b,c)].
\]
\end{definition}

The right lifting is characterized by the following universal diagram:
\[
\begin{tikzcd}[column sep=huge,row sep=large]
B
    \ar[r,dashed,"\Theta"]
    \ar[rr,bend right=18,dashed,"\Psi"']
&
A
    \ar[r,dashed,"\Phi"]
&
C
\ar[phantom,from=1-2,to=1-3,"\Downarrow" description]
\end{tikzcd}
\]
where comparisons
\[
    \Theta\to[\Phi,\Psi]
\]
correspond to comparisons
\[
    \Phi\circ\Theta\to\Psi.
\]
In the locally posetal reflection, this says that \([\Phi,\Psi]\) is the largest
\(\Theta\) satisfying
\[
    \Phi\circ\Theta\leq\Psi.
\]

\begin{proposition}[Universal property of right liftings]
\label{prop:right-lifting-universal-property}
For every fuzzy distributor
\[
    \Theta:B\nrightarrow A,
\]
there is a natural isomorphism
\[
    \operatorname{Dist}_{\mathbb V}(B,A)(\Theta,[\Phi,\Psi])
    \cong
    \operatorname{Dist}_{\mathbb V}(B,C)(\Phi\circ\Theta,\Psi).
\]
Equivalently, postcomposition by \(\Phi\) has right adjoint
\[
    [\Phi,-].
\]
\end{proposition}

\begin{proof}
Using the enriched hom-object in the distributor category, the left-hand side is
\[
\begin{aligned}
\operatorname{Dist}_{\mathbb V}(B,A)(\Theta,[\Phi,\Psi])
&=
\int_{b,a}
[
    \Theta(b,a),
    [\Phi,\Psi](b,a)
]
\\
&=
\int_{b,a}
\left[
    \Theta(b,a),
    \int_c[\Phi(a,c),\Psi(b,c)]
\right]
\\
&\cong
\int_{b,a,c}
[
    \Theta(b,a),
    [\Phi(a,c),\Psi(b,c)]
]
\\
&\cong
\int_{b,a,c}
[
    \Theta(b,a)\tensor\Phi(a,c),
    \Psi(b,c)
].
\end{aligned}
\]
By the end--coend adjunction, this is isomorphic to
\[
    \int_{b,c}
    \left[
        \int^{a}
        \Theta(b,a)\tensor\Phi(a,c),
        \Psi(b,c)
    \right].
\]
The coend is
\[
    (\Phi\circ\Theta)(b,c).
\]
Hence the displayed object is
\[
    \operatorname{Dist}_{\mathbb V}(B,C)(\Phi\circ\Theta,\Psi).
\]
\end{proof}

\subsection{The posetal reflection}
\label{subsec:right-liftings-posetal}

In the posetal reflection, the right lifting is
\[
    [\Phi,\Psi](b,a)
    =
    \bigmeet_{c\in C}
    \left(
        \Phi(a,c)\multimap\Psi(b,c)
    \right).
\]
Its universal property becomes
\[
    \Theta\leq[\Phi,\Psi]
    \quad\Longleftrightarrow\quad
    \Phi\circ\Theta\leq\Psi.
\]
Thus \([\Phi,\Psi](b,a)\) is the degree to which every instance of
\(\Phi(a,c)\) entails the corresponding instance of \(\Psi(b,c)\).

The posetal lifting adjunction can be displayed as
\[
\begin{tikzcd}[column sep=huge,row sep=large]
\operatorname{Dist}_{\Lambda}(B,A)
    \ar[r,bend left=35,"{\Phi\circ -}"]
&
\operatorname{Dist}_{\Lambda}(B,C)
    \ar[l,bend left=35,"{[\Phi,-]}"]
\end{tikzcd}
\]
with order-enriched adjunction condition
\[
    \Phi\circ\Theta\leq\Psi
    \quad\Longleftrightarrow\quad
    \Theta\leq[\Phi,\Psi].
\]

\subsection{Examples}
\label{subsec:right-liftings-examples}

\begin{example}[The crisp case]
\label{ex:right-lifting-crisp}
For \(\Lambda=2\),
\[
    [\Phi,\Psi](b,a)
\]
holds precisely when
\[
    \forall c\in C,\;
    \Phi(a,c)\Rightarrow \Psi(b,c).
\]
The universal property is
\[
    \Theta\subseteq[\Phi,\Psi]
    \quad\Longleftrightarrow\quad
    \Phi\circ\Theta\subseteq\Psi.
\]
\end{example}

\begin{example}[G\"odel right liftings]
\label{ex:right-lifting-godel}
For G\"odel truth values,
\[
    [\Phi,\Psi](b,a)
    =
    \bigmeet_{c\in C}
    \left(
        \Phi(a,c)\multimap_{\mathrm G}\Psi(b,c)
    \right),
\]
where
\[
    u\multimap_{\mathrm G}v
    =
    \begin{cases}
    1,&u\leq v,\\
    v,&u>v.
    \end{cases}
\]
\end{example}

\begin{example}[{\L}ukasiewicz right liftings]
\label{ex:right-lifting-lukasiewicz}
For {\L}ukasiewicz truth values,
\[
    [\Phi,\Psi](b,a)
    =
    \bigmeet_{c\in C}
    \min(1,1-\Phi(a,c)+\Psi(b,c)).
\]
\end{example}

\begin{example}[Product right liftings]
\label{ex:right-lifting-product}
For product truth values,
\[
    [\Phi,\Psi](b,a)
    =
    \bigmeet_{c\in C}
    \left(
        \Phi(a,c)\multimap_\Pi\Psi(b,c)
    \right),
\]
where
\[
    u\multimap_\Pi v
    =
    \begin{cases}
    1,&u\leq v,\\
    v/u,&u>v.
    \end{cases}
\]
\end{example}

\begin{example}[Lawvere right liftings]
\label{ex:right-lifting-lawvere}
For the Lawvere quantale, the internal hom is truncated subtraction. Written in
the usual numerical order,
\[
    [\Phi,\Psi](b,a)
    =
    \sup_{c\in C}
    \bigl(\Psi(b,c)\ominus \Phi(a,c)\bigr),
\]
where
\[
    v\ominus u=\max(0,v-u)
\]
with the usual extended-value conventions. Thus a Lawvere right lifting measures
the largest residual excess needed to transform \(\Phi(a,-)\) into
\(\Psi(b,-)\).
\end{example}

\section{Ideals as distributors and approximable kernels}
\label{sec:ideals-as-distributors}

We next relate the ideal monad from the previous chapter to the distributor
equipment.

\subsection{The enriched construction}
\label{subsec:ideals-as-distributors-enriched}

A fuzzy ideal of \(X\) is a finite-flat weight
\[
    I:X^{\op}\to\underline{\mathbb V}.
\]
Equivalently, it is a fuzzy distributor
\[
    I:X\nrightarrow \mathbf 1,
\]
because
\[
    X^{\op}\tensor\mathbf 1
    \cong
    X^{\op}.
\]
Thus
\[
    \operatorname{Idl}_{\mathbb V}(X)
\]
is the full \(\mathbb V\)-subcategory of
\[
    \operatorname{Dist}_{\mathbb V}(X,\mathbf 1)
    =
    [X^{\op},\underline{\mathbb V}]
\]
spanned by finite-flat distributors. This inclusion is represented by
\[
\begin{tikzcd}[column sep=huge,row sep=large]
{\operatorname{Idl}_{\mathbb V}(X)}
    \ar[r,hook,"i_X"]
&
{\operatorname{Dist}_{\mathbb V}(X,\mathbf{1})}
    \ar[r,"\cong"]
&
{[X^{\op},\underline{\mathbb V}]}.
\end{tikzcd}
\]

\begin{definition}[Fuzzy approximable map]
\label{def:fuzzy-approximable-map-equipment}
Let \(X\) and \(Y\) be small fuzzy categories. A \textbf{fuzzy approximable map}
from \(X\) to \(Y\) is a Kleisli morphism for the fuzzy ideal monad:
\[
    R:X\to \operatorname{Idl}_{\mathbb V}(Y).
\]
Equivalently, it is a kernel
\[
    R:X\tensor Y^{\op}\to\underline{\mathbb V}
\]
such that, for every \(x\in X\), the fibre
\[
    R(x,-):Y^{\op}\to\underline{\mathbb V}
\]
is finite-flat.
\end{definition}

The distributor variance of this kernel is opposite to the variance of a
distributor \(X\nrightarrow Y\). Under the convention
\[
    \Phi:A\nrightarrow B
    \quad\text{means}\quad
    \Phi:A^{\op}\tensor B\to\underline{\mathbb V},
\]
the kernel \(R\) corresponds to a distributor
\[
    R^\dagger:Y\nrightarrow X,
    \qquad
    R^\dagger(y,x)=R(x,y).
\]
The variance conversion is summarized by
\[
\begin{tikzcd}[column sep=huge,row sep=large]
X
    \ar[r,"R"]
&
{\operatorname{Idl}_{\mathbb V}(Y)}
    \ar[r,hook,"i_Y"]
&
{[Y^{\op},\underline{\mathbb V}]}
\end{tikzcd}
\qquad
\begin{tikzcd}[column sep=huge,row sep=large]
Y
    \ar[r,dashed,"R^\dagger"]
&
X.
\end{tikzcd}
\]
The upper row is vertical/Kleisli data, while the lower row is distributor data.

\begin{proposition}[Kleisli extension as distributor convolution]
\label{prop:kleisli-extension-convolution}
Let
\[
    R:X\rightsquigarrow Y
\]
be a fuzzy approximable map. Its Kleisli extension
\[
    \overline R:
    \operatorname{Idl}_{\mathbb V}(X)
    \to
    \operatorname{Idl}_{\mathbb V}(Y)
\]
is given by
\[
    \overline R(I)(y)
    =
    \int^{x\in X}
    I(x)\tensor R(x,y).
\]
Equivalently, if
\[
    R^\dagger:Y\nrightarrow X
\]
is the associated oppositely oriented distributor, then
\[
    \overline R(I)
    =
    I\circ R^\dagger.
\]
\end{proposition}

\begin{proof}
The Kleisli extension is
\[
\begin{tikzcd}[column sep=huge,row sep=large]
\operatorname{Idl}_{\mathbb V}(X)
    \ar[r,"\operatorname{Idl}_{\mathbb V}(R)"]
&
\operatorname{Idl}_{\mathbb V}
    (\operatorname{Idl}_{\mathbb V}(Y))
    \ar[r,"\mu_Y"]
&
\operatorname{Idl}_{\mathbb V}(Y).
\end{tikzcd}
\]
The multiplication \(\mu_Y\) is flattening, so
\[
    \overline R(I)(y)
    =
    \int^{x\in X}
    I(x)\tensor R(x,y).
\]
On the other hand,
\[
    (I\circ R^\dagger)(y)
    =
    \int^{x\in X}
    R^\dagger(y,x)\tensor I(x)
    =
    \int^{x\in X}
    R(x,y)\tensor I(x),
\]
which is canonically isomorphic to the same coend by symmetry of \(\tensor\).
\end{proof}

The convolution picture is
\[
\begin{tikzcd}[column sep=huge,row sep=large]
Y
    \ar[r,dashed,"R^\dagger"]
    \ar[rr,bend right=18,dashed,"I\circ R^\dagger"']
&
X
    \ar[r,dashed,"I"]
&
\mathbf 1.
\end{tikzcd}
\]

\subsection{The posetal reflection}
\label{subsec:ideals-as-distributors-posetal}

In the posetal specialization, a fuzzy approximable map
\[
    R:X\rightsquigarrow Y
\]
is a map
\[
    R:X\times Y\to L
\]
such that each fibre
\[
    R(x,-):Y^{\op}\to \Lambda
\]
is a fuzzy ideal of \(Y\), and such that the assignment
\[
    x\longmapsto R(x,-)
\]
is a \(\Lambda\)-functor
\[
    X\to \operatorname{Idl}_\Lambda(Y).
\]
Its extension is
\[
    \overline R(I)(y)
    =
    \bigjoin_{x\in X}
    I(x)\tensor R(x,y).
\]
Kleisli composition is therefore
\[
    (S\circ R)(x,z)
    =
    \bigjoin_{y\in Y}
    R(x,y)\tensor S(y,z).
\]

\subsection{Examples}
\label{subsec:ideals-as-distributors-examples}

\begin{example}[The crisp case]
\label{ex:ideals-as-distributors-crisp}
For \(\Lambda=2\), a fuzzy approximable map
\[
    R:X\rightsquigarrow Y
\]
is a monotone map
\[
    X\to\operatorname{Idl}(Y).
\]
Equivalently, it is a relation
\[
    R\subseteq X\times Y
\]
such that every fibre
\[
    R[x]=\{y\in Y\mid R(x,y)\}
\]
is an ideal of \(Y\), and
\[
    x\leq x'
    \quad\Longrightarrow\quad
    R[x]\subseteq R[x'].
\]
The extension formula becomes
\[
    \overline R(I)=\bigcup_{x\in I}R[x].
\]
\end{example}

\begin{example}[G\"odel approximable maps]
\label{ex:ideals-as-distributors-godel}
For G\"odel truth values,
\[
    \overline R(I)(y)
    =
    \bigjoin_{x\in X}
    \min(I(x),R(x,y)).
\]
\end{example}

\begin{example}[{\L}ukasiewicz approximable maps]
\label{ex:ideals-as-distributors-lukasiewicz}
For {\L}ukasiewicz truth values,
\[
    \overline R(I)(y)
    =
    \bigjoin_{x\in X}
    \max(0,I(x)+R(x,y)-1).
\]
\end{example}

\begin{example}[Product approximable maps]
\label{ex:ideals-as-distributors-product}
For product truth values,
\[
    \overline R(I)(y)
    =
    \bigjoin_{x\in X}
    I(x)R(x,y).
\]
\end{example}

\begin{example}[Lawvere approximable maps]
\label{ex:ideals-as-distributors-lawvere}
For the Lawvere quantale, the extension formula is infimal convolution:
\[
    \overline R(I)(y)
    =
    \inf_{x\in X}
    \bigl(I(x)+R(x,y)\bigr),
\]
written in the usual numerical order.
\end{example}

\section{The ideal algebra map inside the distributor equipment}
\label{sec:ideal-algebra-map-equipment}

We now place the ideal-algebra map and the principal-ideal embedding inside the
distributor equipment.

\subsection{The enriched construction}
\label{subsec:ideal-algebra-map-equipment-enriched}

Let \(X\) be a fuzzy domain. We have the algebra map
\[
    \alpha_X:\operatorname{Idl}_{\mathbb V}(X)\to X
\]
and the principal-ideal embedding
\[
    \eta_X:X\to\operatorname{Idl}_{\mathbb V}(X).
\]

The conjoint of \(\alpha_X\) is the distributor
\[
    \alpha_X^\ast:
    X\nrightarrow \operatorname{Idl}_{\mathbb V}(X)
\]
given by
\[
    \alpha_X^\ast(b,I)
    =
    X(b,\alpha_X I).
\]
The companion of \(\eta_X\) is the distributor
\[
    (\eta_X)_\ast:
    X\nrightarrow \operatorname{Idl}_{\mathbb V}(X)
\]
given by
\[
    (\eta_X)_\ast(a,I)
    =
    \operatorname{Idl}_{\mathbb V}(X)(\eta_X a,I).
\]

The two constructions arise from opposite sides of the algebra triangle:
\[
\begin{tikzcd}[column sep=huge,row sep=large]
X
    \ar[r,"\eta_X"]
    \ar[dr,"1_X"']
&
\operatorname{Idl}_{\mathbb V}(X)
    \ar[d,"\alpha_X"]
\\
&
X
\end{tikzcd}
\qquad
\begin{tikzcd}[column sep=huge,row sep=large]
&
\operatorname{Idl}_{\mathbb V}(X)
\\
X
    \ar[ur,dashed,bend left=14,"\alpha_X^\ast"]
    \ar[ur,dashed,bend right=14,"{(\eta_X)_\ast}"']
\end{tikzcd}
\]

Thus
\[
    \alpha_X^\ast
    \qquad\text{and}\qquad
    (\eta_X)_\ast
\]
are two fuzzy distributors with the same source and the same target:
\[
    X\nrightarrow \operatorname{Idl}_{\mathbb V}(X).
\]
Consequently, the right lifting
\[
    [\alpha_X^\ast,(\eta_X)_\ast]
\]
is a fuzzy distributor
\[
    X\nrightarrow X.
\]
This is the distributor that will be defined as the fuzzy way-below relation.

\subsection{The posetal reflection}
\label{subsec:ideal-algebra-map-equipment-posetal}

For a fuzzy domain over \(\Lambda\), the two distributors are
\[
    \alpha_X^\ast(b,I)
    =
    X(b,\alpha_X I),
\]
and
\[
    (\eta_X)_\ast(a,I)
    =
    \operatorname{Idl}_\Lambda(X)(\eta_X a,I).
\]
By the enriched Yoneda lemma in the posetal setting,
\[
    \operatorname{Idl}_\Lambda(X)(\eta_X a,I)
    =
    I(a).
\]
Hence
\[
    (\eta_X)_\ast(a,I)=I(a).
\]

\subsection{Examples}
\label{subsec:ideal-algebra-map-equipment-examples}

\begin{example}[The crisp case]
\label{ex:ideal-algebra-map-crisp}
For \(\Lambda=2\),
\[
    \alpha_X^\ast(b,I)
    \Longleftrightarrow
    b\leq \bigvee I,
\]
and
\[
    (\eta_X)_\ast(a,I)
    \Longleftrightarrow
    \downarrow a\subseteq I.
\]
Since \(I\) is lower,
\[
    \downarrow a\subseteq I
    \quad\Longleftrightarrow\quad
    a\in I.
\]
Thus the right lifting expresses the ordinary implication
\[
    b\leq\bigvee I
    \quad\Longrightarrow\quad
    a\in I
\]
uniformly over all ideals \(I\).
\end{example}

\begin{example}[Standard fuzzy truth values]
\label{ex:ideal-algebra-map-standard}
For G\"odel, {\L}ukasiewicz, and product truth values, the formulas are
\[
    \alpha_X^\ast(b,I)=X(b,\alpha_X I),
    \qquad
    (\eta_X)_\ast(a,I)=I(a).
\]
The later right-lifting formula compares these two degrees using the residual
appropriate to the chosen tensor.
\end{example}

\begin{example}[Lawvere truth values]
\label{ex:ideal-algebra-map-lawvere}
For the Lawvere quantale,
\[
    \alpha_X^\ast(b,I)=d_X(b,\alpha_X I)
\]
is the cost from \(b\) to the ideal colimit of \(I\), while
\[
    (\eta_X)_\ast(a,I)=I(a)
\]
is the cost-degree with which \(a\) belongs to the weight \(I\).
\end{example}

\section{The fuzzy way-below distributor}
\label{sec:fuzzy-way-below-distributor}

The fuzzy way-below relation is now obtained by applying the right-lifting
construction to the two distributors built from \(\alpha_X\) and \(\eta_X\).

\subsection{The enriched construction}
\label{subsec:fuzzy-way-below-enriched}

\begin{definition}[Fuzzy way-below distributor]
\label{def:fuzzy-way-below-distributor}
Let \(X\) be a fuzzy domain. The \textbf{fuzzy way-below distributor} of \(X\)
is the right lifting
\[
    \ll_X
    \coloneqq
    [\alpha_X^\ast,(\eta_X)_\ast]
    :
    X\nrightarrow X.
\]
We write
\[
    \ll_X(a,b)
\]
for its value at \((a,b)\), where \(a\) is the approximant variable and \(b\) is
the element being approximated.
\end{definition}

With the right-lifting convention
\[
    [\Phi,\Psi](b,a)=\int_c[\Phi(a,c),\Psi(b,c)],
\]
the first variable of \(\ll_X(a,b)\) comes from the source variable of
\((\eta_X)_\ast\), while the second variable comes from the source variable of
\(\alpha_X^\ast\). Thus
\[
    \ll_X(a,b)
    =
    [\alpha_X^\ast,(\eta_X)_\ast](a,b)
    =
    \int_I
    [
        \alpha_X^\ast(b,I),
        (\eta_X)_\ast(a,I)
    ].
\]

The universal right-lifting triangle is
\[
\begin{tikzcd}[column sep=huge,row sep=large]
X
    \ar[r,dashed,"\ll_X"]
    \ar[rr,bend right=18,dashed,"{(\eta_X)_\ast}"']
&
X
    \ar[r,dashed,"\alpha_X^\ast"]
&
\operatorname{Idl}_{\mathbb V}(X)
\ar[phantom,from=1-2,to=1-3,"\Downarrow" description]
\end{tikzcd}
\]
and its universal property says that every compatible triangle factors through
\(\ll_X\):
\[
    \operatorname{Dist}_{\mathbb V}(X,X)(\Theta,\ll_X)
    \cong
    \operatorname{Dist}_{\mathbb V}
    \bigl(X,\operatorname{Idl}_{\mathbb V}(X)\bigr)
    \bigl(\alpha_X^\ast\circ\Theta,(\eta_X)_\ast\bigr).
\]
Equivalently, a comparison
\[
    \alpha_X^\ast\circ\Theta\to(\eta_X)_\ast
\]
is the same as a comparison
\[
    \Theta\to\ll_X.
\]

\begin{proposition}[Formula for the fuzzy way-below distributor]
\label{prop:fuzzy-way-below-formula}
For every fuzzy domain \(X\),
\[
    \ll_X(a,b)
    \cong
    \int_{I\in \operatorname{Idl}_{\mathbb V}(X)}
    \left[
        X(b,\alpha_X I),
        \operatorname{Idl}_{\mathbb V}(X)(\eta_X a,I)
    \right].
\]
By enriched Yoneda, this is equivalently
\[
    \ll_X(a,b)
    \cong
    \int_{I\in \operatorname{Idl}_{\mathbb V}(X)}
    \left[
        X(b,\alpha_X I),
        I(a)
    \right].
\]
\end{proposition}

\begin{proof}
This is the right-lifting formula applied to
\[
    \Phi=\alpha_X^\ast,
    \qquad
    \Psi=(\eta_X)_\ast.
\]
Thus
\[
    \ll_X(a,b)
    =
    [\alpha_X^\ast,(\eta_X)_\ast](a,b)
\]
is
\[
    \int_{I\in \operatorname{Idl}_{\mathbb V}(X)}
    [
        \alpha_X^\ast(b,I),
        (\eta_X)_\ast(a,I)
    ].
\]
Using the definitions of \(\alpha_X^\ast\) and \((\eta_X)_\ast\), this becomes
\[
    \int_{I\in \operatorname{Idl}_{\mathbb V}(X)}
    \left[
        X(b,\alpha_X I),
        \operatorname{Idl}_{\mathbb V}(X)(\eta_X a,I)
    \right].
\]
Finally,
\[
    \operatorname{Idl}_{\mathbb V}(X)(\eta_X a,I)
    \cong
    I(a)
\]
by enriched Yoneda.
\end{proof}

\begin{proposition}[Universal property of way-below]
\label{prop:way-below-universal-property}
For every fuzzy distributor
\[
    \Theta:X\nrightarrow X,
\]
there is a natural isomorphism
\[
    \operatorname{Dist}_{\mathbb V}(X,X)(\Theta,\ll_X)
    \cong
    \operatorname{Dist}_{\mathbb V}
    \bigl(X,\operatorname{Idl}_{\mathbb V}(X)\bigr)
    \bigl(\alpha_X^\ast\circ\Theta,(\eta_X)_\ast\bigr).
\]
In the locally posetal reflection this becomes
\[
    \Theta\leq\ll_X
    \quad\Longleftrightarrow\quad
    \alpha_X^\ast\circ\Theta\leq(\eta_X)_\ast.
\]
\end{proposition}

\begin{proof}
This is the universal property of the right lifting
\[
    [\alpha_X^\ast,(\eta_X)_\ast].
\]
\end{proof}

The universal property can be visualized as the adjunction
\[
\begin{tikzcd}[column sep=huge,row sep=large]
\operatorname{Dist}_{\mathbb V}(X,X)
    \ar[r,bend left=35,"{\alpha_X^\ast\circ -}"]
&
\operatorname{Dist}_{\mathbb V}
    \bigl(X,\operatorname{Idl}_{\mathbb V}(X)\bigr)
    \ar[l,bend left=35,"{[\alpha_X^\ast,-]}"]
\end{tikzcd}
\]
evaluated at \((\eta_X)_\ast\).

\begin{proposition}[Way-below implies below]
\label{prop:fuzzy-way-below-implies-below}
For every fuzzy domain \(X\), there is a canonical comparison
\[
    \ll_X(a,b)\to X(a,b).
\]
In the posetal reflection this says
\[
    \ll_X(a,b)\leq X(a,b).
\]
\end{proposition}

\begin{proof}
By Proposition~\ref{prop:fuzzy-way-below-formula}, the end defining
\(\ll_X(a,b)\) maps to its component at the principal ideal \(\eta_X b\):
\[
    \ll_X(a,b)
    \to
    \left[
        X(b,\alpha_X\eta_X b),
        \operatorname{Idl}_{\mathbb V}(X)(\eta_X a,\eta_X b)
    \right].
\]
The strict algebra unit law gives
\[
    \alpha_X\eta_X b=b.
\]
The Yoneda embedding is fully faithful, so
\[
    \operatorname{Idl}_{\mathbb V}(X)(\eta_X a,\eta_X b)
    \cong
    X(a,b).
\]
Thus we have a map
\[
    \ll_X(a,b)
    \to
    [X(b,b),X(a,b)].
\]
The unit map
\[
    I\to X(b,b)
\]
induces, by contravariance of the first variable of the internal hom, a morphism
\[
    [X(b,b),X(a,b)]
    \to
    [I,X(a,b)]
    \cong
    X(a,b).
\]
The composite gives
\[
    \ll_X(a,b)\to X(a,b).
\]
\end{proof}

The proof follows the diagram
\[
\begin{tikzcd}[column sep=huge,row sep=large]
{\ll_X(a,b)}
    \ar[r]
    \ar[rr,bend right=18]
&
{\left[
    X(b,\alpha_X\eta_X b),
    \operatorname{Idl}_{\mathbb V}(X)(\eta_X a,\eta_X b)
\right]}
    \ar[r,"{\alpha_X\eta_X b=b}"]
&
{\left[
    X(b,b),
    X(a,b)
\right]}
    \ar[r]
&
{X(a,b)}.
\end{tikzcd}
\]

\subsection{The posetal reflection}
\label{subsec:fuzzy-way-below-posetal}

For a fuzzy domain over \(\Lambda\),
\[
    \ll_X(a,b)
    =
    \bigmeet_{I\in\operatorname{Idl}_\Lambda(X)}
    \left(
        X(b,\alpha_X I)\multimap I(a)
    \right).
\]
The distributor law for
\[
    \ll_X:X\nrightarrow X
\]
gives fuzzy monotonicity:
\[
    X(a',a)\tensor \ll_X(a,b)\tensor X(b,b')
    \leq
    \ll_X(a',b').
\]
The comparison of Proposition~\ref{prop:fuzzy-way-below-implies-below} becomes
\[
    \ll_X(a,b)\leq X(a,b).
\]
Thus the way-below degree is always bounded above by the ordinary fuzzy comparison
degree.

In the Lawvere case this inequality is in the truth-value order. Since the
truth-value order is the reverse of the usual numerical order, it reads
numerically as
\[
    \ll_X(a,b)\geq d_X(a,b).
\]

\subsection{Examples}
\label{subsec:fuzzy-way-below-examples}

\begin{example}[The crisp case]
\label{ex:fuzzy-way-below-crisp}
For \(\Lambda=2\), the formula becomes
\[
    a\ll_X b
    \quad\Longleftrightarrow\quad
    \forall I\in\operatorname{Idl}(X),\;
    b\leq \bigvee I\Rightarrow a\in I.
\]
This is the ordinary way-below relation on a dcpo.
\end{example}

\begin{example}[G\"odel way-below degrees]
\label{ex:fuzzy-way-below-godel}
For G\"odel truth values,
\[
    \ll_X(a,b)
    =
    \bigmeet_{I}
    \left(
        X(b,\alpha_X I)\multimap_{\mathrm G} I(a)
    \right).
\]
This is the degree to which every ideal whose join lies above \(b\) contains
\(a\).
\end{example}

\begin{example}[{\L}ukasiewicz way-below degrees]
\label{ex:fuzzy-way-below-lukasiewicz}
For {\L}ukasiewicz truth values,
\[
    \ll_X(a,b)
    =
    \bigmeet_{I}
    \min(1,1-X(b,\alpha_X I)+I(a)).
\]
\end{example}

\begin{example}[Product way-below degrees]
\label{ex:fuzzy-way-below-product}
For product truth values,
\[
    \ll_X(a,b)
    =
    \bigmeet_I
    \left(
        X(b,\alpha_X I)\multimap_\Pi I(a)
    \right).
\]
\end{example}

\begin{example}[Lawvere way-below costs]
\label{ex:fuzzy-way-below-lawvere}
For the Lawvere quantale, the meet in the truth-value order is a supremum in the
usual numerical order. Hence
\[
    \ll_X(a,b)
    =
    \sup_{I\in\operatorname{Idl}_\Lambda(X)}
    \bigl(I(a)\ominus X(b,\alpha_X I)\bigr),
\]
where
\[
    v\ominus u=\max(0,v-u)
\]
with the extended-value conventions. Thus \(\ll_X(a,b)\) measures the
worst-case residual cost needed to turn evidence that \(b\) lies below an ideal
colimit into evidence that \(a\) belongs to that ideal.
\end{example}

\section{Continuity as representability of way-below}
\label{sec:continuity-representability-way-below}

We now show that the previous chapter's way-below formula is precisely the
representable form of the right-lifted distributor.

\subsection{The enriched construction}
\label{subsec:continuity-representability-enriched}

\begin{definition}[Continuous fuzzy domain, recalled]
\label{def:continuous-fuzzy-domain-recalled-equipment}
A fuzzy domain \(X\) is \textbf{continuous} if its ideal-algebra map
\[
    \alpha_X:\operatorname{Idl}_{\mathbb V}(X)\to X
\]
has an enriched left adjoint
\[
    \beta_X:X\to\operatorname{Idl}_{\mathbb V}(X),
    \qquad
    \beta_X\dashv\alpha_X.
\]
\end{definition}

The continuity data form the adjoint diagram
\[
\begin{tikzcd}[column sep=huge,row sep=large]
X
    \ar[r,bend left=35,"\beta_X"]
&
\operatorname{Idl}_{\mathbb V}(X)
    \ar[l,bend left=35,"\alpha_X"]
\end{tikzcd}
\]
and hence give natural isomorphisms
\[
    X(b,\alpha_X I)
    \cong
    \operatorname{Idl}_{\mathbb V}(X)(\beta_X b,I).
\]

\begin{theorem}[Continuity represents the way-below distributor]
\label{thm:continuity-represents-way-below}
Let \(X\) be a continuous fuzzy domain. Then
\[
    \ll_X(a,b)
    \cong
    \operatorname{Idl}_{\mathbb V}(X)(\eta_X a,\beta_X b).
\]
Consequently,
\[
    \ll_X(a,b)\cong \beta_X(b)(a).
\]
Equivalently, as distributors,
\[
    \ll_X
    \cong
    \beta_X^\ast\circ(\eta_X)_\ast.
\]
\end{theorem}

\begin{proof}
By Proposition~\ref{prop:fuzzy-way-below-formula},
\[
    \ll_X(a,b)
    \cong
    \int_I
    \left[
        X(b,\alpha_X I),
        \operatorname{Idl}_{\mathbb V}(X)(\eta_X a,I)
    \right].
\]
Since
\[
    \beta_X\dashv\alpha_X,
\]
we have
\[
    X(b,\alpha_X I)
    \cong
    \operatorname{Idl}_{\mathbb V}(X)(\beta_X b,I).
\]
Therefore
\[
    \ll_X(a,b)
    \cong
    \int_I
    \left[
        \operatorname{Idl}_{\mathbb V}(X)(\beta_X b,I),
        \operatorname{Idl}_{\mathbb V}(X)(\eta_X a,I)
    \right].
\]
By enriched Yoneda in \(\operatorname{Idl}_{\mathbb V}(X)\),
\[
    \int_I
    \left[
        \operatorname{Idl}_{\mathbb V}(X)(\beta_X b,I),
        \operatorname{Idl}_{\mathbb V}(X)(\eta_X a,I)
    \right]
    \cong
    \operatorname{Idl}_{\mathbb V}(X)(\eta_X a,\beta_X b).
\]
Applying Yoneda again to the ideal
\[
    \beta_X b:X^{\op}\to\underline{\mathbb V}
\]
gives
\[
    \operatorname{Idl}_{\mathbb V}(X)(\eta_X a,\beta_X b)
    \cong
    \beta_X(b)(a).
\]

For the distributor-level statement, compute
\[
\begin{aligned}
(\beta_X^\ast\circ(\eta_X)_\ast)(a,b)
&=
\int^{I\in\operatorname{Idl}_{\mathbb V}(X)}
\operatorname{Idl}_{\mathbb V}(X)(\eta_Xa,I)
\tensor
\operatorname{Idl}_{\mathbb V}(X)(I,\beta_Xb)
\\
&\cong
\operatorname{Idl}_{\mathbb V}(X)(\eta_Xa,\beta_Xb)
\\
&\cong
\ll_X(a,b),
\end{aligned}
\]
by enriched co-Yoneda and the calculation above.
\end{proof}

Thus the approximant map
\[
    \beta_X:X\to\operatorname{Idl}_{\mathbb V}(X)
\]
represents the right-lifted way-below distributor.

\subsection{The posetal reflection}
\label{subsec:continuity-representability-posetal}

For a continuous fuzzy domain over \(\Lambda\),
\[
    \ll_X(a,b)
    =
    \operatorname{Idl}_\Lambda(X)(\eta_X a,\beta_X b)
    =
    \beta_X(b)(a).
\]
Thus the way-below degree
\[
    \operatorname{wb}_X(a,b)=\beta_X(b)(a)
\]
from the previous chapter is exactly the value of the right-lifted distributor
constructed in this chapter.

If \(X\) is separated, then the reconstruction theorem from the previous chapter
gives
\[
    \alpha_X\beta_X=1_X.
\]
For completeness, we recall the argument. The unit of
\[
    \beta_X\dashv\alpha_X
\]
gives
\[
    1_X\leq \alpha_X\beta_X.
\]
The counit gives
\[
    \beta_X\alpha_X\leq 1_{\operatorname{Idl}_\Lambda(X)}.
\]
Applying this to the principal ideal \(\eta_X b\), and using the algebra law
\[
    \alpha_X\eta_X b=b,
\]
gives
\[
    \beta_X b\leq \eta_X b.
\]
Applying \(\alpha_X\) gives
\[
    \alpha_X\beta_X b
    \leq
    \alpha_X\eta_X b
    =
    b.
\]
Together with
\[
    b\leq \alpha_X\beta_X b,
\]
this gives equality in the induced crisp preorder. Since \(X\) is separated,
\[
    \alpha_X\beta_X b=b.
\]

The reconstruction is captured by
\[
\begin{tikzcd}[column sep=huge,row sep=large]
X
    \ar[r,"\beta_X"]
    \ar[dr,"1_X"']
&
\operatorname{Idl}_\Lambda(X)
    \ar[d,"\alpha_X"]
\\
&
X
\end{tikzcd}
\]
in the separated posetal reflection.

Hence every object \(b\) of a separated continuous fuzzy domain is recovered as
the fuzzy ideal colimit of its approximant ideal:
\[
    b=\alpha_X(\beta_X b).
\]

\subsection{Examples}
\label{subsec:continuity-representability-examples}

\begin{example}[The crisp case]
\label{ex:continuity-representability-crisp}
For \(\Lambda=2\), continuity means that
\[
    \alpha_X:\operatorname{Idl}(X)\to X
\]
has a left adjoint
\[
    \beta_X:X\to\operatorname{Idl}(X).
\]
Then
\[
    \beta_X(b)=\{a\in X\mid a\ll b\},
\]
and
\[
    a\ll b
    \quad\Longleftrightarrow\quad
    a\in \beta_X(b).
\]
The reconstruction equation
\[
    \alpha_X\beta_X(b)=b
\]
says that every element is the directed supremum of its way-below approximants.
\end{example}

\begin{example}[G\"odel continuous fuzzy domains]
\label{ex:continuity-representability-godel}
For G\"odel truth values,
\[
    \operatorname{wb}_X(a,b)
    =
    \beta_X(b)(a)
\]
is the G\"odel-valued degree to which \(a\) approximates \(b\).
\end{example}

\begin{example}[{\L}ukasiewicz continuous fuzzy domains]
\label{ex:continuity-representability-lukasiewicz}
For {\L}ukasiewicz truth values,
\[
    \operatorname{wb}_X(a,b)
    =
    \beta_X(b)(a)
\]
is an MV-valued approximation degree. The right-lifting formula computes it
using the {\L}ukasiewicz residual.
\end{example}

\begin{example}[Product continuous fuzzy domains]
\label{ex:continuity-representability-product}
For product truth values,
\[
    \operatorname{wb}_X(a,b)
    =
    \beta_X(b)(a),
\]
and the comparison between ideal-colimit evidence and ideal-membership evidence
is governed by the product residual.
\end{example}

\begin{example}[Lawvere continuous fuzzy domains]
\label{ex:continuity-representability-lawvere}
For the Lawvere quantale,
\[
    \beta_X(b)(a)
\]
is the cost with which \(a\) approximates \(b\). Continuity says that the
right-lifted approximation cost is represented by the finite-flat weight
\[
    \beta_X(b).
\]
\end{example}

\section{Scott opens through way-below approximants}
\label{sec:scott-opens-through-way-below}

The previous chapter defined fuzzy Scott opens as upper fuzzy predicates
preserving fuzzy ideal joins. The way-below distributor reconstructs Scott opens
from their values on approximants.

\subsection{The enriched construction}
\label{subsec:scott-opens-way-below-enriched}

Let \(X\) be a weak fuzzy domain with structure map
\[
    \alpha_X:\operatorname{Idl}_{\mathbb V}(X)\to X.
\]
A fuzzy Scott open is an upper fuzzy predicate
\[
    U:X\to\underline{\mathbb V}
\]
satisfying
\[
    U(\alpha_X I)
    \cong
    \int^{x\in X}
    I(x)\tensor U(x)
\]
for every fuzzy ideal \(I\).

The right-hand side is the enriched pairing of the lower weight
\[
    I:X^{\op}\to\underline{\mathbb V}
\]
with the upper predicate
\[
    U:X\to\underline{\mathbb V}.
\]
Thus Scott openness says that \(U\) sends the chosen ideal colimit of \(I\) to
the corresponding weighted colimit of its values. The pairing is represented by
\[
\begin{tikzcd}[column sep=huge,row sep=large]
X^{\op}
    \ar[r,"I"]
&
\underline{\mathbb V}
&
X
    \ar[l,"U"']
\\
&
\int^{x\in X} I(x)\tensor U(x)
    \ar[u,dashed]
\end{tikzcd}
\]
and the Scott condition by
\[
\begin{tikzcd}[column sep=huge,row sep=large]
\operatorname{Idl}_{\mathbb V}(X)
    \ar[r,"\alpha_X"]
    \ar[dr,"{I\mapsto\int^{x}I(x)\tensor U(x)}"']
&
X
    \ar[d,"U"]
\\
&
\underline{\mathbb V}.
\end{tikzcd}
\]

If an enriched reconstruction equivalence
\[
    \alpha_X\beta_X\cong 1_X
\]
is available, then the Scott-open condition gives
\[
    U(b)
    \cong
    U(\alpha_X(\beta_X b))
    \cong
    \int^{a\in X}
        \beta_X(b)(a)\tensor U(a)
    \cong
    \int^{a\in X}
        \ll_X(a,b)\tensor U(a).
\]
In the posetal separated case, the reconstruction equality was obtained in the
previous section, so the formula becomes an equality of truth values.

The reconstruction path is
\[
\begin{tikzcd}[column sep=huge,row sep=large]
X
    \ar[r,"\beta_X"]
    \ar[rr,bend left=22,"1_X"]
&
\operatorname{Idl}_{\mathbb V}(X)
    \ar[r,"\alpha_X"]
    \ar[d,"{I\mapsto\int^{a}I(a)\tensor U(a)}"']
&
X
    \ar[d,"U"]
\\
&
\underline{\mathbb V}
    \ar[r,equal]
&
\underline{\mathbb V}.
\end{tikzcd}
\]

\subsection{The posetal reflection}
\label{subsec:scott-opens-way-below-posetal}

In the posetal specialization, a fuzzy Scott open is an upper fuzzy predicate
\[
    U:X\to L
\]
satisfying
\[
    U(\alpha_X I)
    =
    \bigjoin_{x\in X}
    I(x)\tensor U(x)
\]
for every fuzzy ideal \(I\).

\begin{proposition}[Scott opens are reconstructed from way-below approximants]
\label{prop:scott-open-reconstruction-way-below}
Let \(X\) be a separated continuous fuzzy domain over \(\Lambda\), and let
\[
    U:X\to L
\]
be a fuzzy Scott open. Then for every \(b\in X\),
\[
    U(b)
    =
    \bigjoin_{a\in X}
    \ll_X(a,b)\tensor U(a).
\]
\end{proposition}

\begin{proof}
Since \(X\) is separated and continuous,
\[
    b=\alpha_X(\beta_X b).
\]
Since \(U\) is Scott open,
\[
    U(b)
    =
    U(\alpha_X(\beta_X b))
    =
    \bigjoin_{a\in X}
    \beta_X(b)(a)\tensor U(a).
\]
By Theorem~\ref{thm:continuity-represents-way-below},
\[
    \beta_X(b)(a)=\ll_X(a,b).
\]
Substituting gives
\[
    U(b)
    =
    \bigjoin_{a\in X}
    \ll_X(a,b)\tensor U(a).
\]
\end{proof}

Equivalently,
\[
    U(b)
    =
    \bigjoin_{a\in X}
    \beta_X(b)(a)\tensor U(a).
\]

\subsection{Examples}
\label{subsec:scott-opens-way-below-examples}

\begin{example}[The crisp case]
\label{ex:scott-opens-way-below-crisp}
For \(\Lambda=2\), if \(X\) is a continuous dcpo and \(U\subseteq X\) is Scott
open, then
\[
    b\in U
    \quad\Longleftrightarrow\quad
    \exists a\in X,\; a\ll b\text{ and }a\in U.
\]
Thus Scott opens are determined by way-below approximants.
\end{example}

\begin{example}[G\"odel Scott opens]
\label{ex:scott-opens-way-below-godel}
For G\"odel truth values, the reconstruction formula is
\[
    U(b)
    =
    \bigjoin_{a\in X}
    \min(\ll_X(a,b),U(a)).
\]
\end{example}

\begin{example}[{\L}ukasiewicz Scott opens]
\label{ex:scott-opens-way-below-lukasiewicz}
For {\L}ukasiewicz truth values,
\[
    U(b)
    =
    \bigjoin_{a\in X}
    \max(0,\ll_X(a,b)+U(a)-1).
\]
\end{example}

\begin{example}[Product Scott opens]
\label{ex:scott-opens-way-below-product}
For product truth values,
\[
    U(b)
    =
    \bigjoin_{a\in X}
    \ll_X(a,b)U(a).
\]
\end{example}

\begin{example}[Lawvere Scott opens]
\label{ex:scott-opens-way-below-lawvere}
For the Lawvere quantale, the Scott reconstruction formula becomes
\[
    U(b)
    =
    \inf_{a\in X}
    \bigl(\ll_X(a,b)+U(a)\bigr),
\]
written in the usual numerical order. Thus the cost of satisfying the Scott
predicate at \(b\) is obtained by optimizing over approximants \(a\) of \(b\).
\end{example}

\section{Summary}
\label{sec:summary-fuzzy-distributors-way-below}

This chapter constructed the fuzzy distributor equipment and used it to explain
the origin of the fuzzy way-below relation.

The construction starts with the data already present in the fuzzy ideal monad:
\[
    \eta_X:X\to\operatorname{Idl}_{\mathbb V}(X),
    \qquad
    \alpha_X:\operatorname{Idl}_{\mathbb V}(X)\to X.
\]
The first is the principal-ideal embedding, and the second is the algebra map of a
fuzzy domain. Their strict algebra interaction is
\[
\begin{tikzcd}[column sep=huge,row sep=large]
X
    \ar[r,"\eta_X"]
    \ar[dr,"1_X"']
&
\operatorname{Idl}_{\mathbb V}(X)
    \ar[d,"\alpha_X"]
\\
&
X.
\end{tikzcd}
\]

Fuzzy distributors
\[
    \Phi:X\nrightarrow Y
\]
are enriched profunctors
\[
    X^{\op}\tensor Y\to\underline{\mathbb V}.
\]
They compose by coends, admit companions and conjoints for fuzzy functors, and
have right liftings
\[
    [\Phi,\Psi](b,a)
    =
    \int_c
    [\Phi(a,c),\Psi(b,c)].
\]
The distributor equipment has the schematic form
\[
\begin{tikzcd}[column sep=huge,row sep=large]
X
    \ar[d,"f"']
    \ar[r,dashed,"\Phi"]
&
Y
    \ar[d,"g"]
\\
X'
    \ar[r,dashed,"\Psi"']
&
Y'
\ar[phantom,from=1-2,to=2-1,"\Downarrow" description]
\end{tikzcd}
\]
where vertical arrows are fuzzy functors, horizontal arrows are fuzzy
distributors, and cells are distributor transformations.

For a fuzzy domain \(X\), the conjoint of \(\alpha_X\) and the companion of
\(\eta_X\) are distributors
\[
    \alpha_X^\ast,
    \qquad
    (\eta_X)_\ast
    :
    X\nrightarrow \operatorname{Idl}_{\mathbb V}(X).
\]
Their right lifting defines the fuzzy way-below distributor:
\[
    \ll_X
    =
    [\alpha_X^\ast,(\eta_X)_\ast]
    :
    X\nrightarrow X.
\]
The defining triangle is
\[
\begin{tikzcd}[column sep=huge,row sep=large]
X
    \ar[r,dashed,"\ll_X"]
    \ar[rr,bend right=18,dashed,"{(\eta_X)_\ast}"']
&
X
    \ar[r,dashed,"\alpha_X^\ast"]
&
\operatorname{Idl}_{\mathbb V}(X).
\end{tikzcd}
\]
Its objectwise formula is
\[
    \ll_X(a,b)
    \cong
    \int_{I\in \operatorname{Idl}_{\mathbb V}(X)}
    \left[
        X(b,\alpha_X I),
        I(a)
    \right].
\]
In the posetal reflection this becomes
\[
    \ll_X(a,b)
    =
    \bigmeet_{I\in\operatorname{Idl}_\Lambda(X)}
    \left(
        X(b,\alpha_X I)\multimap I(a)
    \right).
\]

Finally, if \(X\) is continuous, with
\[
    \beta_X\dashv\alpha_X,
\]
then the right-lifted way-below distributor is represented by \(\beta_X\):
\[
    \ll_X(a,b)
    \cong
    \operatorname{Idl}_{\mathbb V}(X)(\eta_X a,\beta_X b)
    \cong
    \beta_X(b)(a).
\]
Equivalently, as distributors,
\[
    \ll_X
    \cong
    \beta_X^\ast\circ(\eta_X)_\ast.
\]
Thus the way-below degree introduced in the previous chapter is not extra
structure: it is the value of the right-lifted distributor represented by the
approximant map of a continuous fuzzy domain.
    %\chapter{Recursive Domain Equations by Finite-Flat Cofinal Approximation}
\label{ch:recursive-domain-equations-finite-flat}

In the preceding chapters we constructed the fuzzy ideal monad
\[
    T_{\mathbb V}X
    =
    \operatorname{Idl}_{\mathbb V}(X),
\]
where \(\operatorname{Idl}_{\mathbb V}(X)\) is the full
\(\mathbb V\)-subcategory of
\[
    [X^{\op},\underline{\mathbb V}]
\]
spanned by the finite-flat weights
\[
    I:X^{\op}\to\underline{\mathbb V}.
\]
Here \(\mathbb V\) denotes the fuzzy cosmos and
\(\underline{\mathbb V}\) denotes its self-enriched
\(\mathbb V\)-category of truth values.

Strict algebras for this monad are the fuzzy domains.  Equivalently, a fuzzy
domain is a small \(\mathbb V\)-category equipped with a strict coherent choice
of finite-flat weighted colimits, and a fuzzy-domain morphism is a
finite-flat-colimit-preserving \(\mathbb V\)-functor.

We also use the symmetric monoidal closed structure on fuzzy domains:
\[
    A\boxtimes_{\mathrm{ff}}B
    \dashv
    [B,-]_{\mathrm{ff}}.
\]
Thus
\[
    \mathbf{FDom}_{\mathbb V}
    (A\boxtimes_{\mathrm{ff}}B,C)
    \cong
    \mathbf{FDom}_{\mathbb V}
    \bigl(A,[B,C]_{\mathrm{ff}}\bigr),
\]
where \([B,C]_{\mathrm{ff}}\) is the fuzzy domain of
finite-flat-colimit-preserving \(\mathbb V\)-functors \(B\to C\).

The purpose of this chapter is to formulate recursive domain equation
constructions inside the finite-flat doctrine.  The classical construction by
\(\omega\)-chain bilimits in dcpos is recovered as the crisp separated
specialization.  The general replacement for an ordinary directed indexing
chain is a finite-flat approximation shape
\[
    (A,W),
    \qquad
    W:A^{\op}\to\underline{\mathbb V}.
\]
Cofinality is expressed by lower direct image of weights.  Tails and common
successors are not chosen arbitrarily; they are given by canonical
comma-style constructions, and the relevant exactness conditions say that
these canonical constructions are cofinal in the finite-flat sense.

The semantic target in the general resource-sensitive setting is not an
ordinary cartesian reflexive object
\[
    D\simeq [D,D]_{\mathrm{ff}}.
\]
In an affine enriched setting there is no canonical diagonal functor
\[
    A\to A\tensor A,
\]
since such a functor would require canonical contraction morphisms
\[
    A(a,a')\to A(a,a')\tensor A(a,a').
\]
Affinity gives weakening morphisms
\[
    A(a,a')\to I,
\]
not contraction into tensor squares.

Consequently the general internal-hom approximation theorem is
external-product-indexed:
\[
    (A\tensor A,W\boxtimes W).
\]
The same-index collapse is recovered only in the crisp separated/cartesian
specialization, where the ordinary diagonal
\[
    \omega\to\omega\times\omega
\]
exists and is cofinal for the ordinary directed ideal.

\section{Finite-flat approximation shapes}
\label{sec:finite-flat-approximation-shapes}

Throughout this chapter, finite-flatness is understood relative to the
finite-limit doctrine used in the definition of categorical fuzzy frames.  Thus
in a predicate category \([X,\underline{\mathbb V}]\), the specified finite
limits are computed pointwise.  In the posetal scalar-frame specialization this
means preservation of the terminal predicate and binary meets.

\begin{definition}[Finite-flat approximation shape]
\label{def:finite-flat-approximation-shape}
A \emph{finite-flat approximation shape} is a pair
\[
    (A,W),
\]
where \(A\) is a small \(\mathbb V\)-category and
\[
    W:A^{\op}\to\underline{\mathbb V}
\]
is a finite-flat weight.  Equivalently,
\[
    W\in\operatorname{Idl}_{\mathbb V}(A).
\]

If \(X\) is a fuzzy domain and
\[
    D:A\to X
\]
is a \(\mathbb V\)-functor, the corresponding approximation object is the
finite-flat weighted colimit
\[
    W*D.
\]
\end{definition}

Thus the passage from ordinary directed joins to fuzzy-domain approximation is
the passage
\[
    \bigvee_{a\in A}D_a
    \quad\leadsto\quad
    W*D.
\]

\begin{definition}[Lower direct image of a weight]
\label{def:lower-direct-image-weight-rde}
Let
\[
    h:A\to B
\]
be a \(\mathbb V\)-functor, and let
\[
    W:A^{\op}\to\underline{\mathbb V}
\]
be a weight.  The lower direct image of \(W\) along \(h\) is the weight
\[
    h_!W:B^{\op}\to\underline{\mathbb V}
\]
defined by
\[
    (h_!W)(b)
    =
    \int^{a\in A}
        W(a)\tensor B(b,ha).
\]
Equivalently, \(h_!W\) is the weighted left Kan extension of \(W\) along
\[
    h^{\op}:A^{\op}\to B^{\op}.
\]
\end{definition}

By the finite-flat direct-image theorem for the ideal doctrine, if \(W\) is
finite-flat, then \(h_!W\) is finite-flat.

\begin{definition}[Cofinal map of approximation shapes]
\label{def:cofinal-map-approximation-shapes}
Let
\[
    (A,W),
    \qquad
    (B,U)
\]
be finite-flat approximation shapes.  A \(\mathbb V\)-functor
\[
    h:A\to B
\]
is \emph{cofinal from \((A,W)\) to \((B,U)\)} if the canonical comparison
exhibits an isomorphism of weights
\[
    h_!W\cong U.
\]
When \(A=B\) and \(U=W\), we say that \(h\) is \(W\)-cofinal.
\end{definition}

Thus cofinality is equality of approximation weights after lower direct image,
not merely a pointwise finality condition on objects.

\begin{proposition}[Cofinality for finite-flat weighted colimits]
\label{prop:finite-flat-cofinality}
Let
\[
    h:(A,W)\to(B,U)
\]
be a cofinal map of finite-flat approximation shapes.  Then for every fuzzy
domain \(X\) and every diagram
\[
    D:B\to X,
\]
there is a canonical isomorphism
\[
    U*D
    \cong
    W*(Dh).
\]
\end{proposition}

\begin{proof}
Since \(h\) is cofinal,
\[
    U\cong h_!W.
\]
For every \(x\in X\), the universal property of weighted colimits and the
weighted left-Kan-extension adjunction give
\[
\begin{aligned}
    X((h_!W)*D,x)
    &\cong
    [B^{\op},\underline{\mathbb V}]
    (h_!W,X(D-,x))
    \\
    &\cong
    [A^{\op},\underline{\mathbb V}]
    (W,X(Dh-,x))
    \\
    &\cong
    X(W*(Dh),x).
\end{aligned}
\]
By enriched Yoneda,
\[
    (h_!W)*D\cong W*(Dh).
\]
Using \(U\cong h_!W\), this gives
\[
    U*D\cong W*(Dh).
\]
\end{proof}

\begin{definition}[External product of approximation shapes]
\label{def:external-product-approximation-shapes-rde}
Let
\[
    (A,W),
    \qquad
    (B,Q)
\]
be finite-flat approximation shapes.  Their external product is
\[
    (A,W)\boxtimes(B,Q)
    =
    (A\tensor B,W\boxtimes Q),
\]
where
\[
    (W\boxtimes Q)(a,b)
    =
    W(a)\tensor Q(b).
\]
By closure of finite-flat weights under external products,
\[
    W\boxtimes Q
\]
is finite-flat.
\end{definition}

\begin{lemma}[Constant diagrams]
\label{lem:constant-diagrams-finite-flat}
Let
\[
    (A,W)
\]
be a finite-flat approximation shape.  For every fuzzy domain \(X\) and every
object \(x\in X\), the \(W\)-weighted colimit of the constant diagram at \(x\)
is canonically \(x\):
\[
    W*\Delta x\cong x.
\]
\end{lemma}

\begin{proof}
Let
\[
    !:A\to\mathbb 1
\]
be the terminal functor, and let
\[
    \bar x:\mathbb 1\to X
\]
classify the object \(x\).  The constant diagram at \(x\) is
\[
    \Delta x=\bar x!.
\]

Since \(W\) is finite-flat, its character
\[
    \chi_W:[A,\underline{\mathbb V}]\to\underline{\mathbb V}
\]
preserves the terminal predicate.  Let
\[
    \top_A:A\to\underline{\mathbb V}
\]
be the constant upper predicate at \(I\).  Then
\[
\begin{aligned}
    (!_!W)(*)
    &=
    \int^{a\in A}W(a)\tensor \mathbb 1(*,!a)
    \\
    &\cong
    \int^{a\in A}W(a)\tensor I
    \\
    &=
    \chi_W(\top_A)
    \\
    &\cong
    I.
\end{aligned}
\]
Thus
\[
    !_!W\cong \mathbb 1(-,*).
\]
By Proposition~\ref{prop:finite-flat-cofinality},
\[
    W*(\bar x!)
    \cong
    (!_!W)*\bar x
    \cong
    \mathbb 1(-,*)*\bar x
    \cong
    x.
\]
\end{proof}

\begin{lemma}[Functoriality of finite-flat colimits in diagrams]
\label{lem:ff-colimits-functorial-in-diagrams}
Let \(C\) be a fuzzy domain, let
\[
    R:P^{\op}\to\underline{\mathbb V}
\]
be finite-flat, and let
\[
    H,H':P\to C
\]
be \(\mathbb V\)-functors.  A \(\mathbb V\)-natural transformation
\[
    \gamma:H\Rightarrow H'
\]
induces a canonical morphism
\[
    R*H\longrightarrow R*H'
\]
in \(C\).  This construction is functorial in \(\gamma\).
\end{lemma}

\begin{proof}
For each object \(c\in C\), the universal property of the weighted colimit gives
natural isomorphisms
\[
    C(R*H,c)
    \cong
    [P^{\op},\underline{\mathbb V}]
    (R,C(H-,c))
\]
and
\[
    C(R*H',c)
    \cong
    [P^{\op},\underline{\mathbb V}]
    (R,C(H'-,c)).
\]
The natural transformation
\[
    \gamma:H\Rightarrow H'
\]
induces by precomposition a \(\mathbb V\)-natural transformation
\[
    C(H'-,c)\Rightarrow C(H-,c).
\]
Transporting across the representing isomorphisms gives a morphism
\[
    C(R*H',c)\longrightarrow C(R*H,c)
\]
natural in \(c\).  By enriched Yoneda this family is represented by a unique
morphism
\[
    R*H\longrightarrow R*H'.
\]
Functoriality follows from functoriality of precomposition and uniqueness in
the representing property.
\end{proof}

\section{Support categories and projection-pair systems}
\label{sec:support-categories-split-systems}

Projection-pair systems indexed by an enriched category use global support
arrows for the actual transition projection pairs.  The remaining enriched
hom-structure enters through transition-strength morphisms.

\begin{definition}[Support category of an enriched index]
\label{def:support-category-enriched-index}
Let \(A\) be a small \(\mathbb V\)-category.  Its support category
\[
    A_0
\]
is the ordinary category with the same objects as \(A\), and hom-set
\[
    A_0(a,a')=\mathbb V(I,A(a,a')).
\]
If
\[
    u:I\to A(a,a')
    \qquad\text{and}\qquad
    v:I\to A(a',a'')
\]
are arrows in \(A_0\), their composite is the global element
\[
    I\cong I\tensor I
    \xrightarrow{\,v\tensor u\,}
    A(a',a'')\tensor A(a,a')
    \longrightarrow
    A(a,a'')
\]
induced by composition in \(A\).  Identities are the enriched identity maps
\[
    I\to A(a,a).
\]

In the posetal specialization, \(A_0(a,a')\) has a unique arrow precisely when
\[
    A(a,a')=\top.
\]
\end{definition}

\begin{lemma}[Support projections from an enriched tensor product]
\label{lem:support-projections-tensor-product}
Let \(A\) and \(B\) be small \(\mathbb V\)-categories.  There is a canonical
ordinary functor
\[
    (A\tensor B)_0\to A_0\times B_0.
\]
On objects it sends
\[
    (a,b)\longmapsto(a,b).
\]
On arrows, a global element
\[
    \xi:I\to A(a,a')\tensor B(b,b')
\]
is sent to the pair
\[
    \xi_A:I\to A(a,a'),
    \qquad
    \xi_B:I\to B(b,b'),
\]
where \(\xi_A\) is obtained by composing \(\xi\) with
\[
    A(a,a')\tensor B(b,b')
    \xrightarrow{1\tensor !}
    A(a,a')\tensor I
    \cong
    A(a,a'),
\]
and \(\xi_B\) is obtained similarly using
\[
    A(a,a')\tensor B(b,b')
    \xrightarrow{!\tensor 1}
    I\tensor B(b,b')
    \cong
    B(b,b').
\]
\end{lemma}

\begin{proof}
The terminality of \(I\) gives the weakening maps
\[
    B(b,b')\to I,
    \qquad
    A(a,a')\to I.
\]
Thus the two displayed projections are canonical.

Compatibility with identities is immediate from the unit constraints.  For
composition, the composition in \(A\tensor B\) is induced by the tensor of the
compositions in \(A\) and \(B\).  Applying the two weakening projections sends
this composite to the ordinary composites in \(A_0\) and \(B_0\).  Hence the
assignment defines an ordinary functor.
\end{proof}

\begin{definition}[The locally ordinary \(2\)-category of fuzzy domains]
\label{def:locally-ordinary-2-category-fdom}
Let
\[
    \mathbf{FDom}^{(2)}_{\mathbb V}
\]
be the locally ordinary \(2\)-category defined as follows.

Its objects are fuzzy domains.  Its \(1\)-cells are fuzzy-domain morphisms.
For fuzzy domains \(X,Y\), the hom-category is the support category of the
internal-hom fuzzy domain:
\[
    \mathbf{FDom}^{(2)}_{\mathbb V}(X,Y)
    =
    \bigl([X,Y]_{\mathrm{ff}}\bigr)_0.
\]
Thus a \(2\)-cell
\[
    \alpha:f\Rightarrow g
\]
is a global element
\[
    I\to [X,Y]_{\mathrm{ff}}(f,g).
\]

Horizontal composition is induced by the finite-flat-continuous composition
bimorphism
\[
    [Y,Z]_{\mathrm{ff}}\boxtimes_{\mathrm{ff}}[X,Y]_{\mathrm{ff}}
    \longrightarrow
    [X,Z]_{\mathrm{ff}}.
\]
Vertical composition is inherited from the support categories of the internal
homs.
\end{definition}

\begin{definition}[Fuzzy projection pair]
\label{def:fuzzy-projection-pair-ff}
Let \(X\) and \(Y\) be fuzzy domains.  A \emph{fuzzy projection pair}
\[
    (e,p):X\triangleleft Y
\]
consists of fuzzy-domain morphisms
\[
    e:X\to Y,
    \qquad
    p:Y\to X,
\]
together with an adjunction
\[
    e\dashv p
\]
in \(\mathbf{FDom}^{(2)}_{\mathbb V}\), whose unit
\[
    \eta:\id_X\Rightarrow pe
\]
is invertible.

Equivalently, after choosing the inverse of the unit, one has
\[
    pe\cong \id_X
\]
and a counit comparison
\[
    ep\Rightarrow \id_Y
\]
satisfying the triangle identities.  The morphism \(e\) is called the
embedding and \(p\) the projection.
\end{definition}

\begin{lemma}[Composition of fuzzy projection pairs]
\label{lem:composition-fuzzy-projection-pairs-ff}
Fuzzy projection pairs compose.  If
\[
    (e,p):X\triangleleft Y
\]
and
\[
    (e',p'):Y\triangleleft Z
\]
are fuzzy projection pairs, then
\[
    (e'e,pp'):X\triangleleft Z
\]
is a fuzzy projection pair.
\end{lemma}

\begin{proof}
The composite of adjunctions
\[
    e\dashv p,
    \qquad
    e'\dashv p'
\]
is an adjunction
\[
    e'e\dashv pp'.
\]
Its unit is the composite
\[
    \id_X
    \Rightarrow
    pe
    \Rightarrow
    pp'e'e.
\]
Both displayed unit maps are invertible, hence the composite unit is invertible.
The counit and the triangle identities are those of the composite adjunction.
\end{proof}

\begin{definition}[Projection-pair 2-cells]
\label{def:projection-pair-2-cells}
Let
\[
    (e,p),(e',p'):X\triangleleft Y
\]
be fuzzy projection pairs.  A 2-cell
\[
    (\alpha,\beta):(e,p)\Rightarrow(e',p')
\]
consists of enriched natural transformations
\[
    \alpha:e\Rightarrow e',
    \qquad
    \beta:p'\Rightarrow p
\]
which are mates under the adjunctions
\[
    e\dashv p,
    \qquad
    e'\dashv p'.
\]
\end{definition}

\begin{definition}[Projection-pair 2-category]
\label{def:projection-pair-2-category-ff}
The locally ordinary \(2\)-category
\[
    \mathbf{FDom}_{\mathbb V}^{\triangleleft}
\]
has fuzzy domains as objects, fuzzy projection pairs as \(1\)-cells, and
mate-compatible pairs of enriched natural transformations as \(2\)-cells.
The identity \(1\)-cell on \(X\) is
\[
    (\id_X,\id_X):X\triangleleft X,
\]
and composition is
\[
    (e',p')\circ(e,p)=(e'e,pp').
\]
Vertical and horizontal composition of \(2\)-cells are inherited from
\(\mathbf{FDom}^{(2)}_{\mathbb V}\); mate-compatibility is preserved by the
calculus of mates.
\end{definition}

\begin{definition}[Split \(A\)-indexed projection-pair system]
\label{def:split-A-indexed-projection-pair-system}
Let \(A\) be a small \(\mathbb V\)-category.  A
\emph{split \(A\)-indexed projection-pair system}
\[
    D:A\rightsquigarrow\mathbf{FDom}_{\mathbb V}^{\triangleleft}
\]
consists of:

\begin{enumerate}
    \item for every object \(a\in A\), a fuzzy domain \(D_a\);

    \item for every arrow
    \[
        u:a\to a'
    \]
    in the support category \(A_0\), a fuzzy projection pair
    \[
        (e_u,p_u):D_a\triangleleft D_{a'};
    \]

    \item normality isomorphisms
    \[
        (e_{1_a},p_{1_a})\cong(\id_{D_a},\id_{D_a});
    \]

    \item composition isomorphisms
    \[
        e_v e_u\cong e_{vu},
        \qquad
        p_u p_v\cong p_{vu}
    \]
    for composable support arrows
    \[
        a\xrightarrow{u}a'\xrightarrow{v}a'';
    \]

    \item the coherence equations for identities, composition, and
    associativity.
\end{enumerate}
\end{definition}

If
\[
    h:A\to B
\]
is a \(\mathbb V\)-functor, it induces an ordinary functor
\[
    h_0:A_0\to B_0.
\]
Thus a split \(B\)-indexed system
\[
    D:B\rightsquigarrow\mathbf{FDom}_{\mathbb V}^{\triangleleft}
\]
has a reindexing
\[
    Dh:A\rightsquigarrow\mathbf{FDom}_{\mathbb V}^{\triangleleft}
\]
defined on objects by
\[
    (Dh)_a=D_{ha}
\]
and on support arrows by applying \(D\) to their images under \(h_0\).

\begin{definition}[Compatible projection-pair cone]
\label{def:compatible-projection-pair-cone-split}
Let
\[
    D:A\rightsquigarrow\mathbf{FDom}_{\mathbb V}^{\triangleleft}
\]
be a split \(A\)-indexed projection-pair system, and let \(Z\) be a fuzzy
domain.

A \emph{compatible projection-pair cone} from \(D\) to \(Z\) consists of
projection pairs
\[
    (\varepsilon_a,\pi_a):D_a\triangleleft Z
    \qquad(a\in A),
\]
together with:

\begin{enumerate}
    \item for every support arrow
    \[
        u:a\to a'
    \]
    in \(A_0\), coherent isomorphisms
    \[
        \varepsilon_{a'}e_u\cong \varepsilon_a,
        \qquad
        p_u\pi_{a'}\cong \pi_a;
    \]

    \item writing
    \[
        r_a=\varepsilon_a\pi_a:Z\to Z,
    \]
    enriched transition-strength morphisms
    \[
        \theta^Z_{a,a'}:
        A(a,a')
        \longrightarrow
        [Z,Z]_{\mathrm{ff}}(r_a,r_{a'})
    \]
    satisfying the identity and composition axioms for a \(\mathbb V\)-functor
    \[
        r:A\to[Z,Z]_{\mathrm{ff}};
    \]

    \item for every support arrow
    \[
        u:I\to A(a,a'),
    \]
    the composite
    \[
        I
        \xrightarrow{u}
        A(a,a')
        \xrightarrow{\theta^Z_{a,a'}}
        [Z,Z]_{\mathrm{ff}}(r_a,r_{a'})
    \]
    is the canonical comparison
    \[
        \varepsilon_a\pi_a
        \cong
        \varepsilon_{a'}e_u p_u\pi_{a'}
        \Rightarrow
        \varepsilon_{a'}\pi_{a'}
    \]
    induced by the counit
    \[
        e_u p_u\Rightarrow \id_{D_{a'}}.
    \]
\end{enumerate}

A morphism of compatible cones is a family of mate-compatible enriched natural
transformations between the corresponding projection pairs, compatible with the
support-arrow cone isomorphisms and with the transition-strength morphisms.  The
resulting category of compatible cones is denoted
\[
    \operatorname{Cone}_{\triangleleft}(D;Z).
\]
\end{definition}

\begin{lemma}[The approximant functor of a compatible cone]
\label{lem:approximants-are-functorial}
Let
\[
    D:A\rightsquigarrow\mathbf{FDom}_{\mathbb V}^{\triangleleft}
\]
be a split \(A\)-indexed projection-pair system, and let
\[
    (\varepsilon_a,\pi_a):D_a\triangleleft Z
\]
be a compatible projection-pair cone.  Define
\[
    r_a=\varepsilon_a\pi_a:Z\to Z.
\]
Then the objects \(r_a\) assemble canonically to a \(\mathbb V\)-functor
\[
    r:A\to[Z,Z]_{\mathrm{ff}}.
\]
\end{lemma}

\begin{proof}
This is precisely the content of the transition-strength data in
Definition~\ref{def:compatible-projection-pair-cone-split}.
\end{proof}

\section{Canonical tails and common successors}
\label{sec:canonical-tails-common-successors}

The ordinary \(\omega\)-chain construction uses two elementary cofinality
facts: every tail of \(\omega\) is cofinal, and any two indices have a common
later index.  In the finite-flat enriched setting these facts are represented
by canonical support-comma constructions.

\subsection{Canonical tails}

\begin{definition}[Canonical support tail]
\label{def:canonical-support-tail}
Let \(A\) be a small \(\mathbb V\)-category and let \(a_0\in A\).  The
\emph{canonical support tail} of \(A\) at \(a_0\) is the enriched category
\[
    \operatorname{Tail}_A(a_0)
\]
defined as follows.

Its objects are pairs
\[
    (a,u),
    \qquad
    u:a_0\to a
\]
where \(u\) is an arrow in the support category \(A_0\).

The hom-object from
\[
    (a,u)
\]
to
\[
    (a',u')
\]
is the equalizer in \(\mathbb V\)
\[
    \operatorname{Tail}_A(a_0)((a,u),(a',u'))
    \longrightarrow
    A(a,a')
    \rightrightarrows
    A(a_0,a').
\]
The first parallel map is composition with the global element
\[
    u:I\to A(a_0,a).
\]
The second is the constant map determined by the global element
\[
    u':I\to A(a_0,a'),
\]
using the terminality of \(I\).  Composition and identities are inherited from
\(A\), by the universal property of the equalizers.

The evident projection
\[
    \tau_{a_0}:\operatorname{Tail}_A(a_0)\to A,
    \qquad
    (a,u)\longmapsto a,
\]
is a \(\mathbb V\)-functor.
\end{definition}

In the crisp preorder case, \(\operatorname{Tail}_A(a_0)\) is the full tail
suborder on the elements \(a\) satisfying \(a_0\leq a\).

\begin{definition}[Canonical tail weight]
\label{def:canonical-tail-weight}
Let
\[
    (A,W)
\]
be a finite-flat approximation shape and let \(a_0\in A\).  The
\emph{canonical tail weight} at \(a_0\) is
\[
    W^{a_0}
    =
    \tau_{a_0}^{*}W
    :
    \operatorname{Tail}_A(a_0)^{\op}
    \to
    \underline{\mathbb V}.
\]
Thus
\[
    W^{a_0}(a,u)=W(a).
\]
\end{definition}

\begin{definition}[Tail-cofinal approximation shape]
\label{def:tail-cofinal-approximation-shape}
A finite-flat approximation shape
\[
    (A,W)
\]
is \emph{tail-cofinal} if, for every object \(a_0\in A\), the canonical tail
weight
\[
    W^{a_0}=\tau_{a_0}^{*}W
\]
is finite-flat and the canonical direct-image comparison is an isomorphism:
\[
    (\tau_{a_0})_!W^{a_0}\cong W.
\]
\end{definition}

\begin{lemma}[Constant diagrams over canonical tails]
\label{lem:constant-diagrams-canonical-tail}
Let
\[
    (A,W)
\]
be tail-cofinal, and let \(a_0\in A\).  For every fuzzy domain \(X\) and every
object \(x\in X\),
\[
    W^{a_0}*\Delta x\cong x.
\]
\end{lemma}

\begin{proof}
By tail-cofinality, \(W^{a_0}\) is finite-flat.  Hence the result is
Lemma~\ref{lem:constant-diagrams-finite-flat} applied to
\[
    (\operatorname{Tail}_A(a_0),W^{a_0}).
\]
\end{proof}

\subsection{Canonical common successors}

\begin{definition}[Canonical common-successor category]
\label{def:canonical-common-successor-category}
Let \(P\) be a small \(\mathbb V\)-category.  The
\emph{canonical common-successor category}
\[
    \operatorname{Com}(P)
\]
is defined as follows.

An object of \(\operatorname{Com}(P)\) is a tuple
\[
    (i,j,k,u,v),
\]
where \(i,j,k\in P\), and
\[
    u:i\to k,
    \qquad
    v:j\to k
\]
are arrows in the support category \(P_0\).

For two objects
\[
    c=(i,j,k,u,v),
    \qquad
    c'=(i',j',k',u',v'),
\]
the hom-object
\[
    \operatorname{Com}(P)(c,c')
\]
is the equalizer in \(\mathbb V\) of the compatibility maps expressing
\[
    u'r=tu,
    \qquad
    v's=tv
\]
for triples
\[
    r:i\to i',
    \qquad
    s:j\to j',
    \qquad
    t:k\to k'.
\]
Explicitly, it is the equalizer of the two pairs of maps from
\[
    P(i,i')\tensor P(j,j')\tensor P(k,k')
\]
to
\[
    P(i,k')\times P(j,k')
\]
obtained respectively from the composites
\[
    u'r,\quad tu,
    \qquad
    v's,\quad tv.
\]
Composition and identities are inherited from those of \(P\).

There are canonical \(\mathbb V\)-functors
\[
    c_P:\operatorname{Com}(P)\to P\tensor P,
    \qquad
    m_P:\operatorname{Com}(P)\to P
\]
defined on objects by
\[
    c_P(i,j,k,u,v)=(i,j),
    \qquad
    m_P(i,j,k,u,v)=k.
\]
\end{definition}

In the crisp preorder case, an object of \(\operatorname{Com}(P)\) is exactly
a common upper bound
\[
    i\leq k,
    \qquad
    j\leq k.
\]

\begin{definition}[Canonical common-successor weight]
\label{def:canonical-common-successor-weight}
Let
\[
    (P,R)
\]
be a finite-flat approximation shape.  The
\emph{canonical common-successor weight} is
\[
    R^{\mathrm{com}}
    =
    m_P^{*}R:
    \operatorname{Com}(P)^{\op}
    \to
    \underline{\mathbb V}.
\]
Thus, on objects,
\[
    R^{\mathrm{com}}(i,j,k,u,v)=R(k).
\]
\end{definition}

\begin{definition}[Common-successor-cofinal approximation shape]
\label{def:common-successor-cofinal-shape}
A finite-flat approximation shape
\[
    (P,R)
\]
is \emph{common-successor-cofinal} if the canonical common-successor weight
\[
    R^{\mathrm{com}}=m_P^*R
\]
is finite-flat and the canonical maps
\[
    c_P:\operatorname{Com}(P)\to P\tensor P,
    \qquad
    m_P:\operatorname{Com}(P)\to P
\]
satisfy
\[
    (c_P)_!R^{\mathrm{com}}\cong R\boxtimes R,
    \qquad
    (m_P)_!R^{\mathrm{com}}\cong R.
\]
\end{definition}

Thus common-successor-cofinality is the finite-flat analogue of directedness:
double approximation indexed by \(P\tensor P\) can be compared with
approximation over canonical common successors, and then with single
approximation over \(P\).

\begin{definition}[Linear-cofinal finite-flat approximation shape]
\label{def:linear-cofinal-ff-approximation-shape}
A finite-flat approximation shape
\[
    (A,W)
\]
is \emph{linear-cofinal} if:

\begin{enumerate}
    \item it is tail-cofinal;

    \item its external product
    \[
        (A\tensor A,W\boxtimes W)
    \]
    is common-successor-cofinal.
\end{enumerate}
\end{definition}

This deliberately does not include a diagonal comparison
\[
    A\to A\tensor A.
\]
The diagonal is a cartesian feature and appears only in the crisp separated
collapse.

\section{Weight-cofinal iteration shapes}
\label{sec:weight-cofinal-iteration-shapes}

\begin{definition}[Weight-cofinal iteration shape]
\label{def:weight-cofinal-iteration-shape}
A \emph{weight-cofinal iteration shape} is a triple
\[
    (A,W,s),
\]
where \((A,W)\) is a finite-flat approximation shape and
\[
    s:A\to A
\]
is a \(W\)-cofinal endofunctor:
\[
    s_!W\cong W.
\]
\end{definition}

\begin{lemma}[Shift invariance for weighted approximation]
\label{lem:shift-invariance-rde}
Let
\[
    (A,W,s)
\]
be a weight-cofinal iteration shape.  For every fuzzy domain \(X\) and every
diagram
\[
    D:A\to X,
\]
there is a canonical isomorphism
\[
    W*(Ds)
    \cong
    W*D.
\]
\end{lemma}

\begin{proof}
Apply Proposition~\ref{prop:finite-flat-cofinality} to the \(W\)-cofinal map
\[
    s:(A,W)\to(A,W).
\]
\end{proof}

\section{Finite-flat cofinal bilimits}
\label{sec:finite-flat-cofinal-bilimits}

\begin{definition}[\(W\)-bilimit cone]
\label{def:W-bilimit-cone}
Let
\[
    (A,W)
\]
be a finite-flat approximation shape, and let
\[
    D:A\rightsquigarrow\mathbf{FDom}_{\mathbb V}^{\triangleleft}
\]
be a split \(A\)-indexed projection-pair system.

A \emph{\(W\)-bilimit cone} over \(D\) consists of a fuzzy domain
\[
    D_\infty
\]
and a compatible projection-pair cone
\[
    (\varepsilon_a,\pi_a):D_a\triangleleft D_\infty
    \qquad(a\in A),
\]
such that:

\begin{enumerate}
    \item the cone represents compatible projection-pair cones: for every fuzzy
    domain \(Z\), composition with the cone induces an equivalence
    \[
        \mathbf{FDom}_{\mathbb V}^{\triangleleft}(D_\infty,Z)
        \simeq
        \operatorname{Cone}_{\triangleleft}(D;Z);
    \]

    \item if
    \[
        r:A\to[D_\infty,D_\infty]_{\mathrm{ff}}
    \]
    is the approximant functor
    \[
        r_a=\varepsilon_a\pi_a,
    \]
    then the finite-flat approximation identity holds:
    \[
        W*r
        \cong
        \id_{D_\infty}
    \]
    in
    \[
        [D_\infty,D_\infty]_{\mathrm{ff}}.
    \]
\end{enumerate}
\end{definition}

\begin{lemma}[Uniqueness of \(W\)-bilimits]
\label{lem:uniqueness-W-bilimits}
Any two \(W\)-bilimit cones over the same split projection-pair system are
equivalent, and the equivalence is unique up to invertible compatible \(2\)-cell.
\end{lemma}

\begin{proof}
This is the uniqueness part of the representing equivalence in
Definition~\ref{def:W-bilimit-cone}.  The approximation identity is transported
along the resulting compatible equivalence because it is expressed in the
internal hom of endomorphisms.
\end{proof}

\begin{definition}[Bilimit-cofinal reindexing]
\label{def:bilimit-cofinal-reindexing}
Let
\[
    h:(A,W)\to(B,U)
\]
be a cofinal map of finite-flat approximation shapes, and let
\[
    D:B\rightsquigarrow\mathbf{FDom}_{\mathbb V}^{\triangleleft}
\]
be a split \(B\)-indexed projection-pair system.

We say that \(h\) is \emph{bilimit-cofinal for \(D\)} if restriction along
\(h\) induces, for every fuzzy domain \(Z\), an equivalence of compatible cone
categories
\[
    \operatorname{Cone}_{\triangleleft}(D;Z)
    \simeq
    \operatorname{Cone}_{\triangleleft}(Dh;Z).
\]
Under this equivalence, the derived approximant functor for a reindexed cone is
the precomposition of the original derived approximant functor with \(h\).
\end{definition}

\begin{lemma}[Cofinal reindexing of \(W\)-bilimits]
\label{lem:cofinal-reindexing-W-bilimits}
Let
\[
    h:(A,W)\to(B,U)
\]
be cofinal and bilimit-cofinal for
\[
    D:B\rightsquigarrow\mathbf{FDom}_{\mathbb V}^{\triangleleft}.
\]
Then \(D\) has a \(U\)-bilimit with vertex \(D_\infty\) if and only if the
reindexed system
\[
    Dh:A\rightsquigarrow\mathbf{FDom}_{\mathbb V}^{\triangleleft}
\]
has a \(W\)-bilimit with the corresponding transported cone on the same object
\(D_\infty\).
\end{lemma}

\begin{proof}
The projection-pair universal property is transported along the equivalence of
cone categories.

It remains to compare approximation identities.  Let
\[
    r:B\to[D_\infty,D_\infty]_{\mathrm{ff}}
\]
be the approximant functor for the \(B\)-indexed cone.  The approximant functor
for the reindexed cone is
\[
    rh:A\to[D_\infty,D_\infty]_{\mathrm{ff}}.
\]
Since
\[
    h_!W\cong U,
\]
Proposition~\ref{prop:finite-flat-cofinality} gives
\[
    W*(rh)
    \cong
    U*r.
\]
Thus
\[
    U*r\cong\id_{D_\infty}
\]
if and only if
\[
    W*(rh)\cong\id_{D_\infty}.
\]
\end{proof}

\begin{theorem}[Shifted bilimits]
\label{thm:shifted-bilimits-ff}
Let
\[
    (A,W,s)
\]
be a weight-cofinal iteration shape.  Let
\[
    D:A\rightsquigarrow\mathbf{FDom}_{\mathbb V}^{\triangleleft}
\]
be a split \(A\)-indexed projection-pair system with \(W\)-bilimit
\(D_\infty\).  If
\[
    s:(A,W)\to(A,W)
\]
is bilimit-cofinal for \(D\), then \(D_\infty\) is also a \(W\)-bilimit of the
shifted system
\[
    Ds:A\rightsquigarrow\mathbf{FDom}_{\mathbb V}^{\triangleleft}.
\]
\end{theorem}

\begin{proof}
Apply Lemma~\ref{lem:cofinal-reindexing-W-bilimits} to the map
\[
    s:(A,W)\to(A,W).
\]
Equivalently, if
\[
    r:A\to[D_\infty,D_\infty]_{\mathrm{ff}}
\]
is the approximant functor for the original cone, then the approximant functor
for the shifted cone is
\[
    rs:A\to[D_\infty,D_\infty]_{\mathrm{ff}},
\]
and
\[
    W*(rs)
    \cong
    W*r
    \cong
    \id_{D_\infty}
\]
by Lemma~\ref{lem:shift-invariance-rde}.
\end{proof}

\section{Recursive domain equations by cofinal approximation}
\label{sec:rde-by-cofinal-approximation}

\begin{definition}[\(W\)-bilimit-preserving endopseudofunctor]
\label{def:W-bilimit-preserving-endofunctor}
Let
\[
    (A,W)
\]
be a finite-flat approximation shape.  A projection-pair endopseudofunctor
\[
    F:\mathbf{FDom}_{\mathbb V}^{\triangleleft}
      \to
      \mathbf{FDom}_{\mathbb V}^{\triangleleft}
\]
is \emph{\(W\)-bilimit-preserving} if it sends \(W\)-bilimit cones of split
\(A\)-indexed projection-pair systems to \(W\)-bilimit cones.
\end{definition}

\begin{theorem}[Finite-flat cofinal recursive-domain-equation theorem]
\label{thm:finite-flat-cofinal-rde}
Let
\[
    (A,W,s)
\]
be a weight-cofinal iteration shape.  Let
\[
    F:\mathbf{FDom}_{\mathbb V}^{\triangleleft}
      \to
      \mathbf{FDom}_{\mathbb V}^{\triangleleft}
\]
be a \(W\)-bilimit-preserving projection-pair endopseudofunctor.

Let
\[
    D:A\rightsquigarrow\mathbf{FDom}_{\mathbb V}^{\triangleleft}
\]
be a split \(A\)-indexed projection-pair system equipped with an isomorphism of
split projection-pair systems
\[
    Ds\cong FD.
\]
If \(s\) is bilimit-cofinal for \(D\) and \(D_\infty\) is a \(W\)-bilimit of
\(D\), then there is an equivalence of fuzzy domains
\[
    D_\infty\simeq F(D_\infty).
\]
\end{theorem}

\begin{proof}
By Theorem~\ref{thm:shifted-bilimits-ff}, \(D_\infty\) is a \(W\)-bilimit of
the shifted system \(Ds\).  Since
\[
    Ds\cong FD
\]
as split projection-pair systems, \(D_\infty\) is also a \(W\)-bilimit of
\(FD\).

Since \(F\) preserves \(W\)-bilimits,
\[
    F(D_\infty)
\]
is also a \(W\)-bilimit of \(FD\).  By
Lemma~\ref{lem:uniqueness-W-bilimits},
\[
    D_\infty\simeq F(D_\infty).
\]
\end{proof}

\section{Composition and external-product internal-hom approximation}
\label{sec:composition-and-internal-homs-rde}

\begin{lemma}[Separate continuity of composition]
\label{lem:separate-continuity-composition-rde}
For fuzzy domains \(A,B,C\), composition determines a finite-flat bimorphism
\[
    [B,C]_{\mathrm{ff}}\tensor [A,B]_{\mathrm{ff}}
    \longrightarrow
    [A,C]_{\mathrm{ff}},
    \qquad
    (g,f)\longmapsto gf.
\]
Equivalently, it induces a fuzzy-domain morphism
\[
    [B,C]_{\mathrm{ff}}
    \boxtimes_{\mathrm{ff}}
    [A,B]_{\mathrm{ff}}
    \longrightarrow
    [A,C]_{\mathrm{ff}}.
\]
\end{lemma}

\begin{proof}
Finite-flat colimits in internal homs are computed pointwise.

Fix
\[
    f:A\to B.
\]
For a finite-flat weight
\[
    W:K^{\op}\to\underline{\mathbb V}
\]
and a diagram
\[
    H:K\to[B,C]_{\mathrm{ff}},
\]
one has pointwise
\[
\begin{aligned}
    \bigl(W*H\bigr)\circ f
    &=
    a\longmapsto (W*H)(fa)
    \\
    &\cong
    a\longmapsto W*(k\mapsto Hk(fa))
    \\
    &\cong
    W*(k\mapsto Hk\circ f).
\end{aligned}
\]
Hence precomposition by \(f\) preserves finite-flat colimits.

Fix
\[
    g:B\to C.
\]
Since \(g\) is a fuzzy-domain morphism, it preserves finite-flat colimits in
\(B\).  Therefore, for every finite-flat
\[
    W:K^{\op}\to\underline{\mathbb V}
\]
and every diagram
\[
    H:K\to[A,B]_{\mathrm{ff}},
\]
one has pointwise
\[
    g\bigl(W*(k\mapsto Hk(a))\bigr)
    \cong
    W*(k\mapsto g(Hk(a))).
\]
Thus postcomposition by \(g\) preserves finite-flat colimits.  Therefore
composition is finite-flat-continuous separately in both variables.
\end{proof}

Let
\[
    X:A\rightsquigarrow\mathbf{FDom}_{\mathbb V}^{\triangleleft},
    \qquad
    Y:A\rightsquigarrow\mathbf{FDom}_{\mathbb V}^{\triangleleft}
\]
be split \(A\)-indexed projection-pair systems.  Their associated
\emph{external-product internal-hom system}
\[
    [X,Y]_{\mathrm{ext}}
    :
    A\tensor A
    \rightsquigarrow
    \mathbf{FDom}_{\mathbb V}^{\triangleleft}
\]
is defined on objects by
\[
    [X,Y]_{\mathrm{ext}}(a,a')
    =
    [X_{a'},Y_a]_{\mathrm{ff}}.
\]

For a support arrow
\[
    \xi:(a,a')\to(b,b')
\]
in \((A\tensor A)_0\), let
\[
    \xi_1:a\to b,
    \qquad
    \xi_2:a'\to b'
\]
be its two support projections.  The transition projection pair
\[
    [X_{a'},Y_a]_{\mathrm{ff}}
    \triangleleft
    [X_{b'},Y_b]_{\mathrm{ff}}
\]
has embedding
\[
    \mathsf E_\xi(f)
    =
    e^Y_{\xi_1}\, f\, p^X_{\xi_2},
\]
and projection
\[
    \mathsf P_\xi(g)
    =
    p^Y_{\xi_1}\, g\, e^X_{\xi_2}.
\]

\begin{lemma}[Internal-hom transitions are projection pairs]
\label{lem:internal-hom-transitions-projection-pairs}
For every support arrow \(\xi\) in \((A\tensor A)_0\), the pair
\[
    (\mathsf E_\xi,\mathsf P_\xi)
\]
defined above is a fuzzy projection pair.
\end{lemma}

\begin{proof}
The two maps preserve finite-flat colimits by
Lemma~\ref{lem:separate-continuity-composition-rde}.

For
\[
    f:X_{a'}\to Y_a,
\]
one has
\[
\begin{aligned}
    \mathsf P_\xi\mathsf E_\xi(f)
    &=
    p^Y_{\xi_1}e^Y_{\xi_1}\, f\,
    p^X_{\xi_2}e^X_{\xi_2}
    \\
    &\cong
    f,
\end{aligned}
\]
using the unit isomorphisms of the projection pairs for \(X\) and \(Y\).

For
\[
    g:X_{b'}\to Y_b,
\]
one has
\[
    \mathsf E_\xi\mathsf P_\xi(g)
    =
    e^Y_{\xi_1}p^Y_{\xi_1}\, g\,
    e^X_{\xi_2}p^X_{\xi_2}.
\]
The counit comparisons
\[
    e^Y_{\xi_1}p^Y_{\xi_1}\Rightarrow \id_{Y_b},
    \qquad
    e^X_{\xi_2}p^X_{\xi_2}\Rightarrow \id_{X_{b'}}
\]
induce
\[
    \mathsf E_\xi\mathsf P_\xi(g)
    \Rightarrow
    g.
\]
The triangle identities follow from the triangle identities of the projection
pairs in \(X\) and \(Y\).
\end{proof}

Now suppose
\[
    (\varepsilon_a^X,\pi_a^X):X_a\triangleleft X_\infty,
    \qquad
    (\varepsilon_a^Y,\pi_a^Y):Y_a\triangleleft Y_\infty
\]
are compatible projection-pair cones.  Define
\[
    E_{a,a'}:[X_{a'},Y_a]_{\mathrm{ff}}
    \to [X_\infty,Y_\infty]_{\mathrm{ff}}
\]
and
\[
    \Pi_{a,a'}:[X_\infty,Y_\infty]_{\mathrm{ff}}
    \to [X_{a'},Y_a]_{\mathrm{ff}}
\]
by
\[
    E_{a,a'}(f)=\varepsilon_a^Y f\pi_{a'}^X,
\]
and
\[
    \Pi_{a,a'}(g)=\pi_a^Y g\varepsilon_{a'}^X.
\]

\begin{lemma}[Internal-hom cone components are projection pairs]
\label{lem:internal-hom-cone-components-projection-pairs}
For every \(a,a'\), the pair
\[
    (E_{a,a'},\Pi_{a,a'}):
    [X_{a'},Y_a]_{\mathrm{ff}}
    \triangleleft
    [X_\infty,Y_\infty]_{\mathrm{ff}}
\]
is a fuzzy projection pair.
\end{lemma}

\begin{proof}
The maps \(E_{a,a'}\) and \(\Pi_{a,a'}\) preserve finite-flat colimits by
Lemma~\ref{lem:separate-continuity-composition-rde}.

For \(f:X_{a'}\to Y_a\),
\[
\begin{aligned}
    \Pi_{a,a'}E_{a,a'}(f)
    &=
    \pi_a^Y\varepsilon_a^Y f\pi_{a'}^X\varepsilon_{a'}^X
    \\
    &\cong
    f,
\end{aligned}
\]
using the unit isomorphisms of the projection pairs in the two bilimit cones.

For \(g:X_\infty\to Y_\infty\),
\[
    E_{a,a'}\Pi_{a,a'}(g)
    =
    \varepsilon_a^Y\pi_a^Y g\varepsilon_{a'}^X\pi_{a'}^X.
\]
The counit comparisons
\[
    \varepsilon_a^Y\pi_a^Y\Rightarrow\id_{Y_\infty},
    \qquad
    \varepsilon_{a'}^X\pi_{a'}^X\Rightarrow\id_{X_\infty}
\]
induce
\[
    E_{a,a'}\Pi_{a,a'}(g)\Rightarrow g.
\]
The triangle identities follow from the triangle identities of the two
projection-pair cones.
\end{proof}

\begin{theorem}[External-product internal-hom approximation theorem]
\label{thm:external-product-internal-hom-approximation}
Let
\[
    (A,W)
\]
be a finite-flat approximation shape.  Suppose
\[
    X_\infty
\]
and
\[
    Y_\infty
\]
are \(W\)-bilimits of split \(A\)-indexed projection-pair systems
\[
    X
    \qquad\text{and}\qquad
    Y.
\]
Then the external-product internal-hom cone
\[
    (E_{a,a'},\Pi_{a,a'}):
    [X_{a'},Y_a]_{\mathrm{ff}}
    \triangleleft
    [X_\infty,Y_\infty]_{\mathrm{ff}}
\]
has approximants satisfying
\[
    (W\boxtimes W)*
    \bigl((a,a')\mapsto E_{a,a'}\Pi_{a,a'}\bigr)
    \cong
    \id_{[X_\infty,Y_\infty]_{\mathrm{ff}}}.
\]
\end{theorem}

\begin{proof}
Let
\[
    x_{a'}=\varepsilon_{a'}^X\pi_{a'}^X
    \in [X_\infty,X_\infty]_{\mathrm{ff}},
    \qquad
    y_a=\varepsilon_a^Y\pi_a^Y
    \in [Y_\infty,Y_\infty]_{\mathrm{ff}}.
\]
Then
\[
    E_{a,a'}\Pi_{a,a'}(g)=y_a\,g\,x_{a'}.
\]
By separate finite-flat-continuity of composition and finite-flat Fubini, the
\((W\boxtimes W)\)-weighted colimit of these endofunctors is
\[
    g
    \longmapsto
    \bigl(W*(a\mapsto y_a)\bigr)\,
    g\,
    \bigl(W*(a'\mapsto x_{a'})\bigr).
\]
Since \(X_\infty\) and \(Y_\infty\) are \(W\)-bilimits,
\[
    W*(a'\mapsto x_{a'})
    \cong
    \id_{X_\infty},
    \qquad
    W*(a\mapsto y_a)
    \cong
    \id_{Y_\infty}.
\]
Thus the displayed endofunctor is naturally isomorphic to the identity on
\[
    [X_\infty,Y_\infty]_{\mathrm{ff}}.
\]
\end{proof}

\section{Canonical common-successor comparisons}
\label{sec:canonical-common-successor-comparisons}

The general enriched setting permits canonical comparison maps from double
approximations to single approximations.  These comparisons are not generally
invertible; this is the precise point at which the absence of contraction
matters.

Let
\[
    (P,R)
\]
be common-successor-cofinal.  Let
\[
    H:P\rightsquigarrow\mathbf{FDom}_{\mathbb V}^{\triangleleft}
\]
be a split \(P\)-indexed projection-pair system.  Let
\[
    (E_i,\Pi_i):H_i\triangleleft H_\infty
    \qquad(i\in P)
\]
be a compatible cone, and let
\[
    (\delta_i,\rho_i):H_i\triangleleft Z
    \qquad(i\in P)
\]
be another compatible cone.

Write
\[
    L_i=\delta_i\Pi_i:H_\infty\to Z,
    \qquad
    M_i=E_i\rho_i:Z\to H_\infty.
\]

\begin{lemma}[Canonical tail restrictions]
\label{lem:canonical-tail-restrictions}
Assume \((P,R)\) is tail-cofinal.  For every \(i\in P\),
\[
    R*(j\mapsto L_jE_i)
    \cong
    \delta_i,
\]
and
\[
    R*(j\mapsto \Pi_iM_j)
    \cong
    \rho_i.
\]
\end{lemma}

\begin{proof}
Let
\[
    \tau_i:\operatorname{Tail}_P(i)\to P
\]
be the canonical support tail.  Since \((P,R)\) is tail-cofinal,
\[
    (\tau_i)_!R^i\cong R.
\]
Therefore, by Proposition~\ref{prop:finite-flat-cofinality},
\[
    R*(j\mapsto L_jE_i)
    \cong
    R^i*((j,u)\mapsto L_jE_i).
\]

For an object \((j,u:i\to j)\) of \(\operatorname{Tail}_P(i)\), cone
compatibility gives
\[
    E_j e_u\cong E_i,
    \qquad
    p_u\Pi_j\cong \Pi_i,
    \qquad
    \delta_j e_u\cong \delta_i.
\]
Hence
\[
\begin{aligned}
    L_jE_i
    &=
    \delta_j\Pi_jE_i
    \\
    &\cong
    \delta_j\Pi_jE_j e_u
    \\
    &\cong
    \delta_j e_u
    \\
    &\cong
    \delta_i.
\end{aligned}
\]
Thus the restricted diagram is constant at \(\delta_i\), and
\[
    R^i*((j,u)\mapsto L_jE_i)
    \cong
    R^i*\Delta(\delta_i)
    \cong
    \delta_i
\]
by Lemma~\ref{lem:constant-diagrams-canonical-tail}.

The proof of
\[
    R*(j\mapsto \Pi_iM_j)
    \cong
    \rho_i
\]
is analogous.  For \((j,u:i\to j)\), cone compatibility gives
\[
    p_u\rho_j\cong \rho_i
\]
and
\[
\begin{aligned}
    \Pi_iM_j
    &=
    \Pi_iE_j\rho_j
    \\
    &\cong
    p_u\Pi_jE_j\rho_j
    \\
    &\cong
    p_u\rho_j
    \\
    &\cong
    \rho_i.
\end{aligned}
\]
\end{proof}

\begin{lemma}[Canonical double-to-single comparison]
\label{lem:canonical-double-to-single-comparison}
Assume \((P,R)\) is common-successor-cofinal.  Then there is a canonical
comparison
\[
    (R\boxtimes R)*
    \bigl((i,j)\mapsto M_iL_j\bigr)
    \Rightarrow
    R*(k\mapsto E_k\Pi_k).
\]
Moreover, there is a canonical counit comparison
\[
    (R\boxtimes R)*
    \bigl((i,j)\mapsto L_iM_j\bigr)
    \Rightarrow
    \id_Z.
\]
\end{lemma}

\begin{proof}
Let
\[
    R^{\mathrm{com}}=m_P^*R
\]
be the canonical common-successor weight.  By common-successor-cofinality,
\[
    (c_P)_!R^{\mathrm{com}}\cong R\boxtimes R,
    \qquad
    (m_P)_!R^{\mathrm{com}}\cong R.
\]
Hence
\[
\begin{aligned}
    (R\boxtimes R)*
    \bigl((i,j)\mapsto M_iL_j\bigr)
    &\cong
    R^{\mathrm{com}}*
    \bigl((i,j,k,u,v)\mapsto M_iL_j\bigr).
\end{aligned}
\]

For an object
\[
    (i,j,k,u:i\to k,v:j\to k)
\]
of \(\operatorname{Com}(P)\), cone compatibility gives
\[
    E_k e_u\cong E_i,
    \qquad
    p_u\rho_k\cong \rho_i,
\]
and
\[
    \delta_k e_v\cong \delta_j,
    \qquad
    p_v\Pi_k\cong \Pi_j.
\]
Therefore
\[
\begin{aligned}
    M_iL_j
    &=
    E_i\rho_i\delta_j\Pi_j
    \\
    &\cong
    E_k e_u p_u\rho_k\delta_k e_vp_v\Pi_k
    \\
    &\cong
    E_k e_u p_u e_vp_v\Pi_k
    \\
    &\Rightarrow
    E_k\Pi_k.
\end{aligned}
\]
The final comparison is induced by the counits
\[
    e_up_u\Rightarrow \id_{H_k},
    \qquad
    e_vp_v\Rightarrow \id_{H_k}.
\]
These comparisons are coherent in \(\operatorname{Com}(P)\).  Taking the
\(R^{\mathrm{com}}\)-weighted colimit and using
\[
    (m_P)_!R^{\mathrm{com}}\cong R
\]
gives the first comparison.

For the second comparison, again use cofinality of \(c_P\).  For an object
\[
    (i,j,k,u:i\to k,v:j\to k)
\]
there is a canonical comparison
\[
    L_iM_j\Rightarrow \id_Z.
\]
Indeed,
\[
\begin{aligned}
    L_iM_j
    &=
    \delta_i\Pi_iE_j\rho_j
    \\
    &\cong
    \delta_k e_u p_u\Pi_kE_k e_vp_v\rho_k
    \\
    &\Rightarrow
    \delta_k\rho_k
    \\
    &\Rightarrow
    \id_Z.
\end{aligned}
\]
The final comparison is the counit of
\[
    (\delta_k,\rho_k):H_k\triangleleft Z.
\]
Taking the \(R^{\mathrm{com}}\)-weighted colimit and using
Lemma~\ref{lem:constant-diagrams-finite-flat} for the constant diagram at
\(\id_Z\) gives the second comparison.
\end{proof}

\begin{theorem}[Lax reconstruction of external-product cones]
\label{thm:lax-reconstruction-internal-hom-cones}
Let
\[
    (A,W)
\]
be linear-cofinal.  Put
\[
    P=A\tensor A,
    \qquad
    R=W\boxtimes W.
\]
Let
\[
    H_\infty=[X_\infty,Y_\infty]_{\mathrm{ff}}
\]
with external-product approximants
\[
    (E_i,\Pi_i):H_i\triangleleft H_\infty
    \qquad(i\in P).
\]
Let
\[
    (\delta_i,\rho_i):H_i\triangleleft Z
\]
be a compatible cone.  Define
\[
    \delta
    =
    R*(i\mapsto \delta_i\Pi_i):
    H_\infty\to Z
\]
and
\[
    \rho
    =
    R*(i\mapsto E_i\rho_i):
    Z\to H_\infty.
\]
Then
\[
    \delta E_i\cong \delta_i,
    \qquad
    \Pi_i\rho\cong \rho_i,
\]
and there are canonical comparisons
\[
    \rho\delta\Rightarrow \id_{H_\infty},
    \qquad
    \delta\rho\Rightarrow \id_Z.
\]
\end{theorem}

\begin{proof}
The restriction equations follow from
Lemma~\ref{lem:canonical-tail-restrictions}.

For the first comparison, separate finite-flat-continuity of composition and
finite-flat Fubini give
\[
    \rho\delta
    \cong
    (R\boxtimes R)*
    \bigl((i,j)\mapsto M_iL_j\bigr).
\]
By Lemma~\ref{lem:canonical-double-to-single-comparison},
\[
    (R\boxtimes R)*
    \bigl((i,j)\mapsto M_iL_j\bigr)
    \Rightarrow
    R*(k\mapsto E_k\Pi_k).
\]
By Theorem~\ref{thm:external-product-internal-hom-approximation},
\[
    R*(k\mapsto E_k\Pi_k)
    \cong
    \id_{H_\infty}.
\]
Thus
\[
    \rho\delta\Rightarrow\id_{H_\infty}.
\]
The comparison
\[
    \delta\rho\Rightarrow \id_Z
\]
is the second comparison of
Lemma~\ref{lem:canonical-double-to-single-comparison}.
\end{proof}

\section{Linear application in the general finite-flat setting}
\label{sec:linear-application-general-ff}

The preceding theorem reconstructs canonical lax application data from
external-product finite-stage cones.  It does not produce a projection-pair
retract in full generality; the missing ingredient is precisely a diagonal or
contraction comparison.

\begin{definition}[Linear reflexive fuzzy domain]
\label{def:linear-reflexive-fuzzy-domain-ff}
A \emph{linear reflexive fuzzy domain} is a fuzzy domain \(D\) equipped with a
fuzzy projection pair
\[
    [D,D]_{\mathrm{ff}}\triangleleft D.
\]
Equivalently, it consists of fuzzy-domain morphisms
\[
    \operatorname{encode}:[D,D]_{\mathrm{ff}}\to D,
    \qquad
    \operatorname{decode}:D\to [D,D]_{\mathrm{ff}},
\]
together with an adjunction
\[
    \operatorname{encode}\dashv \operatorname{decode}
\]
whose unit is invertible.  Thus
\[
    \operatorname{decode}\operatorname{encode}
    \cong
    \id_{[D,D]_{\mathrm{ff}}},
\]
and there is a counit comparison
\[
    \operatorname{encode}\operatorname{decode}
    \Rightarrow
    \id_D
\]
satisfying the triangle identities.

It is \emph{extensional} if the counit is also invertible.
\end{definition}

The associated application morphism is
\[
    D\boxtimes_{\mathrm{ff}}D
    \xrightarrow{\operatorname{decode}\boxtimes 1_D}
    [D,D]_{\mathrm{ff}}\boxtimes_{\mathrm{ff}}D
    \xrightarrow{\operatorname{ev}}
    D.
\]

\begin{definition}[\((G,T)\)-linear reflexive fuzzy domain]
\label{def:mixed-linear-reflexive-fuzzy-domain-ff}
Let
\[
    G,T:\mathbf{FDom}_{\mathbb V}^{\triangleleft}
      \to
      \mathbf{FDom}_{\mathbb V}^{\triangleleft}
\]
be projection-pair endopseudofunctors.  A
\emph{\((G,T)\)-linear reflexive fuzzy domain} is a fuzzy domain \(D\) equipped
with a fuzzy projection pair
\[
    [GD,TD]_{\mathrm{ff}}\triangleleft D.
\]
The associated application morphism has type
\[
    D\boxtimes_{\mathrm{ff}}GD\to TD
\]
and is the composite
\[
    D\boxtimes_{\mathrm{ff}}GD
    \xrightarrow{\operatorname{decode}\boxtimes 1_{GD}}
    [GD,TD]_{\mathrm{ff}}\boxtimes_{\mathrm{ff}}GD
    \xrightarrow{\operatorname{ev}}
    TD.
\]
\end{definition}

\begin{definition}[External finite-stage endomorphism system]
\label{def:external-finite-stage-endomorphism-system}
Let
\[
    D:A\rightsquigarrow\mathbf{FDom}_{\mathbb V}^{\triangleleft}
\]
be a split \(A\)-indexed projection-pair system.  Its
\emph{external finite-stage endomorphism system} is the split
\((A\tensor A)\)-indexed projection-pair system
\[
    \operatorname{End}_{\mathrm{ext}}(D)
    :
    A\tensor A
    \rightsquigarrow
    \mathbf{FDom}_{\mathbb V}^{\triangleleft}
\]
defined by
\[
    \operatorname{End}_{\mathrm{ext}}(D)(a,a')
    =
    [D_{a'},D_a]_{\mathrm{ff}}.
\]
\end{definition}

\begin{definition}[Linear self-application cone]
\label{def:linear-self-application-cone}
Let
\[
    (A,W)
\]
be linear-cofinal, and let
\[
    D:A\rightsquigarrow\mathbf{FDom}_{\mathbb V}^{\triangleleft}
\]
be a split \(A\)-indexed projection-pair system with \(W\)-bilimit
\[
    D_\infty.
\]
A \emph{linear self-application cone} on \(D\) is a compatible projection-pair
cone
\[
    (\kappa_{a,a'},\lambda_{a,a'}):
    [D_{a'},D_a]_{\mathrm{ff}}
    \triangleleft
    D_\infty
\]
over the external finite-stage endomorphism system
\[
    \operatorname{End}_{\mathrm{ext}}(D).
\]
\end{definition}

\begin{theorem}[Pure lax linear application theorem]
\label{thm:pure-lax-linear-application}
Let
\[
    (A,W)
\]
be linear-cofinal, and suppose
\[
    D_\infty
\]
is a \(W\)-bilimit of a split \(A\)-indexed projection-pair system \(D\).
Let
\[
    H_\infty=[D_\infty,D_\infty]_{\mathrm{ff}}.
\]
If \(D\) carries a linear self-application cone, then there are canonical
fuzzy-domain morphisms
\[
    \operatorname{encode}:H_\infty\to D_\infty,
    \qquad
    \operatorname{decode}:D_\infty\to H_\infty
\]
defined by
\[
    \operatorname{encode}
    =
    (W\boxtimes W)*
    \bigl((a,a')\mapsto
        \kappa_{a,a'}\Pi_{a,a'}
    \bigr),
\]
and
\[
    \operatorname{decode}
    =
    (W\boxtimes W)*
    \bigl((a,a')\mapsto
        E_{a,a'}\lambda_{a,a'}
    \bigr).
\]
They satisfy canonical comparisons
\[
    \operatorname{decode}\operatorname{encode}
    \Rightarrow
    \id_{H_\infty},
    \qquad
    \operatorname{encode}\operatorname{decode}
    \Rightarrow
    \id_{D_\infty}.
\]
Consequently there is a canonical resource-sensitive application morphism
\[
    D_\infty\boxtimes_{\mathrm{ff}}D_\infty
    \longrightarrow
    D_\infty.
\]
\end{theorem}

\begin{proof}
Apply Theorem~\ref{thm:lax-reconstruction-internal-hom-cones} to the external
endomorphism system
\[
    (a,a')\longmapsto [D_{a'},D_a]_{\mathrm{ff}}
\]
and to the compatible cone
\[
    (\kappa_{a,a'},\lambda_{a,a'}):
    [D_{a'},D_a]_{\mathrm{ff}}
    \triangleleft
    D_\infty.
\]
The resulting maps are precisely the displayed
\(\operatorname{encode}\) and \(\operatorname{decode}\).  The application
morphism is
\[
    D_\infty\boxtimes_{\mathrm{ff}}D_\infty
    \xrightarrow{\operatorname{decode}\boxtimes 1}
    [D_\infty,D_\infty]_{\mathrm{ff}}
    \boxtimes_{\mathrm{ff}}
    D_\infty
    \xrightarrow{\operatorname{ev}}
    D_\infty.
\]
\end{proof}

\begin{definition}[Mixed external finite-stage function-space system]
\label{def:mixed-external-finite-stage-function-space-system}
Let
\[
    G,T:\mathbf{FDom}_{\mathbb V}^{\triangleleft}
      \to
      \mathbf{FDom}_{\mathbb V}^{\triangleleft}
\]
be projection-pair endopseudofunctors, and let
\[
    D:A\rightsquigarrow\mathbf{FDom}_{\mathbb V}^{\triangleleft}
\]
be a split \(A\)-indexed projection-pair system.  The
\emph{mixed external finite-stage function-space system} is the split
\((A\tensor A)\)-indexed system
\[
    (a,a')\longmapsto [GD_{a'},TD_a]_{\mathrm{ff}}.
\]
\end{definition}

\begin{definition}[Mixed linear self-application cone]
\label{def:mixed-linear-self-application-cone}
Let
\[
    (A,W)
\]
be linear-cofinal.  Let
\[
    G,T:\mathbf{FDom}_{\mathbb V}^{\triangleleft}
      \to
      \mathbf{FDom}_{\mathbb V}^{\triangleleft}
\]
be projection-pair endopseudofunctors preserving \(W\)-bilimits.  Let
\[
    D:A\rightsquigarrow\mathbf{FDom}_{\mathbb V}^{\triangleleft}
\]
be a split \(A\)-indexed projection-pair system with \(W\)-bilimit
\[
    D_\infty.
\]
A \emph{mixed linear self-application cone} is a compatible projection-pair cone
\[
    [GD_{a'},TD_a]_{\mathrm{ff}}
    \triangleleft
    D_\infty
\]
over the mixed external finite-stage function-space system.
\end{definition}

\begin{theorem}[Mixed lax linear application theorem]
\label{thm:mixed-lax-linear-application}
Let
\[
    (A,W)
\]
be linear-cofinal.  Let
\[
    G,T:\mathbf{FDom}_{\mathbb V}^{\triangleleft}
      \to
      \mathbf{FDom}_{\mathbb V}^{\triangleleft}
\]
be projection-pair endopseudofunctors preserving \(W\)-bilimits.  Suppose
\[
    D_\infty
\]
is a \(W\)-bilimit of a split \(A\)-indexed projection-pair system \(D\).
If \(D\) carries a mixed linear self-application cone, then there are canonical
morphisms
\[
    \operatorname{encode}:[GD_\infty,TD_\infty]_{\mathrm{ff}}\to D_\infty,
\]
and
\[
    \operatorname{decode}:D_\infty\to[GD_\infty,TD_\infty]_{\mathrm{ff}},
\]
with canonical comparisons
\[
    \operatorname{decode}\operatorname{encode}
    \Rightarrow
    \id_{[GD_\infty,TD_\infty]_{\mathrm{ff}}},
    \qquad
    \operatorname{encode}\operatorname{decode}
    \Rightarrow
    \id_{D_\infty}.
\]
The associated application morphism has type
\[
    D_\infty\boxtimes_{\mathrm{ff}}GD_\infty\to TD_\infty.
\]
\end{theorem}

\begin{proof}
Since \(G\) and \(T\) preserve \(W\)-bilimits,
\[
    GD_\infty
\]
is a \(W\)-bilimit of \(GD\), and
\[
    TD_\infty
\]
is a \(W\)-bilimit of \(TD\).  Apply
Theorem~\ref{thm:lax-reconstruction-internal-hom-cones} to the mixed external
finite-stage function-space system.
\end{proof}

\subsection{Finite-stage generation of self-application cones}

The preceding lax application theorem requires a compatible cone from the
finite-stage function-space system into the bilimit.  The following construction
supplies such data over the canonical common-successor category.

\begin{definition}[Finite-stage linear self-application data]
\label{def:finite-stage-linear-self-application-data}
Let
\[
    (A,W)
\]
be linear-cofinal, and let
\[
    D:A\rightsquigarrow\mathbf{FDom}_{\mathbb V}^{\triangleleft}
\]
be a split \(A\)-indexed projection-pair system.  Finite-stage linear
self-application data consist of:

\begin{enumerate}
    \item a support endofunctor
    \[
        s:A\to A;
    \]

    \item for every \(k\in A\), a fuzzy projection pair
    \[
        (\sigma_k,\tau_k):
        [D_k,D_k]_{\mathrm{ff}}
        \triangleleft
        D_{sk};
    \]

    \item coherence isomorphisms saying that, for every support arrow
    \[
        w:k\to k',
    \]
    the square of projection pairs comparing
    \[
        [D_k,D_k]_{\mathrm{ff}}
        \triangleleft
        D_{sk}
    \]
    with
    \[
        [D_{k'},D_{k'}]_{\mathrm{ff}}
        \triangleleft
        D_{sk'}
    \]
    is compatible with the internal-hom transition maps and with the transition
    projection pair
    \[
        D_{sk}\triangleleft D_{sk'}.
    \]
\end{enumerate}
\end{definition}

\begin{construction}[Common-successor finite-stage encoding]
\label{con:common-successor-finite-stage-encoding}
Let
\[
    (i,j,k,u:i\to k,v:j\to k)
\]
be an object of \(\operatorname{Com}(A)\).  The transition projection pairs in
\(D\) induce a projection pair
\[
    [D_j,D_i]_{\mathrm{ff}}
    \triangleleft
    [D_k,D_k]_{\mathrm{ff}}.
\]
Its embedding sends
\[
    f:D_j\to D_i
\]
to
\[
    e_u f p_v:D_k\to D_k,
\]
and its projection sends
\[
    g:D_k\to D_k
\]
to
\[
    p_u g e_v:D_j\to D_i.
\]
Composing with
\[
    [D_k,D_k]_{\mathrm{ff}}
    \triangleleft
    D_{sk}
\]
and then with the bilimit cone component
\[
    D_{sk}\triangleleft D_\infty
\]
gives a projection pair
\[
    [D_j,D_i]_{\mathrm{ff}}
    \triangleleft
    D_\infty.
\]
The coherence axioms for finite-stage self-application data make these
projection pairs compatible over \(\operatorname{Com}(A)\).
\end{construction}

\begin{theorem}[Common-successor finite-stage encodings over \(\operatorname{Com}(A)\)]
\label{thm:common-successor-encodings-over-com}
Let
\[
    (A,W)
\]
be linear-cofinal, and let
\[
    D:A\rightsquigarrow\mathbf{FDom}_{\mathbb V}^{\triangleleft}
\]
be a split \(A\)-indexed system with \(W\)-bilimit \(D_\infty\).  Suppose \(D\)
is equipped with finite-stage linear self-application data.  Then the
construction above gives a compatible projection-pair cone over the pullback of
the external endomorphism system along
\[
    c_A:\operatorname{Com}(A)\to A\tensor A.
\]
\end{theorem}

\begin{proof}
For each object
\[
    (i,j,k,u,v)
\]
of \(\operatorname{Com}(A)\), Construction~\ref{con:common-successor-finite-stage-encoding}
gives a projection pair
\[
    [D_j,D_i]_{\mathrm{ff}}\triangleleft D_\infty.
\]
Compatibility with arrows of \(\operatorname{Com}(A)\) follows from the
coherence of the split system \(D\), the coherence of the finite-stage
self-application data, and the calculus of mates.
\end{proof}

\section{Constant-map seeds}
\label{sec:constant-map-seeds-ff}

The following construction supplies standard finite-stage seeds.

\begin{lemma}[Constant finite-flat-continuous maps]
\label{lem:constant-ff-continuous}
Let \(X\) and \(Y\) be fuzzy domains, and let \(y\in Y\).  Since
\(\mathbb V\) is affine, \(y\) determines a constant \(\mathbb V\)-functor
\[
    \operatorname{const}_y:X\to Y.
\]
This constant functor preserves finite-flat weighted colimits.
\end{lemma}

\begin{proof}
The constant assignment is a \(\mathbb V\)-functor because, for every
\(x,x'\in X\), affinity gives a canonical morphism
\[
    X(x,x')\to I,
\]
and the identity of \(y\) gives
\[
    I\to Y(y,y).
\]
Their composite supplies the hom-map
\[
    X(x,x')\to Y(y,y).
\]

Let
\[
    W:K^{\op}\to\underline{\mathbb V}
\]
be finite-flat and let
\[
    D:K\to X
\]
be a diagram.  Since \(W\) is finite-flat,
Lemma~\ref{lem:constant-diagrams-finite-flat} gives
\[
    W*(\operatorname{const}_yD)
    \cong
    y.
\]
Therefore
\[
    \operatorname{const}_y(W*D)
    =
    y
    \cong
    W*(\operatorname{const}_yD).
\]
Thus \(\operatorname{const}_y\) preserves finite-flat weighted colimits.
\end{proof}

\begin{proposition}[Pure constant-map seed]
\label{prop:pure-constant-map-seed-ff}
Suppose a fuzzy domain \(D_0\) has an enriched initial object
\[
    \bot\in D_0.
\]
Define
\[
    e:D_0\to [D_0,D_0]_{\mathrm{ff}}
\]
by
\[
    e(a)=\operatorname{const}_a,
\]
and define
\[
    p:[D_0,D_0]_{\mathrm{ff}}\to D_0
\]
by
\[
    p(f)=f(\bot).
\]
Then
\[
    (e,p):D_0\triangleleft [D_0,D_0]_{\mathrm{ff}}
\]
is a fuzzy projection pair.
\end{proposition}

\begin{proof}
By Lemma~\ref{lem:constant-ff-continuous}, each \(\operatorname{const}_a\)
lies in \([D_0,D_0]_{\mathrm{ff}}\).  The assignment
\[
    a\longmapsto\operatorname{const}_a
\]
is a \(\mathbb V\)-functor because the end defining the internal hom admits the
evident pointwise wedge
\[
    D_0(a,a')\to D_0(a,a').
\]
It preserves finite-flat colimits by Lemma~\ref{lem:constant-ff-continuous} and
Lemma~\ref{lem:constant-diagrams-finite-flat}.

The functor \(p\) is evaluation at \(\bot\).  Evaluation is a
\(\mathbb V\)-functor and preserves pointwise finite-flat colimits.  Thus \(p\)
is a fuzzy-domain morphism.

Moreover,
\[
    p(e(a))
    =
    \operatorname{const}_a(\bot)
    =
    a.
\]
Thus
\[
    pe=\id_{D_0}.
\]

For \(f:D_0\to D_0\), the composite \(ep(f)\) is
\[
    \operatorname{const}_{f(\bot)}.
\]
The enriched initiality of \(\bot\) gives a canonical enriched natural
transformation
\[
    \Delta\bot\Rightarrow 1_{D_0}.
\]
Applying \(f\) yields
\[
    \Delta f(\bot)\Rightarrow f.
\]
This is precisely
\[
    ep(f)\Rightarrow f.
\]
The triangle identities follow from \(pe=\id_{D_0}\) and from functoriality of
the canonical initial-object comparison.  Therefore \(e\dashv p\) with
invertible unit.
\end{proof}

\begin{proposition}[Mixed constant-map seed]
\label{prop:mixed-constant-map-seed-ff}
Let
\[
    G,T:\mathbf{FDom}_{\mathbb V}^{\triangleleft}
      \to
      \mathbf{FDom}_{\mathbb V}^{\triangleleft}
\]
be projection-pair endopseudofunctors.  Suppose \(GD_0\) has an enriched initial
object
\[
    \bot_{GD_0}\in GD_0,
\]
and suppose there is a fuzzy projection pair
\[
    (u,v):D_0\triangleleft TD_0.
\]
Define
\[
    e:D_0\to[GD_0,TD_0]_{\mathrm{ff}}
\]
by
\[
    e(a)=\operatorname{const}_{u(a)},
\]
and define
\[
    p:[GD_0,TD_0]_{\mathrm{ff}}\to D_0
\]
by
\[
    p(f)=v(f(\bot_{GD_0})).
\]
Then
\[
    (e,p):D_0\triangleleft[GD_0,TD_0]_{\mathrm{ff}}
\]
is a fuzzy projection pair.
\end{proposition}

\begin{proof}
Constant maps preserve finite-flat colimits by
Lemma~\ref{lem:constant-ff-continuous}.  Since \(u\) preserves finite-flat
colimits, the functor
\[
    e(a)=\operatorname{const}_{u(a)}
\]
preserves finite-flat colimits.

The functor \(p\) is the composite of evaluation at \(\bot_{GD_0}\) with
\[
    v:TD_0\to D_0.
\]
Evaluation preserves pointwise finite-flat colimits, and \(v\) is a
fuzzy-domain morphism.  Hence \(p\) is a fuzzy-domain morphism.

For \(a\in D_0\),
\[
    p(e(a))
    =
    v(u(a))
    \cong
    a
\]
by the unit isomorphism of the projection pair \((u,v)\).

For \(f:GD_0\to TD_0\), the composite \(ep(f)\) is the constant map at
\[
    u(v(f(\bot_{GD_0}))).
\]
The counit comparison
\[
    uv\Rightarrow \id_{TD_0}
\]
gives
\[
    u(v(f(\bot_{GD_0})))\to f(\bot_{GD_0}).
\]
The enriched initiality of \(\bot_{GD_0}\) gives
\[
    \Delta\bot_{GD_0}\Rightarrow 1_{GD_0}.
\]
Applying \(f\) gives
\[
    \Delta f(\bot_{GD_0})\Rightarrow f.
\]
Composing these transformations gives
\[
    ep(f)\Rightarrow f.
\]
The triangle identities follow from the triangle identities of \((u,v)\) and
from functoriality of the canonical initial-object comparison.
\end{proof}

\section{The posetal reflection}
\label{sec:posetal-reflection-ff-rde}

Let
\[
    \Lambda=(L,\leq,\meet,\join,\tensor,\multimap,\bot_L,\top_L)
\]
be a fuzzy truth-value object.  In the posetal reflection, finite-flat
approximation shapes are fuzzy ideal weights
\[
    W:A^{\op}\to\Lambda.
\]
For a \(\Lambda\)-functor
\[
    h:A\to B,
\]
the lower direct image is computed by
\[
    (h_!W)(b)
    =
    \bigjoin_{a\in A}
        W(a)\tensor B(b,ha).
\]
Cofinality means
\[
    h_!W\cong U.
\]
Since the truth-value fuzzy preorder \(\Lambda\) is separated, the presheaf
fuzzy preorder
\[
    [B^{\op},\Lambda]
\]
is separated.  Hence an isomorphism of weights may be written as equality:
\[
    h_!W=U.
\]

For every finite-flat weight \(W\), the terminal direct image satisfies
\[
    !_!W\cong \top_L.
\]
Equivalently,
\[
    \bigjoin_{a\in A}W(a)=\top_L.
\]
This is the posetal reflection of preservation of the terminal finite-limit
predicate by the character \(\chi_W\).

A projection pair
\[
    (e,p):X\triangleleft Y
\]
is a pair of fuzzy-domain morphisms satisfying
\[
    pe\cong \id_X,
    \qquad
    ep\leq \id_Y
\]
in the induced crisp preorder on fuzzy-domain morphisms, where
\[
    f\leq g
    \quad:\Longleftrightarrow\quad
    [X,Y]_{\mathrm{ff}}(f,g)=\top_L.
\]

A \(W\)-bilimit cone has approximation identity
\[
    W*(a\mapsto \varepsilon_a\pi_a)
    \cong
    \id_{D_\infty}
\]
inside
\[
    [D_\infty,D_\infty]_{\mathrm{ff}}.
\]
The finite-flat cofinal recursive-domain-equation theorem therefore gives:
if
\[
    Ds\cong FD,
\]
if \(s\) is bilimit-cofinal for \(D\), and if \(F\) preserves the corresponding
\(W\)-bilimit, then
\[
    D_\infty\simeq F(D_\infty).
\]

The general enriched internal-hom theorem gives external-product approximation
\[
    (W\boxtimes W)*
    \bigl((a,a')\mapsto E_{a,a'}\Pi_{a,a'}\bigr)
    \cong
    \id_{[X_\infty,Y_\infty]_{\mathrm{ff}}}.
\]
It does not generally collapse to a same-index equation, because there is no
canonical diagonal
\[
    A\to A\tensor A.
\]

\section{Pointed finite-flat minimal invariants}
\label{sec:pointed-ff-minimal-invariants}

The preceding recursive-domain-equation theorem produces solutions from
finite-flat bilimits.  We now isolate the pointed form of the construction,
which is the finite-flat analogue of the classical pointed-dcpo minimal
solution theorem.

The point of pointedness is not that every solution of
\[
    X\simeq F(X)
\]
is unique.  That statement is false already in ordinary pointed dcpos: for
\(F=\id\), every pointed dcpo is a solution.  Rather, pointedness supplies a
canonical initial approximation process.  The resulting bilimit is the
canonical generated invariant of \(F\).

\begin{definition}[Pointed fuzzy domain]
\label{def:pointed-fuzzy-domain}
A \emph{pointed fuzzy domain} is a fuzzy domain \(X\) equipped with an enriched
initial object
\[
    \bot_X\in X.
\]
Thus
\[
    X(\bot_X,x)\cong I
\]
naturally in \(x\in X\).

A morphism of pointed fuzzy domains is a fuzzy-domain morphism
\[
    f:X\to Y.
\]
It is called \emph{strict} if
\[
    f(\bot_X)\cong \bot_Y.
\]
\end{definition}

\begin{lemma}[Projection pairs between pointed domains are strict]
\label{lem:projection-pairs-between-pointed-domains-are-strict}
Let
\[
    (e,p):X\triangleleft Y
\]
be a fuzzy projection pair between pointed fuzzy domains.  Then both
\[
    e:X\to Y
    \qquad
    p:Y\to X
\]
preserve the chosen initial objects:
\[
    e(\bot_X)\cong \bot_Y,
    \qquad
    p(\bot_Y)\cong \bot_X.
\]
\end{lemma}

\begin{proof}
Since \(e\dashv p\), for every \(y\in Y\) there are natural isomorphisms
\[
    Y(e\bot_X,y)
    \cong
    X(\bot_X,py)
    \cong
    I.
\]
Thus \(e\bot_X\) is an enriched initial object of \(Y\), so
\[
    e(\bot_X)\cong\bot_Y.
\]
Using the invertible unit of the projection pair,
\[
    pe\cong\id_X,
\]
we obtain
\[
    p(\bot_Y)
    \cong
    p(e\bot_X)
    \cong
    \bot_X.
\]
\end{proof}

\begin{definition}[Pointed projection-pair category]
\label{def:pointed-projection-pair-category}
Let
\[
    \mathbf{FDom}^{\triangleleft}_{\mathbb V,\bot}
\]
denote the locally ordinary \(2\)-category whose objects are pointed fuzzy
domains, whose \(1\)-cells are fuzzy projection pairs between pointed fuzzy
domains, and whose \(2\)-cells are mate-compatible pairs of enriched natural
transformations.

By Lemma~\ref{lem:projection-pairs-between-pointed-domains-are-strict}, every
\(1\)-cell in
\[
    \mathbf{FDom}^{\triangleleft}_{\mathbb V,\bot}
\]
is automatically strict on the chosen initial objects.
\end{definition}

\begin{definition}[The one-point pointed fuzzy domain]
\label{def:one-point-pointed-fuzzy-domain}
Let
\[
    \mathbf 1_\bot
\]
denote the one-object fuzzy domain.  Its unique object is written
\[
    \bot.
\]
It is pointed by its unique object.
\end{definition}

\begin{lemma}[The one-point domain embeds into every pointed fuzzy domain]
\label{lem:one-point-domain-embeds-into-pointed-domain}
For every pointed fuzzy domain \(X\), there is a canonical fuzzy projection pair
\[
    \mathbf 1_\bot\triangleleft X.
\]
Its embedding sends the unique object of \(\mathbf 1_\bot\) to
\[
    \bot_X,
\]
and its projection is the unique functor
\[
    X\to\mathbf 1_\bot.
\]
\end{lemma}

\begin{proof}
Let
\[
    e:\mathbf 1_\bot\to X
\]
send the unique object to \(\bot_X\), and let
\[
    p:X\to\mathbf 1_\bot
\]
be the unique functor.  Since \(\bot_X\) is enriched initial, for every
\(x\in X\) one has
\[
    X(e*,x)=X(\bot_X,x)\cong I
    \cong
    \mathbf 1_\bot(*,px).
\]
Thus
\[
    e\dashv p.
\]
The unit is the identity, and the counit is the enriched natural transformation
\[
    \Delta\bot_X\Rightarrow\id_X
\]
induced by initiality of \(\bot_X\).
\end{proof}

\begin{definition}[Sequential finite-flat approximation shape]
\label{def:sequential-ff-approximation-shape}
A \emph{sequential finite-flat approximation shape} is a linear-cofinal
finite-flat approximation shape
\[
    (\omega_{\mathbb V},W_\omega)
\]
whose support category is the ordinary preorder \(\omega\), together with a
\(\mathbb V\)-endofunctor
\[
    s:\omega_{\mathbb V}\to\omega_{\mathbb V}
\]
whose action on support objects is
\[
    n\longmapsto n+1,
\]
such that
\[
    s_!W_\omega\cong W_\omega.
\]
\end{definition}

\begin{definition}[Sequential finite-flat cofinality]
\label{def:sequential-ff-cofinality}
The fuzzy cosmos \(\mathbb V\) is \emph{sequentially finite-flat cofinal} if
there is a sequential finite-flat approximation shape
\[
    (\omega_{\mathbb V},W_\omega,s)
\]
such that:

\begin{enumerate}
    \item \(s\) is \(W_\omega\)-cofinal:
    \[
        s_!W_\omega\cong W_\omega;
    \]

    \item \(s\) is bilimit-cofinal for the pointed sequential projection-pair
    systems considered below;

    \item the canonical tails and canonical external-product common-successor
    comparisons are cofinal in the senses of
    Definitions~\ref{def:tail-cofinal-approximation-shape},
    \ref{def:common-successor-cofinal-shape}, and
    \ref{def:linear-cofinal-ff-approximation-shape}.
\end{enumerate}
\end{definition}

\begin{definition}[Pointed projection-pair endopseudofunctor]
\label{def:pointed-projection-pair-endopseudofunctor}
A \emph{pointed projection-pair endopseudofunctor} is a pseudofunctor
\[
    F:
    \mathbf{FDom}^{\triangleleft}_{\mathbb V,\bot}
    \longrightarrow
    \mathbf{FDom}^{\triangleleft}_{\mathbb V,\bot}.
\]
\end{definition}

\begin{definition}[Pointed initial iteration]
\label{def:pointed-initial-iteration}
Let
\[
    F:
    \mathbf{FDom}^{\triangleleft}_{\mathbb V,\bot}
    \to
    \mathbf{FDom}^{\triangleleft}_{\mathbb V,\bot}
\]
be a pointed projection-pair endopseudofunctor.

The \emph{pointed initial \(F\)-iteration} is the sequential split
projection-pair system
\[
    D:\omega_{\mathbb V}
      \rightsquigarrow
      \mathbf{FDom}^{\triangleleft}_{\mathbb V,\bot}
\]
defined on objects by
\[
    D_0=\mathbf 1_\bot,
    \qquad
    D_{n+1}=F(D_n).
\]
The adjacent transition projection pair
\[
    D_n\triangleleft D_{n+1}
\]
is obtained by applying \(F^n\) to the canonical projection pair
\[
    \mathbf 1_\bot\triangleleft F(\mathbf 1_\bot).
\]
For \(n\leq m\), the transition projection pair
\[
    D_n\triangleleft D_m
\]
is the composite of the adjacent transition projection pairs.
\end{definition}

\begin{definition}[Pointed generated invariant]
\label{def:pointed-generated-invariant}
Let
\[
    F:
    \mathbf{FDom}^{\triangleleft}_{\mathbb V,\bot}
    \to
    \mathbf{FDom}^{\triangleleft}_{\mathbb V,\bot}
\]
be a pointed projection-pair endopseudofunctor, and let
\[
    D:\omega_{\mathbb V}
      \rightsquigarrow
      \mathbf{FDom}^{\triangleleft}_{\mathbb V,\bot}
\]
be its pointed initial iteration.

A \emph{pointed generated invariant} of \(F\) consists of a pointed fuzzy domain
\(X\), an equivalence
\[
    X\simeq F(X),
\]
and a compatible projection-pair cone
\[
    (\varepsilon_n,\pi_n):D_n\triangleleft X
    \qquad(n\in\omega_{\mathbb V})
\]
which exhibits \(X\) as the \(W_\omega\)-bilimit of the pointed initial
iteration \(D\).
\end{definition}

\begin{theorem}[Pointed finite-flat minimal-invariant theorem]
\label{thm:pointed-ff-minimal-invariant}
Assume that \(\mathbb V\) is sequentially finite-flat cofinal.  Let
\[
    F:
    \mathbf{FDom}^{\triangleleft}_{\mathbb V,\bot}
    \to
    \mathbf{FDom}^{\triangleleft}_{\mathbb V,\bot}
\]
be a pointed projection-pair endopseudofunctor preserving \(W_\omega\)-bilimits
of pointed sequential projection-pair systems.

Let
\[
    D:\omega_{\mathbb V}
      \rightsquigarrow
      \mathbf{FDom}^{\triangleleft}_{\mathbb V,\bot}
\]
be the pointed initial \(F\)-iteration.  Suppose \(D\) has a
\(W_\omega\)-bilimit
\[
    D_\infty.
\]
Then there is a canonical equivalence of pointed fuzzy domains
\[
    D_\infty\simeq F(D_\infty).
\]

Moreover, \(D_\infty\) is unique as a pointed generated invariant: if \(X\) is
any other pointed generated invariant of \(F\), then there is an equivalence
\[
    X\simeq D_\infty
\]
unique up to invertible compatible \(2\)-cell.
\end{theorem}

\begin{proof}
Let
\[
    s:\omega_{\mathbb V}\to\omega_{\mathbb V}
\]
be the successor endofunctor.  By construction of the pointed initial
iteration,
\[
    Ds\cong FD
\]
as split projection-pair systems.

Since \(s\) is \(W_\omega\)-cofinal and bilimit-cofinal for \(D\),
Theorem~\ref{thm:shifted-bilimits-ff} shows that
\[
    D_\infty
\]
is a \(W_\omega\)-bilimit of
\[
    Ds.
\]
Since \(F\) preserves \(W_\omega\)-bilimits,
\[
    F(D_\infty)
\]
is a \(W_\omega\)-bilimit of
\[
    FD.
\]
Using
\[
    Ds\cong FD,
\]
both \(D_\infty\) and \(F(D_\infty)\) are \(W_\omega\)-bilimits of the same
system.  By uniqueness,
\[
    D_\infty\simeq F(D_\infty).
\]

If \(X\) is another pointed generated invariant, then \(X\) is also a
\(W_\omega\)-bilimit of the pointed initial iteration \(D\).  Hence uniqueness
of \(W_\omega\)-bilimits gives
\[
    X\simeq D_\infty
\]
uniquely up to invertible compatible \(2\)-cell.
\end{proof}

\section{Pointed finite-flat generated algebraic compactness}
\label{sec:pointed-ff-algebraic-compactness}

The minimal-invariant theorem gives an equivalence
\[
    D_\infty\simeq F(D_\infty).
\]
To express the algebraic-compactness content, we record the associated
initial-algebra and final-coalgebra behavior inside the generated
projection-pair fragment.

Throughout this section \(F\) is assumed to act locally on fuzzy-domain
morphisms and to preserve finite-flat colimits on hom fuzzy domains.

\begin{definition}[Locally finite-flat-continuous pointed endopseudofunctor]
\label{def:locally-ff-continuous-endopseudofunctor}
A pointed projection-pair endopseudofunctor
\[
    F:
    \mathbf{FDom}^{\triangleleft}_{\mathbb V,\bot}
    \to
    \mathbf{FDom}^{\triangleleft}_{\mathbb V,\bot}
\]
is \emph{locally finite-flat-continuous} if it is the restriction to projection
pairs of a pseudofunctor
\[
    \widetilde F:
    \mathbf{FDom}^{(2)}_{\mathbb V,\bot}
    \to
    \mathbf{FDom}^{(2)}_{\mathbb V,\bot}
\]
on the locally ordinary \(2\)-category of pointed fuzzy domains and
fuzzy-domain morphisms, and for every pair of pointed fuzzy domains \(X,Y\),
the induced hom-action
\[
    \widetilde F_{X,Y}:
    [X,Y]_{\mathrm{ff}}
    \to
    [FX,FY]_{\mathrm{ff}}
\]
preserves finite-flat weighted colimits.
\end{definition}

\begin{definition}[Generated projection-pair \(F\)-algebra]
\label{def:generated-projection-pair-F-algebra}
Let
\[
    F:
    \mathbf{FDom}^{\triangleleft}_{\mathbb V,\bot}
    \to
    \mathbf{FDom}^{\triangleleft}_{\mathbb V,\bot}
\]
be locally finite-flat-continuous, and let
\[
    D:\omega_{\mathbb V}\rightsquigarrow
    \mathbf{FDom}^{\triangleleft}_{\mathbb V,\bot}
\]
be its pointed initial iteration.

A \emph{generated projection-pair \(F\)-algebra} is a pointed fuzzy domain
\(X\) equipped with a projection pair
\[
    a=(a_e,a_p):FX\triangleleft X.
\]
The embedding
\[
    a_e:FX\to X
\]
is the algebra structure, and the projection
\[
    a_p:X\to FX
\]
is the associated coalgebra structure.

The projection pair \(a\) generates a compatible cone
\[
    \lambda^a_n:D_n\triangleleft X
\]
by recursion:
\[
    \lambda^a_0:\mathbf 1_\bot\triangleleft X
\]
is the canonical projection pair, and
\[
    \lambda^a_{n+1}
    =
    a\circ F(\lambda^a_n).
\]
\end{definition}

\begin{definition}[Generated algebra and coalgebra categories]
\label{def:generated-algebra-coalgebra-categories}
The locally ordinary \(2\)-category of generated projection-pair \(F\)-algebras
is denoted
\[
    \operatorname{Alg}^{\mathrm{gen}}_{\triangleleft}(F).
\]
The same data may be read coalgebraically: the projection half
\[
    a_p:X\to FX
\]
of
\[
    a=(a_e,a_p):FX\triangleleft X
\]
is an \(F\)-coalgebra structure.  With the same objects, but with the
coalgebra-half compatibility emphasized, we write
\[
    \operatorname{Coalg}^{\mathrm{gen}}_{\triangleleft}(F).
\]
\end{definition}

\begin{construction}[The canonical generated algebra-coalgebra]
\label{con:canonical-generated-algebra-coalgebra}
Let
\[
    D_\infty
\]
be the \(W_\omega\)-bilimit of the pointed initial iteration.  By
Theorem~\ref{thm:pointed-ff-minimal-invariant}, there is an equivalence
\[
    D_\infty\simeq F(D_\infty).
\]
Choosing the corresponding projection-pair representative gives a canonical
generated projection-pair algebra
\[
    a_\infty=(a_{\infty,e},a_{\infty,p}):
    F(D_\infty)\triangleleft D_\infty.
\]
\end{construction}

\begin{lemma}[Generated folds from the bilimit]
\label{lem:generated-folds-from-bilimit}
Let
\[
    a=(a_e,a_p):FX\triangleleft X
\]
be a generated projection-pair \(F\)-algebra.  Then the generated cone
\[
    \lambda^a_n:D_n\triangleleft X
\]
induces a projection pair
\[
    \operatorname{fold}_a:D_\infty\triangleleft X.
\]
Moreover, its embedding half is an algebra morphism
\[
    (D_\infty,a_\infty)\to(X,a),
\]
and its projection half is a coalgebra morphism
\[
    (X,a_p)\to(D_\infty,a_{\infty,p}).
\]
\end{lemma}

\begin{proof}
Since \(D_\infty\) represents compatible cones over the pointed initial
iteration, the generated compatible cone
\[
    \lambda^a_n:D_n\triangleleft X
\]
is represented by a projection pair
\[
    \operatorname{fold}_a:D_\infty\triangleleft X.
\]
The algebra and coalgebra compatibility equations are checked on finite
approximants.  At stage \(n+1\), both sides are obtained from stage \(n\) by
applying \(F\) and then composing with \(a\).  Detection by the bilimit
representation gives the required compatible \(2\)-cells.
\end{proof}

\begin{lemma}[Uniqueness of generated folds]
\label{lem:uniqueness-generated-folds}
Let
\[
    a=(a_e,a_p):FX\triangleleft X
\]
be a generated projection-pair \(F\)-algebra.  Any two generated algebra
morphisms
\[
    (D_\infty,a_\infty)\to(X,a)
\]
are uniquely isomorphic by an invertible compatible \(2\)-cell.  Dually, any
two generated coalgebra morphisms
\[
    (X,a_p)\to(D_\infty,a_{\infty,p})
\]
are uniquely isomorphic by an invertible compatible \(2\)-cell.
\end{lemma}

\begin{proof}
The restrictions of a generated algebra morphism to finite approximants are
forced recursively.  At stage \(0\) they are the canonical map from
\(\mathbf 1_\bot\).  If the restrictions agree at stage \(n\), algebra
compatibility forces agreement at stage \(n+1\).  Thus any two such morphisms
represent the same compatible cone, and uniqueness follows from the
representing property of the bilimit.  The coalgebra statement is the mate
dual.
\end{proof}

\begin{theorem}[Pointed finite-flat generated algebraic compactness]
\label{thm:pointed-ff-algebraic-compactness}
Assume that \(\mathbb V\) is sequentially finite-flat cofinal.  Let
\[
    F:
    \mathbf{FDom}^{\triangleleft}_{\mathbb V,\bot}
    \to
    \mathbf{FDom}^{\triangleleft}_{\mathbb V,\bot}
\]
be a pointed projection-pair endopseudofunctor which preserves
\(W_\omega\)-bilimits of pointed sequential projection-pair systems and is
locally finite-flat-continuous.

Let
\[
    D:\omega_{\mathbb V}
      \rightsquigarrow
      \mathbf{FDom}^{\triangleleft}_{\mathbb V,\bot}
\]
be the pointed initial \(F\)-iteration, and suppose it has a
\(W_\omega\)-bilimit
\[
    D_\infty.
\]
Then the canonical projection pair
\[
    a_\infty=(a_{\infty,e},a_{\infty,p}):
    F(D_\infty)\triangleleft D_\infty
\]
exhibits \(D_\infty\) as the generated algebraically compact solution of \(F\):

\begin{enumerate}
    \item the embedding half
    \[
        a_{\infty,e}:F(D_\infty)\to D_\infty
    \]
    is initial in
    \[
        \operatorname{Alg}^{\mathrm{gen}}_{\triangleleft}(F);
    \]

    \item the projection half
    \[
        a_{\infty,p}:D_\infty\to F(D_\infty)
    \]
    is final in
    \[
        \operatorname{Coalg}^{\mathrm{gen}}_{\triangleleft}(F);
    \]

    \item the algebra and coalgebra structures are inverse halves of the same
    projection-pair equivalence
    \[
        F(D_\infty)\triangleleft D_\infty.
    \]
\end{enumerate}
\end{theorem}

\begin{proof}
Theorem~\ref{thm:pointed-ff-minimal-invariant} gives
\[
    D_\infty\simeq F(D_\infty).
\]
Choosing the corresponding projection-pair representative gives
\[
    a_\infty:F(D_\infty)\triangleleft D_\infty.
\]

For any generated projection-pair \(F\)-algebra
\[
    a:FX\triangleleft X,
\]
Lemma~\ref{lem:generated-folds-from-bilimit} gives a generated fold
\[
    \operatorname{fold}_a:D_\infty\triangleleft X.
\]
Its embedding half is an algebra morphism, and its projection half is a
coalgebra morphism.  Lemma~\ref{lem:uniqueness-generated-folds} gives
uniqueness up to invertible compatible \(2\)-cell.  Hence the embedding half of
\(a_\infty\) is initial in the generated algebra category, and the projection
half is final in the generated coalgebra category.

The two structures are inverse halves of the same projection-pair equivalence
by construction.
\end{proof}

\section{The crisp separated case}
\label{sec:crisp-separated-ff-rde}

For
\[
    \Lambda=2,
\]
finite-flat weights are ordinary ideals.  On separated fuzzy domains, strict
algebras for the ideal monad are precisely dcpos, and fuzzy-domain morphisms
are precisely Scott-continuous maps.

Let
\[
    \omega=0\leq1\leq2\leq\cdots
\]
and let
\[
    W_\omega:\omega^{\op}\to 2
\]
be the constant-true weight:
\[
    W_\omega(n)=1.
\]
This is the ordinary directed ideal \(\omega\subseteq\omega\).

For a dcpo system
\[
    D_0\triangleleft D_1\triangleleft D_2\triangleleft\cdots
\]
of projection pairs, a \(W_\omega\)-bilimit is exactly the usual bilimit:
there are projection pairs
\[
    (\varepsilon_n,\pi_n):D_n\triangleleft D_\infty
\]
such that the approximants satisfy
\[
    \bigvee_{n<\omega}
        \varepsilon_n\pi_n
    =
    \id_{D_\infty}
\]
in the dcpo of Scott-continuous endomorphisms of \(D_\infty\).

In the crisp case, the ordinary diagonal
\[
    \Delta_\omega:\omega\to\omega\times\omega,
    \qquad
    n\longmapsto(n,n),
\]
is a monotone map.  It is cofinal for the ideal weight because for every pair
\[
    (m,n)\in\omega\times\omega
\]
there is
\[
    k\geq m,n.
\]
Hence
\[
    (\Delta_\omega)_!W_\omega
    =
    W_\omega\boxtimes W_\omega.
\]
This diagonal is the cartesian feature that permits same-index function-space
equations.

\begin{theorem}[Crisp separated same-index function-space collapse]
\label{thm:crisp-separated-same-index-collapse}
Let
\[
    \Lambda=2.
\]
Let
\[
    D_0\triangleleft D_1\triangleleft D_2\triangleleft\cdots
\]
be a sequential dcpo projection-pair system with bilimit \(D_\infty\), and
suppose
\[
    D_{n+1}\cong [D_n,D_n]
\]
as projection-pair systems.  Then
\[
    D_\infty\cong [D_\infty,D_\infty].
\]
\end{theorem}

\begin{proof}
For \(\Lambda=2\), the finite-flat weight
\[
    W_\omega:\omega^{\op}\to 2
\]
is the ordinary directed ideal.  The product weight
\[
    W_\omega\boxtimes W_\omega
\]
is the ordinary directed ideal on \(\omega\times\omega\).

The diagonal
\[
    \Delta_\omega:\omega\to\omega\times\omega
\]
is cofinal, so restriction along \(\Delta_\omega\) preserves the relevant
bilimit.  The external function-space system is
\[
    (m,n)\longmapsto [D_n,D_m].
\]
Restricting along the diagonal gives
\[
    n\longmapsto [D_n,D_n]\cong D_{n+1}.
\]
The successor shift is cofinal, so the shifted chain has the same bilimit
\(D_\infty\).  On the other hand, the standard bilimit theorem for dcpo
projection-pair chains gives that
\[
    [D_\infty,D_\infty]
\]
is the bilimit of the same external function-space system.  By uniqueness of
bilimits,
\[
    D_\infty\cong [D_\infty,D_\infty].
\]
\end{proof}

Similarly, for a locally Scott-continuous bilimit-preserving endofunctor \(T\)
on dcpos, the effectful constructor
\[
    \Phi_T(D)=[D,TD]
\]
gives, after diagonal collapse, the equation
\[
    D_\infty\cong [D_\infty,TD_\infty].
\]
For a bilimit-preserving coeffect functor \(G\), the coeffectful form is
\[
    D_\infty\cong [GD_\infty,D_\infty],
\]
and for mixed effect/coeffect structure,
\[
    D_\infty\cong [GD_\infty,TD_\infty].
\]

In the crisp separated case the Kelly tensor product of fuzzy domains is the
cartesian product of dcpos, so the application maps have their usual cartesian
types:
\[
    D_\infty\times D_\infty\to D_\infty,
\]
\[
    D_\infty\times D_\infty\to TD_\infty,
\]
and
\[
    D_\infty\times GD_\infty\to TD_\infty
\]
in the pure, effectful, and mixed cases respectively.

\begin{proposition}[Crisp separated pointed minimal invariants]
\label{prop:crisp-separated-pointed-minimal-invariants}
For
\[
    \Lambda=2,
\]
the pointed finite-flat minimal-invariant and generated algebraic-compactness
theorems specialize to the classical pointed-dcpo bilimit construction.

More explicitly:

\begin{enumerate}
    \item pointed fuzzy domains are precisely pointed dcpos;

    \item fuzzy projection pairs are precisely embedding-projection pairs
    \[
        e:D\to E,
        \qquad
        p:E\to D,
        \qquad
        pe=\id_D,
        \qquad
        ep\leq\id_E;
    \]

    \item the finite-flat weight
    \[
        W_\omega:\omega^{\op}\to 2
    \]
    is the ordinary directed ideal on \(\omega\);

    \item the \(W_\omega\)-bilimit approximation identity becomes
    \[
        \bigvee_{n<\omega}\varepsilon_n\pi_n
        =
        \id_{D_\infty}
    \]
    in the dcpo of Scott-continuous endomorphisms of \(D_\infty\);

    \item a locally finite-flat-continuous endopseudofunctor becomes a locally
    Scott-continuous endofunctor on the projection-pair category of pointed
    dcpos;

    \item the canonical projection pair
    \[
        F(D_\infty)\triangleleft D_\infty
    \]
    gives the generated coincidence of the initial algebra and final coalgebra.
\end{enumerate}
\end{proposition}

\begin{proof}
For \(\Lambda=2\), finite-flat weights are ordinary ideals and strict algebras
for the ideal monad on separated preorders are precisely dcpos.  An enriched
initial object is exactly a least element, so pointed fuzzy domains are pointed
dcpos.

A fuzzy projection pair
\[
    (e,p):D\triangleleft E
\]
is an adjunction
\[
    e\dashv p
\]
with invertible unit.  In order-enriched notation this is precisely
\[
    pe=\id_D,
    \qquad
    ep\leq\id_E.
\]
Thus projection pairs are the usual embedding-projection pairs.

The weight
\[
    W_\omega:\omega^{\op}\to 2
\]
is constant true, hence is the ordinary directed ideal on \(\omega\).  Therefore
a \(W_\omega\)-weighted colimit in an internal hom dcpo is an ordinary directed
join.  The approximation identity
\[
    W_\omega*(n\mapsto\varepsilon_n\pi_n)
    \cong
    \id_{D_\infty}
\]
therefore becomes
\[
    \bigvee_{n<\omega}\varepsilon_n\pi_n
    =
    \id_{D_\infty}.
\]

The generated projection-pair algebra
\[
    F(D_\infty)\triangleleft D_\infty
\]
has embedding half
\[
    F(D_\infty)\to D_\infty
\]
and projection half
\[
    D_\infty\to F(D_\infty).
\]
These are the classical inverse algebra and coalgebra maps obtained from the
bilimit of the initial embedding-projection chain.
\end{proof}

\section{Summary}
\label{sec:summary-ff-rde}

Recursive domain equations in fuzzy domains are governed by the finite-flat
doctrine.  The primitive approximation objects are finite-flat approximation
shapes
\[
    (A,W),
    \qquad
    W\in\operatorname{Idl}_{\mathbb V}(A).
\]
Cofinality of weights is expressed by lower direct image:
\[
    h:(A,W)\to(B,U)
    \quad\Longleftrightarrow\quad
    h_!W\cong U.
\]
An iteration shape is a finite-flat approximation shape equipped with a
weight-cofinal self-map
\[
    s:A\to A,
    \qquad
    s_!W\cong W.
\]

Tails are canonical support-comma categories
\[
    \operatorname{Tail}_A(a_0),
\]
equipped with their canonical projection
\[
    \tau_{a_0}:\operatorname{Tail}_A(a_0)\to A.
\]
Tail-cofinality says that
\[
    \tau_{a_0}^*W
\]
is finite-flat and
\[
    (\tau_{a_0})_!\tau_{a_0}^{*}W\cong W.
\]
Common successors are encoded by
\[
    \operatorname{Com}(P),
\]
whose objects are tuples
\[
    (i,j,k,u:i\to k,v:j\to k).
\]
The canonical maps
\[
    c_P:\operatorname{Com}(P)\to P\tensor P,
    \qquad
    m_P:\operatorname{Com}(P)\to P
\]
are cofinal when
\[
    (c_P)_!m_P^*R\cong R\boxtimes R,
    \qquad
    (m_P)_!m_P^*R\cong R.
\]

For bilimits, weight-cofinality is paired with bilimit-cofinality: restriction
along the indexing map must preserve compatible projection-pair cone
categories.  With this distinction in place, a split projection-pair system
\[
    D:A\rightsquigarrow\mathbf{FDom}_{\mathbb V}^{\triangleleft}
\]
satisfying
\[
    Ds\cong FD
\]
has a generated solution whenever \(D\) has a \(W\)-bilimit, \(s\) is
bilimit-cofinal for \(D\), and \(F\) preserves that bilimit:
\[
    D_\infty\simeq F(D_\infty).
\]

The general internal-hom theorem is external-product-indexed.  If
\[
    X_\infty
    \qquad\text{and}\qquad
    Y_\infty
\]
are \(W\)-bilimits of split systems \(X\) and \(Y\), then
\[
    [X_\infty,Y_\infty]_{\mathrm{ff}}
\]
has external-product approximants
\[
    (a,a')\longmapsto [X_{a'},Y_a]_{\mathrm{ff}},
\]
and these satisfy
\[
    (W\boxtimes W)*
    \bigl((a,a')\mapsto E_{a,a'}\Pi_{a,a'}\bigr)
    \cong
    \id_{[X_\infty,Y_\infty]_{\mathrm{ff}}}.
\]
Compatible cones over the external-product system reconstruct canonically to
lax application data.  They do not generally reconstruct to projection pairs,
because the common-successor comparison supplies
\[
    \rho\delta\Rightarrow\id
\]
rather than an invertible unit
\[
    \id\Rightarrow\rho\delta.
\]
This is exactly the resource-sensitive obstruction caused by the absence of a
canonical diagonal
\[
    A\to A\tensor A.
\]

In the pointed case, a sequential finite-flat approximation shape
\[
    (\omega_{\mathbb V},W_\omega,s)
\]
replaces the ordinary chain
\[
    0\leq1\leq2\leq\cdots.
\]
The pointed initial iteration
\[
    D_0=\mathbf 1_\bot,
    \qquad
    D_{n+1}=F(D_n)
\]
has a \(W_\omega\)-bilimit \(D_\infty\).  If \(F\) preserves this bilimit, then
\[
    D_\infty\simeq F(D_\infty).
\]
If \(F\) is locally finite-flat-continuous, the resulting projection pair
\[
    F(D_\infty)\triangleleft D_\infty
\]
is the generated algebraically compact solution: its embedding half is the
initial generated algebra, its projection half is the final generated
coalgebra, and the two structures are inverse halves of the same projection-pair
equivalence.

In the crisp separated case, taking the ordinary ideal
\[
    W_\omega=\omega\subseteq\omega
\]
recovers the standard bilimit method for solving recursive domain equations of
dcpos.  The ordinary reflexive-domain equation
\[
    D_\infty\cong [D_\infty,D_\infty]
\]
is recovered precisely because the cartesian diagonal
\[
    \omega\to\omega\times\omega
\]
collapses the external-product function-space approximation to a same-index
chain.

    \printbibliography
\end{document}

